\documentclass[11pt]{article}
\usepackage{fontspec}
\directlua{luaotfload.add_fallback("hebfb",{"DejaVu Sans:mode=harf"})}
\usepackage{polyglossia}
\setdefaultlanguage{hebrew}
\setotherlanguage{english}
\setmainfont{Noto Sans Hebrew}[Script=Hebrew,RawFeature={fallback=hebfb}]
\newfontfamily\codefont{DejaVu Sans Mono}[RawFeature={fallback=hebfb}]
\usepackage[a4paper,margin=18mm]{geometry}
\usepackage[table]{xcolor}
\usepackage{mdframed}
\usepackage{tabularx}
\usepackage{xltabular}
\usepackage{array}
\usepackage{enumitem}
\usepackage{fancyhdr}
\setlist{leftmargin=1.6em,itemsep=1pt,topsep=2pt}
% colors
\definecolor{h1}{HTML}{1E3A8A}\definecolor{h2}{HTML}{1E40AF}\definecolor{h3}{HTML}{0F172A}
\definecolor{subgray}{HTML}{64748B}\definecolor{hdrbg}{HTML}{E0E7FF}
\definecolor{cfg}{HTML}{E2E8F0}\definecolor{cbg}{HTML}{0F172A}\definecolor{cbold}{HTML}{7DD3FC}\definecolor{ccom}{HTML}{94A3B8}
\definecolor{metabg}{HTML}{F1F5F9}
\newcommand{\enLR}[1]{\textenglish{#1}}
\newcommand{\code}[1]{\textenglish{\codefont #1}}
\newcommand{\chapterheading}[1]{\vspace{2pt}{\LARGE\bfseries\color{h1}#1\par}\vskip2pt{\color{h1}\hrule height 2pt}\vskip6pt}
\newcommand{\sectionheading}[1]{\vskip10pt\penalty-200{\large\bfseries\color{h2}#1\par}\vskip3pt}
\newcommand{\subheading}[1]{\vskip6pt{\bfseries\color{h3}#1\par}\vskip2pt}
\newcommand{\boxtitle}[1]{{\bfseries\color{boxti}#1}}
% box environments (right border, tinted bg)
\newmdenv[topline=false,bottomline=false,leftline=false,rightline=true,linewidth=2pt,innermargin=0pt,skipabove=6pt,skipbelow=6pt,innertopmargin=6pt,innerbottommargin=6pt,innerleftmargin=8pt,innerrightmargin=8pt]{genericbox}
\definecolor{faqbg}{HTML}{ECFDF5}\definecolor{faqfr}{HTML}{10B981}\definecolor{faqti}{HTML}{065F46}
\definecolor{misbg}{HTML}{FEF2F2}\definecolor{misfr}{HTML}{EF4444}\definecolor{misti}{HTML}{991B1B}
\definecolor{tipbg}{HTML}{FFFBEB}\definecolor{tipfr}{HTML}{F59E0B}\definecolor{tipti}{HTML}{92400E}
\definecolor{exabg}{HTML}{EFF6FF}\definecolor{exafr}{HTML}{3B82F6}\definecolor{exati}{HTML}{1E3A8A}
\definecolor{defbg}{HTML}{F5F3FF}\definecolor{deffr}{HTML}{8B5CF6}\definecolor{defti}{HTML}{5B21B6}
\definecolor{stobg}{HTML}{FDF4FF}\definecolor{stofr}{HTML}{D946EF}\definecolor{stoti}{HTML}{86198F}
\definecolor{exbg}{HTML}{F0FDFA}\definecolor{exfr}{HTML}{14B8A6}\definecolor{exti}{HTML}{115E59}
\newenvironment{faqbox}{\colorlet{boxti}{faqti}\begin{mdframed}[backgroundcolor=faqbg,linecolor=faqfr,rightline=true,leftline=false,topline=false,bottomline=false,linewidth=2pt,innertopmargin=6pt,innerbottommargin=6pt,innerleftmargin=8pt,innerrightmargin=8pt,skipabove=6pt,skipbelow=6pt]}{\end{mdframed}}
\newenvironment{mistakebox}{\colorlet{boxti}{misti}\begin{mdframed}[backgroundcolor=misbg,linecolor=misfr,rightline=true,leftline=false,topline=false,bottomline=false,linewidth=2pt,innertopmargin=6pt,innerbottommargin=6pt,innerleftmargin=8pt,innerrightmargin=8pt,skipabove=6pt,skipbelow=6pt]}{\end{mdframed}}
\newenvironment{tipbox}{\colorlet{boxti}{tipti}\begin{mdframed}[backgroundcolor=tipbg,linecolor=tipfr,rightline=true,leftline=false,topline=false,bottomline=false,linewidth=2pt,innertopmargin=6pt,innerbottommargin=6pt,innerleftmargin=8pt,innerrightmargin=8pt,skipabove=6pt,skipbelow=6pt]}{\end{mdframed}}
\newenvironment{exambox}{\colorlet{boxti}{exati}\begin{mdframed}[backgroundcolor=exabg,linecolor=exafr,rightline=true,leftline=false,topline=false,bottomline=false,linewidth=2pt,innertopmargin=6pt,innerbottommargin=6pt,innerleftmargin=8pt,innerrightmargin=8pt,skipabove=6pt,skipbelow=6pt]}{\end{mdframed}}
\newenvironment{defbox}{\colorlet{boxti}{defti}\begin{mdframed}[backgroundcolor=defbg,linecolor=deffr,rightline=true,leftline=false,topline=false,bottomline=false,linewidth=2pt,innertopmargin=6pt,innerbottommargin=6pt,innerleftmargin=8pt,innerrightmargin=8pt,skipabove=6pt,skipbelow=6pt]}{\end{mdframed}}
\newenvironment{storybox}{\colorlet{boxti}{stoti}\begin{mdframed}[backgroundcolor=stobg,linecolor=stofr,rightline=true,leftline=false,topline=false,bottomline=false,linewidth=2pt,innertopmargin=6pt,innerbottommargin=6pt,innerleftmargin=8pt,innerrightmargin=8pt,skipabove=6pt,skipbelow=6pt]}{\end{mdframed}}
\newenvironment{exbox}{\colorlet{boxti}{exti}\begin{mdframed}[backgroundcolor=exbg,linecolor=exfr,rightline=true,leftline=false,topline=false,bottomline=false,linewidth=2pt,innertopmargin=6pt,innerbottommargin=6pt,innerleftmargin=8pt,innerrightmargin=8pt,skipabove=6pt,skipbelow=6pt]}{\end{mdframed}}
\newenvironment{metabox}{\begin{mdframed}[backgroundcolor=metabg,linewidth=0pt,innertopmargin=5pt,innerbottommargin=5pt,innerleftmargin=8pt,innerrightmargin=8pt,skipabove=4pt,skipbelow=8pt]}{\end{mdframed}}
\newmdenv[backgroundcolor=cbg,linewidth=0pt,innertopmargin=6pt,innerbottommargin=6pt,innerleftmargin=8pt,innerrightmargin=8pt,skipabove=6pt,skipbelow=8pt]{codebox}
\setlength{\parindent}{0pt}\setlength{\parskip}{4pt}
\renewcommand{\thepage}{\enLR{\arabic{page}}}
\pagestyle{fancy}\fancyhf{}\fancyfoot[C]{\thepage}\renewcommand{\headrulewidth}{0pt}
\usepackage[hidelinks,unicode]{hyperref}
\usepackage{bookmark}
\begin{document}
\sloppy
\begin{titlepage}\centering\vspace*{4cm}
{\Huge\bfseries\color{h1} אבטחת מידע\par}\vskip10pt
{\Huge\bfseries\color{h1} ספר הקורס המלא\par}\vskip20pt
{\Large תכנית הלימודים – מגמת תקשוב, משרד החינוך\par}\vskip6pt
{\large\color{subgray} 9 פרקים \enLR{$\cdot$} חומר עיון, סיפורים, תרגילים ושאלות בסגנון בגרות (שאלון \enLR{735001})\par}
\vfill{\large\color{subgray} מדריך למורה\par}\end{titlepage}
\tableofcontents\clearpage
\clearpage
\phantomsection\addcontentsline{toc}{section}{מבוא: איך להשתמש בחומר הזה}

\sloppy


\chapterheading{מבוא: איך להשתמש בחומר הזה}{\small\color{subgray}קראו את הפרק הזה לפני שמתחילים ללמד\par}\vskip4pt

\sectionheading{מה יש כאן}

החומר בנוי משני חלקים שמשלימים זה את זה:\par

\begin{itemize}
\item \textbf{מדריך למורה (ה-\enLR{PDF} הזה)} – פרק לכל פרק בתכנית הלימודים. כל פרק כולל: מטרות ותכנון שעות, את החומר העיוני עצמו בעומק שמאפשר לענות על שאלות של תלמידים, מדריך פקודות \enLR{CLI}, תיבות "שאלת תלמיד נפוצה" עם תשובה מוכנה, תיבות "טעות נפוצה", תיבות "טיפ להוראה", ושאלות בסגנון בגרות עם פתרון מלא.
\item \textbf{מצגות (תיקיית \enLR{slides})} – קובץ \enLR{HTML} אחד לכל פרק, פותחים בדפדפן ומקרינים. עובד ללא אינטרנט. לחיצה על \enLR{N} פותחת הערות למורה בתחתית המסך (מה להגיד, איזו שאלה לשאול את הכיתה). לחיצה על \enLR{F} – מסך מלא. חיצים / רווח – מעבר שקף. בשקפי "שאלה" התשובה מופיעה רק בלחיצה נוספת – כך אפשר לתת לכיתה לנסות קודם.
\end{itemize}

\begin{tipbox}\boxtitle{טיפ: סדר עבודה מומלץ לכל שבוע}\par  \enLR{1.} קראו את פרק המדריך הרלוונטי \enLR{(30}–\enLR{60} דקות). \enLR{2.} עברו על המצגת עם הערות המורה פתוחות \enLR{(N). 3.} הריצו בעצמכם את פקודות ה-\enLR{CLI} ב-\enLR{Packet Tracer} לפני השיעור – תלמידים ישאלו "מה קורה אם…" ועדיף שכבר ראיתם. \enLR{4.} אחרי השיעור – תנו לתלמידים את שאלות "סגנון בגרות" מהמדריך כתרגול.\end{tipbox}

\sectionheading{על תכנית הלימודים ועל שאלון הבגרות}

תכנית הלימודים הרשמית \enLR{(120} שעות: \enLR{96} עיוני + \enLR{24} מעשי) מבוססת על קורס \enLR{CCNA Security} של סיסקו מהדור הקודם. שימו לב לשלושה דברים:\par

\begin{enumerate}
\item \textbf{\enLR{SDM / CCP} הם כלים שיצאו משימוש.} התכנית מזכירה הרבה "הגדרה על ידי \enLR{SDM"} (ממשק גרפי של סיסקו לנתבים). \enLR{SDM} הוחלף ב-\enLR{CCP} וגם הוא הופסק. \textbf{שאלון הבגרות בודק אך ורק \enLR{CLI.}} לכן במדריך הזה ההגדרות הן ב-\enLR{CLI}, ואת \enLR{SDM} אנו מזכירים רק כמושג ("קיים גם ממשק גרפי").
\item \textbf{בתכנית הרשמית יש שגיאת העתקה:} רשימת הנושאים של פרק \enLR{3 (AAA)} זהה לרשימה של פרק \enLR{2.} במדריך הזה פרק \enLR{3} בנוי לפי \emph{המטרות} של פרק \enLR{3} (מודל \enLR{AAA, ACS, TACACS+, RADIUS)} – זה גם מה שנשאל בבגרות.
\item \textbf{מחצית משאלון הבגרות \enLR{(50} נקודות) היא "רשתות תקשורת"} – סאבנטים, \enLR{OSPF, STP, DHCP}, ניתוב סטטי, \enLR{VLAN} וכו'. החומר הזה \underline{לא} נמצא בתכנית "אבטחת מידע" והוא נלמד בקורס הרשתות. במדריך הזה נגענו בו רק היכן שהוא נוגע ישירות לאבטחה (למשל \enLR{STP} בפרק \enLR{6, ACL} בפרק \enLR{4}, פקודות \enLR{line vty} בפרק \enLR{2).} מומלץ לוודא שהתלמידים מגיעים עם רקע ב-\enLR{CCNA 1}–\enLR{2.}
\end{enumerate}

\sectionheading{תכנון שנתי – \enLR{30} שבועות, \enLR{4} שעות שבועיות}

\begin{xltabular}{\linewidth}{|>{\setRTL\raggedleft\arraybackslash}X|>{\setRTL\raggedleft\arraybackslash}X|>{\setRTL\raggedleft\arraybackslash}X|>{\setRTL\raggedleft\arraybackslash}X|>{\setRTL\raggedleft\arraybackslash}X|}
\hline
\rowcolor{hdrbg}\textbf{מעשי} & \textbf{מה מכסים} & \textbf{שעות בתכנית} & \textbf{פרק} & \textbf{שבועות} \\
\hline
\endhead
\enLR{Wireshark} / דמו \enLR{phishing} & מושגי יסוד \enLR{(CIA)}, אבולוציית האיומים, סוגי תוקפים, נוזקות, שלבי תקיפה, \enLR{DoS/DDoS}, הנדסה חברתית, ברוט-פורס. שבוע \enLR{1}: היכרות עם הקורס ומבחן רמה קצר ברשתות. & \enLR{10} & פרק \enLR{1} – מבוא לאיומי רשת & \enLR{1}–\enLR{3} \\
\hline
\enLR{Packet Tracer}: הקשחת נתב + \enLR{SSH} & הקשחת נתב: סיסמאות, \enLR{enable secret, SSH}, רמות הרשאה, \enLR{banner, syslog, NTP, SNMP, AutoSecure, OOB.} & \enLR{12} & פרק \enLR{2} – אבטחת אביזרי רשת & \enLR{4}–\enLR{6} \\
\hline
\enLR{Packet Tracer}: שרת \enLR{RADIUS/TACACS}+ & \enLR{Authentication/Authorization/Accounting, AAA} מקומי, \enLR{RADIUS} מול \enLR{TACACS+, ACS/ISE}, פקודות \enLR{AAA.} & \enLR{14} & פרק \enLR{3} – מודל \enLR{AAA} & \enLR{7}–\enLR{9} \\
\hline
\enLR{Packet Tracer: ACL + ZPF} & \enLR{ACL} סטנדרטי/מורחב/\enLR{named, wildcard}, מיקום \enLR{ACL, ACL} מבוסס זמן/\enLR{reflexive/dynamic}, סוגי חומות אש, \enLR{CBAC, ZPF.} & \enLR{13} & פרק \enLR{4} – חומות אש & \enLR{10}–\enLR{12} \\
\hline
\enLR{Packet Tracer: IOS IPS} & \enLR{IDS} מול \enLR{IPS}, חתימות, \enLR{HIPS} מול \enLR{NIPS}, אזעקות \enLR{(true/false positive), IOS IPS}, ניטור ופתרון תקלות. שבוע \enLR{16}: חזרה + מבחן מחצית. & \enLR{14} & פרק \enLR{5} – \enLR{IDS / IPS} & \enLR{13}–\enLR{16} \\
\hline
\enLR{Packet Tracer: port security + BPDU guard} & התקפות שכבה \enLR{2: MAC flooding, VLAN hopping/DTP, STP, DHCP/ARP spoofing.} הגנות: \enLR{port security, BPDU guard, DHCP snooping, DAI, storm control, SPAN. NAC, IronPort, VoIP/SAN, Evil Twin.} & \enLR{13} & פרק \enLR{6} – אבטחת הרשת המקומית & \enLR{17}–\enLR{19} \\
\hline
תרגול: צופן קיסר, \enLR{hash} בפייתון/אונליין & היסטוריה (קיסר, ויז'נר), \enLR{CIA} + הצפנה, גיבוב \enLR{(MD5/SHA), HMAC}, סימטרי \enLR{(DES/3DES/AES)}, א-סימטרי \enLR{(RSA/DH)}, חתימה דיגיטלית, \enLR{PKI.} & \enLR{14} & פרק \enLR{7} – הצפנה וקריפטולוגיה & \enLR{20}–\enLR{23} \\
\hline
\enLR{Packet Tracer: GRE + IPsec site-to-site} & סוגי \enLR{VPN, GRE, IPsec (AH/ESP, IKE, DH), site-to-site} ב-\enLR{CLI, SSL VPN}, ציוד \enLR{(ASA).} & \enLR{13} & פרק \enLR{8} – \enLR{VPN} & \enLR{24}–\enLR{26} \\
\hline
\enLR{Nmap} במעבדה מבודדת & ניהול סיכונים, סוגי בקרות (מנהלי/פיזי/לוגי), בדיקות חדירה, \enLR{Nmap/Nessus, BCP/DRP, SDLC}, מדיניות ואתיקה. & \enLR{14} & פרק \enLR{9} – ניהול אבטחה מתקדם & \enLR{27}–\enLR{29} \\
\hline
 & פתרון שאלון לדוגמה בתנאי אמת, חזרה על פקודות. & \enLR{3} & חזרה לבגרות & \enLR{30} \\
\hline
\end{xltabular}

\sectionheading{מיפוי שאלון הבגרות (חלק ב' – אבטחת מידע) לפרקי הקורס}

ניתחנו את שאלון \enLR{735001} (אביב תשפ"ג). כך מתפלגים נושאי חלק ב' \enLR{(50} נקודות) בין הפרקים – זה אומר לכם על מה לשים דגש:\par

\begin{xltabular}{\linewidth}{|>{\setRTL\raggedleft\arraybackslash}X|>{\setRTL\raggedleft\arraybackslash}X|>{\setRTL\raggedleft\arraybackslash}X|}
\hline
\rowcolor{hdrbg}\textbf{סוג שאלה} & \textbf{פרק} & \textbf{נושא שנשאל} \\
\hline
\endhead
רב-ברירה על תמונה & \enLR{1} & זיהוי הודעת \enLR{phishing} / הנדסה חברתית \\
\hline
התאמה & \enLR{1} & סוגי תוקפים: \enLR{Script kiddies, Hacktivist, Insider}, פשע מאורגן \\
\hline
רב-ברירה & \enLR{1} & \enLR{Brute force} (מילון), \enLR{Buffer overflow, DoS, Port redirection} \\
\hline
רב-ברירה & \enLR{2} & הפקת מפתחות \enLR{RSA} – משמשים ל-\enLR{SSH} \\
\hline
התאמה & \enLR{3} & שלושת שירותי \enLR{AAA} (הגדרות) \\
\hline
נכון/לא נכון & \enLR{3} & \enLR{RADIUS} מול \enLR{TACACS}+ (הצפנה, \enLR{UDP/TCP}, הרשאות) \\
\hline
השלמת פקודה & \enLR{3} & הפקודה \code{aaa new-model} \\
\hline
רב-ברירה & \enLR{4} & \enLR{ACL} סטנדרטי לחסימת רשת; מיקום \enLR{ACL} וכיוון \enLR{(in/out)} \\
\hline
התאמה & \enLR{5} & \enLR{IDS} מול \enLR{IPS} \\
\hline
רב-ברירה & \enLR{6} & \enLR{Port security} נגד \enLR{MAC table overflow; DTP} ו-\enLR{VLAN hopping; BPDU Guard; Evil Twin} \\
\hline
רב-ברירה & \enLR{7} & \enLR{Diffie-Hellman}, צופן קיסר ו-\enLR{mod 26}, מפתח ציבורי/פרטי, \enLR{AES} מול \enLR{hash} (סודיות) \\
\hline
רב-ברירה & \enLR{8} & \enLR{IPsec} – פועל על \enLR{TCP/IP} \\
\hline
התאמה / רב-ברירה & \enLR{9} & בקרות מנהליות / פיזיות / לוגיות; \enLR{Nmap} \\
\hline
\end{xltabular}

\sectionheading{פתרון מלא – חלק ב' של שאלון הדוגמה \enLR{(735001}, תשפ"ג)}

כדי שתוכלו לפתור עם התלמידים בסוף השנה. ההסברים מופיעים בהרחבה בפרקים.\par

\subheading{שאלה \enLR{4 (20} נק')}

\begin{enumerate}
\item \textbf{א.} הודעת "\enLR{UPS} – \enLR{Last Call} – \enLR{Confirm Now"} – תשובה \textbf{\enLR{4: Social Engineering}} \enLR{(phishing).} סימנים: שולח מ-\enLR{hotmail.com} ולא מדומיין \enLR{UPS}, דחיפות ("\enLR{Last Call")}, כפתור לאישור פרטים.
\item \textbf{ב.} זיהוי ואימות משתמשים = \textbf{\enLR{Authentication}}; ניהול הרשאות = \textbf{\enLR{Authorization}}; רישום למעקב = \textbf{\enLR{Accounting}}.
\item \textbf{ג.} ניטור + חסימה אוטומטית = \textbf{\enLR{IPS}}; ניטור + התראה בלבד = \textbf{\enLR{IDS}}.
\item \textbf{ד.} חסימת \enLR{20.20.20.0/24} ב-\enLR{ACL} סטנדרטי: תשובה \textbf{\enLR{2}: \code{access-list 10 deny 20.20.20.0 0.0.0.255}} (מספר \enLR{1}–\enLR{99} = סטנדרטי; \enLR{wildcard 0.0.0.255).}
\item \textbf{ה.} \enLR{1. RADIUS} מצפין רק את הסיסמה – \textbf{נכון}. \enLR{2. RADIUS} מבצע אימות ולא מתייחס להרשאות – \textbf{נכון} \enLR{(RADIUS} משלב \enLR{authentication+authorization} בתהליך אחד ואינו מאפשר הרשאה נפרדת לכל פקודה). \enLR{3. TACACS}+ משתמש רק ב-\enLR{UDP} – \textbf{לא נכון} \enLR{(TCP 49). 4. TACACS}+ מספק שירות \enLR{AAA} – \textbf{נכון}.
\item \textbf{ו.} הסכמה על מפתח פרטי בערוץ ציבורי – \textbf{\enLR{2: Diffie-Hellman}}.
\end{enumerate}

\subheading{שאלה \enLR{5 (9} נק')}

\begin{enumerate}
\item \textbf{א.} \code{R1(config)\# aaa new-model}
\item \textbf{ב.} מפתחות \enLR{RSA} – \textbf{\enLR{2}: ניתן להשתמש בהם לצורך חיבור \enLR{SSH}} (ללא איפוס).
\item \textbf{ג.} \enLR{modulo 26} – \textbf{\enLR{4}: מספר האפשרויות (האותיות) בשפה האנגלית}; ה-\enLR{mod "}מגלגל" חזרה מ-\enLR{Z} ל-\enLR{A.}
\item \textbf{ד.} הצפנה במפתח ציבורי ⇒ פענוח ב-\textbf{\enLR{1}: מפתח פרטי}.
\item \textbf{ה.} סודיות \enLR{(confidentiality)} של מידע מוצפן – \textbf{\enLR{3: AES}} \enLR{(MD5/HASH} = שלמות; \enLR{WPA-2} הוא פרוטוקול, לא אלגוריתם).
\end{enumerate}

\subheading{שאלה \enLR{6 (12} נק')}

\begin{enumerate}
\item \textbf{א.} \enLR{IPsec} נועד לאבטח את חבילת הפרוטוקולים \textbf{\enLR{2: TCP/IP}} (שכבה \enLR{3).}
\item \textbf{ב.} מנהלי = כתיבת מדיניות, הנחיות, נהלים, תקנים; פיזי = הגנה על חדר שרתים וארונות תקשורת; לוגי = רשימות גישה, סיסמאות, מערכות למניעת גישה.
\item \textbf{ג.} \enLR{Script kiddies} = משתמש בכלי פריצה בלי להיות מודע לתהליכים; \enLR{Hacktivist} = תוקף מתוך אמונות ומניעים אידיאולוגיים; \enLR{Insider} = משתמש בגישה מורשית לצרכים זרים.
\item \textbf{ד.} \enLR{Nmap} – \textbf{\enLR{3}: גילוי, ניתוח ובקרה ברשתות תקשורת} (סריקת פורטים ומארחים).
\item \textbf{ה.} \enLR{Evil Twin} – \textbf{\enLR{1}: נקודת גישה זדונית עם \enLR{SSID} זהה לנקודה חוקית}.
\end{enumerate}

\subheading{שאלה \enLR{7 (9} נק')}

\begin{enumerate}
\item \textbf{א.} ניחוש סיסמה בעזרת מילון – \textbf{\enLR{4: Brute force attack}} (התקפת מילון היא תת-סוג).
\item \textbf{ב.} \enLR{MAC table overflow} – \textbf{\enLR{1: port security}} על כל הממשקים שאינם \enLR{trunk} + הגבלת מספר כתובות \enLR{MAC.}
\item \textbf{ג.} השבתת \enLR{DTP} מונעת \textbf{\enLR{3: VLAN hopping}} \enLR{(switch spoofing).}
\item \textbf{ד.} \enLR{BPDU Guard} גלובלי – \textbf{\enLR{3}: \code{spanning-tree portfast bpduguard default}} (תשובה \enLR{2} גם נכונה תחבירית לממשק בודד – \code{spanning-tree bpduguard enable} – אך התשובה המצופה היא \enLR{3).}
\item \textbf{ה.} חסימת \enLR{192.168.2.0} ל-\enLR{192.168.3.0} בלבד: \textbf{\enLR{2}} – \code{deny 192.168.2.0} ואז \code{permit any}, ומחילים \textbf{\enLR{out}} על \enLR{G0/2} (הממשק שיוצא לרשת המוגבלת). סדר השורות קריטי: \enLR{permit any} קודם היה מבטל את ה-\enLR{deny.}
\end{enumerate}

\sectionheading{סביבת מעבדה מומלצת}

\begin{itemize}
\item \textbf{\enLR{Cisco Packet Tracer}} (חינם דרך \enLR{Cisco Networking Academy)} – תומך ב: \enLR{SSH, AAA} עם \enLR{RADIUS/TACACS+, ACL, ZPF} (חלקית), \enLR{IOS IPS} (בסיסי), \enLR{port security, BPDU guard, DHCP snooping, GRE} ו-\enLR{IPsec site-to-site.} זה מכסה כמעט את כל המעשי בקורס.
\item \textbf{\enLR{Wireshark}} – להדגמת \enLR{Telnet} בטקסט גלוי מול \enLR{SSH} מוצפן (פרק \enLR{2)}, ולהצגת \enLR{handshake} של \enLR{TLS} (פרק \enLR{7).}
\item \textbf{\enLR{Nmap / Zenmap}} – רק ברשת מעבדה מבודדת ובאישור בית הספר (פרק \enLR{9).} אף פעם לא על רשת בית הספר או האינטרנט.
\item אתרי הדגמה: \enLR{cryptii.com} (צופן קיסר/ויז'נר), \enLR{sha256algorithm.com} (הדמיה של גיבוב).
\end{itemize}

\begin{mistakebox}\boxtitle{חשוב – אתיקה ומשמעת}\par  בקורס הזה מלמדים איך התקפות עובדות כדי לדעת להגן. הבהירו לתלמידים כבר בשיעור הראשון: הרצת סורקים, כלי פריצה או ניסיונות "לבדוק" את רשת בית הספר ללא אישור בכתב היא עבירה על חוק המחשבים \enLR{(1995)} – ותאמרו להם זאת במפורש. פרק \enLR{9} מקדיש לכך חלק.\end{mistakebox}


\clearpage
\phantomsection\addcontentsline{toc}{section}{תזכורת רשתות תקשורת}

\sloppy
\chapterheading{תזכורת רשתות תקשורת}{\small\color{subgray}רענון לקראת חלק א' של הבגרות · אינו מחליף את קורס הרשתות\par}\vskip4pt

\begin{metabox}\textbf{מטרה:} לרענן את מושגי הרשתות הנדרשים לחלק א' של שאלון \enLR{735001 (50} נק')\\ \textbf{הערה:} החומר נלמד בהרחבה אצל מורה הרשתות; זהו דף עזר מרוכז לחזרה\end{metabox}

\sectionheading{\enLR{1.} מודל \enLR{OSI} ו-\enLR{TCP/IP} – היכן יושב כל דבר}

\begin{xltabular}{\linewidth}{|>{\setRTL\raggedleft\arraybackslash}X|>{\setRTL\raggedleft\arraybackslash}X|>{\setRTL\raggedleft\arraybackslash}X|>{\setRTL\raggedleft\arraybackslash}X|}
\hline
\rowcolor{hdrbg}\textbf{פרוטוקולים / התקנים} & \textbf{יחידה} & \textbf{\#} & \textbf{שכבה \enLR{OSI}} \\
\hline
\endhead
\enLR{HTTP, HTTPS, DNS, DHCP, FTP, SMTP, IMAP, SSH, Telnet, SNMP} & \enLR{Data} & \enLR{7}–\enLR{5} & \enLR{Application / Presentation / Session} \\
\hline
\textbf{\enLR{TCP}} (אמין, מספרי רצף, לחיצת יד), \textbf{\enLR{UDP}} (מהיר, ללא חיבור) & \enLR{Segment} & \enLR{4} & \enLR{Transport} \\
\hline
\enLR{IP, ICMP, OSPF}; \textbf{נתב \enLR{(Router)}} & \enLR{Packet} & \enLR{3} & \enLR{Network} \\
\hline
\enLR{Ethernet, ARP, STP}; \textbf{מתג \enLR{(Switch)}}, כתובות \enLR{MAC} & \enLR{Frame} & \enLR{2} & \enLR{Data Link} \\
\hline
כבלים, חשמל & \enLR{Bits} & \enLR{1} & \enLR{Physical} \\
\hline
\end{xltabular}

\begin{exambox}\boxtitle{נשאל בבגרות (נכון/לא נכון)}\par  \enLR{TCP} פועל בשכבה \enLR{4} (לא \enLR{3}!) · \enLR{IMAP} מקבל דוא"ל (נכון) · \enLR{UDP} \underline{אינו} יוצר תקשורת אמינה (לא נכון – זה \enLR{TCP).}\end{exambox}

\sectionheading{\enLR{2.} בסיסי ספירה: בינארי · הקסדצימלי · אוקטלי}

\begin{itemize}
\item \enLR{1} ספרה הקסדצימלית = \enLR{4} ביטים; \enLR{1} ספרה אוקטלית = \enLR{3} ביטים.
\item המרה: הקס → בינארי (כל ספרה ל-\enLR{4} ביט), ואז לקבץ ל-\enLR{3} לאוקטלי.
\end{itemize}

\begin{exambox}\boxtitle{דוגמה (כמו בבגרות): \enLR{BA3F} בהקס}\par  \enLR{B=1011, A=1010, 3=0011, F=1111} ⇒ \textbf{בינארי:} \enLR{1011 1010 0011 1111}\newline  לאוקטלי מקבצים \enLR{3} ביטים מימין: \enLR{1 011 101 000 111 111} ⇒ \textbf{אוקטלי:} \enLR{135077}\end{exambox}

\sectionheading{\enLR{3.} כתובות \enLR{IPv4}, סאבנטים ו-\enLR{Wildcard}}

\begin{xltabular}{\linewidth}{|>{\setRTL\raggedleft\arraybackslash}X|>{\setRTL\raggedleft\arraybackslash}X|}
\hline
\enLR{10.0.0.0/8} · \enLR{172.16.0.0/12} · \enLR{192.168.0.0/16} & \textbf{כתובות פרטיות} \\
\hline
\enLR{2}\textsuperscript{\enLR{(32}−\enLR{prefix)}} − \enLR{2} (פחות כתובת רשת ו-\enLR{broadcast)} & \textbf{מספר מארחים שמישים} \\
\hline
\enLR{255.255.255.255} − \enLR{subnet mask} & \textbf{\enLR{Wildcard}} \\
\hline
\end{xltabular}

\subheading{טבלת ייחוס מהירה}

\begin{xltabular}{\linewidth}{|>{\setRTL\raggedleft\arraybackslash}X|>{\setRTL\raggedleft\arraybackslash}X|>{\setRTL\raggedleft\arraybackslash}X|>{\setRTL\raggedleft\arraybackslash}X|>{\setRTL\raggedleft\arraybackslash}X|}
\hline
\rowcolor{hdrbg}\textbf{\enLR{Wildcard}} & \textbf{מארחים שמישים} & \textbf{גודל בלוק} & \textbf{\enLR{Subnet mask}} & \textbf{\enLR{Prefix}} \\
\hline
\endhead
\enLR{0.0.0.255} & \enLR{254} & \enLR{256} & \enLR{255.255.255.0} & /\enLR{24} \\
\hline
\enLR{0.0.0.127} & \enLR{126} & \enLR{128} & \enLR{255.255.255.128} & /\enLR{25} \\
\hline
\enLR{0.0.0.63} & \enLR{62} & \enLR{64} & \enLR{255.255.255.192} & /\enLR{26} \\
\hline
\enLR{0.0.0.31} & \enLR{30} & \enLR{32} & \enLR{255.255.255.224} & /\enLR{27} \\
\hline
\enLR{0.0.0.15} & \enLR{14} & \enLR{16} & \enLR{255.255.255.240} & /\enLR{28} \\
\hline
\enLR{0.0.0.3} & \enLR{2} & \enLR{4} & \enLR{255.255.255.252} & /\enLR{30} \\
\hline
\end{xltabular}

\begin{exambox}\boxtitle{דוגמאות מהבגרות}\par  • מספר מארחים שמישים ב-\code{172.16.16.16/28} ⇒ \enLR{2}\textsuperscript{\enLR{4}}−\enLR{2} = \textbf{\enLR{14}}.\newline  • \code{135.8.x/26} ⇒ \enLR{Subnet mask} = \textbf{\enLR{255.255.255.192}} · \enLR{Wildcard} = \textbf{\enLR{0.0.0.63}}.\newline  • מציאת הרשת: לכתובת עם /\enLR{30}, הבלוקים הם \enLR{0,4,8}… ⇒ \code{10.10.2.20/30} ⟵ רשת .\enLR{20}, מארחים שמישים .\enLR{21}–.\enLR{22, broadcast .23.}\end{exambox}

\sectionheading{\enLR{4. IPv6} – קיצור וסוגי כתובות}

\subheading{כללי קיצור}

\begin{itemize}
\item משמיטים אפסים מובילים בכל בלוק (\code{130F} נשאר, \code{0000}→\code{0}).
\item \code{::} מחליף רצף אחד של בלוקי אפסים (פעם אחת בלבד בכתובת).
\end{itemize}

\begin{exambox}\boxtitle{דוגמה מהבגרות}\par  \code{3001:0000:130F:0000:0000:0000:0001:133B} ⇒ \textbf{\code{3001:0:130F::1:133B}}\end{exambox}

\begin{xltabular}{\linewidth}{|>{\setRTL\raggedleft\arraybackslash}X|>{\setRTL\raggedleft\arraybackslash}X|>{\setRTL\raggedleft\arraybackslash}X|}
\hline
\rowcolor{hdrbg}\textbf{תפקיד} & \textbf{תחילית} & \textbf{סוג כתובת} \\
\hline
\endhead
ניתובה באינטרנט (כמו כתובת ציבורית) & \enLR{2000::/3} & \textbf{\enLR{Global Unicast}} \\
\hline
תקשורת בקו/רשת מקומית בלבד; לא עוברת נתב & \enLR{FE80::/10} & \textbf{\enLR{Link-Local}} \\
\hline
\textbf{תקשורת בין תתי-רשתות פנימיות, לא ניתובה באינטרנט} (מקביל לכתובת פרטית) & \enLR{FC00::/7 (FD00::/8)} & \textbf{\enLR{Unique Local (ULA)}} \\
\hline
שידור לקבוצה & \enLR{FF00::/8} & \textbf{\enLR{Multicast}} \\
\hline
\end{xltabular}

\begin{exambox}\boxtitle{נשאל בבגרות}\par  "סוג כתובת המספק תקשורת בין תתי-רשתות אך שאינו ניתוב באינטרנט" ⇒ \textbf{\enLR{Unique Local}}.\end{exambox}

\sectionheading{\enLR{5.} פרוטוקולים ופורטים חשובים}

\begin{xltabular}{\linewidth}{|>{\setRTL\raggedleft\arraybackslash}X|>{\setRTL\raggedleft\arraybackslash}X|>{\setRTL\raggedleft\arraybackslash}X|>{\setRTL\raggedleft\arraybackslash}X|}
\hline
\rowcolor{hdrbg}\textbf{תפקיד} & \textbf{\enLR{TCP/UDP}} & \textbf{פורט} & \textbf{פרוטוקול} \\
\hline
\endhead
העברת קבצים & \enLR{TCP} & \enLR{20/21} & \enLR{FTP} \\
\hline
ניהול מרחוק מוצפן & \enLR{TCP} & \enLR{22} & \enLR{SSH} \\
\hline
ניהול מרחוק (לא מוצפן!) & \enLR{TCP} & \enLR{23} & \enLR{Telnet} \\
\hline
שליחת דוא"ל & \enLR{TCP} & \enLR{25} & \enLR{SMTP} \\
\hline
תרגום שם ⇄ כתובת & \enLR{UDP/TCP} & \enLR{53} & \enLR{DNS} \\
\hline
הקצאת כתובות אוטומטית & \enLR{UDP} & \enLR{67/68} & \enLR{DHCP} \\
\hline
גלישה & \enLR{TCP} & \enLR{80 / 443} & \enLR{HTTP / HTTPS} \\
\hline
קבלת דוא"ל & \enLR{TCP} & \enLR{143 / 110} & \enLR{IMAP / POP3} \\
\hline
ניטור ציוד & \enLR{UDP} & \enLR{161} & \enLR{SNMP} \\
\hline
יומן אירועים & \enLR{UDP} & \enLR{514} & \enLR{Syslog} \\
\hline
\end{xltabular}

\sectionheading{\enLR{6.} ניתוב: מרחק ניהולי, ניתוב סטטי ו-\enLR{DHCP Relay}}

\subheading{מרחק ניהולי \enLR{(Administrative Distance)} – "עד כמה סומכים על מקור המידע" (נמוך = עדיף)}

\begin{xltabular}{\linewidth}{|>{\setRTL\raggedleft\arraybackslash}X|>{\setRTL\raggedleft\arraybackslash}X|}
\hline
\rowcolor{hdrbg}\textbf{\enLR{AD}} & \textbf{מקור} \\
\hline
\endhead
\enLR{0} & \enLR{Connected} (מחובר ישירות) \\
\hline
\enLR{1} & ניתוב סטטי \enLR{(Static)} \\
\hline
\enLR{20} & \enLR{External BGP} \\
\hline
\enLR{90} & \enLR{Internal EIGRP} \\
\hline
\enLR{110} & \enLR{OSPF} \\
\hline
\enLR{120} & \enLR{RIP} \\
\hline
\end{xltabular}

\begin{codebox}\begin{english}\codefont\footnotesize\color{cfg}\setlength{\parindent}{0pt}
\textcolor{ccom}{\texthebrew{! ניתוב סטטי: אל רשת היעד, דרך ה-\enLR{next-hop}}}\\
R1(config)\# \textcolor{cbold}{\textbf{ip route 10.10.2.20 255.255.255.252 10.10.255.1}}\\
\textcolor{ccom}{\texthebrew{! ניתוב ברירת מחדל (לאינטרנט / כגיבוי)}}\\
R1(config)\# \textcolor{cbold}{\textbf{ip route 0.0.0.0 0.0.0.0 200.1.1.1}}\\
\textcolor{ccom}{\texthebrew{! גיבוי \enLR{(floating static)} – \enLR{AD} גבוה מ-\enLR{OSPF} כדי שלא יבטל את הניתוב הדינמי}}\\
R1(config)\# ip route 10.0.0.0 255.0.0.0 200.1.1.2 \textcolor{cbold}{\textbf{120}}
\end{english}\end{codebox}

\begin{exambox}\boxtitle{\enLR{DHCP Relay} – כשהשרת ברשת אחרת}\par  בקשת \enLR{DHCP} היא \enLR{broadcast} ואינה עוברת נתב. מגדירים על הממשק של הלקוח את כתובת שרת ה-\enLR{DHCP}: \enLR{R1(config)\# interface g0/0 R1(config-if)}\# \textbf{\enLR{ip helper-address 192.50.100.100}} ! כתובת שרת ה-\enLR{DHCP}\end{exambox}

\sectionheading{\enLR{7.} מיתוג: \enLR{VLAN, Trunk} ו-\enLR{Router-on-a-Stick}}

\begin{itemize}
\item \textbf{\enLR{VLAN}} – מפצל מתג לרשתות לוגיות מבודדות. מעבר בין \enLR{VLAN}ים דורש ניתוב.
\item \textbf{\enLR{Trunk (802.1Q)}} – קישור שנושא כמה \enLR{VLAN}ים בין מתגים; מוסיף תג של מספר ה-\enLR{VLAN.}
\item \textbf{\enLR{Router-on-a-Stick}} – נתב עם תת-ממשקים \enLR{(subinterfaces)}, אחד לכל \enLR{VLAN}, לניתוב ביניהם.
\end{itemize}

\begin{codebox}\begin{english}\codefont\footnotesize\color{cfg}\setlength{\parindent}{0pt}
\textcolor{ccom}{\texthebrew{! ניתוב בין \enLR{VLAN}ים על ממשק פיזי אחד}}\\
R1(config)\# interface g0/0/1.10\\
R1(config-subif)\# \textcolor{cbold}{\textbf{encapsulation dot1Q 10}}\\
R1(config-subif)\# \textcolor{cbold}{\textbf{ip address 172.18.1.254 255.255.0.0}}\\
R1(config)\# interface g0/0/1.20\\
R1(config-subif)\# encapsulation dot1Q 20\\
R1(config-subif)\# ip address 172.18.2.254 255.255.0.0\\
\textcolor{ccom}{\texthebrew{! בדיקת טבלת ה-\enLR{MAC} של המתג}}\\
S1\# \textcolor{cbold}{\textbf{show mac address-table}}
\end{english}\end{codebox}

\sectionheading{\enLR{8. STP} ו-\enLR{OSPF} – כללי הבחירה (נשאלים כמעט תמיד)}

\begin{defbox}\boxtitle{\enLR{STP} – בחירת \enLR{Root Bridge}}\par  \enLR{Bridge ID} = \textbf{\enLR{Priority}} (ברירת מחדל \enLR{32768)} + \textbf{\enLR{MAC}}.\newline  מנצח: ה-\enLR{Priority} \underline{הנמוך}; בתיקו – ה-\textbf{\enLR{MAC} הנמוך}.\newline  \enLR{STP} פועל בשכבה \textbf{\enLR{2}} ומונע לולאות ע"י \textbf{חסימת יציאות} \enLR{(port blocking).}\end{defbox}
\begin{defbox}\boxtitle{\enLR{OSPF} – בחירת \enLR{DR / BDR}}\par  מנצח: ה-\textbf{\enLR{Priority} הגבוה}; בתיקו – ה-\textbf{\enLR{Router-ID} הגבוה}.\newline  \underline{שימו לב – הפוך מ-\enLR{STP}:} ב-\enLR{STP} נמוך מנצח, ב-\enLR{OSPF} גבוה מנצח.\end{defbox}

\begin{exambox}\boxtitle{דוגמאות מהבגרות}\par  • \enLR{4} מתגים, \enLR{priority} זהה, \enLR{MAC} הנמוך ⇒ אותו מתג הופך ל-\enLR{root.}\newline  • \enLR{OSPF: R1(pri 2, RID 1.1.1.1), R3(pri 2, RID 3.3.3.3)} ⇒ בין בעלי \enLR{priority 2} מנצח \enLR{R3 (RID} גבוה) כ-\textbf{\enLR{DR}}, ו-\enLR{R1} כ-\textbf{\enLR{BDR}}.\end{exambox}

\sectionheading{\enLR{9.} פקודות \enLR{CLI} ואבחון תקלות}

\begin{xltabular}{\linewidth}{|>{\setRTL\raggedleft\arraybackslash}X|>{\setRTL\raggedleft\arraybackslash}X|}
\hline
\rowcolor{hdrbg}\textbf{הנחיה} & \textbf{מצב} \\
\hline
\endhead
\code{Router>} – פקודות \enLR{show} בסיסיות, \enLR{ping} & \enLR{User EXEC} \\
\hline
\code{Router\#} – \code{enable} & \enLR{Privileged EXEC} \\
\hline
\code{Router(config)\#} – \code{configure terminal} & \enLR{Global config} \\
\hline
\code{Router(config-if)\#} & \enLR{Interface} \\
\hline
\end{xltabular}

\begin{codebox}\begin{english}\codefont\footnotesize\color{cfg}\setlength{\parindent}{0pt}
Switch(config)\# \textcolor{cbold}{\textbf{hostname R1}}\\
R1(config)\# interface g0/1\\
R1(config-if)\# \textcolor{cbold}{\textbf{no switchport}}          \textcolor{ccom}{\texthebrew{! הופך פורט מתג \enLR{L3} לפורט מנותב}}\\
R1(config-if)\# ip address 10.100.20.42 255.255.255.0\\
R1(config-if)\# no shutdown\\
\textcolor{ccom}{\texthebrew{! בדיקות נפוצות}}\\
R1\# show ip interface brief\\
R1\# show ip route\\
R1\# show running-config
\end{english}\end{codebox}

\subheading{אבחון תקלות – שיטת \enLR{Top-Down}}

בודקים מהשכבה העליונה כלפי מטה: גישה לאפליקציה/שירות ⟶ הגדרות רשת \enLR{(IP/gateway/DNS)} ⟶ חיבור פיזי (כבל, נורית). דוגמה: בדיקת גישה לנתב החברה ⟶ בדיקת הגדרות ⟶ בדיקת כבל רשת.\par

\begin{exambox}\boxtitle{פקודות אבחון בתחנת קצה \enLR{(Windows)}}\par  • \code{ipconfig /all} – \enLR{IP, MAC, gateway, DNS.}\newline  • \code{ipconfig /displaydns} – \textbf{הצגת זיכרון המטמון \enLR{(cache)} של ה-\enLR{DNS}} במחשב.\newline  • \code{ipconfig /flushdns} – איפוס המטמון. • \code{nslookup}, \code{ping}, \code{tracert}, \code{arp -a}.\end{exambox}

\begin{exambox}\boxtitle{\enLR{CDP}}\par  פרוטוקול \textbf{\enLR{CDP}} מאפשר לזהות התקנים סמוכים של סיסקו ברשת המקומית (\code{show cdp neighbors}). מטעמי אבטחה מכבים אותו על פורטי קצה (\code{no cdp enable}).\end{exambox}

\sectionheading{\enLR{10.} תרגילי חזרה מהירים}

\begin{exbox}\boxtitle{תרגיל \enLR{1}}\par  א. כמה מארחים שמישים ב-\code{/27}? ב. מהי ה-\enLR{wildcard} של \code{10.20.0.0/16}? ג. באיזו שכבה פועל \enLR{TCP}? \par\vskip2pt\hrule\vskip3pt\textbf{}א. \enLR{30} ב. \enLR{0.0.255.255} ג. שכבה \enLR{4 (Transport)}\end{exbox}

\begin{exbox}\boxtitle{תרגיל \enLR{2}}\par  קצרו: \code{2001:0DB8:0000:0000:0000:00FF:0000:0001} \par\vskip2pt\hrule\vskip3pt\textbf{}\code{2001:DB8::FF:0:1}\end{exbox}

\begin{exbox}\boxtitle{תרגיל \enLR{3}}\par  לנתב יש מסלול \enLR{OSPF (AD 110)} ומסלול סטטי \enLR{(AD 1)} לאותה רשת. איזה ייבחר ולמה? \par\vskip2pt\hrule\vskip3pt\textbf{}הסטטי – ה-\enLR{AD} הנמוך יותר מעיד על מקור מהימן/מועדף.\end{exbox}

\begin{exbox}\boxtitle{תרגיל \enLR{4}}\par  מנהל רוצה ששרת \enLR{DHCP} ברשת אחרת יקצה כתובות ללקוחות ברשת של \enLR{R1.} איזו פקודה נדרשת ועל איזה ממשק? \par\vskip2pt\hrule\vskip3pt\textbf{}\code{ip helper-address <IP של שרת ה-DHCP>} על הממשק של \enLR{R1} הפונה ללקוחות.\end{exbox}


\clearpage
\phantomsection\addcontentsline{toc}{section}{פרק \enLR{1} – מבוא לאיומי רשת חדשים}

\sloppy
\chapterheading{פרק \enLR{1} – מבוא לאיומי רשת חדשים}{\small\color{subgray}\enLR{10} שעות עיוני · שבועות \enLR{1}–\enLR{3}\par}\vskip4pt

\begin{defbox}\boxtitle{איך לקרוא את הפרק הזה}\par  הפרק נכתב כך שאפשר לקרוא אותו ברצף, מההתחלה ועד הסוף, גם אם מעולם לא למדתם רשתות מחשבים. כל מונח מוסבר בפעם הראשונה שהוא מופיע, ואין צורך בידע מוקדם. לאורך הדרך נעצור לסיפורים אמיתיים שקרו בעולם, ובסוף יש תרגילים עם פתרונות ומילון מונחים לחזרה מהירה.\end{defbox}

\sectionheading{\enLR{1.1} על מה בכלל מדברים?}

נתחיל מהשאלה הכי בסיסית: מה זה "אבטחת מידע"? \textbf{אבטחת מידע \enLR{(Information Security)}} היא פשוט ההגנה על מידע – ועל המחשבים שמחזיקים אותו – מפני אנשים שאסור להם לגשת אליו, לשנות אותו או למחוק אותו. המידע יכול להיות תמונות בטלפון שלכם, ציונים במערכת של בית הספר, או פרטי כרטיסי האשראי של לקוחות בחנות. הרעיון זהה בכל המקרים: יש משהו בעל ערך, ואנחנו רוצים שרק האנשים הנכונים יוכלו לגעת בו.\par

בקורס הזה נתמקד בתחום מסוים בתוך אבטחת המידע שנקרא \textbf{אבטחת רשת \enLR{(Network Security)}}. "רשת" היא מה שמחבר מחשבים זה לזה כדי שיוכלו לדבר – למשל הרשת בבית שמחברת את הטלפון, המחשב והטלוויזיה החכמה לאינטרנט. אבטחת רשת עוסקת בהגנה על הצינורות שדרכם עובר המידע ועל המידע עצמו בזמן שהוא בדרך. למה דווקא הרשת? כי היום כמעט כל התקפה מגיעה דרך הרשת – תוקף בצד השני של העולם יכול לנסות להיכנס למחשב שלכם מבלי לקום מהכיסא.\par

\sectionheading{\enLR{1.2} אבני הבניין – המושגים שנשתמש בהם כל הזמן}

לפני שנצלול, בואו נכיר כמה מילים בסיסיות. אל תדאגו – כולן פשוטות, ונחזור אליהן שוב ושוב.\par

\begin{itemize}
\item \textbf{מחשב מארח \enLR{(Host)}} – כל מכשיר שמחובר לרשת: מחשב, טלפון, שרת, מדפסת, ואפילו מצלמת אבטחה חכמה. בקיצור, כל דבר "שיודע לדבר" ברשת.
\item \textbf{כתובת \enLR{IP}} – ה"כתובת" של המכשיר ברשת, בדיוק כמו כתובת של בית (למשל \enLR{192.168.1.10).} לפי הכתובת הזו יודעים לאן לשלוח את המידע. שני מכשירים לא יכולים לתקשר בלי לדעת את הכתובת זה של זה.
\item \textbf{כתובת \enLR{MAC}} – מספר סידורי ייחודי ש"צרוב" בכרטיס הרשת של כל מכשיר עוד מהמפעל (למשל \enLR{0C:0E:15:22:05:97).} אפשר לחשוב עליו כמו מספר השלדה של רכב – הוא מזהה את המכשיר עצמו, בעוד שכתובת ה-\enLR{IP} יכולה להשתנות.
\item \textbf{חבילה \enLR{(Packet)}} – מידע ברשת לא נשלח כמקשה אחת גדולה, אלא נחתך ל"מעטפות" קטנות. כל מעטפה כזו נקראת חבילה, וכתובות עליה כתובת השולח וכתובת היעד – בדיוק כמו מעטפה בדואר.
\item \textbf{שרת \enLR{(Server)} ולקוח \enLR{(Client)}} – ה\textbf{לקוח} הוא מי שמבקש שירות (למשל הדפדפן שלכם), וה\textbf{שרת} הוא מי שנותן אותו (למשל המחשב שמאחסן את אתר האינטרנט). הלקוח שואל, השרת עונה.
\item \textbf{נתב \enLR{(Router)}} – מכשיר שמעביר חבילות \emph{בין רשתות שונות}, למשל בין הרשת בבית לבין האינטרנט. אפשר לחשוב עליו כעל צומת דרכים שמחליט לאן להפנות כל מעטפה.
\item \textbf{מתג \enLR{(Switch)}} – מכשיר שמחבר מכשירים \emph{בתוך אותה רשת מקומית}, למשל את כל המחשבים במשרד אחד.
\item \textbf{פרוטוקול \enLR{(Protocol)}} – "שפה" או אוסף כללים מוסכמים שמכשירים מדברים בהם. למשל, \enLR{HTTP} הוא הפרוטוקול שבו דפדפנים ואתרים מדברים, ו-\enLR{DNS} הוא הפרוטוקול שמתרגם שם של אתר לכתובת \enLR{IP.} כמו כללי הנימוס בשיחת טלפון – שני הצדדים חייבים להסכים עליהם מראש.
\item \textbf{פורט \enLR{(Port)}} – "דלת" ממוספרת בתוך מכשיר, אחת לכל סוג שירות. למשל דלת \enLR{80} לאתרי אינטרנט, דלת \enLR{443} לאתרים מאובטחים, דלת \enLR{22} לניהול מרחוק \enLR{(SSH).} לכל כתובת \enLR{IP} יש עשרות אלפי דלתות כאלה.
\end{itemize}

שימו לב לרעיון החשוב שמלווה את כל הקורס: בכל אחת מנקודות המפגש האלה – מכשיר, כתובת, חבילה, דלת – תוקף יכול לנסות להיכנס, ואנחנו נלמד איך לחסום אותו.\par

\sectionheading{\enLR{1.3} מה בעצם מגנים עליו? שלושת עקרונות ה-\enLR{CIA}}

כשמדברים על אבטחת מידע, מגנים על שלושה דברים. הם כל כך מרכזיים שנוח לזכור אותם בראשי התיבות באנגלית – \textbf{\enLR{CIA}} – ונחזור אליהם בכל פרק בקורס.\par

\textbf{סודיות \enLR{(Confidentiality)}} פירושה שרק מי שמורשה יכול לראות את המידע. כשאתם בודקים את יתרת חשבון הבנק שלכם, רק אתם אמורים לראות אותה – זו סודיות. סודיות נפגעת כשמישהו "מצותת" לתקשורת שלכם או כשמאגר נתונים דולף. מגנים עליה בעזרת \textbf{הצפנה} (ערבוב המידע כך שרק בעל המפתח יבין אותו) ובקרת גישה.\par

\textbf{שלמות \enLR{(Integrity)}} פירושה שהמידע לא שונה בדרך בלי רשות. אם העברתם בבנק \enLR{1,000} ₪, אתם רוצים להיות בטוחים שאף אחד לא שינה את הסכום ל-\enLR{10,000} ₪ באמצע הדרך. שלמות נפגעת כשתוקף מצליח לשנות מידע במעבר. מגנים עליה בעזרת \textbf{גיבוב} (טביעת אצבע דיגיטלית של המידע שנלמד בפרק \enLR{7)} וחתימה דיגיטלית.\par

\textbf{זמינות \enLR{(Availability)}} פירושה שהמידע והשירות זמינים כשצריך אותם. הכספומט צריך לעבוד כשאתם רוצים למשוך מזומן, ואתר בית הספר צריך לעבוד ביום ההרשמה. זמינות נפגעת כשתוקף מפיל שירות (התקפת מניעת שירות, שנכיר בהמשך) או כשכופרה מצפינה את הקבצים.\par

\begin{tipbox}\boxtitle{למה שלושה ולא אחד?}\par  תלמידים רבים חושבים שאם רק נצפין הכול, נהיה מוגנים. אבל הצפנה נותנת רק סודיות. קובץ מוצפן עדיין יכול להימחק (פגיעה בזמינות) או להיות מוחלף (פגיעה בשלמות). למעשה, \textbf{כופרה} – אחת הנוזקות המסוכנות היום – \emph{משתמשת} בהצפנה נגדכם: היא מצפינה את הקבצים שלכם ודורשת כסף כדי לפתוח אותם. אותו כלי בדיוק יכול להיות הגנה או נשק. לכן צריך לחשוב על שלושת העקרונות יחד.\end{tipbox}

לצד שלושת אלה יש שני מושגים נלווים שכדאי להכיר: \textbf{אימות \enLR{(Authentication)}} – להוכיח מי אתה (למשל בעזרת שם משתמש וסיסמה), ו\textbf{אי-הכחשה \enLR{(Non-repudiation)}} – לוודא שמי ששלח הודעה לא יוכל אחר כך להכחיש ששלח אותה (מושג בעזרת חתימה דיגיטלית).\par

\sectionheading{\enLR{1.4} שפת הסיכון: נכס, פגיעות, איום, סיכון ובקרה}

כדי לדבר על אבטחה בצורה מסודרת, אנשי המקצוע משתמשים בכמה מונחים קבועים. הדרך הכי קלה לזכור אותם היא דרך משל פשוט: \textbf{בית}.\par

\begin{itemize}
\item \textbf{נכס \enLR{(Asset)}} – הדבר בעל הערך שרוצים להגן עליו. במשל: הבית עצמו. באבטחת מידע: המידע של הלקוחות, השרתים, ואפילו המוניטין של החברה.
\item \textbf{פגיעות \enLR{(Vulnerability)}} – חולשה שאפשר לנצל. במשל: חלון שנשאר פתוח. באבטחת מידע: סיסמה חלשה, תוכנה עם באג, או דלת (פורט) שנשארה פתוחה. הפגיעות היא מה ש\underline{בשליטתנו} – אנחנו יכולים לסגור את החלון.
\item \textbf{איום \enLR{(Threat)}} – הגורם שעלול לנצל את הפגיעות. במשל: הגנב. באבטחת מידע: האקר, נוזקה, עובד ממורמר, ואפילו שריפה או הצפה. איומים קיימים תמיד ואינם בשליטתנו – גנבים קיימים בעולם בין אם נרצה ובין אם לא.
\item \textbf{ניצול \enLR{(Exploit)}} – הדרך המדויקת שבה האיום מנצל את הפגיעות. במשל: הרגע שבו הגנב מטפס דרך החלון הפתוח.
\item \textbf{סיכון \enLR{(Risk)}} – ההסתברות שהאיום אכן ינצל את הפגיעות, כפול הנזק שייגרם. אפשר לרשום זאת כנוסחה פשוטה: \textbf{סיכון = איום × פגיעות × ערך הנכס}. אם אחד מהם אפס, אין סיכון.
\item \textbf{בקרה \enLR{(Control)}} – כל אמצעי שמקטין את הסיכון. במשל: הסורג על החלון. באבטחת מידע: חומת אש, הצפנה, נוהל בחברה, או מנעול על חדר השרתים.
\end{itemize}

המסקנה המעשית: אנחנו לא יכולים לבטל את האיומים (הגנבים תמיד יהיו שם), אבל אנחנו כן יכולים לצמצם פגיעויות (לסגור את החלונות) ולהוסיף בקרות (להתקין סורגים). זו כל התורה של ניהול סיכונים, שנרחיב עליה בפרק \enLR{9.}\par

\sectionheading{\enLR{1.5} איך הגענו לכאן? קצת היסטוריה}

האיומים על הרשת התפתחו יחד עם האינטרנט. בהתחלה, בשנות ה-\enLR{80} וה-\enLR{90}, האיומים היו סקרנות של חובבנים. עם השנים הם הפכו לעסק של פשע מאורגן ואפילו לכלי מלחמה בין מדינות. במקום להקריא רשימת תאריכים, בואו נבין את המגמה דרך שני סיפורים אמיתיים.\par

\begin{storybox}\boxtitle{\enLR{1988} – הסטודנט שהפיל את האינטרנט (תולעת מוריס)}\par  רוברט מוריס, סטודנט בן \enLR{23}, כתב תוכנה קטנה "כדי למדוד את גודל האינטרנט". התוכנה הייתה אמורה להעתיק את עצמה למחשב פעם אחת בלבד – אבל בגלל באג היא הדביקה כל מחשב שוב ושוב, עד שכל מחשב הריץ מאות עותקים של עצמו וקרס. התוצאה: כ-\enLR{6,000} מחשבים – עשירית מהאינטרנט של אותם ימים – הושבתו. מוריס היה האדם הראשון שהורשע בחוק הפריצה האמריקאי; היום הוא פרופסור מכובד ב-\enLR{MIT.} הסיפור מלמד שני דברים: שקוד יכול להתפשט לבד ברשת, ושלפעמים הנזק הגדול נגרם דווקא מטעות ולא מכוונה זדונית.\end{storybox}

\begin{storybox}\boxtitle{\enLR{2017} – בית החולים שביטל ניתוחים \enLR{(WannaCry)}}\par  בבוקר יום שישי אחד, מסכי המחשבים ב-\enLR{80} בתי חולים ברחבי אנגליה הפכו לאדומים והציגו הודעה: "הקבצים שלך הוצפנו. שלם \enLR{300} דולר בביטקוין". מערכות התורים, תוצאות הבדיקות ומכשירי ההדמיה הפסיקו לעבוד; \enLR{19,000} תורים וניתוחים בוטלו, ואמבולנסים הופנו לבתי חולים אחרים. הנוזקה, כופרה בשם \enLR{WannaCry}, נכנסה דרך חולשה בתוכנת \enLR{Windows} שחברת מיקרוסופט \underline{כבר תיקנה חודשיים קודם} – אבל בתי החולים לא התקינו את העדכון. הלקח: עדכון תוכנה הוא אמצעי אבטחה, לא מטרד, ופגיעה בזמינות יכולה לעלות בחיי אדם.\end{storybox}

\sectionheading{\enLR{1.6} מי הם התוקפים ומה הם רוצים?}

לא כל התוקפים דומים. חשוב להבחין ביניהם לפי שני דברים: מה ה\textbf{מניע} שלהם (למה הם תוקפים) ומה ה\textbf{יכולת} שלהם (כמה הם מתוחכמים).\par

\begin{itemize}
\item \textbf{\enLR{Script Kiddie}} ("ילד סקריפטים") – חובבן שמריץ כלי פריצה שאחרים כתבו, בלי להבין באמת איך הם עובדים. המניע: סקרנות, יוקרה, שעמום. היכולת נמוכה, אבל הכלים שהוא מוריד מהאינטרנט יכולים להיות חזקים מאוד.
\item \textbf{\enLR{Hacktivist}} ("האקטיביסט") – תוקף ממניעים אידיאולוגיים או פוליטיים, למשל כדי למחות נגד ארגון. הקבוצה המפורסמת ביותר היא \enLR{Anonymous}, שהשחיתה אתרים והפילה שרתים מסיבות פוליטיות.
\item \textbf{פשע מאורגן \enLR{(Organized Crime)}} – קבוצה מאורגנת שמטרתה כסף. כנופיות כופרה כמו \enLR{REvil} ו-\enLR{LockBit} סוחטות מארגונים מיליוני דולרים. זהו כיום האיום הנפוץ והמזיק ביותר לעסקים.
\item \textbf{איום פנימי \enLR{(Insider)}} – עובד או קבלן שיש לו כבר גישה מורשית, והוא משתמש בה לרעה – מתוך נקמה, כסף, או סתם רשלנות. זהו האיום המסוכן במיוחד, כי "האויב כבר בפנים ויש לו מפתח". הדוגמה המפורסמת היא אדוארד סנודן, עובד שהדליף מסמכים סודיים של ה-\enLR{NSA.}
\item \textbf{מדינה \enLR{(State-sponsored)}} – יחידות סייבר ממשלתיות שתוקפות למטרות ריגול או חבלה בתשתיות. אלה התוקפים המתוחכמים ביותר, עם משאבים אדירים וסבלנות של שנים. הדוגמה המפורסמת היא \enLR{Stuxnet}, נוזקה שפגעה בצנטריפוגות במתקן גרעיני באיראן.
\end{itemize}

הערה על מילים: לצד החלוקה לפי מניע, נהוג לחלק האקרים לפי אתיקה. \textbf{\enLR{White hat}} ("כובע לבן") הוא האקר שפועל באישור, למשל כדי לבדוק חדירות ולעזור לארגון להתגונן. \textbf{\enLR{Black hat}} ("כובע שחור") הוא פושע. \textbf{\enLR{Gray hat}} ("כובע אפור") פורץ בלי אישור אבל מדווח על מה שמצא. המילה "האקר" עצמה במקור פירושה מי שמבין מערכות לעומק – לא בהכרח פושע.\par

\sectionheading{\enLR{1.7} נוזקות: וירוסים, תולעים, סוסים טרויאנים ועוד}

\textbf{נוזקה \enLR{(Malware)}} היא מילה כללית לכל תוכנה זדונית. שלושת הסוגים הקלאסיים נבדלים זה מזה בעיקר בשאלה אחת: \underline{איך הם מתפשטים}.\par

\begin{itemize}
\item \textbf{וירוס \enLR{(Virus)}} – קוד זדוני שנצמד לקובץ או לתוכנה קיימת. הוא מתפשט רק כאשר \underline{המשתמש} מריץ את הקובץ הנגוע (למשל פותח קובץ \enLR{Excel} נגוע). הוירוס משכפל את עצמו.
\item \textbf{תולעת \enLR{(Worm)}} – תוכנה עצמאית שמשכפלת ומתפשטת \underline{לבד ברשת}, בלי שאף אדם יעשה דבר. תולעת מוריס מ-\enLR{1988} היא דוגמה מושלמת.
\item \textbf{סוס טרויאני \enLR{(Trojan)}} – תוכנה שנראית לגיטימית ומפתה (למשל "משחק חינם"), אבל בתוכה מסתתר קוד זדוני. \underline{המשתמש מתקין אותה מרצונו} כי הוא חושב שהיא בסדר. בניגוד לוירוס ולתולעת, סוס טרויאני \textbf{אינו משכפל את עצמו}.
\end{itemize}

\begin{tipbox}\boxtitle{השאלה שפותרת הכול}\par  כשמתלבטים בין וירוס לתולעת, שאלו את עצמכם: \textbf{האם צריך פעולה של משתמש כדי שהנוזקה תתפשט?} אם כן – זה וירוס. אם היא מתפשטת לבד ברשת – זו תולעת. וסוס טרויאני? אותו המשתמש מתקין במו ידיו, והוא לא משכפל את עצמו.\end{tipbox}

מלבד השלושה הקלאסיים, כדאי להכיר עוד כמה נוזקות בשמן: \textbf{כופרה \enLR{(Ransomware)}} מצפינה את הקבצים ודורשת תשלום (הנוזקה המזיקה ביותר כיום); \textbf{רוגלה \enLR{(Spyware)}} ו-\textbf{\enLR{Keylogger}} מרגלות ומתעדות כל הקשה על המקלדת; \textbf{\enLR{Rootkit}} מתחבאת עמוק בתוך מערכת ההפעלה כדי שיהיה קשה מאוד לגלות אותה; \textbf{\enLR{Botnet}} הוא צבא של מחשבים נגועים ("זומבים") שהתוקף שולט בהם מרחוק ומפעיל אותם יחד; ו-\textbf{\enLR{Backdoor}} ("דלת אחורית") הוא פתח סמוי שעוקף את מנגנון האימות ומאפשר לתוקף להיכנס מתי שירצה.\par

\begin{storybox}\boxtitle{\enLR{2017} – הכופרה שלא הייתה כופרה \enLR{(NotPetya)}}\par  \enLR{NotPetya} נראתה כמו כופרה רגילה, אבל למעשה הייתה נוזקה מסוג \textbf{\enLR{Wiper}} – היא \emph{מחקה} את המידע לצמיתות בלי אפשרות לשחזר, גם אם שילמת. היא התפשטה בצורה מפתיעה: דרך עדכון תוכנה \underline{לגיטימי} של תוכנת הנהלת חשבונות פופולרית באוקראינה (מה שנקרא "מתקפת שרשרת אספקה"). תוך שעות היא שיתקה חברות ענק כמו חברת הספנות \enLR{Maersk}, חברת התרופות \enLR{Merck} ו-\enLR{FedEx}, וגרמה נזק של כ-\enLR{10} מיליארד דולר. הלקח: נוזקה יכולה להיכנס דרך תוכנה שאתם סומכים עליה ומתעדכנת אוטומטית, ולא רק דרך קליק שלכם.\end{storybox}

מול כל אלה, ההגנה הבסיסית מפני נוזקות היא כמה שכבות יחד: \textbf{אנטי-וירוס} (שמזהה נוזקות מוכרות לפי "חתימה" – אבל לא מזהה נוזקה חדשה לגמרי), \textbf{עדכוני מערכת} שסוגרים חולשות, \textbf{גיבוי} שהוא התשובה מספר אחת לכופרה, ומעל הכול – משתמש ערני שלא לוחץ על כל קישור. הרעיון של כמה שכבות עצמאיות נקרא \textbf{הגנה לעומק \enLR{(Defense in Depth)}}: אם שכבה אחת נכשלת, האחרות עדיין מגנות. הוא יחזור לאורך כל הקורס.\par

\sectionheading{\enLR{1.8} איך נראית התקפה מההתחלה עד הסוף}

רוב ההתקפות עוברות שלושה שלבים. אפשר לחשוב עליהם כמו על פורץ שמתכנן שוד: קודם הוא משגיח על הבית, אחר כך מוצא דרך פנימה, ולבסוף לוקח את מה שרצה.\par

\textbf{שלב \enLR{1} – סיור \enLR{(Reconnaissance)}:} התוקף לומד על המטרה. מי היא? אילו כתובות \enLR{IP} יש לה? אילו "דלתות" (פורטים) פתוחות? בשלב הזה תוקף משתמש בכלים כמו \textbf{\enLR{Nmap}} – תוכנה שסורקת רשת ומגלה אילו מכשירים חיים ואילו שירותים פועלים בהם. הסיור יכול להיות גם פסיבי (לחפש מידע גלוי על החברה באינטרנט) וגם אקטיבי (לשלוח חבילות ולראות מה חוזר).\par

\textbf{שלב \enLR{2} – השגת גישה \enLR{(Access)}:} התוקף מנצל מה שגילה כדי להיכנס. הדרכים הנפוצות: ניחוש סיסמה, ניצול חולשה בתוכנה, או \textbf{\enLR{Man-in-the-Middle (MITM)}} – התוקף מכניס את עצמו באמצע התקשורת בין שני צדדים ומקשיב או משנה אותה.\par

\textbf{שלב \enLR{3} – הנזק:} כאן התוקף עושה את מה שרצה – גונב מידע, מצפין אותו בכופרה, או מפיל את השירות כדי שאף אחד לא יוכל להשתמש בו. השלב האחרון מוביל אותנו לסוג התקפה חשוב שכדאי להבין לעומק.\par

\subheading{מניעת שירות – \enLR{DoS} ו-\enLR{DDoS}}

\textbf{מניעת שירות \enLR{(Denial of Service}, בקיצור \enLR{DoS)}} היא התקפה שמטרתה להפיל שירות כדי שלא יהיה זמין. הדרך הפשוטה: להציף את השירות בכל כך הרבה בקשות מזויפות עד שהוא קורס תחת העומס. כשההצפה מגיעה ממקור אחד קוראים לזה \enLR{DoS}; כשהיא מגיעה ב\underline{ו-זמנית מאלפי} מחשבים נגועים \enLR{(botnet)} קוראים לזה \textbf{\enLR{DDoS (Distributed DoS)}}. אנלוגיה: אם מאה אנשים יזמינו במקביל משלוח פיצה לכתובת שלכם, המשלוח האמיתי שהזמנתם לא יוכל להגיע. קשה מאוד לחסום \enLR{DDoS}, כי אין כתובת אחת לחסום – ההתקפה מגיעה מכל מקום.\par

\begin{storybox}\boxtitle{\enLR{2016} – הצבא של מצלמות האבטחה \enLR{(Mirai)}}\par  נוזקה בשם \enLR{Mirai} השתלטה על כ-\enLR{600,000} מכשירים פשוטים – מצלמות אבטחה, ראוטרים ביתיים ומקליטי וידאו. איך? היא פשוט ניסתה \enLR{62} צירופים של שם משתמש וסיסמה מברירת המחדל של היצרנים (כמו \enLR{admin/admin, root/12345)} – ואצל מי שלא שינה את הסיסמה, נכנסה. עם הצבא העצום הזה היא ביצעה \enLR{DDoS} שהפיל את \enLR{Twitter, Netflix, Spotify} ו-\enLR{PayPal} ליום שלם. הסיפור ממחיש בדיוק כמה מסוכנת סיסמת ברירת מחדל שלא שונתה, ומהו \enLR{botnet} בפועל.\end{storybox}

\subheading{סיסמאות והתקפת ניחוש}

אחת הדרכים הנפוצות להיכנס היא פשוט לנחש את הסיסמה. \textbf{התקפת ברוט-פורס \enLR{(Brute Force)}} מנסה באופן שיטתי את כל הצירופים האפשריים עד שהיא מוצאת את הנכון; גרסה נפוצה שלה, \textbf{התקפת מילון}, מנסה קודם רשימה של סיסמאות נפוצות. כאן נכנס עיקרון פשוט וחשוב: \textbf{אורך הסיסמה חשוב יותר ממורכבותה}. סיסמה של \enLR{6} ספרות נפרצת בפחות משנייה; סיסמה של \enLR{8} אותיות קטנות בכ-\enLR{20} שניות; אבל סיסמה של \enLR{12} תווים מכל הסוגים תיקח מיליוני שנים. הסיבה: כל תו נוסף מכפיל את מספר האפשרויות. לכן משפט-סיסמה ארוך עדיף על "\enLR{Aa1}!" קצר. ההגנות הנוספות הן נעילת חשבון אחרי כמה ניסיונות כושלים, ו\textbf{אימות רב-שלבי \enLR{(MFA)}} – שילוב של סיסמה עם קוד מהטלפון, כך שגם אם הסיסמה נפרצה התוקף עדיין חסר את הטלפון שלכם.\par

\sectionheading{\enLR{1.9} ההתקפה שלא צריכה מחשב: הנדסה חברתית}

עד כה דיברנו על תקיפה של מכונות. אבל לעיתים קרובות הדרך הקלה ביותר פנימה היא דרך \underline{האדם}, לא המחשב. \textbf{הנדסה חברתית \enLR{(Social Engineering)}} היא מניפולציה של בני אדם כדי לגרום להם למסור מידע או גישה. אין כאן פריצה טכנית – יש כאן פסיכולוגיה: שילוב של סמכות, דחיפות וסקרנות.\par

הצורה הנפוצה ביותר היא \textbf{\enLR{Phishing}} ("דיוג") – דוא"ל או הודעת \enLR{SMS} שמתחזים לגוף לגיטימי (בנק, חברת שליחויות, בית הספר) ומבקשים מכם ללחוץ על קישור או "לאשר פרטים". יש גרסאות ממוקדות יותר: \textbf{\enLR{Spear Phishing}} מכוון לאדם מסוים עם פרטים אמיתיים עליו, ו-\textbf{\enLR{Whaling}} ("ציד לווייתנים") מכוון למנהלים בכירים. צורות נוספות: \textbf{\enLR{Pretexting}} – סיפור כיסוי ("אני מהתמיכה הטכנית, אני צריך את הסיסמה שלך כדי לתקן"), ו-\textbf{\enLR{Baiting}} ("פיתיון") – דיסק-און-קי "אבוד" עם תווית מפתה כמו "משכורות", בתקווה שמישהו יחבר אותו למחשב מסקרנות.\par

\begin{storybox}\boxtitle{\enLR{2020} – פריצת טוויטר של המפורסמים}\par  קבוצה של בני נוער התקשרה לעובדי חברת טוויטר, התחזתה לצוות ה-\enLR{IT} הפנימי, ושכנעה אותם למסור גישה לכלי הניהול הפנימי. עם הגישה הזו הם השתלטו על חשבונות של ברק אובמה, אילון מאסק, ביל גייטס ואפל, ופרסמו מהם הונאת ביטקוין. הם גרפו כ-\enLR{118,000} דולר תוך שעות – \textbf{בלי לפרוץ אף מחשב}, רק בכוח השכנוע. הסיפור מוכיח שאפילו אחת מחברות הטכנולוגיה הגדולות בעולם נופלת להנדסה חברתית. האדם הוא כמעט תמיד החוליה החלשה.\end{storybox}

\begin{storybox}\boxtitle{\enLR{2016} – המייל ששווה \enLR{70} מיליון \enLR{(Crelan)}}\par  עובד במחלקת הכספים של הבנק הבלגי \enLR{Crelan} קיבל מייל "מהמנכ"ל": "אנחנו באמצע עסקה סודית, העבר בבקשה כסף בדחיפות לחשבון הזה, ואל תדבר על זה עם אף אחד". כתובת השולח הייתה כמעט זהה לזו של המנכ"ל האמיתי, בהבדל של אות אחת. הכסף הועבר, והבנק הפסיד כ-\enLR{70} מיליון אירו. שוב – בלי נוזקה ובלי פריצה, רק שילוב של סמכות, דחיפות וסודיות. שאלה טובה לכיתה: איזה נוהל פשוט היה מונע את זה? (אישור כפול לכל העברה גדולה, ואימות טלפוני בערוץ נפרד.)\end{storybox}

\begin{tipbox}\boxtitle{איך מזהים הודעת \enLR{Phishing} – סימנים שכדאי ללמד}\par  \enLR{1.} כתובת השולח לא תואמת לארגון (חברת \enLR{UPS} ששולחת מכתובת @\enLR{hotmail.com). 2.} תחושת דחיפות או איום ("הצעה אחרונה!", "החשבון ייחסם"). \enLR{3.} בקשה לאשר פרטים או ללחוץ על כפתור. \enLR{4.} שגיאות ניסוח או לוגו מטושטש. \enLR{5.} קישור שהטקסט שלו שונה מהיעד האמיתי (בודקים בריחוף עם העכבר).\end{tipbox}

\sectionheading{\enLR{1.10} איך מתגוננים: סוגי הבקרות}

ראינו הרבה איומים; עכשיו נסדר את ההגנות. כל בקרה (אמצעי הגנה) שייכת לאחת משלוש קטגוריות, ונחזור אליהן בהרחבה בפרק \enLR{9}:\par

\begin{itemize}
\item \textbf{בקרה מנהלית \enLR{(Administrative)}} – מדיניות, נהלים והדרכות. למשל נוהל שאוסר להעביר כסף בלי אישור כפול, או הדרכת עובדים לזהות \enLR{Phishing.} אלה דברים שאנשים עושים.
\item \textbf{בקרה פיזית \enLR{(Physical)}} – הגנה על החומרה עצמה: מנעול על חדר השרתים, מצלמות אבטחה, שומר בכניסה.
\item \textbf{בקרה לוגית / טכנית \enLR{(Logical)}} – אמצעים בתוך התוכנה והחומרה: סיסמאות, הצפנה, חומת אש, רשימות גישה \enLR{(ACL).}
\end{itemize}

שימו לב איך שלוש הקטגוריות עובדות יחד: בסיפור של \enLR{Crelan}, בקרה \emph{מנהלית} (נוהל אישור כפול) הייתה מונעת את ההונאה, גם בלי שום טכנולוגיה. אבטחה טובה משלבת את שלושתן.\par

כדי לחבר את התיאוריה למעשה, הנה הצצה קטנה קדימה. גם בציוד רשת של סיסקו יש הגנה מובנית מפני ניחוש סיסמאות, שנלמד לעומק בפרק \enLR{2}:\par

\begin{codebox}\begin{english}\codefont\footnotesize\color{cfg}\setlength{\parindent}{0pt}
R1(config)\# \textcolor{cbold}{\textbf{login block-for 120 attempts 3 within 60}}\\
\textcolor{ccom}{\texthebrew{! אם היו \enLR{3} ניסיונות כושלים תוך \enLR{60} שניות – חסום כניסות למשך \enLR{120} שניות}}
\end{english}\end{codebox}

\sectionheading{\enLR{1.11} אתיקה וחוק – חשוב מאוד}

בקורס הזה נלמד איך התקפות עובדות – אבל אך ורק כדי לדעת להתגונן מפניהן. חשוב להבהיר כבר עכשיו: הרצת כלי סריקה או פריצה על רשת שאינה שלכם, וללא אישור בכתב, היא \underline{עבירה פלילית} לפי \textbf{חוק המחשבים, התשנ"ה-\enLR{1995}} – גם אם עשיתם זאת "רק כדי לבדוק". הכלים עצמם (כמו \enLR{Nmap)} חוקיים לגמרי; מה שאסור הוא השימוש בהם בלי רשות. אם מצאתם חולשה במערכת של מישהו אחר, הדבר הנכון הוא לדווח עליה בלבד (מה שנקרא "גילוי אחראי"), ולא לנצל אותה.\par

\sectionheading{\enLR{1.12} תרגילים – עם פתרונות}

\begin{exbox}\boxtitle{תרגיל \enLR{1} – סיווג לפי \enLR{CIA}}\par  לכל אירוע, ציינו איזה עיקרון נפגע: (א) תלמיד מצא דרך לשנות ציון במערכת בית הספר. (ב) אתר ההרשמה של בית הספר קרס ביום ההרשמה. (ג) רשימת הטלפונים של ההורים דלפה לקבוצת וואטסאפ. (ד) כופרה הצפינה את שרת הקבצים. \par\vskip2pt\hrule\vskip3pt\textbf{}(א) שלמות (ב) זמינות (ג) סודיות (ד) זמינות (ואם התוקף גם מאיים לפרסם את הקבצים – גם סודיות).\end{exbox}

\begin{exbox}\boxtitle{תרגיל \enLR{2} – וירוס, תולעת או טרויאני?}\par  (א) משחק "מיינקראפט חינם.\enLR{exe"} שמתקין תוכנת ריגול. (ב) קובץ \enLR{Excel} עם מאקרו שמדביק את שאר קובצי ה-\enLR{Excel} כשפותחים אותו. (ג) תוכנה שסורקת את הרשת, מוצאת מחשבים פגיעים ומעתיקה את עצמה אליהם לבד. \par\vskip2pt\hrule\vskip3pt\textbf{}(א) סוס טרויאני (המשתמש מתקין מרצון). (ב) וירוס (צריך שיפתחו את הקובץ). (ג) תולעת (מתפשטת לבד).\end{exbox}

\begin{exbox}\boxtitle{תרגיל \enLR{3} – חוזק סיסמה}\par  מדוע סיסמה של \enLR{12} תווים חזקה בהרבה מסיסמה של \enLR{8} תווים, גם אם שתיהן משתמשות באותם סוגי תווים? \par\vskip2pt\hrule\vskip3pt\textbf{}כי כל תו נוסף מכפיל את מספר הצירופים האפשריים. ההפרש בין \enLR{8} ל-\enLR{12} תווים הוא לא "קצת יותר" אלא פי עשרות מיליונים – מה שהופך ניחוש שיטתי (ברוט-פורס) מבלתי-מעשי לחלוטין.\end{exbox}

\begin{exbox}\boxtitle{תרגיל \enLR{4} – חשיבה על סיפורים אמיתיים}\par  לכל סיפור, איזו בקרה (מנהלית / פיזית / לוגית) הייתה מקטינה את הנזק? (א) \enLR{WannaCry.} (ב) פריצת טוויטר \enLR{2020.} (ג) \enLR{Mirai.} \par\vskip2pt\hrule\vskip3pt\textbf{}(א) בקרה מנהלית – נוהל עדכונים תקופתי (ולוגית – גיבוי לשחזור). (ב) מנהלית – נוהל שאוסר מסירת גישה בטלפון; לוגית – \enLR{MFA} והרשאות מוגבלות בכלי הניהול. (ג) לוגית – שינוי סיסמת ברירת המחדל של המכשירים.\end{exbox}

\sectionheading{\enLR{1.13} מפת הקורס – כל ההגנות שנלמד}

סיימנו בסקירת האיומים. שאר הקורס הוא למעשה סדרה של תשובות לאיומים האלה: נתב מוקשח (פרק \enLR{2)}, אימות מרכזי עם מודל \enLR{AAA} (פרק \enLR{3)}, חומות אש ורשימות גישה (פרק \enLR{4)}, מערכות לזיהוי ומניעת חדירה \enLR{IDS/IPS} (פרק \enLR{5)}, אבטחת הרשת המקומית והמתגים (פרק \enLR{6)}, הצפנה (פרק \enLR{7)}, רשתות פרטיות וירטואליות \enLR{VPN} (פרק \enLR{8)}, וניהול סיכונים כולל (פרק \enLR{9).} כל פרק עונה על איום שראינו כאן.\par

\sectionheading{בנק שאלות שתלמידים שואלים – ותשובות מוכנות}

אוסף שאלות שתלמידים סקרנים נוטים לשאול בפרק זה, עם תשובות מוכנות. מסודר לפי נושא – סרקו לפני השיעור.\par

\subheading{מוטיבציה ותפיסות שגויות}

\textbf{ש: "אין לי מה להסתיר / אני לא מעניין אף אחד, אז למה שיתקפו אותי?"}\newline ת: רוב ההתקפות אינן ממוקדות אלא \underline{אוטומטיות וסיטונאיות}. גם מחשב "לא מעניין" שווה לתוקף: אפשר להפוך אותו ל-\enLR{bot} ברשת זומבים (פרק \enLR{1.4)}, לגנוב סיסמאות שחוזרות על עצמן באתרים אחרים \enLR{(credential stuffing)}, להשתמש בו כקרש קפיצה לרשת של בית הספר/ההורים, או להצפין אותו בכופרה ולדרוש כסף. כופרה לא בוררת מטרות.\par

\textbf{ש: "אנטי-וירוס מספיק כדי להיות מוגן?"}\newline ת: לא. אנטי-וירוס קלאסי מזהה לפי \textbf{חתימות} – תבניות של נוזקות \underline{ידועות}. נוזקה חדשה \enLR{(zero-day)} או כזו ששונתה מעט \enLR{(polymorphic)} לא תזוהה. לכן צריך "הגנה לעומק": עדכוני מערכת, הגנת רשת (חומת אש/\enLR{IPS)}, גיבוי, והכי חשוב – משתמש שלא לוחץ על כל דבר.\par

\textbf{ש: "מה זה \enLR{zero-day}?"}\newline ת: פגיעות שהיצרן עדיין \underline{לא יודע עליה} או שאין לה תיקון – "יום אפס" מאז שהתגלתה. תוקף שמחזיק \enLR{zero-day} יכול לנצל אותה בלי שאף הגנה מבוססת-חתימות תזהה אותה. אלה הפגיעויות היקרות והמסוכנות ביותר.\par

\textbf{ש: "האקר = פושע?"}\newline ת: לא בהכרח. \textbf{\enLR{White hat}} – בודק חדירה חוקי, באישור. \textbf{\enLR{Black hat}} – פושע. \textbf{\enLR{Gray hat}} – פורץ ללא אישור אך מדווח. המילה "\enLR{hacker"} במקורה היא מי שמבין מערכות לעומק; התקשורת הצמידה לה משמעות שלילית.\par

\subheading{נוזקות}

\textbf{ש: "אז מה בדיוק ההבדל בין וירוס לתולעת? זה מבלבל."}\newline ת: שאלת המפתח: \underline{האם צריך פעולה של משתמש כדי שיתפשט?} וירוס נדבק לקובץ וזקוק למשתמש שיריץ אותו. תולעת מתפשטת \underline{לבד} דרך הרשת בלי אף פעולת אדם (למשל תולעת מוריס, \enLR{Slammer).} טרויאני לא משכפל את עצמו כלל – המשתמש מתקין אותו מרצון כי הוא נראה לגיטימי.\par

\textbf{ש: "האם מק / לינוקס / אייפון יכולים להידבק בנוזקה?"}\newline ת: כן, כל מערכת יכולה. \enLR{Windows} הוא היעד הנפוץ ביותר כי הוא הכי נפוץ, אבל יש נוזקות ל-\enLR{macOS}, לינוקס, אנדרואיד ו-\enLR{iOS.} אף מערכת אינה "חסינה" – זו תפיסה שגויה מסוכנת.\par

\textbf{ש: "כדאי לשלם כופר בכופרה?"}\newline ת: הגישה הרשמית (גם של ה-\enLR{FBI} ומערך הסייבר) היא \underline{לא לשלם}: אין ערובה שתקבלו מפתח, זה מממן פשע, והופך אתכם ליעד חוזר. הפתרון האמיתי הוא \textbf{גיבוי} תקין (פרק \enLR{9)} שמאפשר לשחזר בלי לשלם.\par

\subheading{התקפות והגנה}

\textbf{ש: "אם תוקף מפיל אותנו ב-\enLR{DDoS}, למה שלא פשוט נחסום את הכתובת שלו?"}\newline ת: כי ב-\enLR{DDoS} אין כתובת אחת – ההתקפה מגיעה בו-זמנית מאלפי מכשירים נגועים \enLR{(botnet)} בכל העולם, ולעיתים עם כתובות מקור \underline{מזויפות}. חסימה ידנית לא מספיקה; צריך סינון בקנה מידה גדול, שירותי ניקוי בענן, ו-\enLR{rate limiting.}\par

\textbf{ש: "אדם אחד יכול לבצע \enLR{DDoS}?"}\newline ת: כן – אם הוא שולט ב-\enLR{botnet} (רשת מחשבים נגועים) או שוכר אחד בדארקנט ("\enLR{DDoS} כשירות"). ההבחנה היא לא מספר האנשים אלא מספר \underline{מקורות} התעבורה: \enLR{DoS} = מקור אחד, \enLR{DDoS} = מקורות רבים.\par

\textbf{ש: "איך \enLR{phishing} עוקף אנטי-וירוס?"}\newline ת: כי לרוב אין בו נוזקה לסרוק – יש רק קישור לאתר מזויף או בקשה להזין פרטים. זו מניפולציה של \underline{האדם}, לא של המחשב. אנטי-וירוס לא "רואה" שאתם מקלידים סיסמה באתר מתחזה.\par

\textbf{ש: "אם יש נעילה אחרי \enLR{3} ניסיונות, למה צריך בכלל סיסמה חזקה?"}\newline ת: כי לא תמיד התוקף עובר דרך מסך ההתחברות. אם דלף \underline{מסד הנתונים} של הסיסמאות (גיבובים), התוקף מריץ מיליארדי ניחושים אופליין על החומרה שלו – שם אין נעילה. סיסמה ארוכה הופכת את זה לבלתי מעשי.\par

\textbf{ש: "האם זה חוקי להריץ \enLR{Nmap} או להתקין \enLR{Kali Linux}?"}\newline ת: הכלים עצמם חוקיים לחלוטין (זה כמו פטיש). מה שאסור הוא \underline{השימוש} בהם על רשת שאינה שלכם וללא אישור בכתב – זו עבירה על חוק המחשבים התשנ"ה-\enLR{1995.} ברשת מעבדה מבודדת שלכם – מותר ומומלץ ללמוד.\par

\textbf{ש: "מה ההבדל המדויק בין איום, פגיעות וסיכון?"}\newline ת: \textbf{איום} = הגורם המסוכן (האקר, שריפה) – קיים תמיד, לא בשליטתנו. \textbf{פגיעות} = החולשה שהאיום מנצל (סיסמה חלשה, פורט פתוח) – זה מה שבשליטתנו. \textbf{סיכון} = ההסתברות שהאיום ינצל את הפגיעות, כפול הנזק. מנהלים סיכון על ידי צמצום פגיעויות.\par

\sectionheading{מילון מונחים – הגדרות}

הגדרות תמציתיות של המונחים המרכזיים בפרק. שימושי גם כדף עזר לבחינה (מותר מילון מונחים).\par

\begin{xltabular}{\linewidth}{|>{\setRTL\raggedleft\arraybackslash}X|>{\setRTL\raggedleft\arraybackslash}X|}
\hline
\rowcolor{hdrbg}\textbf{הגדרה} & \textbf{מונח} \\
\hline
\endhead
הגנה על מידע ומערכות מפני גישה, שינוי, חשיפה או השמדה בלתי מורשים. & \textbf{אבטחת מידע \enLR{(Information Security)}} \\
\hline
תת-תחום: הגנה על תשתית הרשת ועל התעבורה שעוברת בה. & \textbf{אבטחת רשת \enLR{(Network Security)}} \\
\hline
הבטחה שרק גורמים מורשים רואים את המידע. מושגת בהצפנה ובבקרת גישה. & \textbf{סודיות \enLR{(Confidentiality)}} \\
\hline
הבטחה שהמידע לא שונה ללא הרשאה. מושגת בגיבוב וחתימה דיגיטלית. & \textbf{שלמות \enLR{(Integrity)}} \\
\hline
הבטחה שהמידע והשירות נגישים כשצריך. נפגעת ב-\enLR{DoS} ובכופרה. & \textbf{זמינות \enLR{(Availability)}} \\
\hline
הוכחת זהות – מי אתה (שם משתמש וסיסמה, ביומטריה). & \textbf{אימות \enLR{(Authentication)}} \\
\hline
מניעת אפשרות של השולח להכחיש ששלח; מושגת בחתימה דיגיטלית. & \textbf{אי-הכחשה \enLR{(Non-repudiation)}} \\
\hline
כל דבר בעל ערך שמגנים עליו: מידע, שרת, מוניטין, זמינות. & \textbf{נכס \enLR{(Asset)}} \\
\hline
חולשה שניתן לנצל: באג, סיסמה חלשה, פורט פתוח, עובד לא מודרך. & \textbf{פגיעות \enLR{(Vulnerability)}} \\
\hline
גורם שעלול לנצל פגיעות: האקר, נוזקה, עובד ממורמר, אסון טבע. & \textbf{איום \enLR{(Threat)}} \\
\hline
הסתברות שאיום ינצל פגיעות, כפול הנזק. = איום × פגיעות × ערך הנכס. & \textbf{סיכון \enLR{(Risk)}} \\
\hline
אמצעי שמקטין סיכון. מנהלית (נהלים), פיזית (מנעולים), לוגית (הצפנה, \enLR{ACL).} & \textbf{בקרה \enLR{(Control)}} \\
\hline
שם כולל לתוכנה זדונית: וירוס, תולעת, טרויאני, כופרה ועוד. & \textbf{נוזקה \enLR{(Malware)}} \\
\hline
קוד זדוני שנצמד לקובץ ומתפשט כשמשתמש מריץ אותו. & \textbf{וירוס \enLR{(Virus)}} \\
\hline
תוכנה שמשכפלת ומתפשטת לבד ברשת, ללא פעולת משתמש. & \textbf{תולעת \enLR{(Worm)}} \\
\hline
תוכנה שנראית לגיטימית אך מכילה קוד זדוני; אינה משכפלת את עצמה. & \textbf{סוס טרויאני \enLR{(Trojan)}} \\
\hline
נוזקה שמצפינה קבצים ודורשת תשלום; פוגעת בזמינות. & \textbf{כופרה \enLR{(Ransomware)}} \\
\hline
נוזקה שמסתתרת עמוק במערכת ההפעלה ומסווה נוזקות אחרות. & \textbf{\enLR{Rootkit}} \\
\hline
רשת מחשבים נגועים (זומבים) הנשלטים מרחוק לביצוע התקפות. & \textbf{\enLR{Botnet}} \\
\hline
דלת אחורית שעוקפת את מנגנון האימות ומאפשרת גישה סמויה. & \textbf{\enLR{Backdoor}} \\
\hline
מניפולציה של בני אדם כדי להשיג מידע או גישה (במקום פריצה טכנית). & \textbf{הנדסה חברתית \enLR{(Social Engineering)}} \\
\hline
דוא"ל/הודעה מתחזים המבקשים ללחוץ על קישור או למסור פרטים. & \textbf{\enLR{Phishing}} \\
\hline
מניעת שירות – הצפת שירות עד קריסה, ממקור אחד \enLR{(DoS)} או ממקורות רבים \enLR{(DDoS).} & \textbf{\enLR{DoS / DDoS}} \\
\hline
הצפת שרת בבקשות \enLR{SYN} חצי-פתוחות עד שטבלת החיבורים מתמלאת. & \textbf{\enLR{SYN flood}} \\
\hline
שליחת קלט ארוך מהחוצץ, שגולש ודורס זיכרון ומריץ קוד תוקף. & \textbf{\enLR{Buffer overflow}} \\
\hline
ניסוי שיטתי של כל הסיסמאות האפשריות (או רשימת מילון) עד למציאה. & \textbf{\enLR{Brute force}} \\
\hline
שלב הסיור – איסוף מידע על המטרה \enLR{(whois}, סריקת פורטים, \enLR{sniffing).} & \textbf{\enLR{Reconnaissance}} \\
\hline
פגיעות שאין לה עדיין תיקון או שאינה ידועה ליצרן. & \textbf{\enLR{Zero-day}} \\
\hline
\end{xltabular}


\clearpage
\phantomsection\addcontentsline{toc}{section}{פרק \enLR{2} – אבטחת אביזרי רשת}

\sloppy
\chapterheading{פרק \enLR{2} – אבטחת אביזרי רשת}{\small\color{subgray}\enLR{10} שעות עיוני + \enLR{2} מעשי · שבועות \enLR{4}–\enLR{6}\par}\vskip4pt

\begin{defbox}\boxtitle{איך לקרוא את הפרק הזה}\par  בפרק הקודם הבנו מפני מה מגנים ומי התוקפים. בפרק הזה מתחילים להגן בפועל – על הרכיב הראשון והחשוב ביותר ברשת: הנתב. הפרק כתוב לקריאה רציפה גם למי שלא ראה מעולם ציוד רשת. נתחיל מהבסיס ממש (מה זה נתב, איך "מדברים" איתו), ורק אז נראה פקודות. כל פקודה מוסברת במילים לפני שהיא מופיעה.\end{defbox}

\sectionheading{\enLR{2.1} למה מתחילים דווקא מהנתב?}

נזכיר מהפרק הקודם: \textbf{נתב \enLR{(Router)}} הוא המכשיר שמעביר את התעבורה בין רשתות שונות – למשל בין הרשת של בית הספר לבין האינטרנט. הנתב שיושב על הגבול הזה נקרא \textbf{נתב קצה \enLR{(Edge Router)}}, והוא נקודת המפגש בין הרשת שלנו לעולם החיצון.\par

למה הוא המטרה הראשונה של תוקפים? כי מי ששולט בנתב שולט בכל מה שעובר דרכו. הוא יכול לצותת לכל התעבורה, לחסום אותה, להטות אותה למקום אחר, או לפתוח דרכה דלת אל תוך הרשת. במילים אחרות, הנתב הוא כמו השומר בשער: אם הוא נפרץ, כל שאר ההגנות שמאחוריו כבר לא עוזרות. לכן הצעד הראשון בכל רשת מאובטחת הוא "להקשיח" את הנתב.\par

\textbf{הקשחה \enLR{(Hardening)}} פירושה להפוך מכשיר מהגדרת ברירת המחדל ה"פתוחה" שלו למצב מאובטח: לתת סיסמאות חזקות, לאפשר גישה מוצפנת בלבד, לכבות שירותים מיותרים, ולהפעיל ניטור. את כל אלה נלמד בפרק. נהוג לחלק את אבטחת הנתב לשלושה תחומים: אבטחה \textbf{פיזית} (מי בכלל יכול לגעת בו – חדר נעול), אבטחת \textbf{מערכת ההפעלה} (גרסה עדכנית וגיבוי), ו\textbf{הקשחה} של ההגדרות (לב הפרק).\par

\begin{storybox}\boxtitle{מהעולם האמיתי: חצי מיליון נתבים נחטפו \enLR{(VPNFilter, 2018)}}\par  נוזקה בשם \enLR{VPNFilter} הדביקה מעל \enLR{500,000} נתבים ביתיים וקטנים בעולם. היכולות שלה היו מפחידות: לרגל אחר כל התעבורה שעוברת בנתב, לגנוב סיסמאות, ואף "להרוג" את הנתב בפקודה אחת. ה-\enLR{FBI} ייחס אותה לתקיפה מדינתית והשתלט על שרת השליטה שלה כדי לעצור אותה. איך היא נכנסה? דרך חולשות ידועות שלא עודכנו וסיסמאות ברירת מחדל שלא שונו. הלקח: נתב הוא מחשב מלא לכל דבר, ובלי הקשחה הוא הדלת הקלה ביותר פנימה.\end{storybox}

\sectionheading{\enLR{2.2} איך בכלל "מדברים" עם נתב?}

בניגוד למחשב רגיל, לנתב אין מסך, עכבר וחלונות. מנהלים אותו בעזרת \textbf{\enLR{CLI (Command Line Interface)}} – "ממשק שורת פקודה", כלומר מקלידים פקודות טקסט והנתב מגיב. מערכת ההפעלה של ציוד סיסקו נקראת \textbf{\enLR{IOS}} (כמו ש-\enLR{Windows} היא מערכת ההפעלה של המחשב, \enLR{IOS} היא מערכת ההפעלה של הנתב), וה-\enLR{CLI} הוא הדרך לדבר איתה.\par

\subheading{מצבי העבודה – והסימן שאומר לכם איפה אתם}

ל-\enLR{CLI} יש כמה "מצבים" (רמות), וכל אחד מאפשר דברים אחרים. הדבר החשוב: \underline{הסימן שמופיע על המסך תמיד אומר לכם באיזה מצב אתם נמצאים}. הסימן בנוי משלושה חלקים: שם הנתב, ואחריו – בסוגריים – החלק שאתם מגדירים כרגע, ואחריו הסימן \code{>} או \code{\#}. הסימן \code{>} אומר "אתם רק יכולים להסתכל", והסימן \code{\#} אומר "אתם מנהלים ומותר לכם לשנות" (חשבו על \code{\#} כמו על מנהל/\enLR{root).}\par

\begin{itemize}
\item \code{R1>} – \textbf{מצב משתמש \enLR{(User)}}. זה מה שרואים מיד בכניסה. אפשר רק להסתכל: פקודות \code{show} בסיסיות, \code{ping}.
\item \code{R1\#} – \textbf{מצב מנהל \enLR{(Privileged)}}. מגיעים אליו בהקלדת \code{enable}. כאן מותר הכול, כולל לצפות בכל ההגדרות ולאתחל את הנתב.
\item \code{R1(config)\#} – \textbf{מצב הגדרות \enLR{(Global config)}}. מגיעים אליו בהקלדת \code{configure terminal}. כאן משנים הגדרות של הנתב כולו (שם, משתמשים, ניתוב).
\item \code{R1(config-if)\#} – \textbf{מצב ממשק \enLR{(Interface)}}. מגיעים אליו בהקלדת למשל \code{interface g0/0}, ומגדירים פורט מסוים. (\textbf{ממשק / \enLR{Interface}} הוא פשוט פורט/כרטיס רשת בנתב, למשל \enLR{GigabitEthernet0/0}, בקיצור \enLR{g0/0.)}
\end{itemize}

\begin{tipbox}\boxtitle{מושג חשוב: איפה נשמרות ההגדרות?}\par  לנתב יש שני "עותקים" של ההגדרות. \textbf{\enLR{running-config}} הוא הקונפיגורציה הפעילה כרגע – אבל היא נמצאת בזיכרון זמני ו\underline{נמחקת בכיבוי}. \textbf{\enLR{startup-config}} היא הקונפיגורציה השמורה, שנטענת מחדש בכל הדלקה. לכן אחרי שמסיימים להגדיר – חייבים לשמור, בפקודה \code{write} (או \code{copy running-config startup-config}). מי ששוכח לשמור מגלה שכל עבודתו נעלמה אחרי אתחול.\end{tipbox}

הנה "סיפור" קצר של פקודות. קראו את הסימן בתחילת כל שורה כדי לעקוב איפה אנחנו נמצאים:\par

\begin{codebox}\begin{english}\codefont\footnotesize\color{cfg}\setlength{\parindent}{0pt}
Router> \textcolor{cbold}{\textbf{enable}}                 \textcolor{ccom}{\texthebrew{! ממצב "רק להסתכל" למצב "מנהל"}}\\
Router\# \textcolor{cbold}{\textbf{configure terminal}}     \textcolor{ccom}{\texthebrew{! נכנסים למצב הגדרות}}\\
Router(config)\# \textcolor{cbold}{\textbf{hostname R1}}    \textcolor{ccom}{\texthebrew{! נותנים לנתב שם – שימו לב שהסימן משתנה מיד}}\\
R1(config)\# \textcolor{cbold}{\textbf{interface g0/0}}     \textcolor{ccom}{\texthebrew{! נכנסים לפורט מסוים}}\\
R1(config-if)\# \textcolor{cbold}{\textbf{no shutdown}}     \textcolor{ccom}{\texthebrew{! מדליקים את הפורט}}\\
R1(config-if)\# \textcolor{cbold}{\textbf{end}}             \textcolor{ccom}{\texthebrew{! חוזרים ישר למצב מנהל}}\\
R1\# \textcolor{cbold}{\textbf{write}}                       \textcolor{ccom}{\texthebrew{! שומרים! (אחרת ההגדרות ייעלמו בכיבוי)}}
\end{english}\end{codebox}

\sectionheading{\enLR{2.3} סיסמאות: מה חזק, מה חלש, ולמה}

הצעד הבסיסי בהקשחה הוא סיסמאות טובות. ב-\enLR{IOS} יש כמה סוגים של סיסמאות, וההבדל ביניהן קריטי. הכלל הפשוט: תמיד להשתמש במילה \code{secret}, אף פעם לא במילה \code{password}.\par

\begin{codebox}\begin{english}\codefont\footnotesize\color{cfg}\setlength{\parindent}{0pt}
R1(config)\# \textcolor{cbold}{\textbf{enable secret}} Str0ng-P@ss   \textcolor{ccom}{\texthebrew{! סיסמת מצב מנהל – מאוחסנת מגובבת (בטוח)}}\\
R1(config)\# enable password weak         \textcolor{ccom}{\texthebrew{! מאוחסנת בטקסט גלוי – לא להשתמש!}}\\
R1(config)\# \textcolor{cbold}{\textbf{service password-encryption}}  \textcolor{ccom}{\texthebrew{! מסתיר את הסיסמאות הגלויות בקונפיג}}\\
R1(config)\# username admin \textcolor{cbold}{\textbf{secret}} P@ss   \textcolor{ccom}{\texthebrew{! משתמש מקומי עם סיסמה מגובבת}}
\end{english}\end{codebox}

למה \code{secret} עדיף? כי הוא נשמר כ\textbf{גיבוב \enLR{(hash)}} – טביעת אצבע חד-כיוונית של הסיסמה, שממנה אי אפשר לשחזר את הסיסמה המקורית (נרחיב על גיבוב בפרק \enLR{7).} לעומת זאת \code{password} נשמר בטקסט גלוי, וגם הפקודה \code{service password-encryption} שמסתירה אותו משתמשת בהצפנה חלשה מאוד (מה שנקרא "\enLR{type 7")} שנשברת בשנייה באתרים חינמיים באינטרנט. לכן \code{service password-encryption} מגן רק מפני מישהו שמציץ מעבר לכתף בזמן שאתם צופים בקונפיג – לא מפני תוקף אמיתי. ההגנה האמיתית היא ה-\code{secret}.\par

\sectionheading{\enLR{2.4 Telnet} מול \enLR{SSH} – ההבדל שמציל אתכם}

כדי לנהל נתב מרחוק (מבלי להתחבר אליו פיזית בכבל), מתחברים אליו דרך הרשת. יש לכך שני פרוטוקולים, ואחד מהם מסוכן מאוד.\par

\textbf{\enLR{Telnet}} (פורט \enLR{23)} הוא הפרוטוקול הישן. הבעיה שלו חמורה: הוא מעביר את \underline{הכול} בטקסט גלוי, כולל את שם המשתמש והסיסמה. כל מי ש"מקשיב" לרשת יכול לראות אותם. \textbf{\enLR{SSH (Secure Shell)}} (פורט \enLR{22)} הוא הפרוטוקול המאובטח: הוא \underline{מצפין} את כל התקשורת, כך שמאזין רואה רק ג'יבריש, וגם מוודא שאתם מדברים באמת עם הנתב הנכון. משתמשים ב-\enLR{SSH} גרסה \enLR{2} בלבד.\par

\begin{storybox}\boxtitle{מהעולם האמיתי: \enLR{Telnet} שנלכד בבית חולים}\par  בבדיקת חדירה מוזמנת, הבודק חיבר מחשב נייד לשקע רשת בחדר המתנה והריץ תוכנה שמקשיבה לתעבורה \enLR{(Wireshark).} תוך \enLR{40} דקות הוא לכד תקשורת \enLR{Telnet} של טכנאי שהתחבר למתג הראשי – ושם המשתמש והסיסמה השתקפו על מסכו בטקסט גלוי. עם הסיסמה הזו הוא קיבל שליטה בכל המתגים בבניין. הטכנאי לא עשה שום דבר "לא נכון" – הוא פשוט השתמש בכלי הישן שהיה זמין. ההגנה מפני כל התרחיש הזה היא שורה אחת שחוסמת \enLR{Telnet} ומחייבת \enLR{SSH.}\end{storybox}

\subheading{הגדרת \enLR{SSH} – חמישה שלבים, כל אחד מוסבר}

כדי ש-\enLR{SSH} יעבוד, הנתב צריך לייצר לעצמו \textbf{זוג מפתחות \enLR{RSA}} – מפתח ציבורי ומפתח פרטי (נבין אותם לעומק בפרק \enLR{7).} המפתחות האלה הם מה שמצפין את החיבור. כדי לייצר אותם, הנתב חייב קודם שם מלא, שמורכב מהשם שלו \enLR{(hostname)} ומשם דומיין. לכן שני אלה חייבים לבוא ראשונים:\par

\begin{codebox}\begin{english}\codefont\footnotesize\color{cfg}\setlength{\parindent}{0pt}
R1(config)\# \textcolor{cbold}{\textbf{hostname R1}}                       \textcolor{ccom}{\texthebrew{! \enLR{1.} שם לנתב}}\\
R1(config)\# \textcolor{cbold}{\textbf{ip domain-name school.local}}       \textcolor{ccom}{\texthebrew{! \enLR{2.} שם דומיין (חובה לפני יצירת מפתחות)}}\\
R1(config)\# \textcolor{cbold}{\textbf{crypto key generate rsa}} modulus 2048  \textcolor{ccom}{\texthebrew{! \enLR{3.} ייצור זוג מפתחות \enLR{RSA}}}\\
R1(config)\# username admin privilege 15 secret P@ss \textcolor{ccom}{\texthebrew{! \enLR{4.} משתמש שאיתו נתחבר}}\\
R1(config)\# ip ssh version 2                         \textcolor{ccom}{\texthebrew{! \enLR{SSHv2} בלבד}}\\
R1(config)\# \textcolor{cbold}{\textbf{line vty 0 4}}                          \textcolor{ccom}{\texthebrew{! \enLR{5.} הגדרת קווי הגישה מרחוק}}\\
R1(config-line)\# \textcolor{cbold}{\textbf{login local}}                      \textcolor{ccom}{\texthebrew{! אימות מול המשתמש המקומי}}\\
R1(config-line)\# \textcolor{cbold}{\textbf{transport input ssh}}              \textcolor{ccom}{\texthebrew{! רק \enLR{SSH} – חוסם \enLR{Telnet}}}
\end{english}\end{codebox}

\begin{mistakebox}\boxtitle{הטעות הנפוצה ביותר}\par  אם שוכחים את שורת \code{ip domain-name} (או משאירים את שם הנתב כברירת המחדל "\enLR{Router")}, הפקודה \code{crypto key generate rsa} נכשלת עם ההודעה "\enLR{Please define a domain-name first".} זכרו: השם והדומיין חייבים לבוא \underline{לפני} יצירת המפתחות, כי שם המפתח נגזר מהם.\end{mistakebox}

הערה על \code{line vty 0 4}: "\enLR{VTY"} הם קווי הגישה הווירטואליים שדרכם מתחברים מרחוק (יש חמישה כאלה, ממוספרים \enLR{0} עד \enLR{4).} עליהם אנחנו קובעים שני דברים: \code{login local} – שהאימות ייעשה מול מסד המשתמשים המקומי (כלומר יידרש גם שם וגם סיסמה), ו-\code{transport input ssh} – שמותר להתחבר רק ב-\enLR{SSH} ולא ב-\enLR{Telnet.}\par

\begin{tipbox}\boxtitle{\enLR{login} מול \enLR{login local} – הבחנה שנשאלת בבגרות}\par  \code{login} (לבד) מבקש רק סיסמה אחת משותפת של הקו, בלי שם משתמש. \code{login local} מבקש גם שם משתמש וגם סיסמה, ומאמת אותם מול מסד המשתמשים המקומי (\code{username ... secret}). מכיוון ש-\enLR{SSH} דורש שם משתמש, הוא חייב \code{login local} – ולכן זו התשובה הנכונה בשאלות בגרות על התחברות מאובטחת.\end{tipbox}

\sectionheading{\enLR{2.5} עוד כלי הקשחה: רמות הרשאה ו-\enLR{Banner}}

\textbf{רמות הרשאה \enLR{(Privilege Levels)}:} לפעמים רוצים לתת למישהו גישה חלקית – למשל למתמחה שיוכל רק לצפות, בלי לשנות. ל-\enLR{IOS} יש \enLR{16} רמות הרשאה \enLR{(0} עד \enLR{15).} רמה \enLR{1} היא ברירת המחדל (רק צפייה), רמה \enLR{15} היא הרשאת מנהל מלאה, ואת הרמות שביניהן אפשר להתאים אישית ולהעביר אליהן פקודות מסוימות. כך אפשר לתת לכל תפקיד בדיוק את מה שהוא צריך – עיקרון שנקרא "מינימום הרשאות".\par

\textbf{\enLR{Banner}:} זוהי הודעת אזהרה שמופיעה בכניסה לנתב. מטרתה משפטית – להבהיר שהגישה מיועדת למורשים בלבד ושהפעילות מנוטרת.\par

\begin{codebox}\begin{english}\codefont\footnotesize\color{cfg}\setlength{\parindent}{0pt}
R1(config)\# \textcolor{cbold}{\textbf{banner motd \#}}\\
*** גישה בלתי מורשית אסורה. הפעילות מנוטרת ומתועדת. ***\\
\#
\end{english}\end{codebox}

\begin{mistakebox}\boxtitle{מה לא לכתוב בבאנר}\par  לא לכתוב "\enLR{Welcome"} (ברוכים הבאים) – במקרה משפטי בארה"ב פורץ טען שהבאנר "הזמין" אותו להיכנס. לא לכתוב את שם החברה ולא את דגם הנתב – זה רק עוזר לתוקף. הבאנר צריך רק להזהיר, לא לקבל את פני האורח.\end{mistakebox}

\begin{storybox}\boxtitle{מהעולם האמיתי: הדגל האמריקאי על \enLR{168,000} נתבים \enLR{(2018)}}\par  קבוצת האקרים גילתה שירות ניהול של סיסקו בשם \enLR{Smart Install} שנשאר פתוח וחשוף לאינטרנט אצל עשרות אלפי ארגונים. הם מצאו את הנתבים החשופים דרך מנוע חיפוש בשם \enLR{Shodan}, איפסו את הקונפיגורציה של כ-\enLR{168,000} נתבים, ובחלקם השאירו הודעה עם דגל אמריקאי וכיתוב "אל תתעסקו עם הבחירות שלנו". נפגעו ספקיות אינטרנט באיראן וברוסיה. הלקח ישיר: שירות ניהול מיותר שנשאר פתוח הוא הזמנה – ולכן מכבים כל שירות שלא נחוץ (הנושא של סעיף \enLR{2.7).}\end{storybox}

\sectionheading{\enLR{2.6} ניטור: לדעת מה קורה בנתב}

אבטחה טובה היא לא רק חסימה – היא גם לדעת מה קורה. לשם כך יש שלושה שירותים שכדאי להפעיל.\par

\subheading{\enLR{Syslog} – יומן האירועים}

הנתב מייצר כל הזמן הודעות על אירועים: ממשק שעלה או נפל, ניסיון כניסה כושל, חבילה שנחסמה. \textbf{\enLR{Syslog}} הוא המנגנון ששולח את ההודעות האלה ל\textbf{שרת מרכזי} (על פורט \enLR{UDP 514).} למה שרת מרכזי ולא רק בנתב עצמו? כי תוקף חכם שמצליח להיכנס לנתב ימחק את היומן המקומי כדי לטשטש עקבות – אבל אין לו גישה לשרת החיצוני, ושם הראיות נשמרות.\par

\begin{codebox}\begin{english}\codefont\footnotesize\color{cfg}\setlength{\parindent}{0pt}
R1(config)\# \textcolor{cbold}{\textbf{logging host 10.0.0.5}}       \textcolor{ccom}{\texthebrew{! שולח את היומן לשרת}}\\
R1(config)\# \textcolor{cbold}{\textbf{logging trap warnings}}       \textcolor{ccom}{\texthebrew{! שולח רמות \enLR{0} עד \enLR{4}}}
\end{english}\end{codebox}

להודעות היומן יש \textbf{\enLR{8} רמות חומרה}, מ-\enLR{0 (Emergency}, החמור ביותר) עד \enLR{7 (Debug}, מידע פנימי מפורט). כשבוחרים רמה מקבלים אותה וכל מה שחמור ממנה – למשל אם בוחרים רמה \enLR{4 (warnings)} מקבלים את רמות \enLR{0} עד \enLR{4.}\par

\subheading{\enLR{NTP} – למה זמן מדויק הוא עניין של אבטחה}

\textbf{\enLR{NTP (Network Time Protocol)}} מסנכרן את השעון של כל הציוד ברשת למקור זמן אחד. זה נשמע טכני ומשעמם, אבל הוא קריטי לאבטחה: בלי זמן מסונכרן, אי אפשר לשחזר את סדר האירועים בזמן חקירה. אם הנתב מראה שעה אחת והשרת שעה אחרת, לא תדעו אם הכניסה הכושלת לנתב קרתה לפני החדירה לשרת או אחריה. חוץ מזה, גם תעודות דיגיטליות ומפתחות \enLR{VPN} תלויים בשעון נכון.\par

\begin{codebox}\begin{english}\codefont\footnotesize\color{cfg}\setlength{\parindent}{0pt}
R1(config)\# \textcolor{cbold}{\textbf{ntp server 10.0.0.5}}
\end{english}\end{codebox}

\subheading{\enLR{SNMP} – ניטור מצב הציוד}

\textbf{\enLR{SNMP}} הוא פרוטוקול שמאפשר לתחנת ניהול מרכזית לאסוף נתונים מהנתב (כמה תעבורה עוברת, עומס המעבד). הבעיה: בגרסאות הישנות \enLR{(v1, v2c)} הגישה מאובטחת רק ב"מחרוזת קהילה" \enLR{(community string)} שעוברת בטקסט גלוי, ולעיתים קרובות נשארת ברירת המחדל "\enLR{public".} רק \textbf{גרסה \enLR{3 (SNMPv3)}} מוסיפה אימות והצפנה אמיתיים. הכלל: אם משתמשים ב-\enLR{SNMP}, גרסה \enLR{3} בלבד.\par

\begin{storybox}\boxtitle{מהעולם האמיתי: \enLR{Shodan} – "גוגל של המכשירים המחוברים"}\par  \enLR{Shodan} הוא מנוע חיפוש שסורק כל הזמן את האינטרנט ומקטלג מכשירים חשופים – נתבים, מתגים, מצלמות אבטחה. אפשר לחפש בו מכשירים עם \enLR{SNMP "public"}, עם \enLR{Telnet} פתוח, או עם סיסמת ברירת מחדל, ולקבל רשימה של אלפים. תוקפים משתמשים בו לשלב הסיור (פרק \enLR{1)}; אנשי אבטחה משתמשים בו כדי לגלות מה הארגון שלהם חושף בטעות. כל פעולת הקשחה שאנחנו לומדים בפרק הזה מורידה אתכם מהרשימה של \enLR{Shodan.}\end{storybox}

\sectionheading{\enLR{2.7} כיבוי שירותים מיותרים ו-\enLR{AutoSecure}}

נתב מגיע מהמפעל עם הרבה שירותים דלוקים שהיו שימושיים פעם והיום הם רק שטח תקיפה מיותר. עיקרון ההקשחה פשוט: \textbf{כל שירות שלא צריך – כבה אותו}. כל שירות פתוח הוא דלת אפשרית (זכרו את סיפור \enLR{Cisco Smart Install}!).\par

\begin{codebox}\begin{english}\codefont\footnotesize\color{cfg}\setlength{\parindent}{0pt}
R1(config)\# no ip http server      \textcolor{ccom}{\texthebrew{! ממשק ניהול \enLR{web} מיותר}}\\
R1(config)\# no cdp run             \textcolor{ccom}{\texthebrew{! פרוטוקול שחושף דגם וגרסה לשכנים}}
\end{english}\end{codebox}

כדי לחסוך זמן, קיימת פקודה אחת שמבצעת חלק גדול מההקשחה אוטומטית: \textbf{\enLR{AutoSecure}}. היא מכבה שירותים מסוכנים, מפעילה ניטור, ומקשיחה את הנתב בכמה שאלות פשוטות.\par

\begin{codebox}\begin{english}\codefont\footnotesize\color{cfg}\setlength{\parindent}{0pt}
R1\# \textcolor{cbold}{\textbf{auto secure}}
\end{english}\end{codebox}

\begin{mistakebox}\boxtitle{זהירות עם \enLR{AutoSecure}}\par  היא "עיוורת" – עלולה לכבות דברים שהרשת בכל זאת צריכה (למשל \enLR{CDP} שהמתגים משתמשים בו). לכן מריצים אותה, בודקים מה היא שינתה בעזרת \code{show auto secure config}, ומכווננים לפי הצורך.\end{mistakebox}

\sectionheading{\enLR{2.8} עוד דוגמאות מהעולם האמיתי}

שלושת הסיפורים בפרק – \enLR{VPNFilter}, הדגל האמריקאי \enLR{(Cisco Smart Install)} ו-\enLR{Shodan} – אינם אנקדוטות; הם ההוכחה החיה לכל מה שלמדנו. בכל אחד מהם, נתב או שירות שלא הוקשח היה הדלת פנימה: סיסמת ברירת מחדל שלא שונתה, שירות מיותר שנשאר פתוח, או ניהול לא מוצפן שנחשף. שורת המחץ של הפרק: הקשחה הופכת את הנתב משער פתוח לשער נעול.\par

\sectionheading{\enLR{2.9} תרגילים – עם פתרונות}

\begin{exbox}\boxtitle{תרגיל \enLR{1} – מצאו את הבעיות בקונפיג}\par  \enLR{Router(config)\# enable password cisco Router(config)\# line vty 0 4 Router(config-line)\# password cisco Router(config-line)\# login Router(config-line)\# transport input telnet Router(config)\# banner motd \# Welcome to ACME Corp}! \# \par\vskip2pt\hrule\vskip3pt\textbf{}\enLR{1.} \code{enable password} (גלוי) במקום \code{enable secret}. \enLR{2.} הסיסמה "\enLR{cisco"} חלשה מאוד. \enLR{3.} \code{login} בלי שם משתמש (צריך \code{login local} עם \code{username}). \enLR{4.} \code{transport input telnet} – לא מוצפן, צריך \code{ssh}. \enLR{5.} באנר "\enLR{Welcome"} עם שם החברה. \enLR{6.} השם נשאר "\enLR{Router"} ואין \code{ip domain-name} – כך שאי אפשר בכלל ליצור מפתחות ל-\enLR{SSH.}\end{exbox}

\begin{exbox}\boxtitle{תרגיל \enLR{2} – סדרו את הפקודות}\par  באיזה סדר צריך להריץ כדי ש-\enLR{SSH} יעבוד? \code{crypto key generate rsa} · \code{login local} · \code{hostname R1} · \code{transport input ssh} · \code{ip domain-name lab.local} · \code{line vty 0 4} \par\vskip2pt\hrule\vskip3pt\textbf{}\code{hostname R1} → \code{ip domain-name lab.local} → \code{crypto key generate rsa} → \code{line vty 0 4} → \code{login local} → \code{transport input ssh}. השם והדומיין חייבים לבוא לפני יצירת המפתחות.\end{exbox}

\begin{exbox}\boxtitle{תרגיל \enLR{3} – למה שרת \enLR{syslog} מרכזי?}\par  תוקף הצליח להיכנס לנתב. מדוע עדיף שהיומן יישלח לשרת חיצוני ולא יישמר רק בנתב עצמו? \par\vskip2pt\hrule\vskip3pt\textbf{}כי תוקף עם גישה לנתב יכול למחוק את היומן המקומי כדי לטשטש את עקבותיו. ליומן שנשלח לשרת חיצוני אין לו גישה, ולכן הראיות נשמרות שם – מה שמאפשר לחקור את האירוע.\end{exbox}

\begin{exbox}\boxtitle{תרגיל \enLR{4} – תכנון הקשחה}\par  אתם מנהלי הרשת של בית ספר עם שלושה נתבים. כתבו שישה צעדי הקשחה שתבצעו על כל נתב, לפי סדר עדיפות, ונמקו את הראשון. \par\vskip2pt\hrule\vskip3pt\textbf{}דוגמה: \enLR{1.} \code{enable secret} + \code{username ... secret} – בלי הגנת סיסמה בסיסית כל השאר חסר ערך. \enLR{2. SSH} בלבד (\code{login local} + \code{transport input ssh}), כיבוי \enLR{Telnet. 3.} \code{login block-for} נגד ניחוש סיסמאות. \enLR{4. syslog + NTP} לשרת מרכזי. \enLR{5. SNMPv3} וכיבוי שירותים מיותרים \enLR{(HTTP, CDP). 6.} באנר אזהרה + גיבוי הקונפיג.\end{exbox}

\sectionheading{בנק שאלות שתלמידים שואלים – ותשובות מוכנות}

שאלות נפוצות בפרק הקשחת הנתב, עם תשובות מוכנות.\par

\subheading{סיסמאות וגישה}

\textbf{ש: "אם אני שם סיסמה חזקה ל-\enLR{Telnet}, למה בכל זאת צריך \enLR{SSH}?"}\newline ת: כי \enLR{Telnet} מעביר את \underline{הכול} בטקסט גלוי, כולל את הסיסמה החזקה עצמה. כל מי שמאזין לרשת \enLR{(Wireshark)} רואה אותה. \enLR{SSH} מצפין את כל התקשורת. חוזק הסיסמה לא עוזר אם היא נשלחת גלויה.\par

\textbf{ש: "מה ההבדל בין \code{password} ל-\code{secret}?"}\newline ת: \code{secret} נשמר כ\underline{גיבוב חד-כיווני} (לא ניתן לשחזר). \code{password} נשמר גלוי, ואפילו עם \code{service password-encryption} הוא רק "\enLR{type 7"} שמפוענח באתר אונליין בשנייה. תמיד להשתמש ב-\code{secret}.\par

\textbf{ש: "אם \code{service password-encryption} זו הצפנה, למה היא חלשה?"}\newline ת: כי זה אלגוריתם הפיך וישן (וריאציה על \enLR{Vigen}è\enLR{re).} הוא מגן רק מפני מישהו שמציץ מעבר לכתף בזמן שאתם צופים בקונפיג – לא מפני תוקף אמיתי. ה"חוזק" האמיתי הוא ב-\code{secret} (גיבוב).\par

\textbf{ש: "מה קורה אם שכחתי את ה-\enLR{enable secret}?"}\newline ת: יש הליך \textbf{\enLR{password recovery}} דרך גישה פיזית לקונסול: מאתחלים, נכנסים ל-\enLR{ROMMON}, משנים את ה-\enLR{configuration register (0x2142)} כדי לדלג על ה-\enLR{startup-config}, ומגדירים סיסמה חדשה. המסקנה החשובה: \underline{מי שיש לו גישה פיזית לנתב – שולט בו}, ולכן אבטחה פיזית קריטית.\par

\textbf{ש: "מה ההבדל בין רמות הרשאה \enLR{(privilege levels)} לתצוגות \enLR{(views)}?"}\newline ת: רמות הרשאה גסות – רמה נתונה מקבלת אוטומטית גם את כל מה שמתחתיה. תצוגות \enLR{(parser view)} מדויקות – מגדירים רשימת פקודות מדויקת לכל תפקיד, אך דורשות שה-\enLR{AAA} יופעל.\par

\subheading{\enLR{SSH}}

\textbf{ש: "למה צריך \code{ip domain-name} כדי ליצור מפתחות \enLR{SSH}?"}\newline ת: כי שם המפתח נגזר מהשם המלא של הנתב = \enLR{hostname + domain-name (FQDN).} בלי דומיין הנתב לא יודע איזה שם לתת למפתח, והפקודה נכשלת עם "\enLR{Please define a domain-name first".}\par

\textbf{ש: "אחרי \code{transport input ssh} ניתקתי בטעות את \enLR{Telnet} – זו טעות?"}\newline ת: לא, זה בדיוק מה שרוצים – לחסום את \enLR{Telnet} הלא-מאובטח ולאפשר רק \enLR{SSH.} חיבור \enLR{Telnet} חדש ייכשל, וזו התנהגות רצויה.\par

\subheading{שירותים וניטור}

\textbf{ש: "למה זמן מדויק \enLR{(NTP)} קשור לאבטחה בכלל?"}\newline ת: כי בלי זמן מסונכרן, הלוגים מכל המכשירים לא ניתנים לשחזור – אי אפשר לדעת מה קרה לפני מה כשחוקרים אירוע. בנוסף, תעודות דיגיטליות ומפתחות \enLR{VPN} תלויים בזמן נכון.\par

\textbf{ש: "\enLR{SNMPv2} באמת כל כך גרוע?"}\newline ת: הבעיה שלו היא ה-\enLR{community string} שעובר בטקסט גלוי (ולעיתים נשאר "\enLR{public").} אם חייבים \enLR{v2c} – לפחות \enLR{read-only} + הגבלה ב-\enLR{ACL} לכתובת תחנת הניהול. אבל הנכון הוא \enLR{v3}, שמוסיף אימות \enLR{(HMAC)} והצפנה.\par

\textbf{ש: "אז למה שלא נריץ תמיד \code{auto secure} ונגמור עם זה?"}\newline ת: כי הוא "עיוור" – יכול לכבות \enLR{CDP} שהמתגים צריכים, להפעיל חומת אש שתחסום תעבורה לגיטימית, ואי אפשר לבטלו בפקודה אחת. מריצים אותו, בודקים מה שינה עם \code{show auto secure config}, ומכווננים.\par

\textbf{ש: "מה זה בעצם \enLR{CLI}, ולמה לא ללמוד ממשק גרפי?"}\newline ת: \enLR{CLI} הוא ממשק שורת הפקודה של הנתב. אנחנו לומדים אותו כי הוא אוניברסלי (עובד על כל ציוד סיסקו וגם דרך \enLR{SSH)}, מדויק, וניתן לאוטומציה – \underline{והבגרות בודקת אותו בלבד}. הממשקים הגרפיים \enLR{(SDM/CCP)} יצאו משימוש.\par

\sectionheading{מילון מונחים – הגדרות}

הגדרות תמציתיות של המונחים המרכזיים בפרק. שימושי גם כדף עזר לבחינה (מותר מילון מונחים).\par

\begin{xltabular}{\linewidth}{|>{\setRTL\raggedleft\arraybackslash}X|>{\setRTL\raggedleft\arraybackslash}X|}
\hline
\rowcolor{hdrbg}\textbf{הגדרה} & \textbf{מונח} \\
\hline
\endhead
ממשק שורת פקודה לניהול ציוד רשת בהקלדת פקודות טקסט. & \textbf{\enLR{CLI (Command Line Interface)}} \\
\hline
מערכת ההפעלה של ציוד סיסקו. & \textbf{\enLR{IOS}} \\
\hline
ההגדרות הפעילות בזיכרון; נמחקות בכיבוי אם לא נשמרו. & \textbf{\enLR{running-config}} \\
\hline
ההגדרות השמורות שנטענות באתחול הנתב. & \textbf{\enLR{startup-config}} \\
\hline
דרכי גישה לנתב: קונסול (פיזי), \enLR{VTY} (מרחוק ברשת), \enLR{AUX} (מודם ישן). & \textbf{\enLR{Console / VTY / AUX}} \\
\hline
ניהול דרך רשת הנתונים \enLR{(In-band)} או דרך רשת ניהול נפרדת \enLR{(Out-of-band).} & \textbf{\enLR{In-band / OOB}} \\
\hline
סיסמת מצב מנהל, נשמרת כגיבוב חד-כיווני (בטוח). & \textbf{\enLR{enable secret}} \\
\hline
מסתיר סיסמאות בקונפיג בהצפנה חלשה \enLR{(type 7)} – מגן רק מהצצה. & \textbf{\enLR{service password-encryption}} \\
\hline
פרוטוקול ניהול מרחוק מוצפן (פורט \enLR{22)}; מחליף את \enLR{Telnet} הלא-מאובטח. & \textbf{\enLR{SSH (Secure Shell)}} \\
\hline
פרוטוקול ניהול מרחוק לא מוצפן (פורט \enLR{23)} – מסוכן, לא לשימוש. & \textbf{\enLR{Telnet}} \\
\hline
זוג מפתחות (ציבורי/פרטי) שהנתב מייצר כדי לאפשר הצפנת \enLR{SSH.} & \textbf{מפתחות \enLR{RSA}} \\
\hline
\enLR{16} רמות \enLR{(0}–\enLR{15)} הקובעות אילו פקודות מותרות; \enLR{15} = הכול. & \textbf{רמות הרשאה \enLR{(Privilege levels)}} \\
\hline
תצוגת תפקיד – רשימת פקודות מדויקת למשתמש; דורשת \enLR{AAA.} & \textbf{\enLR{Parser View}} \\
\hline
הודעת אזהרה משפטית המוצגת בכניסה לנתב. & \textbf{\enLR{Banner}} \\
\hline
מנגנון יומן אירועים; שולח הודעות לשרת מרכזי \enLR{(UDP 514).} & \textbf{\enLR{Syslog}} \\
\hline
\enLR{8} רמות חומרה ליומן \enLR{(0=Emergency} החמור ביותר עד \enLR{7=Debug).} & \textbf{\enLR{Severity levels}} \\
\hline
מסנכרן שעונים בכל הציוד \enLR{(UDP 123)}; קריטי ללוגים ולתעודות. & \textbf{\enLR{NTP (Network Time Protocol)}} \\
\hline
פרוטוקול לניטור ציוד \enLR{(UDP 161)}; גרסה \enLR{3} בלבד מאובטחת. & \textbf{\enLR{SNMP}} \\
\hline
פקודה שמבצעת הקשחה אוטומטית של הנתב בצעד אחד. & \textbf{\enLR{AutoSecure}} \\
\hline
סגירת שירותים מיותרים והגדרות אבטחה למזעור שטח התקיפה. & \textbf{\enLR{Hardening} (הקשחה)} \\
\hline
מצב שבו פורט מושבת אוטומטית עקב הפרת אבטחה. & \textbf{\enLR{err-disabled}} \\
\hline
\end{xltabular}


\clearpage
\phantomsection\addcontentsline{toc}{section}{פרק \enLR{3} – מודל ה-\enLR{AAA}}

\sloppy
\chapterheading{פרק \enLR{3} – מודל ה-\enLR{AAA}}{\small\color{subgray}\enLR{11} שעות עיוני + \enLR{3} מעשי · שבועות \enLR{7}–\enLR{9}\par}\vskip4pt

\begin{metabox}\textbf{מטרות:} מודל \enLR{AAA} · שרתי \enLR{AAA (ACS/ISE)} · \enLR{TACACS}+ ו-\enLR{RADIUS} עם מפתחות מוצפנים\\ \textbf{בבגרות:} הגדרת שלושת השירותים, \code{aaa new-model}, \enLR{RADIUS} מול \enLR{TACACS}+\end{metabox}

\begin{defbox}\boxtitle{לפני שמתחילים – מה זו בכלל "הזדהות"? (למי שאין רקע)}\par  כשאתם נכנסים לחשבון (מייל, בנק, נתב) קורים שלושה דברים שאנשים נוטים לבלבל – וזה בדיוק מה שמודל \enLR{AAA} מסדר:  \textbf{הזדהות \enLR{(Authentication)}} – להוכיח \underline{מי אתה} (שם משתמש + סיסמה). כמו להראות תעודת זהות. \textbf{הרשאה \enLR{(Authorization)}} – לקבוע \underline{מה מותר לך} אחרי שנכנסת (לצפות בלבד? לשנות הכול?). כמו דרגת גישה בעבודה. \textbf{רישום \enLR{(Accounting)}} – לתעד \underline{מה עשית} (מתי נכנסת, אילו פקודות הקלדת). כמו יומן מעקב.  \textbf{למה מרכזי?} בלי מערכת מרכזית, כל נתב מחזיק רשימת משתמשים משלו. אם יש \enLR{50} נתבים ועובד עוזב – צריך למחוק אותו מ-\enLR{50} מקומות. מודל \enLR{AAA} מאפשר "מקום אחד" שכל הנתבים שואלים אותו (שרת). זה כל הרעיון. \newline \textbf{מושג שיחזור:} \textbf{שרת} = מחשב מרכזי שנותן שירות (כאן: שירות "מי מורשה"). \textbf{לקוח} של השרת הזה = הנתב עצמו, ששואל את השרת "האם להכניס את המשתמש הזה?".\end{defbox}

\begin{tipbox}\boxtitle{הערה על תכנית הלימודים}\par  רשימת הנושאים של פרק \enLR{3} בתכנית הרשמית הועתקה בטעות מפרק \enLR{2.} הפרק כאן בנוי לפי \emph{המטרות} של פרק \enLR{3} – שהן גם בדיוק מה שהבגרות שואלת.\end{tipbox}

\sectionheading{\enLR{3.1} שלוש ה-\enLR{A}}

עד עכשיו כל האימות היה \textbf{מקומי}: שם משתמש וסיסמה בקונפיג של כל נתב. זה לא מתרחב: \enLR{50} נתבים × \enLR{10} מנהלים = \enLR{500} חשבונות לתחזק, ועובד שעוזב צריך להימחק מ-\enLR{50} מקומות. \textbf{\enLR{AAA}} הוא המסגרת שפותרת את זה – בין אם המשתמשים מקומיים ובין אם בשרת מרכזי.\par

\begin{xltabular}{\linewidth}{|>{\setRTL\raggedleft\arraybackslash}X|>{\setRTL\raggedleft\arraybackslash}X|>{\setRTL\raggedleft\arraybackslash}X|>{\setRTL\raggedleft\arraybackslash}X|}
\hline
\rowcolor{hdrbg}\textbf{דוגמה} & \textbf{הגדרה (כפי שנדרש בבגרות)} & \textbf{השאלה שהוא עונה עליה} & \textbf{שירות} \\
\hline
\endhead
שם משתמש + סיסמה, כרטיס חכם, טביעת אצבע & \textbf{זיהוי ואימות משתמשים} & מי אתה? & \textbf{\enLR{Authentication}} (אימות) \\
\hline
רמה \enLR{15} או רק \enLR{show}; גישה ל-\enLR{VLAN 10} בלבד; מותר \enLR{reload}? & \textbf{ניהול הרשאות} & מה מותר לך לעשות? & \textbf{\enLR{Authorization}} (הרשאה) \\
\hline
נכנס ב-\enLR{08:12}, הקליד \code{reload} ב-\enLR{08:14}, ניתק ב-\enLR{08:15; 2GB} תעבורה & \textbf{רישום למעקב של פעולות המשתמש} & מה עשית, מתי וכמה? & \textbf{\enLR{Accounting}} (רישום/חיוב) \\
\hline
\end{xltabular}

\begin{faqbox}\boxtitle{שאלת תלמיד: "מה ההבדל בין \enLR{Authentication} ל-\enLR{Authorization}? זה נשמע אותו דבר"}\par  המשל שעובד: שדה תעופה. \textbf{אימות} = בדיקת הדרכון בכניסה – מוודאים שאתה מי שאתה טוען. \textbf{הרשאה} = כרטיס הטיסה – גם אחרי שידוע מי אתה, מותר לך לעלות רק לטיסה שלך ולא לכל מטוס. \textbf{\enLR{Accounting}} = רישום הנוסעים – מי עלה, מתי, עם כמה מזוודות. אי אפשר לעשות הרשאה בלי אימות קודם, אבל אימות לבד לא אומר שמותר לך הכול.\end{faqbox}

\subheading{גורמי אימות \enLR{(Authentication factors)}}

\begin{itemize}
\item \textbf{משהו שאתה יודע} – סיסמה, \enLR{PIN.}
\item \textbf{משהו שיש לך} – טלפון \enLR{(OTP/SMS)}, טוקן, כרטיס חכם.
\item \textbf{משהו שאתה} – ביומטריה: טביעת אצבע, פנים.
\end{itemize}

\textbf{\enLR{MFA / 2FA}} – שילוב של שני גורמים \emph{מסוגים שונים}. שתי סיסמאות אינן \enLR{MFA.}\par

\begin{storybox}\boxtitle{סיפור מהחיים: מנהל הרשת שפוטר ביום שלישי}\par  בחברת שירותי אינטרנט בינונית פיטרו מנהל רשת. המנכ"ל ביקש מה-\enLR{IT "}לבטל לו את הגישה". ה-\enLR{IT} מחק את חשבון ה-\enLR{Windows} שלו. אבל ב-\enLR{40} הנתבים והמתגים היו חשבונות \emph{מקומיים} (\code{username} בקונפיג) – שהוא עצמו יצר. שלושה ימים אחר כך, ב-\enLR{2} בלילה, כל הנתבים "אבדו" את הקונפיג. אף אחד לא ידע מי עשה את זה – כי לא היה \enLR{Accounting}, ובלוג היה רק "\enLR{admin".} עם \enLR{AAA} מבוסס שרת: לחיצה אחת ב-\enLR{ISE} הייתה מנתקת אותו מכל הציוד, וה-\enLR{accounting} היה מראה שם מלא, שעה ופקודה. זה בדיוק "למה \enLR{AAA".}\end{storybox}

\sectionheading{\enLR{3.2} שני מצבי \enLR{AAA}: מקומי ומבוסס-שרת}

\begin{xltabular}{\linewidth}{|>{\setRTL\raggedleft\arraybackslash}X|>{\setRTL\raggedleft\arraybackslash}X|>{\setRTL\raggedleft\arraybackslash}X|}
\hline
\rowcolor{hdrbg}\textbf{\enLR{AAA} מבוסס שרת \enLR{(Server-based)}} & \textbf{\enLR{AAA} מקומי \enLR{(Local / Self-contained)}} & \textbf{} \\
\hline
\endhead
בשרת מרכזי \enLR{(Cisco ISE / ACS, FreeRADIUS, Windows NPS)} & בקונפיג של הנתב (\code{username}) & איפה המשתמשים \\
\hline
\enLR{RADIUS} או \enLR{TACACS}+ בין הנתב לשרת & – & פרוטוקול \\
\hline
ארגונים: ניהול מרכזי, מדיניות אחידה, לוג מרכזי & רשתות קטנות, גיבוי \enLR{(fallback)} כשהשרת לא זמין & מתאים ל \\
\hline
\end{xltabular}

בטרמינולוגיה: הנתב/מתג/\enLR{ASA} שפונה לשרת נקרא \textbf{\enLR{NAS}} \enLR{(Network Access Server)} או \textbf{\enLR{AAA client}}. השרת מחזיק את מסד המשתמשים – או פונה בעצמו ל-\enLR{Active Directory / LDAP.}\par

\sectionheading{\enLR{3.3 AAA} מקומי – הפקודות}

\begin{codebox}\begin{english}\codefont\footnotesize\color{cfg}\setlength{\parindent}{0pt}
\textcolor{ccom}{\texthebrew{! \enLR{1.} הפעלת המודל – מרגע זה כל ה-\enLR{lines} משתמשות ב-\enLR{AAA} (ולא ב-\enLR{login / login local)}}}\\
R1(config)\# \textcolor{cbold}{\textbf{aaa new-model}}\\
\textcolor{ccom}{\texthebrew{! \enLR{2.} משתמש מקומי}}\\
R1(config)\# username admin privilege 15 secret Adm1n-P@ss\\
\textcolor{ccom}{\texthebrew{! \enLR{3.} רשימת שיטות אימות ברירת מחדל: \enLR{local}}}\\
R1(config)\# \textcolor{cbold}{\textbf{aaa authentication login default local}}\\
\textcolor{ccom}{\texthebrew{! \enLR{4.} רשימה בשם – למשל לקונסול: ללא אימות (זהירות!) או \enLR{enable password}}}\\
R1(config)\# \textcolor{cbold}{\textbf{aaa authentication login CONSOLE-LIST local-case}}\\
R1(config)\# line console 0\\
R1(config-line)\# \textcolor{cbold}{\textbf{login authentication CONSOLE-LIST}}
\end{english}\end{codebox}

\begin{mistakebox}\boxtitle{הטעות המסוכנת ביותר בפרק}\par  \code{aaa new-model} משנה מיד את התנהגות כל הקווים. אם הקלדתם אותה בלי משתמש מקומי ובלי \code{aaa authentication login default} ואז התנתקתם – \textbf{ננעלתם בחוץ} (ברירת המחדל אחרי \enLR{aaa new-model} היא \enLR{local}, ואם אין \enLR{username} – אין כניסה). כלל ברזל: \textbf{קודם \code{username}, אחר כך \code{aaa new-model}, ותמיד להשאיר חלון \enLR{SSH} אחד פתוח עד שבודקים בחלון שני.}\end{mistakebox}

\textbf{שיטות אימות \enLR{(methods)} אפשריות ברשימה:} \code{local}, \code{local-case} (תלוי רישיות), \code{enable} (סיסמת \enLR{enable)}, \code{group radius}, \code{group tacacs+}, \code{group שם}, \code{none} (ללא אימות – רק לקונסול בבדיקות). ניתן לשרשר: הנתב מנסה את הראשונה; \textbf{רק אם היא לא זמינה} (שרת לא עונה) הוא עובר לבאה. תשובת "סיסמה שגויה" מהשרת \underline{לא} גורמת למעבר לשיטה הבאה.\par

\sectionheading{\enLR{3.4 RADIUS} מול \enLR{TACACS}+ – הטבלה שחייבים לדעת}

\begin{xltabular}{\linewidth}{|>{\setRTL\raggedleft\arraybackslash}X|>{\setRTL\raggedleft\arraybackslash}X|>{\setRTL\raggedleft\arraybackslash}X|}
\hline
\rowcolor{hdrbg}\textbf{\enLR{TACACS}+} & \textbf{\enLR{RADIUS}} & \textbf{} \\
\hline
\endhead
\enLR{Cisco} (קנייני, אך מפורסם) & תקן פתוח \enLR{(IETF, RFC 2865)} & \textbf{מי פיתח / תקן} \\
\hline
\textbf{\enLR{TCP}} – פורט \enLR{49} & \textbf{\enLR{UDP}} – פורטים \enLR{1812} (אימות) ו-\enLR{1813 (accounting)}; ישן: \enLR{1645/1646} & \textbf{שכבת תעבורה} \\
\hline
\textbf{כל גוף החבילה} מוצפן (רק הכותרת גלויה) & \textbf{רק שדה הסיסמה} מוצפן בבקשת הגישה \enLR{(Access-Request)}; שם המשתמש ושאר המידע גלויים & \textbf{הצפנה} \\
\hline
שלושת השירותים \textbf{נפרדים} – ניתן לבצע אימות מול \enLR{AD} והרשאה מול \enLR{TACACS}+ & אימות והרשאה \textbf{משולבים} בתגובה אחת \enLR{(Access-Accept} מכיל גם את ההרשאות). לא ניתן לבצע הרשאה ללא אימות & \textbf{הפרדת השירותים} \\
\hline
נתמך – כל פקודה ש-המנהל מקליד נשלחת לאישור בשרת & לא נתמך & \textbf{הרשאה לכל פקודה} \\
\hline
קיים, לרוב לרישום פקודות & מפורט וחזק – פותח לחיוב \enLR{(ISP)} & \textbf{\enLR{Accounting}} \\
\hline
תומך גם ב-\enLR{AppleTalk, NetBIOS} וכו' (היסטורי) & \enLR{IP} בלבד (בעיקר) & \textbf{ריבוי פרוטוקולים} \\
\hline
\textbf{ניהול ציוד} \enLR{(device administration)}: מנהלי רשת שמתחברים לנתבים ולמתגים & \textbf{גישת משתמשים לרשת}: \enLR{Wi-Fi/802.1X, VPN} מרחוק, דיאל-אפ & \textbf{שימוש אופייני} \\
\hline
\end{xltabular}

\begin{exambox}\boxtitle{בבגרות (שאלה \enLR{4}ה) – נכון/לא נכון}\par  \enLR{1. RADIUS} מצפין רק את הסיסמה כשהוא שולח "בקשת גישה" – \textbf{נכון}.\newline  \enLR{2. RADIUS} מבצע אימות משתמשים ולא מתייחס לניהול הרשאות – \textbf{נכון} (במובן שאין הפרדה/הרשאה עצמאית; ההרשאות מגיעות רק כחלק מתשובת האימות).\newline  \enLR{3. TACACS}+ משתמש רק ב-\enLR{UDP} – \textbf{לא נכון} \enLR{(TCP 49).}\newline  \enLR{4. TACACS}+ מספק שירות \enLR{AAA} – \textbf{נכון}.\end{exambox}

\begin{faqbox}\boxtitle{שאלת תלמיד: "למה \enLR{RADIUS} על \enLR{UDP} אם \enLR{TCP} אמין יותר?"}\par  \enLR{RADIUS} תוכנן ב-\enLR{1991} לשרתי דיאל-אפ עם אלפי חיבורים; \enLR{UDP} קל, ללא לחיצת יד, והלקוח פשוט שולח שוב אם לא קיבל תשובה. \enLR{TACACS}+ בחר \enLR{TCP} כי הוא מנהל \emph{שיחה} רב-שלבית (שאלה–תשובה–שאלה) ורוצה לדעת מיד אם השרת נפל. שניהם עובדים היום; ההבדל המעשי החשוב הוא ההצפנה והפרדת השירותים.\end{faqbox}

\subheading{שיחה אופיינית – \enLR{RADIUS}}

\begin{defbox}\boxtitle{משתמש → \enLR{NAS} (נתב) → שרת \enLR{RADIUS}}\par  \enLR{1.} המשתמש מקליד שם וסיסמה. \enLR{2.} הנתב שולח \textbf{\enLR{Access-Request}} (הסיסמה מוצפנת עם ה-\enLR{shared secret + MD5). 3.} השרת עונה \textbf{\enLR{Access-Accept}} (עם \enLR{attributes} – \enLR{VLAN}, רמת הרשאה) / \textbf{\enLR{Access-Reject}} / \textbf{\enLR{Access-Challenge}} (נדרש גורם נוסף). \enLR{4.} הנתב שולח \textbf{\enLR{Accounting-Request (start)}}, ובסוף \enLR{(stop).}\end{defbox}

\subheading{שיחה אופיינית – \enLR{TACACS}+}

\begin{defbox}\boxtitle{שלוש שיחות נפרדות}\par  \enLR{Authentication: START} → \enLR{REPLY (GETUSER)} → \enLR{CONTINUE} (שם) → \enLR{REPLY (GETPASS)} → \enLR{CONTINUE} (סיסמה) → \enLR{REPLY (PASS/FAIL). Authorization: REQUEST} (הפקודה שהוקלדה) → \enLR{RESPONSE (PASS\_ADD/FAIL). Accounting: REQUEST (start/stop/watchdog)} → \enLR{REPLY.}\end{defbox}

\begin{storybox}\boxtitle{סיפור מהחיים: למה ה-\enLR{Wi-Fi} באוניברסיטה שואל שם משתמש – ובבית קפה סיסמה}\par  ברשת \enLR{eduroam} (שפועלת גם באוניברסיטאות בישראל) סטודנט מחיפה מתחבר ל-\enLR{Wi-Fi} באוניברסיטה בברלין עם החשבון של הטכניון. איך? המחשב שולח שם משתמש ל-\enLR{Access Point (802.1X)}, ה-\enLR{AP} הוא רק \enLR{Authenticator} ומעביר את הבקשה לשרת \enLR{RADIUS} בברלין, שרואה "@\enLR{technion.ac.il"} ומעביר לשרת \enLR{RADIUS} בטכניון, שמאמת ומחזיר \enLR{Access-Accept.} שרת \enLR{RADIUS} אחד לא צריך להכיר את כל המשתמשים בעולם – רק לדעת למי להעביר. זו הסיבה ש-\enLR{RADIUS} הוא "פרוטוקול גישת המשתמשים לרשת", ו-\enLR{TACACS}+ נשאר לניהול הציוד.\end{storybox}

\sectionheading{\enLR{3.5 Cisco ACS} ו-\enLR{ISE}}

\textbf{\enLR{ACS}} \enLR{(Access Control Server)} – שרת ה-\enLR{AAA} ההיסטורי של סיסקו (התכנית מזכירה אותו). תמך ב-\enLR{RADIUS} וב-\enLR{TACACS}+, עם \enLR{GUI} לניהול משתמשים, קבוצות ומדיניות. הוחלף ב-\textbf{\enLR{ISE}} \enLR{(Identity Services Engine)} – שרת מדיניות שעושה \enLR{AAA, 802.1X, NAC}, פרופיילינג של מכשירים ו-\enLR{BYOD.} חלופות חינמיות: \textbf{\enLR{FreeRADIUS}}, \textbf{\enLR{TACACS+ daemon}}, ו-\textbf{\enLR{Windows NPS}} \enLR{(RADIUS} על \enLR{Windows Server).} ב-\enLR{Packet Tracer} יש שרת \enLR{AAA} מובנה שתומך בשניהם.\par

\sectionheading{\enLR{3.6} הגדרת \enLR{AAA} מבוסס-שרת ב-\enLR{CLI}}

\subheading{\enLR{TACACS}+}

\begin{codebox}\begin{english}\codefont\footnotesize\color{cfg}\setlength{\parindent}{0pt}
R1(config)\# aaa new-model\\
\textcolor{ccom}{\texthebrew{! תחביר חדש \enLR{(IOS 15+)}:}}\\
R1(config)\# \textcolor{cbold}{\textbf{tacacs server ISE1}}\\
R1(config-server-tacacs)\# \textcolor{cbold}{\textbf{address ipv4 10.0.0.5}}\\
R1(config-server-tacacs)\# \textcolor{cbold}{\textbf{key T@cacsK3y}}            \textcolor{ccom}{\texthebrew{! \enLR{shared secret} – חייב להיות זהה בשרת}}\\
R1(config-server-tacacs)\# single-connection          \textcolor{ccom}{\texthebrew{! חיבור \enLR{TCP} אחד לכל השיחות}}\\
R1(config-server-tacacs)\# exit\\
\textcolor{ccom}{\texthebrew{! תחביר ישן (וב-\enLR{Packet Tracer)}:}}\\
R1(config)\# \textcolor{cbold}{\textbf{tacacs-server host 10.0.0.5 key T@cacsK3y}}\\
\textcolor{ccom}{\texthebrew{! אימות: קודם \enLR{TACACS}+, ואם השרת לא זמין – מקומי}}\\
R1(config)\# \textcolor{cbold}{\textbf{aaa authentication login default group tacacs+ local}}
\end{english}\end{codebox}

\subheading{\enLR{RADIUS}}

\begin{codebox}\begin{english}\codefont\footnotesize\color{cfg}\setlength{\parindent}{0pt}
R1(config)\# \textcolor{cbold}{\textbf{radius server NPS1}}\\
R1(config-radius-server)\# \textcolor{cbold}{\textbf{address ipv4 10.0.0.6 auth-port 1812 acct-port 1813}}\\
R1(config-radius-server)\# \textcolor{cbold}{\textbf{key R@diusK3y}}\\
R1(config-radius-server)\# exit\\
\textcolor{ccom}{\texthebrew{! תחביר ישן (וב-\enLR{Packet Tracer)}:}}\\
R1(config)\# \textcolor{cbold}{\textbf{radius-server host 10.0.0.6 key R@diusK3y}}\\
R1(config)\# aaa authentication login default group radius local
\end{english}\end{codebox}

\subheading{קבוצת שרתים \enLR{(server group)} – כמה שרתים, ולפי סדר}

\begin{codebox}\begin{english}\codefont\footnotesize\color{cfg}\setlength{\parindent}{0pt}
R1(config)\# \textcolor{cbold}{\textbf{aaa group server tacacs+ MY-TACACS}}\\
R1(config-sg-tacacs+)\# server name ISE1\\
R1(config-sg-tacacs+)\# server name ISE2\\
R1(config)\# aaa authentication login default \textcolor{cbold}{\textbf{group MY-TACACS}} local
\end{english}\end{codebox}

\begin{faqbox}\boxtitle{שאלת תלמיד: "מה זה ה-\enLR{key}? זה הסיסמה של המשתמש?"}\par  לא. ה-\textbf{\enLR{key (shared secret)}} הוא סוד משותף בין \emph{הנתב} ל\emph{שרת}, שמשמש להצפנת התקשורת ביניהם ולאימות שהנתב הוא באמת לקוח מורשה של השרת. סיסמאות המשתמשים נמצאות בשרת. שני מפתחות שונים לגמרי: המפתח בין הציוד לשרת, והסיסמה של האדם. "מפתחות מוצפנים" בתכנית הלימודים = ה-\enLR{key} הזה.\end{faqbox}

\sectionheading{\enLR{3.7 Authorization} – הרשאות}

\begin{codebox}\begin{english}\codefont\footnotesize\color{cfg}\setlength{\parindent}{0pt}
\textcolor{ccom}{\texthebrew{! מי רשאי לקבל \enLR{shell (EXEC)} ובאיזו רמה}}\\
R1(config)\# \textcolor{cbold}{\textbf{aaa authorization exec default group tacacs+ local}}\\
\textcolor{ccom}{\texthebrew{! אישור כל פקודה ברמה \enLR{15} מול השרת \enLR{(TACACS}+ בלבד)}}\\
R1(config)\# \textcolor{cbold}{\textbf{aaa authorization commands 15 default group tacacs+ local}}\\
\textcolor{ccom}{\texthebrew{! הרשאות לשירותי רשת \enLR{(PPP, VPN)} – בדרך כלל \enLR{RADIUS}}}\\
R1(config)\# \textcolor{cbold}{\textbf{aaa authorization network default group radius}}\\
\textcolor{ccom}{\texthebrew{! שהפקודות בקונפיג עצמו \enLR{(config mode)} גם יאושרו}}\\
R1(config)\# aaa authorization config-commands
\end{english}\end{codebox}

שימו לב: אחרי \code{aaa authorization exec}, אם השרת מחזיר \enLR{privilege 15} – המשתמש נכנס ישירות ל-\code{R1\#} בלי \enLR{enable.} זה נוח ומאובטח (אין סיסמת \enLR{enable} משותפת).\par

\sectionheading{\enLR{3.8 Accounting} – רישום}

\begin{codebox}\begin{english}\codefont\footnotesize\color{cfg}\setlength{\parindent}{0pt}
\textcolor{ccom}{\texthebrew{! רישום התחלה וסיום של כל \enLR{session}}}\\
R1(config)\# \textcolor{cbold}{\textbf{aaa accounting exec default start-stop group tacacs+}}\\
\textcolor{ccom}{\texthebrew{! רישום כל פקודה ברמה \enLR{15} – "מי הקליד \enLR{reload}?"}}\\
R1(config)\# \textcolor{cbold}{\textbf{aaa accounting commands 15 default start-stop group tacacs+}}\\
\textcolor{ccom}{\texthebrew{! חיבורי רשת \enLR{(VPN/PPP)} – חיוב לפי זמן/נפח}}\\
R1(config)\# \textcolor{cbold}{\textbf{aaa accounting network default start-stop group radius}}
\end{english}\end{codebox}

אפשרויות: \code{start-stop} (רשומה בהתחלה ובסוף), \code{stop-only} (רק בסיום, עם סיכום), \code{none}.\par

\sectionheading{\enLR{3.9 802.1X} – \enLR{AAA} לפורט של מתג (הצצה לפרק \enLR{6)}}

\enLR{RADIUS} משמש גם לאימות \emph{מכשירים} לפני שמקבלים גישה לרשת: תקן \textbf{\enLR{IEEE 802.1X}}. שלושה תפקידים: \textbf{\enLR{Supplicant}} (המחשב/הטלפון), \textbf{\enLR{Authenticator}} (המתג או ה-\enLR{Access Point} – חוסם את הפורט עד לאימות), \textbf{\enLR{Authentication Server}} \enLR{(RADIUS).} הפרוטוקול בין ה-\enLR{supplicant} למתג: \enLR{EAP over LAN (EAPOL)}; בין המתג לשרת: \enLR{RADIUS.} זה הבסיס ל-\enLR{NAC} (פרק \enLR{6)} ול-\enLR{WPA2-Enterprise.}\par

\begin{codebox}\begin{english}\codefont\footnotesize\color{cfg}\setlength{\parindent}{0pt}
S1(config)\# aaa new-model\\
S1(config)\# radius-server host 10.0.0.6 key R@diusK3y\\
S1(config)\# \textcolor{cbold}{\textbf{aaa authentication dot1x default group radius}}\\
S1(config)\# \textcolor{cbold}{\textbf{dot1x system-auth-control}}\\
S1(config)\# interface f0/1\\
S1(config-if)\# switchport mode access\\
S1(config-if)\# \textcolor{cbold}{\textbf{authentication port-control auto}}\\
S1(config-if)\# \textcolor{cbold}{\textbf{dot1x pae authenticator}}
\end{english}\end{codebox}

\sectionheading{\enLR{3.10} בדיקה ופתרון תקלות}

\begin{codebox}\begin{english}\codefont\footnotesize\color{cfg}\setlength{\parindent}{0pt}
R1\# \textcolor{cbold}{\textbf{show aaa sessions}}\\
R1\# \textcolor{cbold}{\textbf{show aaa user all}}\\
R1\# \textcolor{cbold}{\textbf{show tacacs}}\\
R1\# \textcolor{cbold}{\textbf{show radius statistics}}\\
R1\# \textcolor{cbold}{\textbf{test aaa group tacacs+ admin Adm1n-P@ss legacy}}   \textcolor{ccom}{\texthebrew{! בודק אימות מול השרת בלי להתחבר}}\\
R1\# \textcolor{cbold}{\textbf{debug aaa authentication}}\\
R1\# \textcolor{cbold}{\textbf{debug tacacs}}\\
R1\# \textcolor{cbold}{\textbf{debug radius}}
\end{english}\end{codebox}

\begin{xltabular}{\linewidth}{|>{\setRTL\raggedleft\arraybackslash}X|>{\setRTL\raggedleft\arraybackslash}X|}
\hline
\rowcolor{hdrbg}\textbf{סיבה נפוצה} & \textbf{תסמין} \\
\hline
\endhead
השרת לא נגיש \enLR{(IP/ACL}/שירות כבוי) – ה-\enLR{fallback} ל-\enLR{local} עבד. בדקו \enLR{ping} לשרת ו-\code{debug tacacs}. & הנתב "נתקע" \enLR{10}–\enLR{20} שניות ואז מקבל את המשתמש המקומי \\
\hline
ה-\enLR{key} לא תואם בין הנתב לשרת; או שהשרת לא מכיר את ה-\enLR{IP} של הנתב כלקוח \enLR{(AAA client).} ב-\enLR{Packet Tracer}: בדקו ב-\enLR{Server} → \enLR{AAA} → \enLR{Client Name/IP/Secret.} & "\enLR{Authentication failed"} מיד, גם עם סיסמה נכונה \\
\hline
חסר \code{aaa authorization exec}, או שהשרת לא מחזיר את ה-\enLR{attribute (shell:priv-lvl=15} ב-\enLR{TACACS+; Service-Type / cisco-avpair} ב-\enLR{RADIUS).} & המשתמש נכנס אבל לא מקבל רמה \enLR{15} \\
\hline
שכחו רשימה נפרדת לקונסול. תמיד: \code{aaa authentication login CONSOLE local} + \code{login authentication CONSOLE} על \enLR{line con 0.} & הקונסול דורש אימות שרת ואין שרת \\
\hline
\end{xltabular}

\sectionheading{\enLR{3.11} מדריך פקודות מרוכז – פרק \enLR{3}}

\begin{xltabular}{\linewidth}{|>{\setRTL\raggedleft\arraybackslash}X|>{\setRTL\raggedleft\arraybackslash}X|}
\hline
\rowcolor{hdrbg}\textbf{תפקיד} & \textbf{פקודה} \\
\hline
\endhead
הפעלת \enLR{AAA} (בבגרות: "הפקודה שתגרום לנתב לעבוד עם מודל \enLR{AAA")} & \code{aaa new-model} \\
\hline
רשימת אימות: שרת ואז מקומי & \code{aaa authentication login default group tacacs+ local} \\
\hline
רשימה בשם לקו ספציפי & \code{aaa authentication login NAME …} + \code{login authentication NAME} (ב-\enLR{line)} \\
\hline
הגדרת שרת \enLR{TACACS}+ & \code{tacacs-server host IP key K} / \code{tacacs server NAME} \\
\hline
הגדרת שרת \enLR{RADIUS} & \code{radius-server host IP key K} / \code{radius server NAME} \\
\hline
קבוצת שרתים & \code{aaa group server tacacs+ G} \\
\hline
הרשאת \enLR{shell} & \code{aaa authorization exec default …} \\
\hline
הרשאה לכל פקודה & \code{aaa authorization commands 15 default …} \\
\hline
רישום & \code{aaa accounting exec/commands/network default start-stop …} \\
\hline
בדיקה & \code{show aaa sessions}, \code{debug aaa authentication}, \code{test aaa} \\
\hline
\end{xltabular}

\sectionheading{\enLR{3.12} תרגילים לתלמידים}

\begin{exbox}\boxtitle{תרגיל \enLR{1} – איזה \enLR{A}?}\par  לכל משפט – \enLR{Authentication, Authorization} או \enLR{Accounting}:\newline  א. "המשתמש \enLR{dana} ניסתה להקליד \code{configure terminal} ונדחתה." ב. "החשבון של \enLR{yossi} נעול אחרי \enLR{3} ניסיונות." ג. "בדוח החודשי: \enLR{4,200} חיבורי \enLR{VPN}, ממוצע \enLR{38} דקות." ד. "המשתמש הזדהה עם טביעת אצבע." ה. "המורה יכולה לראות ציונים אבל לא לשנות." \par\vskip2pt\hrule\vskip3pt\textbf{}א. \enLR{Authorization} ב. \enLR{Authentication} ג. \enLR{Accounting} ד. \enLR{Authentication} ה. \enLR{Authorization}\end{exbox}

\begin{exbox}\boxtitle{תרגיל \enLR{2} – מה יקרה?}\par  נתב עם: \code{aaa authentication login default group tacacs+ group radius local none}. תארו את סדר הפעולות כאשר: (א) שרת ה-\enLR{TACACS}+ עונה "\enLR{FAIL".} (ב) שרת ה-\enLR{TACACS}+ לא עונה, ושרת ה-\enLR{RADIUS} עונה \enLR{Accept.} (ג) שני השרתים לא עונים ואין \enLR{username} בקונפיג. \par\vskip2pt\hrule\vskip3pt\textbf{}(א) המשתמש נדחה – לא עוברים לשיטה הבאה אחרי דחייה. (ב) אחרי \enLR{timeout} של \enLR{TACACS}+ הנתב פונה ל-\enLR{RADIUS} ומקבל את המשתמש. (ג) \enLR{TACACS+ timeout} → \enLR{RADIUS timeout} → \enLR{local} (אין משתמשים ⇒ השיטה "לא זמינה", לא "דחייה") → \enLR{none} ⇒ המשתמש נכנס בלי אימות! לכן \code{none} מסוכן.\end{exbox}

\begin{exbox}\boxtitle{תרגיל \enLR{3} – בחירת פרוטוקול}\par  לכל תרחיש בחרו \enLR{RADIUS} או \enLR{TACACS}+ ונמקו במשפט:\newline  א. אימות \enLR{2,000} תלמידים ל-\enLR{Wi-Fi} של בית הספר. ב. \enLR{5} מנהלי רשת שמנהלים \enLR{60} מתגים, ורוצים שמתמחה יוכל להקליד רק \enLR{show.} ג. ספק אינטרנט שמחייב לקוחות לפי נפח. ד. חברה שהרגולטור דורש ממנה תיעוד של כל פקודת קונפיגורציה. \par\vskip2pt\hrule\vskip3pt\textbf{}א. \enLR{RADIUS} (גישת משתמשים, \enLR{802.1X).} ב. \enLR{TACACS}+ (הרשאה לכל פקודה). ג. \enLR{RADIUS (accounting} מפורט לנפח). ד. \enLR{TACACS+ (command accounting).}\end{exbox}

\begin{exbox}\boxtitle{תרגיל \enLR{4} – תקלה}\par  מנהל הקליד \code{aaa new-model} ואז \code{aaa authentication login default group tacacs+} (בלי \enLR{local).} השרת נפל. המנהל מתנתק. מה קרה, ואיך מתקנים? \par\vskip2pt\hrule\vskip3pt\textbf{}ננעל בחוץ – אין שיטת גיבוי. תיקון: גישה פיזית לקונסול (אם לקונסול לא הוגדרה רשימה אחרת – גם הוא נעול!) ⇒ \enLR{password recovery} דרך \enLR{ROMMON (register 0x2142)}, ואז להוסיף \code{local} לרשימה ומשתמש מקומי. לקח: תמיד \code{… local} בסוף, ורשימה נפרדת לקונסול.\end{exbox}

\sectionheading{\enLR{3.13} שאלות בסגנון בגרות}

\begin{exambox}\boxtitle{שאלה \enLR{1}}\par  התאימו: \enLR{(1)} רישום הפקודות שהקליד המנהל \enLR{(2)} בדיקה אם מותר למשתמש להקליד \enLR{reload (3)} בדיקת שם משתמש וסיסמה. אפשרויות: \enLR{Authentication / Authorization / Accounting.} \par\vskip2pt\hrule\vskip3pt\textbf{}\enLR{(1) Accounting (2) Authorization (3) Authentication}\end{exambox}

\begin{exambox}\boxtitle{שאלה \enLR{2}}\par  מנהל הגדיר: \code{aaa authentication login default group radius local}. שרת ה-\enLR{RADIUS} זמין ומשיב "סיסמה שגויה". מה יקרה? \par\vskip2pt\hrule\vskip3pt\textbf{}המשתמש יידחה. המעבר ל-\enLR{local} קורה רק כששרת ה-\enLR{RADIUS} \emph{לא זמין} \enLR{(timeout)}, לא כשהוא דוחה.\end{exambox}

\begin{exambox}\boxtitle{שאלה \enLR{3}}\par  איזה פרוטוקול מתאים יותר לניהול ציוד רשת \enLR{(device administration)} עם אישור לכל פקודה, ומדוע? \par\vskip2pt\hrule\vskip3pt\textbf{}\enLR{TACACS}+ – מפריד בין אימות להרשאה ותומך בהרשאה לכל פקודה \enLR{(per-command authorization); RADIUS} משלב אימות והרשאה ואינו תומך בכך.\end{exambox}

\begin{exambox}\boxtitle{שאלה \enLR{4}}\par  השלימו: פרוטוקול \_\_\_\_\_\_ משתמש ב-\enLR{TCP} פורט \_\_\_\_\_\_ ומצפין את \_\_\_\_\_\_ החבילה; פרוטוקול \_\_\_\_\_\_ משתמש ב-\enLR{UDP} פורטים \_\_\_\_\_\_ ומצפין רק את \_\_\_\_\_\_. \par\vskip2pt\hrule\vskip3pt\textbf{}\enLR{TACACS}+ · \enLR{49} · כל גוף · \enLR{RADIUS} · \enLR{1812/1813} · הסיסמה\end{exambox}

\begin{exambox}\boxtitle{שאלה \enLR{5}}\par  כתבו את הפקודות שמגדירות שרת \enLR{TACACS}+ בכתובת \enLR{192.168.1.10} עם מפתח \code{Key123}, ואימות ברירת מחדל מולו עם גיבוי מקומי. \par\vskip2pt\hrule\vskip3pt\textbf{}\code{aaa new-model} · \code{tacacs-server host 192.168.1.10 key Key123} · \code{aaa authentication login default group tacacs+ local}\end{exambox}

\begin{exambox}\boxtitle{שאלה \enLR{6}}\par  ב-\enLR{802.1X} – מה תפקיד המתג? \par\vskip2pt\hrule\vskip3pt\textbf{}\enLR{Authenticator} – חוסם את הפורט עד שה-\enLR{supplicant} מאומת מול שרת ה-\enLR{RADIUS}, ומתווך בין \enLR{EAPOL} ל-\enLR{RADIUS.}\end{exambox}

\sectionheading{\enLR{3.14} מעבדה \enLR{(3} שעות) – \enLR{AAA} ב-\enLR{Packet Tracer}}

\begin{enumerate}
\item טופולוגיה: \enLR{PC-Admin} — \enLR{SW} — \enLR{R1} — \enLR{Server (AAA).} בשרת: \enLR{Services} → \enLR{AAA} → \enLR{On; Client Name: R1, Client IP}: כתובת \enLR{R1, Secret}: \code{Key123}, \enLR{ServerType: Tacacs} (ואחר כך \enLR{Radius).} הוסיפו \enLR{User: admin / cisco123.}
\item ב-\enLR{R1}: משתמש מקומי גיבוי, \code{aaa new-model}, \code{tacacs-server host … key Key123}, \code{aaa authentication login default group tacacs+ local}, \enLR{SSH} לפי פרק \enLR{2.}
\item התחברו ב-\enLR{SSH} מה-\enLR{PC} עם \enLR{admin/cisco123} – הצליח? כבו את שירות ה-\enLR{AAA} בשרת ונסו שוב – ה-\enLR{fallback} למשתמש המקומי אמור לעבוד אחרי השהיה.
\item שנו את ה-\enLR{secret} בשרת בלבד – ראו את הכישלון; הריצו \code{debug aaa authentication} וקראו את הפלט.
\item חזרו על התרגיל עם \enLR{RADIUS.} השוו ב-\enLR{Simulation mode} את החבילות: ב-\enLR{RADIUS} רואים \enLR{UDP}; ב-\enLR{TACACS}+ רואים \enLR{TCP.}
\item הוסיפו רשימה בשם לקונסול עם \code{local} בלבד, וודאו שהקונסול לא פונה לשרת.
\end{enumerate}

\sectionheading{בנק שאלות שתלמידים שואלים – ותשובות מוכנות}

שאלות נפוצות בפרק ה-\enLR{AAA}, עם תשובות מוכנות.\par

\subheading{מושגי יסוד}

\textbf{ש: "אימות והרשאה נשמעים אותו דבר – מה ההבדל?"}\newline ת: משל שדה התעופה: \textbf{אימות} = בדיקת הדרכון (מי אתה). \textbf{הרשאה} = כרטיס הטיסה (לאיזה מטוס מותר לך לעלות). קודם מוודאים מי אתה, ורק אז קובעים מה מותר לך. אי אפשר הרשאה בלי אימות.\par

\textbf{ש: "ה-\enLR{key} בין הנתב לשרת – זו הסיסמה של המשתמש?"}\newline ת: לא! ה-\enLR{key} הוא סוד משותף בין ה\underline{נתב} ל\underline{שרת} (מצפין ומאמת את הקשר ביניהם). סיסמאות המשתמשים נמצאות בשרת. שני דברים שונים לגמרי.\par

\textbf{ש: "\enLR{MFA} ו-\enLR{2FA} זה אותו דבר?"}\newline ת: \enLR{2FA} (שני גורמים) הוא מקרה פרטי של \enLR{MFA} (רב-גורמי). העיקרון: שילוב גורמים מ\underline{סוגים שונים} – משהו שאתה יודע (סיסמה) + משהו שיש לך (טלפון) + משהו שאתה (טביעת אצבע). שתי סיסמאות אינן \enLR{MFA.}\par

\subheading{מקומי מול שרת, ותקלות}

\textbf{ש: "אם שרת ה-\enLR{AAA} נופל, ננעלנו החוצה מכל הנתבים?"}\newline ת: לא, אם הגדרתם נכון: רשימת השיטות מסתיימת ב-\code{local} (למשל \code{... group tacacs+ local}). כשהשרת לא זמין, הנתב עובר אוטומטית לאימות מול משתמש מקומי. לכן \underline{תמיד} משאירים גיבוי מקומי.\par

\textbf{ש: "מה קורה אם הקלדתי \code{aaa new-model} בטעות בלי משתמש מקומי?"}\newline ת: אתם עלולים להינעל החוצה – כי מרגע הפקודה כל הקווים עוברים ל-\enLR{AAA}, ואם אין משתמש ואין רשימת ברירת מחדל, אין כניסה. כלל ברזל: קודם \code{username}, אחר כך \code{aaa new-model}, ולהשאיר חלון פתוח עד שבודקים בחלון שני.\par

\textbf{ש: "אם השרת עונה 'סיסמה שגויה', הנתב עובר לאימות מקומי?"}\newline ת: לא! מעבר לשיטה הבאה קורה רק כשהשרת \underline{לא זמין} \enLR{(timeout).} דחייה מפורשת מהשרת = המשתמש נדחה, בלי מעבר. זו נקודה שנשאלת בבגרות.\par

\subheading{\enLR{RADIUS} מול \enLR{TACACS}+}

\textbf{ש: "אז איזה מהם 'יותר טוב', \enLR{RADIUS} או \enLR{TACACS}+?"}\newline ת: תלוי במשימה. ל\underline{ניהול ציוד} (מנהלים שמתחברים לנתבים, עם אישור לכל פקודה) – \enLR{TACACS}+ עדיף. ל\underline{גישת משתמשים לרשת} \enLR{(Wi-Fi/802.1X, VPN)} – \enLR{RADIUS.} בפועל ארגונים משתמשים בשניהם, כל אחד לתפקידו.\par

\textbf{ש: "למה \enLR{RADIUS} על \enLR{UDP} אם \enLR{TCP} אמין יותר?"}\newline ת: \enLR{RADIUS} תוכנן ב-\enLR{1991} לשרתי דיאל-אפ עם אלפי חיבורים; \enLR{UDP} קל ומהיר, והלקוח פשוט שולח שוב אם אין תשובה. \enLR{TACACS}+ בחר \enLR{TCP} כי הוא מנהל שיחה רב-שלבית ורוצה לדעת מיד אם השרת נפל. שניהם עובדים היום.\par

\textbf{ש: "אומרים ש-\enLR{RADIUS '}לא עושה הרשאה' – אבל הוא כן מחזיר הרשאות, לא?"}\newline ת: נכון, זו נקודה עדינה. \enLR{RADIUS} כן נושא \enLR{attributes} של הרשאה \enLR{(VLAN}, רמה) בתוך תשובת ה-\enLR{Accept}, אבל הוא \underline{משלב} אימות והרשאה בצעד אחד ואינו תומך באישור \underline{נפרד לכל פקודה} כמו \enLR{TACACS+.} בבגרות "אינו מתייחס לניהול הרשאות" מתקבל כנכון במובן הזה.\par

\textbf{ש: "מה זה \enLR{802.1X} במשפט אחד?"}\newline ת: תקן שבו המתג/נקודת הגישה \underline{חוסמים את הפורט} עד שהמכשיר מאמת את עצמו מול שרת \enLR{RADIUS.} זה מה שמונע ממישהו לחבר מחשב לשקע ולקבל גישה, וזה הבסיס ל-\enLR{Wi-Fi} ארגוני \enLR{(WPA2-Enterprise).}\par

\textbf{ש: "אפשר לחבר \enLR{AAA} ל-\enLR{Active Directory} של בית הספר?"}\newline ת: כן – זה בדיוק מה שעושים בארגונים. שרת ה-\enLR{AAA (ISE/NPS/FreeRADIUS)} מקבל את הבקשה מהנתב ומאמת מול ה-\enLR{AD/LDAP.} כך המורים משתמשים באותם שם וסיסמה בכל מקום.\par

\sectionheading{מילון מונחים – הגדרות}

הגדרות תמציתיות של המונחים המרכזיים בפרק. שימושי גם כדף עזר לבחינה (מותר מילון מונחים).\par

\begin{xltabular}{\linewidth}{|>{\setRTL\raggedleft\arraybackslash}X|>{\setRTL\raggedleft\arraybackslash}X|}
\hline
\rowcolor{hdrbg}\textbf{הגדרה} & \textbf{מונח} \\
\hline
\endhead
מסגרת של שלושה שירותים: \enLR{Authentication, Authorization, Accounting.} & \textbf{\enLR{AAA}} \\
\hline
קביעת זהות המשתמש – מי אתה. & \textbf{\enLR{Authentication} (אימות)} \\
\hline
קביעת מה מותר למשתמש לעשות אחרי שאומת. & \textbf{\enLR{Authorization} (הרשאה)} \\
\hline
תיעוד פעולות המשתמש: מתי נכנס, אילו פקודות הקליד. & \textbf{\enLR{Accounting} (רישום)} \\
\hline
אימות מול מסד משתמשים בקונפיג של הנתב עצמו. & \textbf{\enLR{AAA} מקומי} \\
\hline
אימות מול שרת מרכזי \enLR{(ISE/ACS)} בפרוטוקול \enLR{RADIUS} או \enLR{TACACS+.} & \textbf{\enLR{AAA} מבוסס-שרת} \\
\hline
ההתקן (נתב/מתג) שפונה לשרת ה-\enLR{AAA} בשם המשתמש. & \textbf{\enLR{NAS / AAA client}} \\
\hline
פרוטוקול \enLR{AAA} פתוח, \enLR{UDP (1812/1813)}; מצפין רק את הסיסמה; לגישת משתמשים. & \textbf{\enLR{RADIUS}} \\
\hline
פרוטוקול \enLR{AAA} של סיסקו, \enLR{TCP 49}; מצפין את כל החבילה; מפריד שירותים; לניהול ציוד. & \textbf{\enLR{TACACS}+} \\
\hline
סוד משותף בין הנתב לשרת המצפין ומאמת את הקשר ביניהם (לא סיסמת המשתמש). & \textbf{\enLR{key / shared secret}} \\
\hline
שרתי ה-\enLR{AAA} של סיסקו \enLR{(ACS} ישן, \enLR{ISE} מודרני). & \textbf{\enLR{ACS / ISE}} \\
\hline
הפקודה שמפעילה את מסגרת ה-\enLR{AAA} בנתב. & \textbf{\enLR{aaa new-model}} \\
\hline
רשימת שיטות אימות מסודרת; עוברים לשיטה הבאה רק אם הקודמת לא זמינה. & \textbf{\enLR{method list}} \\
\hline
תקן שבו המתג/\enLR{AP} חוסם פורט עד לאימות מול שרת \enLR{RADIUS.} & \textbf{\enLR{802.1X}} \\
\hline
ב-\enLR{802.1X}: המכשיר המבקש גישה \enLR{(supplicant)} והמתג שחוסם \enLR{(authenticator).} & \textbf{\enLR{Supplicant / Authenticator}} \\
\hline
אימות רב-שלבי – שילוב גורמים מסוגים שונים (יודע/יש/הוא). & \textbf{\enLR{MFA / 2FA}} \\
\hline
\end{xltabular}


\clearpage
\phantomsection\addcontentsline{toc}{section}{פרק \enLR{4} – חומות אש ברשת הבינונית והרחבה}

\sloppy
\chapterheading{פרק \enLR{4} – חומות אש ברשת הבינונית והרחבה}{\small\color{subgray}\enLR{11} שעות עיוני + \enLR{2} מעשי · שבועות \enLR{10}–\enLR{12}\par}\vskip4pt

\begin{metabox}\textbf{מטרות:} \enLR{ACL} סטנדרטי/מורחב · \enLR{ACL} מבוסס זמן · עצירת התקפות ב-\enLR{ACL} · \enLR{CBAC} · \enLR{ZPF}\\ \textbf{בבגרות:} \enLR{ACL} סטנדרטי (חסימת רשת), מיקום \enLR{ACL} וכיוון \enLR{in/out, ACL} מורחב עם \enLR{host}\end{metabox}

\begin{defbox}\boxtitle{לפני שמתחילים – איך תעבורה "עוברת" בנתב? (למי שאין רקע)}\par  חומת אש מחליטה אילו חבילות לעבור. כדי להבין זאת צריך שני מושגים:  \textbf{כיוון תעבורה} – לכל ממשק (פורט) בנתב יש שני כיוונים: \textbf{\enLR{in}} = חבילות שנכנסות לנתב דרך הפורט, \textbf{\enLR{out}} = חבילות שיוצאות מהנתב דרך הפורט. תמיד חושבים מנקודת המבט של \underline{הנתב}, לא של המשתמש. \textbf{סינון לפי כותרת החבילה} – כל חבילה נושאת בכותרת שלה: כתובת מקור, כתובת יעד, סוג פרוטוקול \enLR{(TCP/UDP/ICMP)}, ופורט. חומת אש קוראת את הכותרת ומחליטה \enLR{permit} (לעבור) או \enLR{deny} (לחסום). \textbf{\enLR{Wildcard mask}} – דרך לומר "התאם רק חלק מהכתובת". נסביר לעומק בהמשך; רק דעו שזה ההפך ממסכת רשת רגילה: \enLR{0 = "}בדוק את הביט הזה", \enLR{1 = "}לא אכפת לי".  \textbf{אנלוגיה:} חומת אש היא כמו שומר בכניסה לבניין עם רשימה. כל אדם (חבילה) שמגיע – השומר בודק ברשימה \enLR{(ACL)} לפי הזהות שלו (כתובת/פורט), ומחליט להכניס או לא. הרשימה נבדקת \underline{מלמעלה למטה}, והשומר עוצר בהתאמה הראשונה.\end{defbox}

\sectionheading{\enLR{4.1} מהי חומת אש}

\textbf{חומת אש} \enLR{(Firewall)} היא מערכת – חומרה, תוכנה או שילוב – שעומדת בין רשתות ברמות אמון שונות ומחליטה לפי מדיניות אילו חבילות לעבור ואילו לחסום. המונח מגיע מקיר בטון שמונע התפשטות שריפה בין דירות. המאפיינים: \textbf{עמידה בפני התקפות} (היא עצמה מוקשחת), \textbf{נקודת מעבר יחידה} (כל התעבורה חייבת לעבור דרכה – אחרת אין טעם), ו\textbf{אכיפת מדיניות}.\par

\begin{storybox}\boxtitle{סיפור מהחיים: הקזינו שנפרץ דרך האקווריום \enLR{(2017)}}\par  קזינו בלאס וגאס התקין אקווריום חכם עם חיישן טמפרטורה המחובר לאינטרנט. התוקפים פרצו לחיישן (סיסמת ברירת מחדל), ומכיוון שהאקווריום היה מחובר לאותה רשת של שאר הקזינו – עברו ממנו למסד הנתונים של "השחקנים הגדולים" ושאבו \enLR{10GB} החוצה דרך… החיישן. אם היה \enLR{ACL} אחד שאומר "מהאקווריום מותר רק לשרת הבקרה שלו" – הסיפור היה נגמר בחיישן. חומת אש היא לא רק בין "הפנים" ל"אינטרנט"; היא גם בין חלקי הפנים.\end{storybox}

\sectionheading{\enLR{4.2} סוגי חומות אש}

\begin{xltabular}{\linewidth}{|>{\setRTL\raggedleft\arraybackslash}X|>{\setRTL\raggedleft\arraybackslash}X|>{\setRTL\raggedleft\arraybackslash}X|>{\setRTL\raggedleft\arraybackslash}X|}
\hline
\rowcolor{hdrbg}\textbf{חסרונות} & \textbf{יתרונות} & \textbf{איך זה עובד} & \textbf{סוג} \\
\hline
\endhead
לא מבינה "שיחה": כדי לאפשר תשובות חייבים לפתוח פורטים לכיוון פנימה; ניתנת לעקיפה ב-\enLR{spoofing} & מהיר, זול, בכל נתב & בודקת כל חבילה לבד לפי כותרות שכבה \enLR{3}–\enLR{4: IP} מקור/יעד, פרוטוקול, פורט & \textbf{סינון חבילות} \enLR{(Packet filtering, stateless)} – \enLR{ACL} \\
\hline
לא מבינה את תוכן האפליקציה; \enLR{UDP/ICMP "}חסרי מצב" מטופלים לפי \enLR{timeout} & מאובטחת הרבה יותר, פשוטה להגדרה; הסטנדרט היום \enLR{(ASA, ZPF, CBAC)} & שומרת טבלת חיבורים \enLR{(state table).} תעבורה שיוצאת מהפנים נרשמת; רק התשובות שתואמות לרשומה מורשות להיכנס & \textbf{\enLR{Stateful}} (בדיקת מצב) \\
\hline
איטית, צריך \enLR{proxy} לכל פרוטוקול & מבינה תוכן, יכולה לסנן \enLR{URL}, וירוסים & מסיימת את החיבור, קוראת את תוכן שכבה \enLR{7 (HTTP, FTP)} ופותחת חיבור חדש ליעד & \textbf{\enLR{Application gateway / Proxy}} \\
\hline
אינה חומת אש אמיתית – רק מסתירה & חינם & מסתירה כתובות פנימיות & \textbf{\enLR{NAT}} כ"חומת אש" \\
\hline
יקרה, דורשת כוונון & מקיפה & \enLR{Stateful} + זיהוי אפליקציה \enLR{(Facebook} מול \enLR{YouTube} על אותו פורט \enLR{443) + IPS} + זהות משתמש + בדיקת \enLR{TLS} & \textbf{\enLR{NGFW}} (דור הבא) \\
\hline
מנוהלת בנפרד בכל מחשב & מגנה גם מפני מחשבים באותה רשת & על המחשב עצמו & \textbf{\enLR{Host-based}} \enLR{(Windows Firewall)} \\
\hline
 & מוסיפים בלי לשנות כתובות & בשכבה \enLR{2} – ללא כתובת \enLR{IP, "}שקופה" & \textbf{\enLR{Transparent}} \\
\hline
\end{xltabular}

\sectionheading{\enLR{4.3 ACL} – רשימות בקרת גישה}

\textbf{\enLR{ACL}} \enLR{(Access Control List)} היא רשימה סדורה של משפטי \enLR{permit/deny} שהנתב עובר עליה \textbf{מלמעלה למטה} עבור כל חבילה; ב\textbf{התאמה הראשונה} הוא מבצע את הפעולה ומפסיק. אם אף שורה לא תואמת – יש \textbf{\enLR{deny any} סמוי \enLR{(implicit deny)}} בסוף כל \enLR{ACL.} שלושת החוקים האלה מסבירים \enLR{90\%} מהטעויות.\par

\begin{defbox}\boxtitle{\enLR{ACL} יכולה לשמש לא רק לסינון}\par  גם לזהות תעבורה עבור \enLR{NAT, QoS, route-map, VPN} (פרק \enLR{8)}, ו-\code{access-class} על קווי \enLR{VTY.} בפרק זה – סינון.\end{defbox}

\subheading{מסכת \enLR{wildcard} – הדבר שתלמידים הכי מתבלבלים בו}

ב-\enLR{ACL} לא כותבים \enLR{subnet mask} אלא \textbf{\enLR{wildcard mask}}: ביט \textbf{\enLR{0} = חייב להתאים}, ביט \textbf{\enLR{1} = לא אכפת לי}. במסכת רשת רגילה זה הפוך. הדרך המהירה: \textbf{\enLR{wildcard = 255.255.255.255} − \enLR{subnet mask}}.\par

\begin{xltabular}{\linewidth}{|>{\setRTL\raggedleft\arraybackslash}X|>{\setRTL\raggedleft\arraybackslash}X|>{\setRTL\raggedleft\arraybackslash}X|>{\setRTL\raggedleft\arraybackslash}X|}
\hline
\rowcolor{hdrbg}\textbf{משמעות} & \textbf{\enLR{Wildcard}} & \textbf{\enLR{Subnet mask}} & \textbf{רשת} \\
\hline
\endhead
\enLR{3} האוקטטים הראשונים חייבים להתאים & \textbf{\enLR{0.0.0.255}} & \enLR{255.255.255.0} & \enLR{192.168.1.0/24} \\
\hline
\enLR{172.16.16.0}–\enLR{172.16.16.15} & \textbf{\enLR{0.0.0.15}} & \enLR{255.255.255.240} & \enLR{172.16.16.0/28} \\
\hline
כל מה שמתחיל ב-\enLR{10} & \textbf{\enLR{0.255.255.255}} & \enLR{255.0.0.0} & \enLR{10.0.0.0/8} \\
\hline
הכול חייב להתאים & \textbf{\enLR{0.0.0.0}} = \code{host 10.1.1.7} & /\enLR{32} & מארח בודד \enLR{10.1.1.7} \\
\hline
\enLR{0.0.0.0 255.255.255.255} & \textbf{\enLR{255.255.255.255}} = \code{any} & – & כל כתובת \\
\hline
ביט אחרון באוקטט השלישי חייב \enLR{0} – טריק לבחינות & \enLR{192.168.0.0} \textbf{\enLR{0.0.254.255}} &  & רשתות זוגיות \enLR{192.168.0.0, .2.0, .4.0}… \\
\hline
\end{xltabular}

\begin{exambox}\boxtitle{בבגרות (חלק א, שאלה \enLR{1}ט)}\par  "נתונה \enLR{135.8.267/26"} – מהי ה-\enLR{subnet mask} ומהי ה-\enLR{wildcard}? \enLR{Subnet: 255.255.255.192} · \enLR{Wildcard: 0.0.0.63.} (שימו לב: \enLR{267} אינו אוקטט חוקי – טעות בשאלון; מתעלמים.)\end{exambox}

\subheading{\enLR{ACL} סטנדרטית – רק לפי מקור}

מספרים \textbf{\enLR{1}–\enLR{99}} ו-\enLR{1300}–\enLR{1999.} בודקת \textbf{רק את כתובת המקור}. לכן ממקמים אותה \textbf{קרוב ליעד} – אחרת היא תחסום את המקור מלהגיע גם ליעדים שלא התכוונו אליהם.\par

\begin{codebox}\begin{english}\codefont\footnotesize\color{cfg}\setlength{\parindent}{0pt}
R1(config)\# \textcolor{cbold}{\textbf{access-list 10 deny 20.20.20.0 0.0.0.255}}    \textcolor{ccom}{\texthebrew{! בבגרות שאלה \enLR{4}ד – חסימת רשת /\enLR{24}}}\\
R1(config)\# \textcolor{cbold}{\textbf{access-list 10 permit any}}                 \textcolor{ccom}{\texthebrew{! בלי זה – הכול נחסם \enLR{(implicit deny)}}}\\
R1(config)\# interface g0/1\\
R1(config-if)\# \textcolor{cbold}{\textbf{ip access-group 10 out}}
\end{english}\end{codebox}

\subheading{\enLR{ACL} מורחבת – מקור, יעד, פרוטוקול ופורט}

מספרים \textbf{\enLR{100}–\enLR{199}} ו-\enLR{2000}–\enLR{2699.} תחביר: \code{access-list N \{permit|deny\} PROTOCOL SRC SRC-WC [op PORT] DST DST-WC [op PORT] [established] [log]}. ממקמים \textbf{קרוב למקור} – כדי לא להעביר תעבורה שתיחסם ממילא דרך כל הרשת.\par

\begin{codebox}\begin{english}\codefont\footnotesize\color{cfg}\setlength{\parindent}{0pt}
\textcolor{ccom}{\texthebrew{! רשת \enLR{IT (172.18.1.0/24)} רשאית ל-\enLR{DNS} ול-\enLR{DHCP}, אך לא ל-\enLR{web 10.0.0.190} – בבגרות חלק א שאלה \enLR{3}ה}}\\
R1(config)\# \textcolor{cbold}{\textbf{access-list 101 deny ip 172.18.1.0 0.0.0.255 host 10.0.0.190}}\\
R1(config)\# \textcolor{cbold}{\textbf{access-list 101 permit ip any any}}\\
R1(config)\# interface g0/0/1\\
R1(config-if)\# \textcolor{cbold}{\textbf{ip access-group 101 in}}\\
\textcolor{ccom}{\texthebrew{! דוגמאות לפורטים}}\\
R1(config)\# access-list 110 permit tcp 192.168.1.0 0.0.0.255 any \textcolor{cbold}{\textbf{eq 80}}\\
R1(config)\# access-list 110 permit tcp 192.168.1.0 0.0.0.255 any eq 443\\
R1(config)\# access-list 110 permit udp 192.168.1.0 0.0.0.255 host 8.8.8.8 eq 53\\
R1(config)\# access-list 110 deny icmp any any \textcolor{cbold}{\textbf{echo}}            \textcolor{ccom}{\texthebrew{! חסימת \enLR{ping}}}\\
R1(config)\# access-list 110 permit tcp any any \textcolor{cbold}{\textbf{established}}      \textcolor{ccom}{\texthebrew{! רק תשובות \enLR{(ACK/RST} דלוק) – "\enLR{stateful} לעניים"}}\\
R1(config)\# access-list 110 deny ip any any \textcolor{cbold}{\textbf{log}}                 \textcolor{ccom}{\texthebrew{! רשום מה נחסם (רק בסוף!)}}
\end{english}\end{codebox}

אופרטורים: \code{eq} שווה, \code{neq}, \code{gt}, \code{lt}, \code{range 20 21}. פורטים ניתן לכתוב בשם (\code{eq www}, \code{eq ftp}, \code{eq domain}).\par

\subheading{\enLR{ACL} בשם \enLR{(Named)} – הדרך המומלצת}

\begin{codebox}\begin{english}\codefont\footnotesize\color{cfg}\setlength{\parindent}{0pt}
R1(config)\# \textcolor{cbold}{\textbf{ip access-list extended BLOCK-WEB}}\\
R1(config-ext-nacl)\# \textcolor{cbold}{\textbf{remark}} IT may not reach the web server\\
R1(config-ext-nacl)\# \textcolor{cbold}{\textbf{10}} deny tcp 172.18.1.0 0.0.0.255 host 10.0.0.190 eq 80\\
R1(config-ext-nacl)\# 20 permit ip any any\\
R1(config-ext-nacl)\# exit\\
R1(config)\# interface g0/0/1\\
R1(config-if)\# ip access-group BLOCK-WEB in\\
\textcolor{ccom}{\texthebrew{! עריכה: מוחקים/מוסיפים שורה לפי מספר רצף – בלי למחוק את כל הרשימה}}\\
R1(config)\# ip access-list extended BLOCK-WEB\\
R1(config-ext-nacl)\# no 10\\
R1(config-ext-nacl)\# 15 deny tcp 172.18.1.0 0.0.0.255 host 10.0.0.190 eq 443\\
\textcolor{ccom}{\texthebrew{! סטנדרטית בשם – זה מה שמופיע בבגרות \enLR{(config-std-nacl)}}}\\
R1(config)\# \textcolor{cbold}{\textbf{ip access-list standard BLOCK\_LAN2}}\\
R1(config-std-nacl)\# deny 192.168.2.0 0.0.0.255\\
R1(config-std-nacl)\# permit any
\end{english}\end{codebox}

\subheading{כיוון: \enLR{in} או \enLR{out}?}

הכיוון הוא \textbf{מנקודת מבטו של הנתב} על ממשק מסוים. \textbf{\enLR{in}} = חבילות שנכנסות לנתב דרך הממשק הזה (לפני ניתוב). \textbf{\enLR{out}} = חבילות שיוצאות מהנתב דרך הממשק הזה (אחרי ניתוב). \enLR{ACL} אחת לכל ממשק, לכל כיוון, לכל פרוטוקול.\par

\begin{exambox}\boxtitle{בבגרות (שאלה \enLR{7}ה) – הדוגמה המושלמת}\par  \enLR{R1} עם שלוש רשתות: \enLR{G0/0}→\enLR{192.168.1.0, G0/1}→\enLR{192.168.2.0, G0/2}→\enLR{192.168.3.0 (Restricted).} חוסמים את \enLR{192.168.2.0} ל-\enLR{192.168.3.0} בלבד. \enLR{ACL} סטנדרטית ⇒ קרוב ליעד ⇒ \textbf{\enLR{G0/2}}. כיוון: התעבורה \emph{יוצאת} מהנתב לרשת \enLR{3} ⇒ \textbf{\enLR{out}}. סדר: \code{deny 192.168.2.0} ואז \code{permit any}. תשובה \enLR{2.} אם היינו שמים \code{in} על \enLR{G0/1} – רשת \enLR{2} הייתה נחסמת גם לרשת \enLR{1} ולאינטרנט. אם \enLR{permit any} היה ראשון – ה-\enLR{deny} לעולם לא היה מגיע.\end{exambox}

\begin{mistakebox}\boxtitle{טעויות נפוצות ב-\enLR{ACL}}\par  • שוכחים \code{permit} בסוף ⇒ הכול נחסם. תלמידים מגדירים "\enLR{deny X"} ומתפלאים שאין אינטרנט.\newline  • סדר הפוך – שורה כללית לפני ספציפית.\newline  • \enLR{wildcard} כ-\enLR{subnet mask (255.255.255.0} במקום \enLR{0.0.0.255).}\newline  • \enLR{ACL} מוגדרת אך לא הוחלה על ממשק (\code{ip access-group}) – ב-\code{show ip interface} רואים "\enLR{Outgoing access list is not set".}\newline  • \enLR{ACL} סטנדרטית שממוקמת קרוב למקור – חוסמת הרבה יותר מהמתוכנן.\newline  • \enLR{ACL} על \enLR{VTY}: לא \code{ip access-group} אלא \code{access-class 10 in} תחת \enLR{line vty.}\newline  • ה-\enLR{ACL} לא בודקת תעבורה \emph{שהנתב עצמו יוצר} \enLR{(ping} מהנתב) – רק תעבורה שעוברת דרכו.\end{mistakebox}

\subheading{וריפיקציה}

\begin{codebox}\begin{english}\codefont\footnotesize\color{cfg}\setlength{\parindent}{0pt}
R1\# \textcolor{cbold}{\textbf{show access-lists}}                \textcolor{ccom}{\texthebrew{! כולל מונה התאמות \enLR{(matches)} לכל שורה – הכלי הכי טוב לדיבוג}}\\
R1\# \textcolor{cbold}{\textbf{show ip access-lists 101}}\\
R1\# \textcolor{cbold}{\textbf{show ip interface g0/1}}           \textcolor{ccom}{\texthebrew{! איזו \enLR{ACL} מוחלת ובאיזה כיוון}}\\
R1\# show running-config | section access-list\\
R1\# \textcolor{cbold}{\textbf{clear access-list counters}}
\end{english}\end{codebox}

\sectionheading{\enLR{4.4 ACL} מורכבות: מבוססות זמן, דינמיות ורפלקסיביות}

\subheading{\enLR{Time-based ACL}}

מאפשרת כלל שתקף רק בזמנים מסוימים. דורש שעון נכון \enLR{(NTP} – פרק \enLR{2).}\par

\begin{codebox}\begin{english}\codefont\footnotesize\color{cfg}\setlength{\parindent}{0pt}
R1(config)\# \textcolor{cbold}{\textbf{time-range WORK-HOURS}}\\
R1(config-time-range)\# \textcolor{cbold}{\textbf{periodic weekdays 8:00 to 17:00}}\\
R1(config-time-range)\# exit\\
R1(config)\# access-list 120 permit tcp 192.168.1.0 0.0.0.255 any eq 80 \textcolor{cbold}{\textbf{time-range WORK-HOURS}}\\
R1(config)\# access-list 120 deny ip any any\\
\textcolor{ccom}{\texthebrew{! אפשרויות: \enLR{periodic Monday Wednesday 9:00 to 12:00 / absolute start 08:00 1 Jan 2026 end 17:00 31 Jan 2026}}}
\end{english}\end{codebox}

\subheading{\enLR{Dynamic ACL (Lock-and-Key)}}

המשתמש מבצע \enLR{Telnet/SSH} לנתב ומזדהה; רק אז נפתחת לו דלת זמנית דרך הנתב. "מנעול ומפתח".\par

\begin{codebox}\begin{english}\codefont\footnotesize\color{cfg}\setlength{\parindent}{0pt}
R1(config)\# username student password 0 pass\\
R1(config)\# access-list 101 permit tcp any host 10.0.0.1 eq telnet\\
R1(config)\# \textcolor{cbold}{\textbf{access-list 101 dynamic OPEN-DOOR timeout 15 permit ip any any}}\\
R1(config)\# interface g0/0\\
R1(config-if)\# ip access-group 101 in\\
R1(config)\# line vty 0 4\\
R1(config-line)\# login local\\
R1(config-line)\# \textcolor{cbold}{\textbf{autocommand access-enable host timeout 5}}
\end{english}\end{codebox}

\subheading{\enLR{Reflexive ACL}}

"מראה": כשתעבורה יוצאת החוצה, הנתב יוצר אוטומטית שורה זמנית שמאפשרת \emph{רק את התשובה} \enLR{(IP}/פורט הפוכים) להיכנס. זה מנגנון \enLR{stateful} פשוט ב-\enLR{ACL}, ועדיף על \code{established} כי עובד גם ל-\enLR{UDP} ו-\enLR{ICMP.}\par

\begin{codebox}\begin{english}\codefont\footnotesize\color{cfg}\setlength{\parindent}{0pt}
R1(config)\# ip access-list extended OUTBOUND\\
R1(config-ext-nacl)\# permit tcp any any \textcolor{cbold}{\textbf{reflect TCP-TRAFFIC}}\\
R1(config-ext-nacl)\# permit udp any any reflect UDP-TRAFFIC\\
R1(config)\# ip access-list extended INBOUND\\
R1(config-ext-nacl)\# \textcolor{cbold}{\textbf{evaluate TCP-TRAFFIC}}\\
R1(config-ext-nacl)\# evaluate UDP-TRAFFIC\\
R1(config-ext-nacl)\# deny ip any any\\
R1(config)\# interface s0/0/0\\
R1(config-if)\# ip access-group OUTBOUND out\\
R1(config-if)\# ip access-group INBOUND in
\end{english}\end{codebox}

\sectionheading{\enLR{4.5} עצירת התקפות בעזרת \enLR{ACL}}

\subheading{\enLR{Anti-spoofing} – מה לחסום בכניסה מהאינטרנט \enLR{(ingress)}}

\begin{codebox}\begin{english}\codefont\footnotesize\color{cfg}\setlength{\parindent}{0pt}
R1(config)\# ip access-list extended ANTI-SPOOF\\
\textcolor{ccom}{\texthebrew{! כתובות פרטיות \enLR{(RFC 1918)} לא אמורות להגיע מהאינטרנט – מישהו מזייף}}\\
R1(config-ext-nacl)\# deny ip 10.0.0.0 0.255.255.255 any\\
R1(config-ext-nacl)\# deny ip 172.16.0.0 0.15.255.255 any\\
R1(config-ext-nacl)\# deny ip 192.168.0.0 0.0.255.255 any\\
\textcolor{ccom}{\texthebrew{! הרשת שלנו כמקור מבחוץ – זיוף ודאי}}\\
R1(config-ext-nacl)\# deny ip 203.0.113.0 0.0.0.255 any\\
R1(config-ext-nacl)\# deny ip 127.0.0.0 0.255.255.255 any     \textcolor{ccom}{\texthebrew{! \enLR{loopback}}}\\
R1(config-ext-nacl)\# deny ip 0.0.0.0 0.255.255.255 any\\
R1(config-ext-nacl)\# deny ip 224.0.0.0 15.255.255.255 any    \textcolor{ccom}{\texthebrew{! \enLR{multicast}}}\\
\textcolor{ccom}{\texthebrew{! \enLR{ICMP} – לאפשר רק מה שצריך, לחסום \enLR{echo} מבחוץ (נגד \enLR{ping sweep} ו-\enLR{Smurf)}}}\\
R1(config-ext-nacl)\# permit icmp any any echo-reply\\
R1(config-ext-nacl)\# permit icmp any any unreachable\\
R1(config-ext-nacl)\# permit icmp any any time-exceeded      \textcolor{ccom}{\texthebrew{! \enLR{traceroute}}}\\
R1(config-ext-nacl)\# deny icmp any any\\
\textcolor{ccom}{\texthebrew{! שירותי ניהול לא נגישים מבחוץ}}\\
R1(config-ext-nacl)\# deny tcp any any eq 23\\
R1(config-ext-nacl)\# deny udp any any eq 161\\
\textcolor{ccom}{\texthebrew{! רק תשובות ל-\enLR{TCP} שהפנים פתח + שירותי \enLR{DMZ}}}\\
R1(config-ext-nacl)\# permit tcp any any established\\
R1(config-ext-nacl)\# permit tcp any host 203.0.113.10 eq 80\\
R1(config-ext-nacl)\# permit tcp any host 203.0.113.10 eq 443\\
R1(config-ext-nacl)\# deny ip any any log\\
R1(config)\# interface g0/0\\
R1(config-if)\# ip access-group ANTI-SPOOF in
\end{english}\end{codebox}

\textbf{\enLR{Egress filtering}} – גם ביציאה: לאפשר רק את כתובות המקור שלנו; כך הרשת שלנו לא תשמש ל-\enLR{DDoS} עם כתובות מזויפות.\par

\sectionheading{\enLR{4.6 CBAC} – \enLR{Context-Based Access Control} (חומת אש "קלאסית" ב-\enLR{IOS)}}

\enLR{ACL} היא \enLR{stateless.} \textbf{\enLR{CBAC}} (נקראת גם \enLR{Classic Firewall)} מוסיפה מצב: היא \emph{בוחנת} \enLR{(inspect)} תעבורה שיוצאת מהפנים, זוכרת את השיחה בטבלת \enLR{state}, ופותחת \textbf{אוטומטית וזמנית} חור ב-\enLR{ACL} הנכנסת לתשובות בלבד. היא גם מבינה פרוטוקולים עם פורטים דינמיים \enLR{(FTP active)} ומגנה מ-\enLR{SYN flood (TCP intercept-like).}\par

\begin{codebox}\begin{english}\codefont\footnotesize\color{cfg}\setlength{\parindent}{0pt}
\textcolor{ccom}{\texthebrew{! \enLR{1. ACL} שחוסמת הכול מבחוץ (התשובות ייפתחו על ידי \enLR{CBAC)}}}\\
R1(config)\# ip access-list extended OUTSIDE-IN\\
R1(config-ext-nacl)\# permit icmp any any echo-reply\\
R1(config-ext-nacl)\# deny ip any any\\
R1(config)\# interface s0/0/0\\
R1(config-if)\# ip access-group OUTSIDE-IN in\\
\textcolor{ccom}{\texthebrew{! \enLR{2.} כלל בדיקה}}\\
R1(config)\# \textcolor{cbold}{\textbf{ip inspect name FW tcp}}\\
R1(config)\# \textcolor{cbold}{\textbf{ip inspect name FW udp}}\\
R1(config)\# ip inspect name FW http\\
R1(config)\# ip inspect name FW ftp\\
\textcolor{ccom}{\texthebrew{! \enLR{3.} החלה בכיוון שבו התעבורה "מתחילה" – יוצאת החוצה}}\\
R1(config)\# interface s0/0/0\\
R1(config-if)\# \textcolor{cbold}{\textbf{ip inspect FW out}}\\
R1\# \textcolor{cbold}{\textbf{show ip inspect sessions}}\\
R1\# show ip inspect config
\end{english}\end{codebox}

\enLR{CBAC} היא "מבוססת ממשק": ככל שיש יותר ממשקים, ההגדרה מסתבכת. לכן סיסקו החליפה אותה ב-\enLR{ZPF.}\par

\sectionheading{\enLR{4.7 ZPF} – \enLR{Zone-Based Policy Firewall}}

הגישה המודרנית ב-\enLR{IOS}: במקום לחשוב על ממשקים, מגדירים \textbf{אזורים} \enLR{(zones)} ומדיניות \textbf{בין זוגות אזורים} \enLR{(zone-pair)}, בכיוון אחד. חוקים:\par

\begin{itemize}
\item ממשק שייך לאזור אחד לכל היותר.
\item תעבורה בין ממשקים \textbf{באותו אזור} – מותרת תמיד.
\item תעבורה בין ממשק באזור לממשק שלא באזור – \textbf{נחסמת}.
\item תעבורה בין שני אזורים – \textbf{נחסמת כברירת מחדל}, אלא אם יש \enLR{zone-pair} עם מדיניות.
\item אזור מיוחד \textbf{\enLR{self}} = הנתב עצמו (תעבורה אל הנתב וממנו – \enLR{SSH, OSPF).} ברירת המחדל ל-\enLR{self}: הכול מותר.
\end{itemize}

\subheading{שלוש פעולות}

\begin{xltabular}{\linewidth}{|>{\setRTL\raggedleft\arraybackslash}X|>{\setRTL\raggedleft\arraybackslash}X|}
\hline
\rowcolor{hdrbg}\textbf{משמעות} & \textbf{פעולה} \\
\hline
\endhead
בדיקה \enLR{stateful} – מאפשר את התעבורה \emph{ואת התשובות} חזרה אוטומטית. זה מה שרוצים כמעט תמיד. & \textbf{\enLR{inspect}} \\
\hline
מאפשר את התעבורה בכיוון הזה בלבד, בלי מעקב (כמו \enLR{permit} ב-\enLR{ACL).} לתשובה צריך \enLR{pass} בזוג ההפוך. & \textbf{\enLR{pass}} \\
\hline
חוסם (ברירת מחדל, אפשר להוסיף \enLR{log).} & \textbf{\enLR{drop}} \\
\hline
\end{xltabular}

\subheading{ההגדרה – \enLR{C3PL (Cisco Common Classification Policy Language)} בחמישה שלבים}

\begin{codebox}\begin{english}\codefont\footnotesize\color{cfg}\setlength{\parindent}{0pt}
\textcolor{ccom}{\texthebrew{! \enLR{1.} אזורים}}\\
R1(config)\# \textcolor{cbold}{\textbf{zone security INSIDE}}\\
R1(config)\# \textcolor{cbold}{\textbf{zone security OUTSIDE}}\\
\textcolor{ccom}{\texthebrew{! \enLR{2. class-map} – איזו תעבורה \enLR{(match-any} = מספיק תנאי אחד; \enLR{match-all} = כולם)}}\\
R1(config)\# \textcolor{cbold}{\textbf{class-map type inspect match-any WEB-DNS}}\\
R1(config-cmap)\# \textcolor{cbold}{\textbf{match protocol http}}\\
R1(config-cmap)\# match protocol https\\
R1(config-cmap)\# match protocol dns\\
R1(config-cmap)\# match protocol icmp\\
R1(config-cmap)\# exit\\
\textcolor{ccom}{\texthebrew{! (אפשר גם \enLR{match access-group 101} – לשלב \enLR{ACL)}}}\\
\textcolor{ccom}{\texthebrew{! \enLR{3. policy-map} – מה עושים עם התעבורה}}\\
R1(config)\# \textcolor{cbold}{\textbf{policy-map type inspect IN-TO-OUT}}\\
R1(config-pmap)\# \textcolor{cbold}{\textbf{class type inspect WEB-DNS}}\\
R1(config-pmap-c)\# \textcolor{cbold}{\textbf{inspect}}\\
R1(config-pmap-c)\# exit\\
R1(config-pmap)\# class class-default          \textcolor{ccom}{\texthebrew{! כל השאר}}\\
R1(config-pmap-c)\# drop log\\
R1(config-pmap-c)\# exit\\
\textcolor{ccom}{\texthebrew{! \enLR{4. zone-pair} – מקור, יעד, ומדיניות}}\\
R1(config)\# \textcolor{cbold}{\textbf{zone-pair security IN-OUT source INSIDE destination OUTSIDE}}\\
R1(config-sec-zone-pair)\# \textcolor{cbold}{\textbf{service-policy type inspect IN-TO-OUT}}\\
R1(config-sec-zone-pair)\# exit\\
\textcolor{ccom}{\texthebrew{! \enLR{5.} שיוך ממשקים לאזורים – ברגע זה החסימה מתחילה!}}\\
R1(config)\# interface g0/1\\
R1(config-if)\# \textcolor{cbold}{\textbf{zone-member security INSIDE}}\\
R1(config)\# interface s0/0/0\\
R1(config-if)\# \textcolor{cbold}{\textbf{zone-member security OUTSIDE}}\\
\textcolor{ccom}{\texthebrew{! וריפיקציה}}\\
R1\# \textcolor{cbold}{\textbf{show zone security}}\\
R1\# \textcolor{cbold}{\textbf{show zone-pair security}}\\
R1\# \textcolor{cbold}{\textbf{show policy-map type inspect zone-pair IN-OUT sessions}}\\
R1\# show class-map type inspect
\end{english}\end{codebox}

\begin{faqbox}\boxtitle{שאלת תלמיד: "הגדרתי \enLR{INSIDE}→\enLR{OUTSIDE} עם \enLR{inspect.} איך התשובות חוזרות, הרי אין \enLR{zone-pair OUTSIDE}→\enLR{INSIDE}?"}\par  זה בדיוק מה ש-\enLR{inspect} עושה: הנתב רושם את השיחה בטבלת \enLR{state}, וכשמגיעה תשובה מ-\enLR{OUTSIDE} שתואמת לשיחה קיימת – היא עוברת, בלי צורך ב-\enLR{zone-pair} הפוך. תעבורה \emph{חדשה} מ-\enLR{OUTSIDE} (למשל מישהו מבחוץ מנסה \enLR{SSH} למחשב פנימי) – נחסמת. אם היינו משתמשים ב-\code{pass} במקום \enLR{inspect}, התשובות היו נחסמות.\end{faqbox}

\begin{faqbox}\boxtitle{שאלת תלמיד: "אחרי שהגדרתי \enLR{ZPF}, אני לא מצליח לעשות \enLR{SSH} לנתב מבחוץ / \enLR{OSPF} נפל"}\par  תעבורה אל הנתב עצמו שייכת לאזור \textbf{\enLR{self}}. ברירת המחדל ל-\enLR{self} היא \enLR{permit}, אז \enLR{SSH} אמור לעבוד. אם הגדרתם \enLR{zone-pair} שמערב \enLR{self} עם מדיניות – עכשיו רק מה שהוגדר עובר. יש להוסיף \enLR{class} ל-\enLR{SSH/OSPF} עם \enLR{pass/inspect.} גם בדקו ש-\enLR{CBAC} (\code{ip inspect}) לא מוגדרת על אותו ממשק – אסור לשלב.\end{faqbox}

\sectionheading{\enLR{4.8} תכנון חומת אש ברשת – שיטות עבודה}

\begin{itemize}
\item \textbf{\enLR{DMZ}}: שלושה אזורים – \enLR{INSIDE, DMZ, OUTSIDE.} מדיניות: \enLR{INSIDE}→\enLR{OUTSIDE inspect; INSIDE}→\enLR{DMZ inspect; OUTSIDE}→\enLR{DMZ} – רק \enLR{http/https/smtp} אל השרתים; \enLR{DMZ}→\enLR{INSIDE} – \textbf{כלום} (אם שרת ה-\enLR{web} נפרץ – התוקף לא ממשיך פנימה). \enLR{OUTSIDE}→\enLR{INSIDE} – כלום.
\item \textbf{מינימום הרשאות} \enLR{(least privilege)}: מתחילים מ-\enLR{deny all} ופותחים רק מה שיש לו הצדקה עסקית.
\item \textbf{תיעוד}: כל שורת \enLR{ACL} עם \code{remark}; מי ביקש, מתי, למה.
\item \textbf{ביקורת}: פעם ברבעון – \code{show access-lists}; שורות עם \enLR{0 matches} – כנראה מיותרות.
\item \textbf{\enLR{Logging}} – \code{log} על \enLR{deny} בסוף, אבל לא על \enLR{permit} (עומס).
\end{itemize}

\begin{storybox}\boxtitle{סיפור מהחיים: \enLR{32} מיליון שורות \enLR{ACL}}\par  בחברת ביטוח גדולה, בדיקת ביקורת מצאה חומת אש עם \enLR{8,000} כללים, מתוכם \enLR{3,200} שלא התאימו לאף חבילה במשך שנה – "כללים זומבי" שאף אחד לא העז למחוק כי "אולי משהו ישבר". בין הזומבים: כלל מ-\enLR{2009} שפתח \enLR{RDP} לכל האינטרנט לשרת שכבר לא קיים – אך הכתובת שלו הוקצתה מחדש לשרת \enLR{HR.} הלקח: חומת אש שלא מתחזקים היא לא חומת אש; והשוואת מוני \enLR{matches} היא לא רק לדיבוג.\end{storybox}

\sectionheading{\enLR{4.9} מדריך פקודות מרוכז – פרק \enLR{4}}

\begin{xltabular}{\linewidth}{|>{\setRTL\raggedleft\arraybackslash}X|>{\setRTL\raggedleft\arraybackslash}X|}
\hline
\rowcolor{hdrbg}\textbf{תפקיד} & \textbf{פקודה} \\
\hline
\endhead
\enLR{ACL} סטנדרטית ממוספרת & \code{access-list 1-99 \{permit|deny\} SRC WC} \\
\hline
\enLR{ACL} מורחבת ממוספרת & \code{access-list 100-199 \{permit|deny\} PROTO SRC WC [eq P] DST WC [eq P]} \\
\hline
\enLR{ACL} בשם, עם מספרי רצף & \code{ip access-list \{standard|extended\} NAME} \\
\hline
החלת \enLR{ACL} & \code{ip access-group N \{in|out\}} (ממשק) \\
\hline
\enLR{ACL} על גישת ניהול & \code{access-class N in} \enLR{(line vty)} \\
\hline
קיצורים ל-\enLR{wildcard 0.0.0.0 / 255.255.255.255} & \code{host A} / \code{any} \\
\hline
מילות מפתח מתקדמות & \code{established}, \code{log}, \code{time-range}, \code{reflect}/\code{evaluate}, \code{dynamic} \\
\hline
וריפיקציה & \code{show access-lists}, \code{show ip interface} \\
\hline
\enLR{CBAC} & \code{ip inspect name N PROTO} + \code{ip inspect N out} \\
\hline
\enLR{ZPF}, בסדר הזה & \code{zone security} → \code{class-map type inspect} → \code{policy-map type inspect} → \code{zone-pair security} → \code{zone-member security} \\
\hline
וריפיקציה \enLR{ZPF} & \code{show zone-pair security}, \code{show policy-map type inspect zone-pair sessions} \\
\hline
\end{xltabular}

\sectionheading{\enLR{4.10} תרגילים לתלמידים}

\begin{exbox}\boxtitle{תרגיל \enLR{1} – \enLR{wildcard}}\par  כתבו את ה-\enLR{wildcard}: א. \enLR{192.168.10.0/25} ב. \enLR{10.10.0.0/16} ג. \enLR{172.16.4.0/22} ד. המארח \enLR{8.8.8.8} ה. כל הכתובות \enLR{192.168.0.0} עד \enLR{192.168.7.255} \par\vskip2pt\hrule\vskip3pt\textbf{}א. \enLR{0.0.0.127} ב. \enLR{0.0.255.255} ג. \enLR{0.0.3.255} ד. \enLR{0.0.0.0 (host 8.8.8.8)} ה. \enLR{192.168.0.0 0.0.7.255}\end{exbox}

\begin{exbox}\boxtitle{תרגיל \enLR{2} – מה עובר?}\par  \enLR{access-list 105 permit tcp 192.168.1.0 0.0.0.255 any eq 443 access-list 105 deny tcp 192.168.1.0 0.0.0.255 any eq 80 access-list 105 permit udp any host 8.8.8.8 eq 53 interface g0/0 (LAN 192.168.1.0/24) ip access-group 105 in} עבור כל חבילה שנכנסת ב-\enLR{g0/0} מ-\enLR{192.168.1.20}: א. \enLR{HTTPS} ל-\enLR{1.1.1.1} ב. \enLR{HTTP} ל-\enLR{1.1.1.1} ג. \enLR{DNS} ל-\enLR{8.8.8.8} ד. \enLR{DNS} ל-\enLR{1.1.1.1} ה. \enLR{ping} ל-\enLR{1.1.1.1} ו. \enLR{SSH} לנתב עצמו \par\vskip2pt\hrule\vskip3pt\textbf{}א. עובר (שורה \enLR{1).} ב. נחסם (שורה \enLR{2).} ג. עובר \enLR{(3).} ד. נחסם \enLR{(implicit deny).} ה. נחסם \enLR{(implicit deny} – אין שורה ל-\enLR{icmp).} ו. נחסם – גם תעבורה אל הנתב עוברת דרך \enLR{ACL} נכנסת. תלמידים רבים חושבים שרק "מה שכתוב \enLR{deny"} נחסם – זו הנקודה.\end{exbox}

\begin{exbox}\boxtitle{תרגיל \enLR{3} – תכנון \enLR{ACL}}\par  נתב עם \enLR{G0/0}→רשת תלמידים \enLR{10.1.0.0/16, G0/1}→רשת מורים \enLR{10.2.0.0/16, G0/2}→שרתים \enLR{10.3.0.0/24} (שרת ציונים \enLR{10.3.0.10}, שרת קבצים \enLR{10.3.0.20).} דרישות: תלמידים לא ניגשים לשרת הציונים אך כן לשרת הקבצים; מורים ניגשים להכול; תלמידים לא רשאים ל-\enLR{Telnet/SSH} לשום מקום. כתבו \enLR{ACL} מורחבת בשם, ובחרו ממשק וכיוון. \par\vskip2pt\hrule\vskip3pt\textbf{}\enLR{ip access-list extended STUDENTS deny ip 10.1.0.0 0.0.255.255 host 10.3.0.10 deny tcp 10.1.0.0 0.0.255.255 any eq 22 deny tcp 10.1.0.0 0.0.255.255 any eq 23 permit ip any any interface g0/0 ip access-group STUDENTS in}מורחבת ⇒ קרוב למקור ⇒ \enLR{G0/0 in.} מורים לא מופיעים – \enLR{permit any} מכסה אותם (הם ממילא לא נכנסים ב-\enLR{G0/0).}\end{exbox}

\begin{exbox}\boxtitle{תרגיל \enLR{4} – \enLR{ZPF}}\par  ציירו טבלה של אזורים \enLR{INSIDE / DMZ / OUTSIDE} וכתבו לכל זוג \enLR{(6} כיוונים) איזו מדיניות תגדירו ומדוע. \par\vskip2pt\hrule\vskip3pt\textbf{}\enLR{IN}→\enLR{OUT: inspect} (הכול/רשימת פרוטוקולים). \enLR{IN}→\enLR{DMZ: inspect. OUT}→\enLR{DMZ: inspect http, https, smtp} בלבד (ל-\enLR{match access-group} עם ה-\enLR{IP} של השרתים). \enLR{OUT}→\enLR{IN}: אין \enLR{zone-pair (drop). DMZ}→\enLR{IN}: אין \enLR{(drop)} – עיקרון "שרת נפרץ לא ממשיך פנימה". \enLR{DMZ}→\enLR{OUT: inspect dns, ntp, http} (עדכונים) בלבד – ולא הכול, כדי ששרת נפרץ לא ישמש ל-\enLR{DDoS.}\end{exbox}

\sectionheading{\enLR{4.11} שאלות בסגנון בגרות}

\begin{exambox}\boxtitle{שאלה \enLR{1}}\par  איזו פקודה חוסמת את המארח \enLR{10.5.5.5} בלבד ב-\enLR{ACL} סטנדרטית?\newline \enLR{1.} \code{access-list 10 deny 10.5.5.5 255.255.255.255} \enLR{2.} \code{access-list 10 deny host 10.5.5.5} \enLR{3.} \code{access-list 100 deny 10.5.5.5 0.0.0.0} \enLR{4.} \code{access-list 10 deny 10.5.5.0 0.0.0.255} \par\vskip2pt\hrule\vskip3pt\textbf{}\textbf{\enLR{2.}} \enLR{(1} – \enLR{wildcard} של \enLR{any; 3} – מספר של מורחבת ותחביר חסר; \enLR{4} – חוסם רשת שלמה.)\end{exambox}

\begin{exambox}\boxtitle{שאלה \enLR{2}}\par  מנהל הגדיר \enLR{ACL} עם שורה אחת: \code{access-list 20 deny 192.168.5.0 0.0.0.255} והחיל אותה. אף אחד לא מצליח לצאת. מדוע? \par\vskip2pt\hrule\vskip3pt\textbf{}\enLR{implicit deny any} בסוף. חסר \code{access-list 20 permit any}.\end{exambox}

\begin{exambox}\boxtitle{שאלה \enLR{3}}\par  היכן נכון למקם \enLR{ACL} מורחבת – קרוב למקור או ליעד? ומדוע? \par\vskip2pt\hrule\vskip3pt\textbf{}קרוב למקור – היא מזהה בדיוק את התעבורה (מקור+יעד+פורט), ולכן אפשר לחסום אותה מוקדם ולא לבזבז רוחב פס. סטנדרטית – קרוב ליעד, כי היא בודקת רק מקור ותחסום יותר מדי אם תמוקם קרוב אליו.\end{exambox}

\begin{exambox}\boxtitle{שאלה \enLR{4}}\par  ב-\enLR{ZPF}, ממשק \enLR{G0/0} באזור \enLR{INSIDE} וממשק \enLR{G0/1} לא שויך לאף אזור. מה יקרה לתעבורה ביניהם? \par\vskip2pt\hrule\vskip3pt\textbf{}תיחסם. ממשק באזור לא מתקשר עם ממשק שאינו באזור.\end{exambox}

\begin{exambox}\boxtitle{שאלה \enLR{5}}\par  מה ההבדל בין הפעולות \enLR{inspect} ו-\enLR{pass} ב-\enLR{ZPF}? \par\vskip2pt\hrule\vskip3pt\textbf{}\enLR{inspect} – מעקב \enLR{stateful}, התשובות מורשות אוטומטית. \enLR{pass} – מעבר בכיוון אחד בלבד ללא מעקב; לתשובות נדרש \enLR{pass} בזוג ההפוך.\end{exambox}

\begin{exambox}\boxtitle{שאלה \enLR{6}}\par  השלימו: \code{R1(config)\# access-list 101 deny ip \_\_\_\_\_\_ host \_\_\_\_\_\_} · \code{R1(config)\# access-list 101 permit \_\_\_\_\_\_} · \code{R1(config-if)\# \_\_\_\_\_\_} – כך שמשתמשי \enLR{IT (172.18.1.0/16} לפי הטבלה) יקבלו \enLR{DNS} ו-\enLR{DHCP} אך לא \enLR{WEB (10.0.0.190).} \par\vskip2pt\hrule\vskip3pt\textbf{}\code{172.18.1.0 0.0.255.255} (או לפי /\enLR{24}: \code{0.0.0.255}) · \code{10.0.0.190} · \code{ip any any} · \code{ip access-group 101 in}\end{exambox}

\sectionheading{\enLR{4.12} מעבדה \enLR{(2} שעות) – \enLR{ACL} ו-\enLR{ZPF} ב-\enLR{Packet Tracer}}

\begin{enumerate}
\item \textbf{\enLR{ACL}:} בנו את הטופולוגיה משאלה \enLR{7}ה בבגרות \enLR{(R1} עם \enLR{3} רשתות). הגדירו את \enLR{BLOCK\_LAN2} והחילו \enLR{out} על \enLR{G0/2.} בדקו \enLR{ping} מרשת \enLR{2} לרשת \enLR{3} (נכשל) ולרשת \enLR{1} (מצליח). הריצו \code{show access-lists} וראו את המונים עולים.
\item הזיזו את אותה \enLR{ACL} ל-\enLR{G0/1 in.} מה נשבר? (רשת \enLR{2} מאבדת הכול.) הסבירו.
\item \textbf{\enLR{ACL} מורחבת:} אפשרו לרשת \enLR{2} רק \enLR{HTTP} לשרת ברשת \enLR{3} וחסמו \enLR{ping.} בדקו בדפדפן של ה-\enLR{PC.}
\item \textbf{\enLR{ZPF}:} טופולוגיה \enLR{INSIDE(PC)}–\enLR{R1}–\enLR{OUTSIDE(Server).} הגדירו \enLR{zones, class-map} ל-\enLR{http+icmp, policy inspect, zone-pair IN-OUT.} מה-\enLR{PC}: דפדפן לשרת עובד, \enLR{ping} עובד. מהשרת: \enLR{ping} ל-\enLR{PC} נכשל. הריצו \code{show policy-map type inspect zone-pair sessions} תוך כדי גלישה.
\item החליפו \enLR{inspect} ב-\enLR{pass.} מה קרה לגלישה? הסבירו בעזרת מושג ה-\enLR{state.}
\end{enumerate}

\sectionheading{בנק שאלות שתלמידים שואלים – ותשובות מוכנות}

שאלות נפוצות בפרק חומות האש וה-\enLR{ACL}, עם תשובות מוכנות.\par

\subheading{הבנת \enLR{ACL}}

\textbf{ש: "הגדרתי רק שורת \code{deny} אחת ופתאום שום דבר לא עובד. למה?"}\newline ת: בגלל ה-\textbf{\enLR{deny any} הסמוי} בסוף כל \enLR{ACL.} ברגע שיש \enLR{ACL}, כל מה שלא הותר במפורש – נחסם. חייבים להוסיף \code{permit} מפורש למה שרוצים לאפשר, אחרת הכול נחסם.\par

\textbf{ש: "למה \enLR{wildcard} הפוך ממסכת רשת רגילה? זה מבלבל."}\newline ת: כי הוא עונה על שאלה אחרת. מסכת רשת אומרת "איזה חלק הוא הרשת". \enLR{Wildcard} אומרת ל-\enLR{ACL "}אילו ביטים לבדוק": \enLR{0} = חייב להתאים, \enLR{1} = לא אכפת לי. טריק: \enLR{wildcard = 255.255.255.255} פחות מסכת הרשת.\par

\textbf{ש: "איך אני יודע אם ה-\enLR{ACL} בכלל עובד?"}\newline ת: \code{show access-lists} – ליד כל שורה יש \underline{מונה התאמות} \enLR{(matches).} אם שולחים תעבורה והמונה לא עולה, ה-\enLR{ACL} לא נתפס (כנראה לא הוחל על ממשק, או בכיוון הלא נכון). זה כלי הדיבוג מספר \enLR{1.}\par

\textbf{ש: "למה סטנדרטי קרוב ליעד, ומורחב קרוב למקור?"}\newline ת: סטנדרטי בודק \underline{רק מקור}, אז אם נמקם אותו קרוב למקור הוא יחסום את המקור מלהגיע לכל יעד – יותר מדי. לכן קרוב ליעד. מורחב מזהה בדיוק מקור+יעד+פורט, אז אפשר לחסום מוקדם (קרוב למקור) בלי לפגוע בתעבורה אחרת.\par

\textbf{ש: "מה ההבדל בין \enLR{in} ל-\enLR{out}? אני מתבלבל."}\newline ת: תמיד מנקודת המבט של \underline{הנתב}: \textbf{\enLR{in}} = נכנס לנתב דרך הממשק (לפני שהנתב מנתב). \textbf{\enLR{out}} = יוצא מהנתב דרך הממשק (אחרי הניתוב). ציירו נתב עם חצים – זה מתבהר מיד.\par

\textbf{ש: "האם \enLR{ACL} חוסם גם תעבורה שהנתב עצמו שולח (למשל \enLR{ping} מהנתב)?"}\newline ת: לא. \enLR{ACL} בודק תעבורה ש\underline{עוברת דרך} הנתב, לא תעבורה שהנתב עצמו יוצר. לכן \enLR{ping} מהנתב עצמו יעבור גם אם יש \enLR{ACL} חוסם.\par

\subheading{חומות אש מתקדמות}

\textbf{ש: "מה ההבדל בין חומת אש \enLR{stateless} ל-\enLR{stateful}?"}\newline ת: \textbf{\enLR{Stateless}} \enLR{(ACL} רגיל) בודקת כל חבילה בנפרד – לא "זוכרת" שיחות, אז צריך לפתוח ידנית פורטים לתשובות. \textbf{\enLR{Stateful}} \enLR{(ZPF/CBAC)} זוכרת שיחות שיצאו, ומאפשרת רק את התשובות התואמות להיכנס – בטוח וקל בהרבה.\par

\textbf{ש: "ב-\enLR{ZPF} הגדרתי \enLR{INSIDE}→\enLR{OUTSIDE} עם \enLR{inspect.} איך התשובות חוזרות אם אין \enLR{zone-pair} בכיוון ההפוך?"}\newline ת: זה בדיוק מה ש-\code{inspect} עושה: הנתב זוכר את השיחה בטבלת \enLR{state}, וכשמגיעה תשובה תואמת – היא עוברת אוטומטית, בלי צורך ב-\enLR{zone-pair} הפוך. תעבורה \underline{חדשה} מבחוץ עדיין נחסמת.\par

\textbf{ש: "מה ההבדל בין \enLR{inspect} ל-\enLR{pass} ב-\enLR{ZPF}?"}\newline ת: \code{inspect} עוקב אחרי השיחה ומאפשר תשובות אוטומטית \enLR{(stateful).} \code{pass} רק מעביר בכיוון אחד בלי מעקב – לתשובה צריך \enLR{pass} נפרד בכיוון ההפוך. כמעט תמיד רוצים \enLR{inspect.}\par

\textbf{ש: "מה זה \enLR{DMZ} ולמה צריך אותו?"}\newline ת: אזור "מפורז" בין הרשת הפנימית לאינטרנט, שבו יושבים השרתים הציבוריים (אתר, מייל). הרעיון: אם שרת ב-\enLR{DMZ} נפרץ, התוקף עדיין לא נמצא ברשת הפנימית – יש חומת אש נוספת בין ה-\enLR{DMZ} לפנים.\par

\textbf{ש: "אם יש לי חומת אש, אני מוגן לגמרי?"}\newline ת: לא. חומת אש בודקת כתובות ופורטים, אבל התקפה יכולה להגיע בפורט שפתחת \enLR{(SQL injection} דרך פורט \enLR{80).} לכן צריך גם \enLR{IPS} (פרק \enLR{5)} שבודק תוכן, וגם הגנות בשכבה \enLR{2} (פרק \enLR{6). "}הגנה לעומק".\par

\sectionheading{מילון מונחים – הגדרות}

הגדרות תמציתיות של המונחים המרכזיים בפרק. שימושי גם כדף עזר לבחינה.\par

\begin{xltabular}{\linewidth}{|>{\setRTL\raggedleft\arraybackslash}X|>{\setRTL\raggedleft\arraybackslash}X|}
\hline
\rowcolor{hdrbg}\textbf{הגדרה} & \textbf{מונח} \\
\hline
\endhead
מערכת שמסננת תעבורה בין רשתות לפי מדיניות \enLR{(permit/deny).} & \textbf{חומת אש \enLR{(Firewall)}} \\
\hline
רשימה סדורה של כללי \enLR{permit/deny}; נסרקת מלמעלה למטה, עצירה בהתאמה ראשונה. & \textbf{\enLR{ACL (Access Control List)}} \\
\hline
חסימה סמויה של כל מה שלא הותר במפורש, בסוף כל \enLR{ACL.} & \textbf{\enLR{implicit deny}} \\
\hline
סינון כל חבילה בנפרד, ללא זכירת שיחות \enLR{(ACL} רגיל). & \textbf{\enLR{Stateless}} \\
\hline
זכירת שיחות שיצאו והתרת התשובות התואמות בלבד \enLR{(ZPF/CBAC/ASA).} & \textbf{\enLR{Stateful}} \\
\hline
מספרים \enLR{1}–\enLR{99}; בודק רק כתובת מקור; ממקמים קרוב ליעד. & \textbf{\enLR{ACL} סטנדרטי} \\
\hline
מספרים \enLR{100}–\enLR{199}; בודק מקור, יעד, פרוטוקול ופורט; ממקמים קרוב למקור. & \textbf{\enLR{ACL} מורחב} \\
\hline
מסכה הפוכה למסכת רשת: \enLR{0}=חייב להתאים, \enLR{1}=לא אכפת. = \enLR{255.255.255.255} פחות המסכה. & \textbf{\enLR{Wildcard mask}} \\
\hline
כיוון החלת \enLR{ACL} מנקודת מבט הנתב: נכנס לממשק \enLR{(in)} או יוצא ממנו \enLR{(out).} & \textbf{\enLR{in / out}} \\
\hline
קיצורי \enLR{wildcard: host}=כתובת בודדת \enLR{(0.0.0.0), any}=כל כתובת \enLR{(255.255.255.255).} & \textbf{\enLR{host / any}} \\
\hline
מילת מפתח ב-\enLR{ACL} המתירה רק חבילות תשובה של חיבור \enLR{TCP} קיים. & \textbf{\enLR{established}} \\
\hline
כלל \enLR{ACL} שתקף רק בזמנים מוגדרים (דורש שעון \enLR{NTP).} & \textbf{\enLR{Time-based ACL}} \\
\hline
\enLR{ACL '}מראה' שיוצר אוטומטית התר זמני לתשובות של תעבורה שיצאה. & \textbf{\enLR{Reflexive ACL}} \\
\hline
דלת זמנית שנפתחת רק לאחר שהמשתמש מזדהה בנתב. & \textbf{\enLR{Dynamic ACL (Lock-and-Key)}} \\
\hline
חומת אש \enLR{stateful} קלאסית ב-\enLR{IOS}; בוחנת תעבורה יוצאת ופותחת התר לתשובות. & \textbf{\enLR{CBAC}} \\
\hline
חומת אש מודרנית ב-\enLR{IOS} מבוססת אזורים ומדיניות בין זוגות אזורים. & \textbf{\enLR{ZPF (Zone-Based Firewall)}} \\
\hline
אזור אבטחה שממשקים משויכים אליו; מדיניות מוגדרת בין זוג אזורים בכיוון אחד. & \textbf{\enLR{Zone / Zone-pair}} \\
\hline
פעולות \enLR{ZPF: inspect=stateful} (תשובות אוטומטית), \enLR{pass}=מעבר חד-כיווני, \enLR{drop}=חסימה. & \textbf{\enLR{inspect / pass / drop}} \\
\hline
אזור מפורז לשרתים ציבוריים, מבודד מהרשת הפנימית. & \textbf{\enLR{DMZ}} \\
\hline
חסימת כתובות מקור מזויפות (פרטיות/פנימיות) בכניסה מהאינטרנט. & \textbf{\enLR{Anti-spoofing}} \\
\hline
\end{xltabular}


\clearpage
\phantomsection\addcontentsline{toc}{section}{פרק \enLR{5} – מניעת חדירה לרשת המקומית והרחבה \enLR{(IDS / IPS)}}

\sloppy
\chapterheading{פרק \enLR{5} – מניעת חדירה לרשת המקומית והרחבה \enLR{(IDS / IPS)}}{\small\color{subgray}\enLR{11} שעות עיוני + \enLR{3} מעשי · שבועות \enLR{13}–\enLR{16}\par}\vskip4pt

\begin{metabox}\textbf{מטרות:} מערכות זיהוי פלישה \enLR{(IDS)} · מערכות מניעת פלישה \enLR{(IPS)} · הגדרה ב-\enLR{CLI} · ניטור ופתרון תקלות\\ \textbf{בבגרות:} ההבדל \enLR{IDS/IPS} (ניטור+התראה מול ניטור+חסימה אוטומטית)\end{metabox}

\begin{defbox}\boxtitle{לפני שמתחילים – במה זה שונה מחומת אש? (למי שאין רקע)}\par   \textbf{חומת אש} (פרק \enLR{4)} בודקת \underline{כתובות ופורטים} – "מי מדבר עם מי ובאיזו דלת". היא לא מסתכלת מה \emph{בתוך} החבילה. \textbf{\enLR{IDS/IPS}} מסתכלות על \underline{התוכן} של החבילה ועל \underline{ההתנהגות} – ומחפשות סימנים של התקפה, גם אם הפורט מותר.  \textbf{אנלוגיה:} חומת אש = שומר שבודק תעודת זהות בכניסה. \enLR{IDS/IPS} = מצלמות אבטחה + מאבטח בתוך הבניין שמזהים התנהגות חשודה גם של מי שנכנס חוקית.  \textbf{\enLR{IDS}} \enLR{(Detection)} = מצלמה ש\underline{מתריעה} אם רואה משהו חשוד (אבל לא עוצרת). \textbf{\enLR{IPS}} \enLR{(Prevention)} = מאבטח ש\underline{גם עוצר} את החשוד במקום. \textbf{חתימה \enLR{(Signature)}} = "תמונת מבוקש" – תבנית ידועה של התקפה שהמערכת מחפשת. \end{defbox}

\sectionheading{\enLR{5.1} למה חומת אש לא מספיקה}

חומת אש (פרק \enLR{4)} מחליטה לפי \emph{כתובות ופורטים}. אבל התקפה יכולה להגיע בדיוק בפורט שפתחנו: \enLR{SQL injection} בתוך בקשת \enLR{HTTP} לפורט \enLR{80}, ניצול חולשה בשרת המייל בפורט \enLR{25}, או תולעת שמתפשטת דרך \enLR{SMB} שאנחנו מאפשרים בין סניפים. חומת האש רואה "\enLR{TCP 80} ל-\enLR{web server} – מותר" ולא מסתכלת מה \emph{בתוך} החבילה. כאן נכנסות מערכות \enLR{IDS/IPS}: הן מסתכלות על \textbf{התוכן} וההתנהגות ומחפשות סימנים של התקפה.\par

\begin{storybox}\boxtitle{סיפור מהחיים: \enLR{10} דקות של \enLR{Slammer (2003)}}\par  תולעת \enLR{SQL Slammer} הייתה חבילת \enLR{UDP} \emph{אחת} בגודל \enLR{376} בתים לפורט \enLR{1434 (Microsoft SQL).} כל שרת נגוע שלח אותה לכתובות אקראיות בקצב מרבי. תוך \enLR{10} דקות נדבקו \enLR{75,000} שרתים; כספומטים של \enLR{Bank of America} הפסיקו לעבוד, חלק מטיסות \enLR{Continental} בוטלו, ו-\enLR{911} בסיאטל נפל. מי שהיה עם \enLR{IPS} מעודכן – חתימה אחת ("\enLR{UDP 1434}, גודל \enLR{376}, המחרוזת הזו") חסמה את זה מיד. מי שהיה רק עם חומת אש שאיפשרה \enLR{SQL} בין הסניפים – נדבק. ומה שמדהים: מיקרוסופט פרסמה תיקון \enLR{6} חודשים קודם.\end{storybox}

\sectionheading{\enLR{5.2 IDS} מול \enLR{IPS} – ההבדל המהותי}

\begin{xltabular}{\linewidth}{|>{\setRTL\raggedleft\arraybackslash}X|>{\setRTL\raggedleft\arraybackslash}X|>{\setRTL\raggedleft\arraybackslash}X|}
\hline
\rowcolor{hdrbg}\textbf{\enLR{IPS} – \enLR{Intrusion Prevention System}} & \textbf{\enLR{IDS} – \enLR{Intrusion Detection System}} & \textbf{} \\
\hline
\endhead
מערכת ניטור שמאפשרת זיהוי \textbf{וחסימה אוטומטית} של ניסיונות חדירה & מערכת ניטור שנועדה \textbf{לספק התראות} בעת ניסיונות חדירה & \textbf{הגדרה (כמו בבגרות)} \\
\hline
\textbf{בתוך הנתיב} \enLR{(inline)}: התעבורה \emph{עוברת דרכה}, חבילה אחר חבילה & \textbf{מחוץ לנתיב} \enLR{(out-of-band, promiscuous)}: מקבלת \emph{עותק} של התעבורה דרך \enLR{SPAN/TAP} & \textbf{מיקום} \\
\hline
מפילה את החבילה הזדונית עוד לפני שהגיעה ליעד; חוסמת את החיבור/התוקף & התראה \enLR{(syslog}, מייל), רישום; יכולה לבקש מחומת אש/נתב לחסום \emph{אחרי} שהחבילה הראשונה כבר עברה, או לשלוח \enLR{TCP reset} & \textbf{תגובה} \\
\hline
מוסיפה השהיה \enLR{(latency)}; אם ה-\enLR{IPS} נופל – או שהרשת נופלת \enLR{(fail-closed)} או שעוברים ללא הגנה \enLR{(fail-open)} & אפס השהיה; אם ה-\enLR{IDS} נופל – הרשת ממשיכה & \textbf{השפעה על הרשת} \\
\hline
תעבורה לגיטימית נחסמת – פוגע בעסק. לכן \enLR{IPS} דורש כוונון קפדני & התראה מיותרת – מטריד & \textbf{מחיר של טעות \enLR{(false positive)}} \\
\hline
רואה הכול, כולל חבילות שהיא עצמה מפילה & יכולה לפספס חבילות בעומס (עותק) & \textbf{מה היא רואה} \\
\hline
הגנה בזמן אמת בקצה ובין אזורים & ניטור, \enLR{forensics}, סביבה שלא סובלת חסימה בטעות & \textbf{מקרי שימוש} \\
\hline
\end{xltabular}

\begin{faqbox}\boxtitle{שאלת תלמיד: "אז למה שמישהו יבחר \enLR{IDS}, אם \enLR{IPS} גם מזהה וגם חוסם?"}\par  שלוש סיבות: \enLR{(1)} \textbf{סיכון של \enLR{false positive}} – בבית חולים או בבורסה, לחסום תעבורה לגיטימית בטעות גרוע יותר מלהחמיץ התקפה; מתחילים ב-\enLR{IDS}, מכווננים, ורק אז עוברים ל-\enLR{IPS. (2)} \textbf{ביצועים} – \enLR{IPS inline} על קו \enLR{100Gbps} הוא ציוד יקר מאוד; \enLR{IDS} על עותק יכול לדגום. \enLR{(3)} \textbf{נראות} – \enLR{IDS} פנימי (על \enLR{SPAN} של המתג המרכזי) רואה תנועה רוחבית בין מחשבים שלא עוברת דרך שום חומת אש. בפועל ארגונים מפעילים את שניהם. ורוב מוצרי ה-\enLR{IPS} ניתנים להגדרה במצב \enLR{IDS (monitor-only).}\end{faqbox}

\sectionheading{\enLR{5.3} מבוסס-מארח \enLR{(HIPS)} מול מבוסס-רשת \enLR{(NIPS)}}

\begin{xltabular}{\linewidth}{|>{\setRTL\raggedleft\arraybackslash}X|>{\setRTL\raggedleft\arraybackslash}X|>{\setRTL\raggedleft\arraybackslash}X|}
\hline
\rowcolor{hdrbg}\textbf{\enLR{NIPS} – \enLR{Network-based}} & \textbf{\enLR{HIPS} – \enLR{Host-based}} & \textbf{} \\
\hline
\endhead
מכשיר/מודול בנתיב הרשת \enLR{(ASA} עם מודול \enLR{IPS, Firepower, IOS IPS, Snort/Suricata)} & תוכנה על השרת/המחשב עצמו \enLR{(agent)} & איפה \\
\hline
תעבורת רשת בלבד; לא רואה לתוך \enLR{TLS} בלי \enLR{SSL inspection} & הכול במחשב: קריאות מערכת, קבצים, \enLR{registry}, זיכרון, \textbf{ותעבורה מוצפנת אחרי הפענוח} & מה רואה \\
\hline
מכשיר אחד מגן על מאות מחשבים; לא צואך מערכת ההפעלה; רואה סריקות ו-\enLR{DoS} & מגן גם כשהמחשב מחוץ לרשת (לפטופ בבית); מבין את ההקשר (איזו אפליקציה) & יתרון \\
\hline
עיוור למה שקורה בתוך המארח ולתעבורה מוצפנת & התקנה ותחזוקה בכל מחשב; צורך משאבים; תלוי במערכת ההפעלה & חיסרון \\
\hline
\enLR{Cisco Firepower NGIPS, Snort} (קוד פתוח, שנרכש על ידי סיסקו), \enLR{Suricata, Palo Alto Threat Prevention} & \textbf{\enLR{Cisco Security Agent (CSA)}} – מוזכר בתכנית; הופסק ב-\enLR{2010.} היום: \enLR{Cisco Secure Endpoint (AMP), CrowdStrike, Microsoft Defender} – התחום נקרא \textbf{\enLR{EDR}} & דוגמה \\
\hline
\end{xltabular}

\textbf{\enLR{CSA}} עבד לפי גישה התנהגותית ולא לפי חתימות: הוא יירט קריאות מערכת \enLR{(system calls)} ועצר התנהגות חשודה – "תוכנה שהגיעה במייל מנסה לשנות את ה-\enLR{registry} ולפתוח פורט" – גם בלי להכיר את הנוזקה. זה הרעיון של \enLR{HIPS} מודרני.\par

\sectionheading{\enLR{5.4} שיטות זיהוי}

\begin{xltabular}{\linewidth}{|>{\setRTL\raggedleft\arraybackslash}X|>{\setRTL\raggedleft\arraybackslash}X|>{\setRTL\raggedleft\arraybackslash}X|>{\setRTL\raggedleft\arraybackslash}X|}
\hline
\rowcolor{hdrbg}\textbf{חיסרון} & \textbf{יתרון} & \textbf{איך} & \textbf{שיטה} \\
\hline
\endhead
לא מזהה התקפה חדשה \enLR{(zero-day)}; דורשת עדכונים & מדויקת, מעט \enLR{false positive}, קלה להבנה & השוואה לתבניות ידועות של התקפות (מחרוזת, רצף חבילות) & \textbf{מבוססת חתימות} \enLR{(Signature / pattern)} \\
\hline
הרבה \enLR{false positive}; אם ה-\enLR{baseline} נלמד בזמן התקפה – היא "נורמלית" & מזהה התקפות לא מוכרות & לומדת "מה נורמלי" \enLR{(baseline)} ומתריעה על חריגה – פתאום \enLR{500} חיבורי \enLR{SSH} בדקה & \textbf{מבוססת אנומליה} \enLR{(Anomaly / profile)} \\
\hline
מוגבל למה שחשבו עליו & פשוט, מדויק לארגון & המנהל מגדיר כללים: "אסור \enLR{Telnet} מרשת \enLR{X"} & \textbf{מבוססת מדיניות} \enLR{(Policy)} \\
\hline
לא מגן ישירות; דורש תחזוקה & אפס \enLR{false positive}; לומדים על התוקף & מערכת "מזמינה" שאין לה שימוש לגיטימי – כל גישה אליה חשודה & \textbf{מבוססת פיתיון} \enLR{(Honeypot)} \\
\hline
תוקף חדש לא ברשימה & מהיר & רשימות של \enLR{IP}/דומיינים זדוניים ידועים \enLR{(Cisco Talos)} & \textbf{מוניטין} \enLR{(Reputation)} \\
\hline
\end{xltabular}

\sectionheading{\enLR{5.5} חתימות \enLR{(Signatures)}}

\textbf{חתימה} היא תיאור של תבנית התקפה שה-\enLR{IPS} מחפש. חתימות מקובצות ב\textbf{קובץ חתימות} \enLR{(Signature package / SDF)} שסיסקו מעדכנת; לכל חתימה מזהה \enLR{(ID)} ותת-מזהה \enLR{(subsig)}, ולכל אחת \textbf{מנוע} \enLR{(micro-engine)} שמתאים לסוג – \enLR{atomic-ip, service-http, string-tcp, multi-string} ועוד.\par

\subheading{אטומית מול מורכבת – נשאל בתכנית}

\begin{xltabular}{\linewidth}{|>{\setRTL\raggedleft\arraybackslash}X|>{\setRTL\raggedleft\arraybackslash}X|>{\setRTL\raggedleft\arraybackslash}X|}
\hline
\rowcolor{hdrbg}\textbf{חתימה מורכבת \enLR{(Composite / Stateful)}} & \textbf{חתימה אטומית \enLR{(Atomic)}} & \textbf{} \\
\hline
\endhead
\textbf{רצף} של חבילות/אירועים לאורך זמן – למשל: "יותר מ-\enLR{100 SYN} לפורטים שונים באותו מארח תוך \enLR{5} שניות" \enLR{(port scan)}, או מחרוזת שמפוצלת בין כמה חבילות & \textbf{חבילה אחת} ובודדת – למשל: "\enLR{ICMP echo} עם גודל > \enLR{1000"} או "\enLR{TCP} עם דגלי \enLR{SYN+FIN} יחד" (לא חוקי) & מה בודקת \\
\hline
חייבת לזכור \enLR{(state)}, עד "חלון אירועים" \enLR{(event horizon)} & לא צריך לשמור מצב – מהיר וזול & זיכרון/מצב \\
\hline
\enLR{SYN flood}, סריקת פורטים, \enLR{brute force} & \enLR{Ping of Death, LAND attack} (מקור = יעד) & דוגמה \\
\hline
\end{xltabular}

\subheading{תכונות של חתימה}

\begin{itemize}
\item \textbf{\enLR{Enabled / Disabled}} – פעילה או לא (מושבתת = טעונה בזיכרון אך לא בודקת).
\item \textbf{\enLR{Retired / Unretired}} – בפנסיה = לא טעונה לזיכרון כלל (חוסך משאבים). כדי להשתמש בחתימה היא צריכה להיות \enLR{unretired} \emph{וגם} \enLR{enabled.}
\item \textbf{\enLR{Severity}} – \enLR{informational / low / medium / high.}
\item \textbf{\enLR{Fidelity rating (SFR)}} – עד כמה החתימה מדויקת \enLR{(0}–\enLR{100).}
\item \textbf{\enLR{Action}} – מה לעשות כשיש התאמה (ראו \enLR{5.6).}
\item \textbf{קטגוריות} – ב-\enLR{IOS IPS}: \code{ios\_ips basic} (חתימות בסיסיות, לנתב עם זיכרון מועט) ו-\code{ios\_ips advanced}.
\end{itemize}

\sectionheading{\enLR{5.6} פעולות \enLR{(Actions)} ואזעקות \enLR{(Alarms)}}

\subheading{מה \enLR{IPS} יכול לעשות כשחתימה תואמת}

\begin{xltabular}{\linewidth}{|>{\setRTL\raggedleft\arraybackslash}X|>{\setRTL\raggedleft\arraybackslash}X|}
\hline
\rowcolor{hdrbg}\textbf{משמעות} & \textbf{פעולה} \\
\hline
\endhead
שלח התראה \enLR{(syslog / SDEE).} ברירת מחדל ברוב החתימות. & \textbf{\enLR{produce-alert}} \\
\hline
התראה + עותק של החבילה (לניתוח). & \textbf{\enLR{produce-verbose-alert}} \\
\hline
הפל את החבילה הזו. & \textbf{\enLR{deny-packet-inline}} \\
\hline
הפל את כל החבילות של החיבור \enLR{(TCP flow)} הזה. & \textbf{\enLR{deny-connection-inline}} \\
\hline
הפל כל תעבורה מכתובת התוקף לזמן מוגדר – החזק ביותר, והמסוכן ביותר ב-\enLR{false positive.} & \textbf{\enLR{deny-attacker-inline}} \\
\hline
שלח \enLR{TCP RST} לשני הצדדים – היחידה שגם \enLR{IDS} (מחוץ לנתיב) יכולה לבצע. & \textbf{\enLR{reset-tcp-connection}} \\
\hline
\end{xltabular}

\subheading{ארבעה סוגי תוצאות – חובה להבין}

\begin{xltabular}{\linewidth}{|>{\setRTL\raggedleft\arraybackslash}X|>{\setRTL\raggedleft\arraybackslash}X|>{\setRTL\raggedleft\arraybackslash}X|}
\hline
\rowcolor{hdrbg}\textbf{לא הייתה התקפה} & \textbf{הייתה התקפה} & \textbf{} \\
\hline
\endhead
\textbf{\enLR{False Positive}} – אזעקת שווא; מטריד, וב-\enLR{IPS} חוסם משהו לגיטימי & \textbf{\enLR{True Positive}} – מצוין & \textbf{המערכת התריעה} \\
\hline
\textbf{\enLR{True Negative}} – מצוין & \textbf{\enLR{False Negative}} – הכי גרוע: התקפה עברה בשקט & \textbf{המערכת שתקה} \\
\hline
\end{xltabular}

\begin{faqbox}\boxtitle{שאלת תלמיד: "מה יותר גרוע – \enLR{false positive} או \enLR{false negative}?"}\par  תלוי במערכת. ב-\textbf{\enLR{IDS}} \enLR{false positive} זה רק רעש; \enLR{false negative} זה פריצה שלא ראינו – גרוע. ב-\textbf{\enLR{IPS}} \enLR{false positive} חוסם עסק (הזמנות לא נכנסות, טלפוניה נופלת) – זה מה שמפחיד מנהלים ולכן הם מעדיפים לפעמים לכבות חתימות; ואז מגיע ה-\enLR{false negative.} האיזון נקרא \textbf{כוונון \enLR{(tuning)}} – מכבים חתימות לא רלוונטיות (חתימות ל-\enLR{IIS} ברשת של \enLR{Linux)}, מעלים סף, ומחריגים כתובות מהימנות. אין תשובה אחת; זה מה שעושים אנשי \enLR{SOC} כל היום.\end{faqbox}

\sectionheading{\enLR{5.7} ניטור \enLR{IPS}}

\begin{itemize}
\item \textbf{\enLR{Syslog}} – ההתראות נשלחות כהודעות (\code{\%IPS-4-SIGNATURE: Sig:2004 Subsig:0 Sev:25 ICMP Echo Request [10.1.1.5:0 -> 10.2.2.2:0]}).
\item \textbf{\enLR{SDEE}} \enLR{(Security Device Event Exchange)} – פרוטוקול מבוסס \enLR{HTTPS/XML} שסיסקו יצרה כדי שתחנת ניהול תשלוף אירועים מה-\enLR{IPS (pull)} בצורה מאובטחת. דורש \code{ip http secure-server}.
\item \textbf{תחנות ניהול:} בעבר \enLR{IEV / IME / MARS} (מוזכרים בחומרים ישנים); היום \enLR{Firepower Management Center (FMC)}, או \enLR{SIEM} כללי \enLR{(Splunk, QRadar, ELK)} שמקבל \enLR{syslog} מכל המערכות ומבצע \textbf{קורלציה} – "אותו \enLR{IP} הופיע ב-\enLR{IPS}, בחומת האש וב-\enLR{AD} תוך דקה".
\end{itemize}

\sectionheading{\enLR{5.8} הגדרת \enLR{IOS IPS} ב-\enLR{CLI}}

נתב סיסקו יכול לתפקד כ-\enLR{IPS} (מוגבל) בעצמו. התהליך המלא בציוד אמיתי:\par

\begin{enumerate}
\item הורדת חבילת חתימות (\code{IOS-Sxxx-CLI.pkg}) ומפתח ציבורי (\code{realm-cisco.pub.key.txt}) מ-\enLR{Cisco.com.}
\item יצירת תיקייה ב-\enLR{flash} לחתימות.
\item הגדרת המפתח הציבורי (כדי שהנתב יוודא שהחבילה חתומה על ידי סיסקו).
\item הגדרת \enLR{IPS}: שם, מיקום, התראות, קטגוריות.
\item החלה על ממשק וטעינת החבילה.
\end{enumerate}

\begin{codebox}\begin{english}\codefont\footnotesize\color{cfg}\setlength{\parindent}{0pt}
\textcolor{ccom}{\texthebrew{! שלב \enLR{2} – תיקייה}}\\
R1\# \textcolor{cbold}{\textbf{mkdir ipsdir}}\\
\textcolor{ccom}{\texthebrew{! שלב \enLR{3} – מפתח ציבורי של סיסקו (מעתיקים מהקובץ; קטע מקוצר)}}\\
R1(config)\# crypto key pubkey-chain rsa\\
R1(config-pubkey-chain)\# \textcolor{cbold}{\textbf{named-key realm-cisco.pub signature}}\\
R1(config-pubkey-key)\# key-string\\
R1(config-pubkey)\# 30820122 300D0609 2A864886 ...\\
R1(config-pubkey)\# quit\\
\textcolor{ccom}{\texthebrew{! שלב \enLR{4} – הגדרת ה-\enLR{IPS}}}\\
R1(config)\# \textcolor{cbold}{\textbf{ip ips config location flash:ipsdir}}\\
R1(config)\# \textcolor{cbold}{\textbf{ip ips name MYIPS}}            \textcolor{ccom}{\texthebrew{! אפשר: \enLR{ip ips name MYIPS list 101} – רק תעבורה שתואמת ל-\enLR{ACL}}}\\
R1(config)\# \textcolor{cbold}{\textbf{ip ips notify log}}            \textcolor{ccom}{\texthebrew{! התראות ל-\enLR{syslog}}}\\
R1(config)\# ip http secure-server\\
R1(config)\# \textcolor{cbold}{\textbf{ip ips notify sdee}}           \textcolor{ccom}{\texthebrew{! ולתחנת ניהול}}\\
R1(config)\# \textcolor{cbold}{\textbf{ip ips signature-category}}\\
R1(config-ips-category)\# \textcolor{cbold}{\textbf{category all}}\\
R1(config-ips-category-action)\# \textcolor{cbold}{\textbf{retired true}}          \textcolor{ccom}{\texthebrew{! קודם – הכול בפנסיה}}\\
R1(config-ips-category-action)\# exit\\
R1(config-ips-category)\# \textcolor{cbold}{\textbf{category ios\_ips basic}}\\
R1(config-ips-category-action)\# \textcolor{cbold}{\textbf{retired false}}         \textcolor{ccom}{\texthebrew{! ואז – רק הבסיסיות פעילות}}\\
R1(config-ips-category-action)\# exit\\
R1(config-ips-category)\# exit\\
\textcolor{ccom}{\texthebrew{! שלב \enLR{5} – החלה על ממשק (בכיוון שממנו מגיעה התעבורה החשודה)}}\\
R1(config)\# interface g0/0\\
R1(config-if)\# \textcolor{cbold}{\textbf{ip ips MYIPS in}}\\
\textcolor{ccom}{\texthebrew{! טעינת חבילת החתימות (מ-\enLR{TFTP/USB)} – ב-\enLR{Packet Tracer} שלב זה לא קיים}}\\
R1\# \textcolor{cbold}{\textbf{copy tftp://10.0.0.5/IOS-S416-CLI.pkg idconf}}
\end{english}\end{codebox}

\subheading{שינוי חתימה בודדת – הדוגמה הקלאסית: התראה + חסימה על \enLR{ping}}

\begin{codebox}\begin{english}\codefont\footnotesize\color{cfg}\setlength{\parindent}{0pt}
R1(config)\# \textcolor{cbold}{\textbf{ip ips signature-definition}}\\
R1(config-sigdef)\# \textcolor{cbold}{\textbf{signature 2004 0}}           \textcolor{ccom}{\texthebrew{! \enLR{2004/0 = ICMP Echo Request}}}\\
R1(config-sigdef-sig)\# \textcolor{cbold}{\textbf{status}}\\
R1(config-sigdef-sig-status)\# \textcolor{cbold}{\textbf{retired false}}\\
R1(config-sigdef-sig-status)\# \textcolor{cbold}{\textbf{enabled true}}\\
R1(config-sigdef-sig-status)\# exit\\
R1(config-sigdef-sig)\# \textcolor{cbold}{\textbf{engine}}\\
R1(config-sigdef-sig-engine)\# \textcolor{cbold}{\textbf{event-action produce-alert}}\\
R1(config-sigdef-sig-engine)\# \textcolor{cbold}{\textbf{event-action deny-packet-inline}}\\
R1(config-sigdef-sig-engine)\# exit\\
R1(config-sigdef-sig)\# exit\\
R1(config-sigdef)\# exit\\
Do you want to accept these changes? [confirm]
\end{english}\end{codebox}

\begin{tipbox}\boxtitle{טיפ: מה עובד ב-\enLR{Packet Tracer}}\par  \enLR{Packet Tracer} תומך בתת-קבוצה: \code{ip ips config location}, \code{ip ips name}, \code{ip ips notify log}, \code{signature-category} \enLR{(all retired / ios\_ips basic)}, \code{ip ips NAME in|out}, ושינוי חתימה \enLR{2004 0 (ICMP echo)} עם \enLR{produce-alert} ו-\enLR{deny-packet-inline.} אין טעינת חבילה ואין \enLR{SDEE.} תרגיל ה-\enLR{ping} הוא התרגיל שעובד – ורואים בשרת ה-\enLR{syslog} את ההתראה ואת ה-\enLR{ping} נכשל. זה מספיק להמחשה.\end{tipbox}

\sectionheading{\enLR{5.9} וריפיקציה ופתרון תקלות}

\begin{codebox}\begin{english}\codefont\footnotesize\color{cfg}\setlength{\parindent}{0pt}
R1\# \textcolor{cbold}{\textbf{show ip ips all}}                   \textcolor{ccom}{\texthebrew{! סיכום: שם, ממשקים, קטגוריות, מספר חתימות טעונות}}\\
R1\# \textcolor{cbold}{\textbf{show ip ips configuration}}\\
R1\# \textcolor{cbold}{\textbf{show ip ips interfaces}}\\
R1\# \textcolor{cbold}{\textbf{show ip ips signatures}} [count]      \textcolor{ccom}{\texthebrew{! אילו חתימות פעילות}}\\
R1\# \textcolor{cbold}{\textbf{show ip ips statistics}}            \textcolor{ccom}{\texthebrew{! כמה חבילות נבדקו, כמה התראות}}\\
R1\# \textcolor{cbold}{\textbf{clear ip ips statistics}}\\
R1\# show ip ips sessions\\
R1\# debug ip ips
\end{english}\end{codebox}

\begin{xltabular}{\linewidth}{|>{\setRTL\raggedleft\arraybackslash}X|>{\setRTL\raggedleft\arraybackslash}X|}
\hline
\rowcolor{hdrbg}\textbf{סיבה נפוצה ופתרון} & \textbf{תסמין} \\
\hline
\endhead
ה-\enLR{IPS} לא מוחל על הממשק / בכיוון הלא נכון. תעבורה נכנסת מהאינטרנט ⇒ \code{in} על הממשק החיצוני. בדקו \code{show ip ips interfaces}. & אין התראות בכלל \\
\hline
קטגוריה \enLR{all retired} ולא \enLR{unretired} כלום; או חבילת חתימות לא נטענה (\code{copy … idconf}); או מפתח ציבורי חסר/שגוי – הנתב דוחה חבילה לא חתומה. & "\enLR{0 signatures loaded"} \\
\hline
יותר מדי חתימות \enLR{(advanced} על נתב קטן). השאירו \enLR{ios\_ips basic}, השתמשו ב-\code{list ACL} כדי לבדוק רק תעבורה רלוונטית. & הנתב איטי / \enLR{CPU 100\%} \\
\hline
ה-\enLR{action} הוא \enLR{produce-alert} בלבד; הוסיפו \enLR{deny-packet-inline.} או שהחתימה \enLR{enabled} אך \enLR{retired.} & יש התראה אך החבילה עוברת \\
\hline
\enLR{false positive.} זהו את החתימה ב-\enLR{syslog} ובטלו/כווננו אותה; או השתמשו ב-\code{ip ips name X list ACL} להחריג. & תעבורה לגיטימית נחסמת \\
\hline
\end{xltabular}

\sectionheading{\enLR{5.10} מדריך פקודות מרוכז – פרק \enLR{5}}

\begin{xltabular}{\linewidth}{|>{\setRTL\raggedleft\arraybackslash}X|>{\setRTL\raggedleft\arraybackslash}X|}
\hline
\rowcolor{hdrbg}\textbf{תפקיד} & \textbf{פקודה} \\
\hline
\endhead
תיקייה לחתימות & \code{mkdir ipsdir} \\
\hline
מפתח לאימות חבילה & \code{crypto key pubkey-chain rsa} → \code{named-key realm-cisco.pub signature} \\
\hline
מיקום & \code{ip ips config location flash:ipsdir} \\
\hline
כלל \enLR{IPS} & \code{ip ips name NAME [list ACL]} \\
\hline
יעד התראות & \code{ip ips notify log | sdee} \\
\hline
כיבוי כל החתימות & \code{ip ips signature-category} → \code{category all} → \code{retired true} \\
\hline
הפעלת הבסיסיות & \code{category ios\_ips basic} → \code{retired false} \\
\hline
החלה & \code{ip ips NAME in|out} (ממשק) \\
\hline
טעינת חתימות & \code{copy tftp://…pkg idconf} \\
\hline
עריכת חתימה & \code{ip ips signature-definition} → \code{signature ID SUB} → \code{status}/\code{engine} \\
\hline
וריפיקציה & \code{show ip ips all | configuration | signatures | statistics} \\
\hline
\end{xltabular}

\sectionheading{\enLR{5.11} תרגילים לתלמידים}

\begin{exbox}\boxtitle{תרגיל \enLR{1} – \enLR{True/False Positive/Negative}}\par  סווגו: א. ה-\enLR{IPS} חסם עדכון \enLR{Windows} לגיטימי כי הוא "נראה כמו הורדת קובץ הרצה חשוד". ב. תולעת עברה דרך ה-\enLR{IDS} בלי התראה כי החתימה לא עודכנה. ג. ה-\enLR{IDS} התריע על סריקת \enLR{Nmap} שאכן בוצעה. ד. תעבורת \enLR{Zoom} עברה כרגיל ללא התראה. \par\vskip2pt\hrule\vskip3pt\textbf{}א. \enLR{False Positive} ב. \enLR{False Negative} ג. \enLR{True Positive} ד. \enLR{True Negative}\end{exbox}

\begin{exbox}\boxtitle{תרגיל \enLR{2} – אטומית או מורכבת?}\par  א. חבילת \enLR{TCP} עם כל \enLR{6} הדגלים דלוקים ("\enLR{Xmas scan").} ב. \enLR{50} ניסיונות כניסה כושלים ל-\enLR{SSH} מאותו \enLR{IP} תוך דקה. ג. חבילת \enLR{IP} שבה כתובת המקור שווה לכתובת היעד. ד. מחרוזת \code{/etc/passwd} שמפוצלת בין שתי חבילות \enLR{HTTP.} \par\vskip2pt\hrule\vskip3pt\textbf{}א. אטומית ב. מורכבת ג. אטומית \enLR{(LAND)} ד. מורכבת (דורשת הרכבת זרם)\end{exbox}

\begin{exbox}\boxtitle{תרגיל \enLR{3} – בחירת פעולה}\par  לכל מצב בחרו את ה-\enLR{action} המתאים ביותר ונמקו: א. חתימה עם \enLR{fidelity 30} (לא מדויקת) לפעילות חשודה קלה. ב. חתימה ל-\enLR{SQL injection} ודאית \enLR{(fidelity 95)} לשרת ה-\enLR{web.} ג. תוקף ששולח \enLR{10,000 SYN} בשנייה. ד. \enLR{IDS} מחוץ לנתיב שזיהה \enLR{session} זדוני. \par\vskip2pt\hrule\vskip3pt\textbf{}א. \enLR{produce-alert} בלבד – סיכון \enLR{false positive} גבוה. ב. \enLR{deny-connection-inline} (או \enLR{deny-packet) + alert.} ג. \enLR{deny-attacker-inline} – לחסום את המקור לזמן. ד. \enLR{reset-tcp-connection} – היחיד שזמין ל-\enLR{IDS} מחוץ לנתיב.\end{exbox}

\begin{exbox}\boxtitle{תרגיל \enLR{4} – תכנון פריסה}\par  ציירו רשת: אינטרנט – נתב קצה – חומת אש – מתג מרכזי – (שרתים, משתמשים). סמנו היכן תמקמו \enLR{IPS} והיכן \enLR{IDS}, ונמקו. היכן תתקינו \enLR{HIPS}? \par\vskip2pt\hrule\vskip3pt\textbf{}\enLR{IPS inline} בין חומת האש למתג המרכזי (או כמודול בחומת האש) – מגן על הכול מהאינטרנט. \enLR{IDS} על \enLR{SPAN} של המתג המרכזי – רואה תנועה רוחבית פנימית. \enLR{HIPS/EDR} על השרתים (בעיקר \enLR{web, DB, DC)} ועל מחשבי המשתמשים (ניידים).\end{exbox}

\sectionheading{\enLR{5.12} שאלות בסגנון בגרות}

\begin{exambox}\boxtitle{שאלה \enLR{1}}\par  התאימו: מערכת \_\_\_\_ היא מערכת ניטור שמאפשרת זיהוי וחסימה אוטומטית של ניסיונות חדירה. מערכת \_\_\_\_ היא מערכת ניטור שנועדה לספק התראות בעת ניסיונות חדירה. \enLR{(IDS / IPS / DNS / RADIUS / TACACS)} \par\vskip2pt\hrule\vskip3pt\textbf{}\enLR{IPS} · \enLR{IDS}\end{exambox}

\begin{exambox}\boxtitle{שאלה \enLR{2}}\par  מהו ההבדל במיקום ברשת בין \enLR{IDS} ל-\enLR{IPS}, ומה המשמעות לגבי השהיה? \par\vskip2pt\hrule\vskip3pt\textbf{}\enLR{IDS} – מחוץ לנתיב \enLR{(promiscuous)}, מקבל עותק ⇒ אין השהיה. \enLR{IPS} – בתוך הנתיב \enLR{(inline)}, התעבורה עוברת דרכו ⇒ מוסיף השהיה.\end{exambox}

\begin{exambox}\boxtitle{שאלה \enLR{3}}\par  מה ההבדל בין חתימה אטומית לחתימה מורכבת? \par\vskip2pt\hrule\vskip3pt\textbf{}אטומית – בודקת חבילה בודדת, ללא שמירת מצב. מורכבת – בודקת רצף חבילות לאורך זמן, שומרת מצב.\end{exambox}

\begin{exambox}\boxtitle{שאלה \enLR{4}}\par  מנהל הגדיר \enLR{IOS IPS} ו-\code{category all retired true} בלבד. אילו חתימות פעילות? \par\vskip2pt\hrule\vskip3pt\textbf{}אף אחת – חייבים לבצע \enLR{unretire} (\code{retired false}) לקטגוריה כלשהי, למשל \enLR{ios\_ips basic.}\end{exambox}

\begin{exambox}\boxtitle{שאלה \enLR{5}}\par  איזו מהפעולות הבאות אינה זמינה ל-\enLR{IDS} שנמצא מחוץ לנתיב התעבורה? \enLR{1. produce-alert 2. deny-packet-inline 3. reset-tcp-connection 4. produce-verbose-alert} \par\vskip2pt\hrule\vskip3pt\textbf{}\textbf{\enLR{2}} – אי אפשר להפיל חבילה שכבר עברה; ה-\enLR{IDS} רואה רק עותק.\end{exambox}

\begin{exambox}\boxtitle{שאלה \enLR{6}}\par  מהו \enLR{HIPS} ומה יתרונו המרכזי על פני \enLR{IPS} רשתי? \par\vskip2pt\hrule\vskip3pt\textbf{}\enLR{Host-based IPS} – תוכנה על המארח עצמו. יתרון: רואה תעבורה מוצפנת אחרי הפענוח ואת ההתנהגות בתוך המחשב (קבצים, זיכרון), ומגן גם כשהמחשב מחוץ לרשת הארגון.\end{exambox}

\sectionheading{\enLR{5.13} מעבדה \enLR{(3} שעות) – \enLR{IOS IPS} ב-\enLR{Packet Tracer}}

\begin{enumerate}
\item טופולוגיה: \enLR{PC-A (10.1.1.10)} — \enLR{R1} — \enLR{PC-B (10.2.2.10)}; שרת \enLR{Syslog} ברשת של \enLR{PC-B.} הגדירו \enLR{logging} לשרת.
\item ב-\enLR{R1}: \code{mkdir ipsdir}, \code{ip ips config location flash:ipsdir}, \code{ip ips name IOSIPS}, \code{ip ips notify log}, קטגוריות: \enLR{all retired, ios\_ips basic unretired.}
\item שנו חתימה \enLR{2004 0: retired false, enabled true, event-action produce-alert + deny-packet-inline.}
\item החילו \code{ip ips IOSIPS out} על הממשק לכיוון \enLR{PC-B} (או \enLR{in} על הממשק מ-\enLR{PC-A).}
\item מ-\enLR{PC-A: ping} ל-\enLR{PC-B} ⇒ נכשל. בשרת ה-\enLR{syslog}: הודעת \code{\%IPS-4-SIGNATURE: Sig:2004}. מ-\enLR{PC-B ping} ל-\enLR{PC-A} ⇒ מצליח (הכיוון!). הסבירו.
\item הסירו את \enLR{deny-packet-inline} והשאירו \enLR{produce-alert. ping} עובר אך יש התראה – זו התנהגות \enLR{IDS.} דונו בהבדל.
\item \code{show ip ips all}, \code{show ip ips statistics} – כמה חבילות נבדקו?
\end{enumerate}

\sectionheading{בנק שאלות שתלמידים שואלים – ותשובות מוכנות}

שאלות נפוצות בפרק \enLR{IDS/IPS}, עם תשובות מוכנות.\par

\textbf{ש: "אם \enLR{IPS} גם מזהה וגם חוסם, למה שמישהו יבחר \enLR{IDS}?"}\newline ת: שלוש סיבות: \enLR{(1)} סיכון של \textbf{\enLR{false positive}} – בבית חולים או בבורסה, לחסום תעבורה לגיטימית בטעות גרוע מלהחמיץ התקפה. \enLR{(2)} ביצועים – \enLR{IPS inline} על קו מהיר יקר מאוד. \enLR{(3) IDS} פנימי (על \enLR{SPAN)} רואה תנועה רוחבית בין מחשבים שלא עוברת בשום חומת אש. בפועל מפעילים את שניהם.\par

\textbf{ש: "מה ההבדל בין \enLR{IDS/IPS} לאנטי-וירוס?"}\newline ת: אנטי-וירוס מגן על \underline{קבצים במחשב אחד}. \enLR{IPS} מגן על \underline{תעבורת הרשת} בין מחשבים. \enLR{HIPS (host-based)} מטשטש את הגבול – הוא \enLR{IPS} שרץ על המחשב עצמו.\par

\textbf{ש: "מה יותר גרוע – \enLR{false positive} או \enLR{false negative}?"}\newline ת: תלוי במערכת. ב-\enLR{IDS, false positive} הוא רק רעש; \enLR{false negative} (התקפה שעברה בשקט) גרוע יותר. ב-\enLR{IPS, false positive} \underline{חוסם עסק} (הזמנות לא נכנסות) – ולכן מנהלים לפעמים מכבים חתימות, וזה מוביל ל-\enLR{false negatives.} האיזון הזה נקרא "כוונון" \enLR{(tuning)} וזו עבודת ה-\enLR{SOC.}\par

\textbf{ש: "מה ההבדל בין חתימה אטומית למורכבת?"}\newline ת: אטומית בודקת \underline{חבילה אחת} (למשל \enLR{Ping of Death} – חבילה גדולה מדי) – מהירה, בלי זיכרון. מורכבת בודקת \underline{רצף} של חבילות לאורך זמן (למשל סריקת פורטים או \enLR{SYN flood)} – חייבת לזכור מצב.\par

\textbf{ש: "אם \enLR{IPS} מבוסס חתימות, איך הוא תופס התקפה חדשה שאין לה חתימה?"}\newline ת: הוא לא, אם הוא רק מבוסס-חתימות – זו החולשה. לכן קיימות שיטות נוספות: זיהוי \textbf{אנומליה} (חריגה מ"נורמלי") ו\textbf{מוניטין} (רשימות \enLR{IP} זדוניים). מערכות מודרניות משלבות כמה שיטות.\par

\textbf{ש: "מה זה \enLR{SPAN} ולמה \enLR{IDS} צריך אותו?"}\newline ת: \enLR{SPAN} הוא "שיקוף פורט" – המתג מעתיק את כל התעבורה לפורט שאליו מחובר ה-\enLR{IDS.} כך ה-\enLR{IDS} מקבל \underline{עותק} של הכול בלי לשבת בנתיב עצמו – ולכן אם ה-\enLR{IDS} נופל, הרשת ממשיכה לעבוד.\par

\textbf{ש: "אילו פעולות \enLR{IPS} יכול לבצע כשהוא מזהה התקפה?"}\newline ת: מ-\enLR{alert} בלבד \enLR{(produce-alert)}, דרך הפלת החבילה \enLR{(deny-packet-inline)}, הפלת כל החיבור \enLR{(deny-connection-inline)}, חסימת התוקף לזמן \enLR{(deny-attacker-inline)}, ועד שליחת \enLR{TCP reset} – הפעולה היחידה שגם \enLR{IDS} מחוץ לנתיב יכול לבצע.\par

\textbf{ש: "מה זה \enLR{CSA} שמופיע בתכנית?"}\newline ת: \enLR{Cisco Security Agent} – תוכנת \enLR{HIPS} ישנה שהותקנה על תחנות קצה וזיהתה התנהגות חשודה (לא רק חתימות). היא הופסקה ב-\enLR{2010}; היום התחום נקרא \enLR{EDR} (כמו \enLR{Microsoft Defender, CrowdStrike).}\par

\sectionheading{מילון מונחים – הגדרות}

הגדרות תמציתיות של המונחים המרכזיים בפרק. שימושי גם כדף עזר לבחינה.\par

\begin{xltabular}{\linewidth}{|>{\setRTL\raggedleft\arraybackslash}X|>{\setRTL\raggedleft\arraybackslash}X|}
\hline
\rowcolor{hdrbg}\textbf{הגדרה} & \textbf{מונח} \\
\hline
\endhead
מערכת שמנטרת ומתריעה על ניסיונות חדירה, אך אינה חוסמת. & \textbf{\enLR{IDS (Intrusion Detection System)}} \\
\hline
מערכת שמזהה וגם חוסמת אוטומטית ניסיונות חדירה. & \textbf{\enLR{IPS (Intrusion Prevention System)}} \\
\hline
\enLR{IPS} יושב בנתיב התעבורה \enLR{(inline); IDS} מקבל עותק מחוץ לנתיב \enLR{(out-of-band).} & \textbf{\enLR{inline / out-of-band}} \\
\hline
\enLR{IPS} מבוסס-מארח – תוכנה על המחשב/שרת עצמו; רואה גם תעבורה מוצפנת. & \textbf{\enLR{HIPS}} \\
\hline
\enLR{IPS} מבוסס-רשת – מכשיר בנתיב שמגן על מארחים רבים. & \textbf{\enLR{NIPS}} \\
\hline
תוכנת \enLR{HIPS} ישנה של סיסקו; היום התחום נקרא \enLR{EDR.} & \textbf{\enLR{CSA (Cisco Security Agent)}} \\
\hline
תבנית ידועה של התקפה שהמערכת מחפשת בתעבורה. & \textbf{חתימה \enLR{(Signature)}} \\
\hline
בודקת חבילה בודדת, ללא שמירת מצב (למשל \enLR{Ping of Death).} & \textbf{חתימה אטומית} \\
\hline
בודקת רצף חבילות לאורך זמן, שומרת מצב (למשל סריקת פורטים). & \textbf{חתימה מורכבת} \\
\hline
זיהוי חריגה מ'התנהגות נורמלית' שנלמדה \enLR{(baseline).} & \textbf{זיהוי מבוסס-אנומליה} \\
\hline
מערכת פיתיון שאין לה שימוש לגיטימי; כל גישה אליה חשודה. & \textbf{\enLR{Honeypot}} \\
\hline
אזעקת שווא – התראה על תעבורה לגיטימית. & \textbf{\enLR{False Positive}} \\
\hline
החמצה – התקפה אמיתית שלא זוהתה (המצב הגרוע ביותר). & \textbf{\enLR{False Negative}} \\
\hline
זיהוי נכון של התקפה \enLR{(positive)} או של תעבורה תקינה \enLR{(negative).} & \textbf{\enLR{True Positive / Negative}} \\
\hline
התאמת חתימות והחרגות כדי לצמצם \enLR{false positives}; עבודת ה-\enLR{SOC.} & \textbf{\enLR{Tuning} (כוונון)} \\
\hline
מה שה-\enLR{IPS} עושה בהתאמה: \enLR{alert, deny-packet/connection/attacker, TCP reset.} & \textbf{\enLR{Action} (פעולה)} \\
\hline
פרוטוקול מאובטח \enLR{(HTTPS/XML)} שדרכו תחנת ניהול שולפת אירועי \enLR{IPS.} & \textbf{\enLR{SDEE}} \\
\hline
שיקוף פורט – העתקת תעבורה לפורט שאליו מחובר \enLR{IDS}/מנתח. & \textbf{\enLR{SPAN}} \\
\hline
\end{xltabular}


\clearpage
\phantomsection\addcontentsline{toc}{section}{פרק \enLR{6} – אבטחת הרשת המקומית \enLR{(LAN Security)}}

\sloppy
\chapterheading{פרק \enLR{6} – אבטחת הרשת המקומית \enLR{(LAN Security)}}{\small\color{subgray}\enLR{11} שעות עיוני + \enLR{2} מעשי · שבועות \enLR{17}–\enLR{19}\par}\vskip4pt

\begin{metabox}\textbf{מטרות:} התקפות שכבה \enLR{2} · \enLR{port security} · \enLR{STP/VLAN} · \enLR{DHCP/ARP} · \enLR{storm control} · \enLR{SPAN} · \enLR{NAC} · \enLR{VoIP/SAN}\\ \textbf{בבגרות:} \enLR{MAC overflow}→\enLR{port security, DTP}→\enLR{VLAN hopping, BPDU Guard, Evil Twin, root bridge}\end{metabox}

\begin{defbox}\boxtitle{לפני שמתחילים – מה קורה ב"שכבה \enLR{2"}? (למי שאין רקע)}\par  עד עכשיו דיברנו על כתובות \enLR{IP} (שכבה \enLR{3).} אבל בתוך רשת מקומית, המכשירים מדברים ביניהם דרך \textbf{מתג} לפי \textbf{כתובות \enLR{MAC}} (שכבה \enLR{2).} כמה מושגים:  \textbf{מתג \enLR{(Switch)}} – מחבר את כל המחשבים ברשת המקומית. הוא לומד איזו כתובת \enLR{MAC} נמצאת באיזה פורט, ושומר זאת ב\textbf{טבלת \enLR{MAC}}. \textbf{\enLR{VLAN}} – "רשת וירטואלית". מפצל מתג פיזי אחד לכמה רשתות נפרדות ומבודדות (למשל \enLR{VLAN} למחלקת כספים ו-\enLR{VLAN} לאורחים) בלי לקנות עוד מתגים. \textbf{\enLR{Trunk}} – קישור בין מתגים שנושא כמה \enLR{VLAN}ים יחד. \textbf{\enLR{Broadcast}} – הודעה שנשלחת ל\underline{כל} המכשירים ברשת בבת אחת ("שידור"). \textbf{\enLR{STP}} – פרוטוקול שמונע "לולאות" בין מתגים (נסביר בהמשך למה לולאה מפילה רשת).  \textbf{למה זה קריטי לאבטחה?} אם תוקף כבר מחובר לרשת המקומית (חיבר מחשב לשקע, פרץ ל-\enLR{Wi-Fi)}, הוא פועל בשכבה \enLR{2} – \underline{מתחת} לכל ההגנות של שכבה \enLR{3} (חומת אש, \enLR{IP).} לכן צריך להגן גם כאן.\end{defbox}

\sectionheading{\enLR{6.1} למה שכבה \enLR{2} היא "הבטן הרכה"}

כל מנגנוני האבטחה שלמדנו עד כה עובדים בשכבה \enLR{3} ומעלה \enLR{(ACL, IPS, VPN).} אבל אם התוקף כבר \textbf{ברשת המקומית} – חיבר לפטופ לשקע, פרץ ל-\enLR{Wi-Fi}, או השתלט על מחשב עובד – הוא פועל בשכבה \enLR{2}, מתחת לכל ההגנות האלה. עיקרון קריטי: \textbf{מודל \enLR{OSI} הוא שרשרת – אם שכבה \enLR{2} נפרצת, כל השכבות מעליה נפגעות.} מתג "בוטח" בכל מה שמחובר אליו כברירת מחדל, וזו הבעיה.\par

\begin{storybox}\boxtitle{סיפור מהחיים: השקע בחדר הישיבות}\par  בבדיקת חדירה נפוצה, הבודק מגיע כ"טכניקאי", מבקש להמתין בחדר ישיבות, ומחבר מכשיר קטן \enLR{(Raspberry Pi)} לשקע רשת מתחת לשולחן. המכשיר מבצע \enLR{ARP spoofing} ומקליט תעבורה, או מריץ \enLR{Responder} לגניבת \enLR{hashes} של סיסמאות. הוא יוצא אחרי חצי שעה, והמכשיר ממשיך לשדר דרך \enLR{LTE.} אף חומת אש לא ראתה כלום – הכול קרה בשכבה \enLR{2}, בתוך ה-\enLR{VLAN.} ההגנות בפרק הזה \enLR{(port security, DHCP snooping, DAI, 802.1X)} הן בדיוק מה שהיה עוצר אותו.\end{storybox}

\sectionheading{\enLR{6.2} מפת ההתקפות וההגנות בשכבה \enLR{2}}

\begin{xltabular}{\linewidth}{|>{\setRTL\raggedleft\arraybackslash}X|>{\setRTL\raggedleft\arraybackslash}X|>{\setRTL\raggedleft\arraybackslash}X|}
\hline
\rowcolor{hdrbg}\textbf{הגנה עיקרית} & \textbf{מה מנצלת} & \textbf{התקפה} \\
\hline
\endhead
\textbf{\enLR{Port Security}} & גודל מוגבל של טבלת ה-\enLR{MAC} & \enLR{MAC table overflow (CAM flooding)} \\
\hline
כיבוי \enLR{DTP}, \code{switchport mode access} & \enLR{DTP} – משא-ומתן \enLR{trunk} אוטומטי & \enLR{VLAN hopping (switch spoofing)} \\
\hline
\enLR{Native VLAN} ייעודי ולא בשימוש & \enLR{Native VLAN} & \enLR{VLAN hopping (double tagging)} \\
\hline
\textbf{\enLR{BPDU Guard, Root Guard}} & בחירת \enLR{root bridge} לפי \enLR{priority} & \enLR{STP manipulation} \\
\hline
\textbf{\enLR{DHCP Snooping}} & אין אימות ל-\enLR{DHCP} & \enLR{DHCP starvation / spoofing (rogue DHCP)} \\
\hline
\textbf{\enLR{Dynamic ARP Inspection (DAI)}} (מסתמך על \enLR{DHCP snooping)} & \enLR{ARP} חסר אימות & \enLR{ARP spoofing / poisoning (MITM)} \\
\hline
\enLR{IP Source Guard} & אין קשירת כתובת לפורט & \enLR{MAC/IP spoofing} \\
\hline
\code{no cdp enable} על פורטי קצה & \enLR{CDP} משדר דגם, \enLR{IOS, IP} & \enLR{CDP reconnaissance} \\
\hline
\end{xltabular}

\sectionheading{\enLR{6.3 MAC Table Overflow} ו-\enLR{Port Security}}

מתג לומד כתובות \enLR{MAC} ורושם אותן בטבלת \enLR{CAM (Content Addressable Memory)} יחד עם הפורט. הטבלה מוגבלת (למשל \enLR{8,000} רשומות). כלי כמו \code{macof} מציף את המתג באלפי כתובות \enLR{MAC} מזויפות בשנייה; הטבלה מתמלאת, והמתג עובר ל\textbf{\enLR{fail-open}} – מתחיל לשדר \emph{כל} מסגרת לכל הפורטים (כמו \enLR{hub).} עכשיו התוקף רואה את תעבורת כולם. זהו "\enLR{MAC flooding".}\par

\subheading{\enLR{Port Security} – ההגנה}

\begin{codebox}\begin{english}\codefont\footnotesize\color{cfg}\setlength{\parindent}{0pt}
S1(config)\# interface f0/1\\
S1(config-if)\# \textcolor{cbold}{\textbf{switchport mode access}}              \textcolor{ccom}{\texthebrew{! חובה – \enLR{port security} לא עובד על \enLR{trunk} דינמי}}\\
S1(config-if)\# \textcolor{cbold}{\textbf{switchport port-security}}\\
S1(config-if)\# \textcolor{cbold}{\textbf{switchport port-security maximum 2}}    \textcolor{ccom}{\texthebrew{! עד \enLR{2 MAC} (מחשב + טלפון \enLR{IP)}}}\\
S1(config-if)\# \textcolor{cbold}{\textbf{switchport port-security mac-address sticky}}  \textcolor{ccom}{\texthebrew{! לומד את הכתובת ושומר בקונפיג}}\\
S1(config-if)\# \textcolor{cbold}{\textbf{switchport port-security violation shutdown}}\\
S1(config-if)\# switchport port-security aging time 60\\
\textcolor{ccom}{\texthebrew{! וריפיקציה}}\\
S1\# \textcolor{cbold}{\textbf{show port-security}}\\
S1\# \textcolor{cbold}{\textbf{show port-security interface f0/1}}\\
S1\# show port-security address
\end{english}\end{codebox}

\subheading{שלושה מצבי הפרה \enLR{(violation)}}

\begin{xltabular}{\linewidth}{|>{\setRTL\raggedleft\arraybackslash}X|>{\setRTL\raggedleft\arraybackslash}X|>{\setRTL\raggedleft\arraybackslash}X|>{\setRTL\raggedleft\arraybackslash}X|}
\hline
\rowcolor{hdrbg}\textbf{מונה} & \textbf{התראה} & \textbf{מה קורה בהפרה} & \textbf{מצב} \\
\hline
\endhead
לא & לא & מפיל מסגרות מ-\enLR{MAC} לא מוכרות; הפורט נשאר \enLR{up} & \textbf{\enLR{protect}} \\
\hline
כן & כן \enLR{(syslog/SNMP)} & מפיל מסגרות; הפורט \enLR{up} & \textbf{\enLR{restrict}} \\
\hline
כן & כן & הפורט עובר ל-\textbf{\enLR{err-disabled}} – מושבת לגמרי & \textbf{\enLR{shutdown}} (ברירת מחדל) \\
\hline
\end{xltabular}

\begin{faqbox}\boxtitle{שאלת תלמיד: "הפורט נכבה \enLR{(err-disabled).} איך מחזירים אותו?"}\par  ידנית: \code{shutdown} ואז \code{no shutdown} על הממשק. אוטומטית אחרי זמן: \code{errdisable recovery cause psecure-violation} + \code{errdisable recovery interval 300}. בדקו סיבה: \code{show interfaces status err-disabled}. הערה: \code{sticky} נדבק ל-\enLR{running-config} – אל תשכחו \code{write} אחרת הכתובת תיעלם ב-\enLR{reload.}\end{faqbox}

\begin{mistakebox}\boxtitle{טעות נפוצה}\par  מגדירים \enLR{port security} ושוכחים \code{switchport mode access}. אם הפורט במצב \enLR{dynamic, port security} לא נכנס לתוקף כראוי. וגם: \enLR{maximum 1} עם טלפון \enLR{IP} + מחשב מאחוריו ⇒ הפרות מיידיות; להגדיר \enLR{maximum 2} או להשתמש ב-\enLR{voice VLAN.}\end{mistakebox}

\sectionheading{\enLR{6.4 VLAN} – והתקפת \enLR{VLAN Hopping}}

\textbf{\enLR{VLAN}} מפצל מתג פיזי אחד לכמה רשתות לוגיות מבודדות. מעבר בין \enLR{VLAN}ים דורש ניתוב \enLR{(router / L3 switch).} קישור בין מתגים שנושא כמה \enLR{VLAN}ים נקרא \textbf{\enLR{trunk}} ומשתמש בתיוג \textbf{\enLR{802.1Q}} (מוסיף תג של \enLR{4} בתים עם מספר ה-\enLR{VLAN} לכל מסגרת). \textbf{\enLR{Native VLAN}} – ה-\enLR{VLAN} היחיד שעובר ב-\enLR{trunk} \emph{בלי} תג (ברירת מחדל \enLR{VLAN 1).}\par

\subheading{התקפה א': \enLR{Switch Spoofing} (ניצול \enLR{DTP)}}

\textbf{\enLR{DTP}} \enLR{(Dynamic Trunking Protocol)} מאפשר לשני מתגים "לסכם" אוטומטית אם הקישור ביניהם יהיה \enLR{trunk.} הבעיה: פורט קצה במצב ברירת מחדל (\code{dynamic auto} / \code{dynamic desirable}) יסכים להפוך ל-\enLR{trunk} גם מול \emph{מחשב של תוקף} שמתחזה למתג ושולח \enLR{DTP.} ברגע שהקישור \enLR{trunk} – התוקף מקבל גישה לכל ה-\enLR{VLAN}ים. זו התשובה בבגרות: "השבתת \enLR{DTP} מונעת \enLR{VLAN hopping".}\par

\subheading{התקפה ב': \enLR{Double Tagging}}

התוקף (ב-\enLR{VLAN} של ה-\enLR{native)} שולח מסגרת עם \emph{שני} תגים: החיצוני = \enLR{native VLAN} (המתג הראשון מסיר אותו), הפנימי = ה-\enLR{VLAN} של הקורבן. המתג השני רואה את התג הפנימי ומעביר ל-\enLR{VLAN} הקורבן. חד-כיווני, אך מסוכן. הגנה: \enLR{native VLAN} שאינו בשימוש ואינו \enLR{VLAN 1.}\par

\subheading{ההגנה על פורטים}

\begin{codebox}\begin{english}\codefont\footnotesize\color{cfg}\setlength{\parindent}{0pt}
\textcolor{ccom}{\texthebrew{! פורט קצה (למחשב) – מפורש \enLR{access}, כיבוי \enLR{DTP}}}\\
S1(config)\# interface range f0/1 - 20\\
S1(config-if-range)\# \textcolor{cbold}{\textbf{switchport mode access}}\\
S1(config-if-range)\# \textcolor{cbold}{\textbf{switchport access vlan 10}}\\
S1(config-if-range)\# \textcolor{cbold}{\textbf{switchport nonegotiate}}          \textcolor{ccom}{\texthebrew{! לא לשלוח \enLR{DTP}}}\\
S1(config-if-range)\# \textcolor{cbold}{\textbf{spanning-tree portfast}}\\
S1(config-if-range)\# \textcolor{cbold}{\textbf{spanning-tree bpduguard enable}}\\
\textcolor{ccom}{\texthebrew{! פורט \enLR{trunk} – מפורש \enLR{trunk, native} ייעודי}}\\
S1(config)\# interface g0/1\\
S1(config-if)\# \textcolor{cbold}{\textbf{switchport mode trunk}}\\
S1(config-if)\# \textcolor{cbold}{\textbf{switchport nonegotiate}}\\
S1(config-if)\# \textcolor{cbold}{\textbf{switchport trunk native vlan 999}}       \textcolor{ccom}{\texthebrew{! \enLR{VLAN "}חור שחור", לא בשימוש}}\\
S1(config-if)\# switchport trunk allowed vlan 10,20,30      \textcolor{ccom}{\texthebrew{! רק מה שצריך}}\\
\textcolor{ccom}{\texthebrew{! פורטים לא בשימוש – לכבות ולשים ב-\enLR{VLAN} מבודד}}\\
S1(config)\# interface range f0/21 - 24\\
S1(config-if-range)\# switchport access vlan 999\\
S1(config-if-range)\# shutdown
\end{english}\end{codebox}

\sectionheading{\enLR{6.5 STP} והתקפות עליו}

\textbf{\enLR{STP}} \enLR{(Spanning Tree Protocol, 802.1D)} מונע \textbf{לולאות} בשכבה \enLR{2.} למה לולאות מסוכנות? למסגרת שכבה \enLR{2} \textbf{אין \enLR{TTL}} (בניגוד ל-\enLR{IP).} אם יש לולאה פיזית בין מתגים, מסגרת \enLR{broadcast} תסתובב \emph{לנצח}, תשוכפל בכל סיבוב, ותייצר \textbf{\enLR{broadcast storm}} שמשתק את הרשת תוך שניות. \enLR{STP} חוסם באופן לוגי פורטים כדי להשאיר מסלול יחיד, ופותח אותם אם קישור נופל.\par

\begin{exambox}\boxtitle{בבגרות (שאלה \enLR{3}ו)}\par  "פרוטוקול \enLR{STP} מונע לולאות מיתוג במודל ה-\enLR{OSI} בשכבה מספר \_\_\_" – \textbf{\enLR{2}} \enLR{(Data Link). "}הפעולה של מניעת הלולאות נעשית על ידי \_\_\_" – \textbf{\enLR{Port blocking} (חסימת יציאות)}.\end{exambox}

\subheading{איך נבחר ה-\enLR{Root Bridge} (נשאל בבגרות!)}

המתג עם \textbf{\enLR{Bridge ID} הנמוך ביותר} הופך ל-\enLR{root. Bridge ID} = \textbf{\enLR{Priority (2} בתים) + \enLR{MAC address}}. ברירת מחדל \enLR{priority = 32768} לכולם. אם ה-\enLR{priority} שווה – מכריע ה-\textbf{\enLR{MAC} הנמוך ביותר}.\par

\begin{exambox}\boxtitle{בבגרות (שאלה \enLR{3}ז)}\par  \enLR{4} מתגים, \enLR{priority} ברירת מחדל זהה. \enLR{MAC}ים: \enLR{SW1=0C:0E:15:22:05:97, SW2=0C:E0:38:00:36:75, SW3=0C:E0:18:A1:B3:19, SW4=0C:0E:15:1A:3C:9D.} מי ה-\enLR{root}? משווים בית-בית: \enLR{SW1} ו-\enLR{SW4} מתחילים ב-\enLR{0C:0E:15}, השאר ב-\enLR{0C:E0.} הבית הרביעי: \enLR{SW1=22, SW4=1A. 1A} < \enLR{22} ⇒ \textbf{\enLR{SW4} הוא ה-\enLR{root bridge}}. הסיבה: "כי כתובת ה-\enLR{MAC} שלו היא הנמוכה ביותר" (וה-\enLR{priority} זהה).\end{exambox}

\subheading{\enLR{OSPF DR/BDR} – רקע לשאלה \enLR{2}ח (שכבה \enLR{3}, אך נשאל)}

בבחירת \enLR{DR/BDR} ב-\enLR{OSPF}: קודם \textbf{\enLR{priority} הגבוה ביותר}, ואם שווה – \textbf{\enLR{Router-ID} הגבוה ביותר}. בשאלת הבגרות \enLR{R1(pri 2, RID 1.1.1.1), R2(pri 1), R3(pri 2, RID 3.3.3.3), R4(pri 1).} בין בעלי \enLR{priority 2: R3 (RID 3.3.3.3) > R1} ⇒ \textbf{\enLR{R3 = DR}}. ה-\enLR{BDR} הוא הבא: בין \enLR{priority 2, R1}; אבל אחרי \enLR{DR} משווים את השאר – \enLR{R1 (RID 1.1.1.1)} הוא הגבוה מבין הנותרים בעלי \enLR{priority 2.} תשובה: \textbf{\enLR{R3 = DR, R1 = BDR}}. (שימו לב: ההיגיון הפוך מ-\enLR{STP} – שם נמוך מנצח, כאן גבוה.)\par

\subheading{הגנות \enLR{STP}}

\begin{xltabular}{\linewidth}{|>{\setRTL\raggedleft\arraybackslash}X|>{\setRTL\raggedleft\arraybackslash}X|>{\setRTL\raggedleft\arraybackslash}X|}
\hline
\rowcolor{hdrbg}\textbf{איפה} & \textbf{מה עושה} & \textbf{מנגנון} \\
\hline
\endhead
פורטי \enLR{access} למחשבים & פורט קצה עובר מיד ל-\enLR{forwarding} (לא מחכה \enLR{30} שניות) & \textbf{\enLR{PortFast}} \\
\hline
פורטי קצה & אם פורט \enLR{PortFast} מקבל \enLR{BPDU} (כלומר חובר אליו מתג/תוקף) – \enLR{err-disable} מיידי & \textbf{\enLR{BPDU Guard}} \\
\hline
פורטים למתגי גישה מתחת & מונע ממתג בפורט זה להפוך ל-\enLR{root} (אם ישלח \enLR{BPDU} עדיף – הפורט נחסם) & \textbf{\enLR{Root Guard}} \\
\hline
קישורים \enLR{redundant} & מגן מלולאה כשפורט מפסיק לקבל \enLR{BPDU} & \textbf{\enLR{Loop Guard}} \\
\hline
\end{xltabular}

\begin{codebox}\begin{english}\codefont\footnotesize\color{cfg}\setlength{\parindent}{0pt}
\textcolor{ccom}{\texthebrew{! פר-ממשק}}\\
S1(config-if)\# spanning-tree portfast\\
S1(config-if)\# \textcolor{cbold}{\textbf{spanning-tree bpduguard enable}}\\
\textcolor{ccom}{\texthebrew{! גלובלי – על כל פורטי \enLR{PortFast} (זו התשובה בבגרות שאלה \enLR{7}ד)}}\\
S1(config)\# \textcolor{cbold}{\textbf{spanning-tree portfast default}}\\
S1(config)\# \textcolor{cbold}{\textbf{spanning-tree portfast bpduguard default}}\\
\textcolor{ccom}{\texthebrew{! שהמתג שלנו יהיה \enLR{root} בוודאות}}\\
S1(config)\# \textcolor{cbold}{\textbf{spanning-tree vlan 10 root primary}}       \textcolor{ccom}{\texthebrew{! מוריד \enLR{priority} ל-\enLR{24576}}}\\
S1(config)\# spanning-tree vlan 10 priority 4096            \textcolor{ccom}{\texthebrew{! ידני – חייב כפולה של \enLR{4096}}}\\
S1\# \textcolor{cbold}{\textbf{show spanning-tree}}
\end{english}\end{codebox}

\sectionheading{\enLR{6.6 DHCP} – \enLR{Starvation, Spoofing} ו-\enLR{DHCP Snooping}}

\textbf{\enLR{DHCP starvation}:} תוקף מבקש את כל הכתובות ב-\enLR{pool} (עם \enLR{MAC}ים מזויפים) – משתמשים אמיתיים לא מקבלים כתובת \enLR{(DoS).} \textbf{\enLR{DHCP spoofing (rogue server)}:} התוקף מקים שרת \enLR{DHCP} משלו שעונה מהר יותר, ומחלק לקורבנות \textbf{\enLR{default gateway} = הכתובת שלו} ו-\textbf{\enLR{DNS} שלו} ⇒ \enLR{MITM} על כל התעבורה.\par

\subheading{\enLR{DHCP Snooping} – ההגנה}

המתג מסמן פורטים כ-\textbf{\enLR{trusted}} (לכיוון שרת ה-\enLR{DHCP} הלגיטימי / \enLR{uplink)} או \textbf{\enLR{untrusted}} (פורטי קצה, ברירת מחדל). הודעות שרת \enLR{(OFFER, ACK)} מפורט \enLR{untrusted} – נחסמות. המתג בונה \textbf{\enLR{DHCP Snooping Binding Table}} \enLR{(MAC}↔\enLR{IP}↔פורט↔\enLR{VLAN)} – בסיס ל-\enLR{DAI} ו-\enLR{IP Source Guard.}\par

\begin{codebox}\begin{english}\codefont\footnotesize\color{cfg}\setlength{\parindent}{0pt}
S1(config)\# \textcolor{cbold}{\textbf{ip dhcp snooping}}\\
S1(config)\# \textcolor{cbold}{\textbf{ip dhcp snooping vlan 10,20}}\\
S1(config)\# interface g0/1                    \textcolor{ccom}{\texthebrew{! לכיוון שרת \enLR{DHCP} הלגיטימי}}\\
S1(config-if)\# \textcolor{cbold}{\textbf{ip dhcp snooping trust}}\\
S1(config)\# interface range f0/1 - 20         \textcolor{ccom}{\texthebrew{! פורטי קצה}}\\
S1(config-if-range)\# \textcolor{cbold}{\textbf{ip dhcp snooping limit rate 6}}   \textcolor{ccom}{\texthebrew{! נגד \enLR{starvation}}}\\
S1\# \textcolor{cbold}{\textbf{show ip dhcp snooping binding}}
\end{english}\end{codebox}

\sectionheading{\enLR{6.7 ARP Spoofing} ו-\enLR{Dynamic ARP Inspection (DAI)}}

\textbf{\enLR{ARP}} ממפה \enLR{IP} ל-\enLR{MAC}, וחסר כל אימות: כל מחשב יכול לשלוח "\enLR{Gratuitous ARP"} שאומר "אני ה-\enLR{gateway"} – והקורבנות יעדכנו את הטבלה שלהם ויתחילו לשלוח את התעבורה לתוקף \enLR{(MITM).} \textbf{\enLR{DAI}} בודק כל הודעת \enLR{ARP} מול טבלת ה-\enLR{DHCP Snooping}: אם ה-\enLR{IP}↔\enLR{MAC} לא תואם לרשומה – ההודעה נזרקת.\par

\begin{codebox}\begin{english}\codefont\footnotesize\color{cfg}\setlength{\parindent}{0pt}
S1(config)\# \textcolor{cbold}{\textbf{ip arp inspection vlan 10,20}}\\
S1(config)\# interface g0/1\\
S1(config-if)\# \textcolor{cbold}{\textbf{ip arp inspection trust}}            \textcolor{ccom}{\texthebrew{! \enLR{uplink} מהימן}}\\
\textcolor{ccom}{\texthebrew{! לפורטים סטטיים בלי \enLR{DHCP} – מגדירים \enLR{ACL} של \enLR{ARP} ידני}}\\
S1(config)\# arp access-list STATIC-HOSTS\\
S1(config-arp-nacl)\# permit ip host 10.0.0.5 mac host aaaa.bbbb.cccc\\
S1\# \textcolor{cbold}{\textbf{show ip arp inspection}}
\end{english}\end{codebox}

המשלים: \textbf{\enLR{IP Source Guard}} – מוודא שכתובת ה-\enLR{IP} במסגרת תואמת לפורט לפי טבלת ה-\enLR{snooping} (נגד \enLR{IP spoofing).}\par

\sectionheading{\enLR{6.8 Storm Control}}

מגביל את אחוז התעבורה מסוג \enLR{broadcast / multicast / unknown-unicast} על פורט; אם עוברים סף – המתג מפיל את העודף (ומתריע). מגן מ-\enLR{broadcast storms} (גם מלולאה וגם מתקלה/התקפה). התכנית מציינת "\enLR{Storm Control (SC)".}\par

\begin{codebox}\begin{english}\codefont\footnotesize\color{cfg}\setlength{\parindent}{0pt}
S1(config-if)\# \textcolor{cbold}{\textbf{storm-control broadcast level 5.00}}       \textcolor{ccom}{\texthebrew{! מעל \enLR{5\%} מרוחב הפס – הגבל}}\\
S1(config-if)\# \textcolor{cbold}{\textbf{storm-control multicast level pps 1k}}\\
S1(config-if)\# \textcolor{cbold}{\textbf{storm-control action shutdown}}            \textcolor{ccom}{\texthebrew{! או \enLR{trap}}}\\
S1\# \textcolor{cbold}{\textbf{show storm-control}}
\end{english}\end{codebox}

\sectionheading{\enLR{6.9 SPAN} – שיקוף פורטים (הבסיס ל-\enLR{IDS)}}

\textbf{\enLR{SPAN}} \enLR{(Switched Port Analyzer, "port mirroring")} מעתיק את כל התעבורה מפורט/\enLR{VLAN} מקור לפורט יעד – שאליו מחובר \enLR{Wireshark} או \textbf{\enLR{IDS}}. זה החיבור לפרק \enLR{5: IDS "}מחוץ לנתיב" מקבל את התעבורה דווקא דרך \enLR{SPAN.} \textbf{\enLR{RSPAN}} – שיקוף בין מתגים דרך \enLR{VLAN} ייעודי; \textbf{\enLR{ERSPAN}} – מעל \enLR{IP (GRE).}\par

\begin{codebox}\begin{english}\codefont\footnotesize\color{cfg}\setlength{\parindent}{0pt}
S1(config)\# \textcolor{cbold}{\textbf{monitor session 1 source interface f0/1 - 10 both}}\\
S1(config)\# \textcolor{cbold}{\textbf{monitor session 1 destination interface f0/24}}   \textcolor{ccom}{\texthebrew{! כאן ה-\enLR{IDS/Wireshark}}}\\
S1\# \textcolor{cbold}{\textbf{show monitor session 1}}
\end{english}\end{codebox}

\begin{faqbox}\boxtitle{שאלת תלמיד: "אם ה-\enLR{IDS} רק מקבל עותק דרך \enLR{SPAN}, למה שלא נחבר אותו \enLR{inline} ונחסום?"}\par  כי אז הוא \enLR{IPS}, לא \enLR{IDS} – והוא מוסיף השהיה ונקודת כשל בנתיב. \enLR{SPAN} מאפשר לנטר בלי לגעת בזרימה: הפורט המשקף יכול אפילו ליפול בלי להשפיע על הרשת. זה בדיוק ה-\enLR{trade-off} מפרק \enLR{5.}\end{faqbox}

\sectionheading{\enLR{6.10 NAC} ו-\enLR{802.1X}}

\textbf{\enLR{NAC}} \enLR{(Network Access Control)} – "בקרת בריאות": לפני שמחשב מקבל גישה מלאה, בודקים שהוא עומד בתנאים (אנטי-וירוס מעודכן, \enLR{patch}, אין תוכנות אסורות). מחשב שנכשל מועבר ל-\enLR{VLAN "}הסגר" \enLR{(quarantine)} לתיקון. הבסיס הטכני: \textbf{\enLR{802.1X}} (פרק \enLR{3)} – המתג/\enLR{AP} חוסם את הפורט עד אימות מול \enLR{RADIUS.} שילוב: \enLR{802.1X} מאמת \emph{מי}, \enLR{NAC} בודק \emph{באיזה מצב} המכשיר.\par

\sectionheading{\enLR{6.11} מערכות קצה: \enLR{IronPort, CSA}, ומיילים/אתרים}

\begin{itemize}
\item \textbf{\enLR{Cisco IronPort}} – משפחת מכשירים ל\textbf{אבטחת תוכן} בקצה: \textbf{\enLR{ESA}} \enLR{(Email Security Appliance)} – סינון \enLR{spam, phishing} ווירוסים במייל; \textbf{\enLR{WSA}} \enLR{(Web Security Appliance)} – \enLR{proxy} שמסנן אתרים זדוניים, \enLR{URL filtering}, בדיקת הורדות. סיסקו רכשה את \enLR{IronPort} ב-\enLR{2007.}
\item \textbf{\enLR{CSA}} \enLR{(Cisco Security Agent)} – \enLR{HIPS} על התחנה (פרק \enLR{5).}
\item היום התחום נקרא \textbf{\enLR{Secure Email / Secure Web / SWG / CASB}}, ולרוב בענן \enLR{(Cisco Umbrella).}
\end{itemize}

\sectionheading{\enLR{6.12} אבטחת \enLR{VoIP} ו-\enLR{SAN}}

\subheading{\enLR{VoIP} (טלפוניה על \enLR{IP)}}

איומים: האזנה לשיחות \enLR{(sniffing)}, \textbf{\enLR{toll fraud}} (שימוש לא מורשה בקווים לחיוב יקר), התחזות \enLR{(caller-ID spoofing), DoS} על ה-\enLR{PBX, SPIT (spam} קולי). הגנות: \textbf{\enLR{Voice VLAN} נפרד} (הפרדה מתעבורת הנתונים), הצפנה (\textbf{\enLR{SRTP}} למדיה, \textbf{\enLR{TLS/SIP-TLS}} לאיתות), אימות מכשירי טלפון, \enLR{ACL} בין \enLR{voice} ל-\enLR{data.}\par

\begin{codebox}\begin{english}\codefont\footnotesize\color{cfg}\setlength{\parindent}{0pt}
S1(config-if)\# switchport access vlan 10        \textcolor{ccom}{\texthebrew{! נתונים}}\\
S1(config-if)\# \textcolor{cbold}{\textbf{switchport voice vlan 20}}       \textcolor{ccom}{\texthebrew{! קול – מתויג בנפרד}}
\end{english}\end{codebox}

\subheading{\enLR{SAN} (רשת אחסון)}

\enLR{SAN} מחברת שרתים לאחסון בלוקים \enLR{(Fibre Channel / iSCSI).} איומים: גישה לא מורשית ל-\enLR{LUN, WWN spoofing.} הגנות: \textbf{\enLR{Zoning}} (מי מדבר עם מי ב-\enLR{fabric} – מקביל ל-\enLR{VLAN)}, \textbf{\enLR{LUN masking}} (איזה שרת רואה איזה נפח), \textbf{\enLR{VSAN}} (בידוד לוגי), אימות \enLR{(FC-SP / CHAP} ל-\enLR{iSCSI)}, והצפנת נתונים במנוחה. חשיבות: כל הנתונים הקריטיים של הארגון נמצאים שם.\par

\sectionheading{\enLR{6.13 Wi-Fi} ו-\enLR{Evil Twin}}

\textbf{\enLR{Evil Twin}:} התוקף מקים \enLR{Access Point} עם \textbf{אותו \enLR{SSID}} כמו הרשת החוקית, בעוצמה חזקה יותר. מכשירים מתחברים אליו אוטומטית (הם זוכרים את השם), והתוקף מבצע \enLR{MITM} – רואה הכול, מציג דף התחברות מזויף \enLR{(captive portal)} לגניבת סיסמה. הגנות: \enLR{WPA2/WPA3-Enterprise (802.1X} – המכשיר מאמת את השרת, לא רק להיפך), \enLR{WIPS (Wireless IPS} שמזהה \enLR{AP} מתחזים), אימות הדדי.\par

\begin{exambox}\boxtitle{בבגרות (שאלה \enLR{6}ה)}\par  "מהם מתכנני התקפת \enLR{Evil Twin}?" – \textbf{לשתול נקודת גישה זדונית, בעלת \enLR{SSID} זהה לנקודת גישה חוקית, לצורך גניבת מידע.}\end{exambox}

\sectionheading{\enLR{6.14} מדריך פקודות מרוכז – פרק \enLR{6}}

\begin{xltabular}{\linewidth}{|>{\setRTL\raggedleft\arraybackslash}X|>{\setRTL\raggedleft\arraybackslash}X|}
\hline
\rowcolor{hdrbg}\textbf{הגנה מפני} & \textbf{פקודה} \\
\hline
\endhead
\enLR{MAC overflow} & \code{switchport port-security} + \code{maximum} + \code{violation} + \code{mac-address sticky} \\
\hline
\enLR{VLAN hopping (DTP)} & \code{switchport mode access} + \code{switchport nonegotiate} \\
\hline
\enLR{double tagging} & \code{switchport trunk native vlan 999} \\
\hline
\enLR{STP manipulation} & \code{spanning-tree portfast} + \code{bpduguard enable} / \code{… default} \\
\hline
קיבוע \enLR{root} & \code{spanning-tree vlan N root primary} / \code{priority} \\
\hline
\enLR{DHCP spoofing/starvation} & \code{ip dhcp snooping} + \code{trust} + \code{limit rate} \\
\hline
\enLR{ARP poisoning} & \code{ip arp inspection vlan} + \code{trust} \\
\hline
\enLR{broadcast storm} & \code{storm-control broadcast level} \\
\hline
\enLR{SPAN} ל-\enLR{IDS} & \code{monitor session} \\
\hline
הפרדת \enLR{VoIP} & \code{switchport voice vlan} \\
\hline
וריפיקציה & \code{show port-security}, \code{show spanning-tree}, \code{show ip dhcp snooping binding} \\
\hline
\end{xltabular}

\sectionheading{\enLR{6.15} תרגילים לתלמידים}

\begin{exbox}\boxtitle{תרגיל \enLR{1} – התאמת התקפה להגנה}\par  חברו: \enLR{(1) MAC flooding (2) rogue DHCP (3) ARP poisoning (4)} מחשב שמתחזה למתג ומקבל את כל ה-\enLR{VLAN}ים \enLR{(5)} מתג מזויף שמנסה להיות \enLR{root.} הגנות: \enLR{DAI / port security / BPDU Guard+Root Guard / switchport nonegotiate / DHCP snooping.} \par\vskip2pt\hrule\vskip3pt\textbf{}\enLR{1}→\enLR{port security} · \enLR{2}→\enLR{DHCP snooping} · \enLR{3}→\enLR{DAI} · \enLR{4}→\enLR{switchport nonegotiate (access)} · \enLR{5}→\enLR{BPDU Guard/Root Guard}\end{exbox}

\begin{exbox}\boxtitle{תרגיל \enLR{2} – מי ה-\enLR{root}?}\par  \enLR{4} מתגים, \enLR{priority: SW1=32768, SW2=32768, SW3}=\textbf{\enLR{4096}}, \enLR{SW4=32768. MAC}ים עולים בסדר \enLR{SW1}\par\vskip2pt\hrule\vskip3pt\textbf{}\enLR{SW3} – ה-\enLR{priority} נבדק ראשון והוא הנמוך \enLR{(4096).} ה-\enLR{MAC} נבדק רק בתיקו. לו כולם היו \enLR{32768} – \enLR{SW1 (MAC} נמוך).\end{exbox}

\begin{exbox}\boxtitle{תרגיל \enLR{3} – ניתוח קונפיג}\par  פורט \enLR{f0/5} הוגדר: \code{switchport port-security}, \code{maximum 1}, \code{violation shutdown}. עובד חיבר אליו מתג קטן \enLR{(switch)} עם \enLR{3} מחשבים. מה יקרה, ומה ההודעה בלוג? \par\vskip2pt\hrule\vskip3pt\textbf{}המתג הקטן יעביר מסגרות מ-\enLR{3 MAC}ים ⇒ בהפרה הראשונה \enLR{(MAC} שני) הפורט עובר ל-\enLR{err-disabled} ונכבה. בלוג: \code{\%PM-4-ERR\_DISABLE: psecure-violation ... Fa0/5}. שחזור: \enLR{shut/no shut.}\end{exbox}

\begin{exbox}\boxtitle{תרגיל \enLR{4} – תכנון הקשחת מתג גישה}\par  כתבו \enLR{8} פקודות הקשחה שתחילו על מתג גישה של קומה בבניין \enLR{(24} פורטים למשתמשים, \enLR{uplink} אחד). \par\vskip2pt\hrule\vskip3pt\textbf{}דוגמה: על פורטי הקצה: \enLR{mode access, access vlan 10, nonegotiate, port-security max 2 sticky violation restrict, portfast, bpduguard enable, ip dhcp snooping (untrusted, limit rate), no cdp enable.} על ה-\enLR{uplink: mode trunk, nonegotiate, native vlan 999, allowed vlan list, dhcp snooping trust, arp inspection trust.} גלובלי: \enLR{ip dhcp snooping + vlan, ip arp inspection vlan, portfast bpduguard default.} כיבוי פורטים לא בשימוש ל-\enLR{vlan 999 + shutdown.}\end{exbox}

\sectionheading{\enLR{6.16} שאלות בסגנון בגרות}

\begin{exambox}\boxtitle{שאלה \enLR{1}}\par  מה הדרך למנוע מתקפת \enLR{MAC Table overflow} במתג? \par\vskip2pt\hrule\vskip3pt\textbf{}\enLR{Port security} על כל הפורטים שאינם \enLR{trunk}, עם הגבלת מספר כתובות \enLR{MAC.}\end{exambox}

\begin{exambox}\boxtitle{שאלה \enLR{2}}\par  איזו התקפת שכבה \enLR{2} אפשר לסכל על ידי השבתת פרוטוקול \enLR{DTP}? \enLR{1. DHCP spoofing 2. ARP spoofing 3. VLAN hopping 4. ARP poisoning} \par\vskip2pt\hrule\vskip3pt\textbf{}\textbf{\enLR{3. VLAN hopping}} \enLR{(switch spoofing).}\end{exambox}

\begin{exambox}\boxtitle{שאלה \enLR{3}}\par  באיזו פקודה מפעילים \enLR{BPDU Guard} על כל פורטי ה-\enLR{PortFast} במתג? \par\vskip2pt\hrule\vskip3pt\textbf{}\code{spanning-tree portfast bpduguard default} (במצב \enLR{config} גלובלי).\end{exambox}

\begin{exambox}\boxtitle{שאלה \enLR{4}}\par  \enLR{STP} פועל בשכבה \_\_\_ של \enLR{OSI} ומונע \_\_\_ על ידי \_\_\_. \par\vskip2pt\hrule\vskip3pt\textbf{}\enLR{2 (Data Link)} · לולאות \enLR{(broadcast storm)} · חסימת פורטים \enLR{(port blocking).}\end{exambox}

\begin{exambox}\boxtitle{שאלה \enLR{5}}\par  מהי מתקפת \enLR{Evil Twin} ואיזו טכנולוגיית \enLR{Wi-Fi} מסייעת למנוע אותה? \par\vskip2pt\hrule\vskip3pt\textbf{}נקודת גישה זדונית עם \enLR{SSID} זהה לחוקית לצורך \enLR{MITM}/גניבת מידע. מניעה: \enLR{WPA2/WPA3-Enterprise} עם \enLR{802.1X} (אימות הדדי – המכשיר מוודא את זהות השרת) ו-\enLR{WIPS.}\end{exambox}

\begin{exambox}\boxtitle{שאלה \enLR{6}}\par  מהו תפקיד \enLR{DHCP Snooping} ואיזה מנגנון אבטחה נוסף מסתמך על הטבלה שהוא בונה? \par\vskip2pt\hrule\vskip3pt\textbf{}חוסם הודעות שרת \enLR{DHCP} מפורטים לא-מהימנים \enLR{(rogue server)} ומגביל קצב \enLR{(starvation). Dynamic ARP Inspection} (ו-\enLR{IP Source Guard)} מסתמכים על טבלת ה-\enLR{binding} שלו.\end{exambox}

\sectionheading{\enLR{6.17} מעבדה \enLR{(2} שעות) – הקשחת מתג ב-\enLR{Packet Tracer}}

\begin{enumerate}
\item \enLR{Port security: 2} מחשבים על פורט אחד (דרך \enLR{hub}/מתג קטן), \enLR{maximum 1, violation shutdown.} הוסיפו מחשב שני ⇒ \enLR{err-disabled.} שחזרו עם \enLR{shut/no shut} והעלו ל-\enLR{maximum 2.}
\item \enLR{BPDU Guard}: הגדירו \enLR{portfast+bpduguard} על פורט קצה; חברו אליו מתג נוסף ⇒ \enLR{err-disabled.}
\item \enLR{Root bridge: 3} מתגים בטבעת; הגדירו את המרכזי כ-\code{root primary}; \code{show spanning-tree} – ודאו מי \enLR{root} ואיזה פורט \enLR{blocking.}
\item \enLR{DTP}: הגדירו \enLR{trunk} מפורש + \enLR{nonegotiate} בין המתגים; פורטי קצה \enLR{access + nonegotiate.}
\item \enLR{DHCP snooping}: הפעילו, סמנו את פורט השרת \enLR{trusted}; הוסיפו "שרת \enLR{DHCP"} תוקף על פורט \enLR{untrusted} וראו שההצעות שלו נחסמות.
\end{enumerate}

\sectionheading{בנק שאלות שתלמידים שואלים – ותשובות מוכנות}

שאלות נפוצות בפרק אבטחת הרשת המקומית, עם תשובות מוכנות.\par

\textbf{ש: "אם יש חומת אש בכניסה לרשת, למה בכלל צריך להגן על המתגים בפנים?"}\newline ת: כי חומת האש מגנה מפני האינטרנט (מבחוץ), אבל אם תוקף כבר \underline{בפנים} (חיבר מחשב לשקע, פרץ ל-\enLR{Wi-Fi}, השתלט על מחשב עובד) – הוא פועל בשכבה \enLR{2}, מתחת לחומת האש. שם צריך \enLR{port security, DHCP snooping} וכו'.\par

\textbf{ש: "הפורט עבר ל-\enLR{err-disabled} ונכבה. איך מחזירים אותו?"}\newline ת: ידנית: \code{shutdown} ואז \code{no shutdown} על הממשק. אוטומטית אחרי זמן: \code{errdisable recovery cause psecure-violation}. בודקים סיבה עם \code{show interfaces status err-disabled}.\par

\textbf{ש: "למה לולאה בין מתגים כל כך מסוכנת?"}\newline ת: כי למסגרת בשכבה \enLR{2} \underline{אין \enLR{TTL}} (בניגוד ל-\enLR{IP).} אם יש לולאה פיזית, הודעת \enLR{broadcast} מסתובבת לנצח ומשוכפלת בכל סיבוב – "סופת \enLR{broadcast"} שמשתקת את הרשת תוך שניות. \enLR{STP} מונע זאת על ידי חסימת פורט.\par

\textbf{ש: "איך בוחרים את ה-\enLR{root bridge} ב-\enLR{STP}?"}\newline ת: לפי ה-\enLR{Bridge ID} הנמוך ביותר = \enLR{Priority} (ברירת מחדל \enLR{32768)} ואז כתובת \enLR{MAC.} אם ה-\enLR{priority} זהה אצל כולם – מנצח ה-\underline{\enLR{MAC} הנמוך}. שימו לב: ב-\enLR{STP} נמוך מנצח, ב-\enLR{OSPF} (פרק רשתות) גבוה מנצח – הפוך!\par

\textbf{ש: "מה זה \enLR{VLAN hopping} ואיך \enLR{DTP} קשור?"}\newline ת: \enLR{DTP} הוא פרוטוקול שמנהל אוטומטית אם קישור יהיה \enLR{trunk.} תוקף יכול להתחזות למתג, "לשכנע" את הפורט להפוך ל-\enLR{trunk}, ואז לקבל גישה לכל ה-\enLR{VLAN}ים. הפתרון: לכבות \enLR{DTP} (\code{switchport nonegotiate}) ולהגדיר פורטי קצה כ-\enLR{access} מפורש.\par

\textbf{ש: "מה זה \enLR{rogue DHCP} ולמה זה מסוכן?"}\newline ת: תוקף מקים שרת \enLR{DHCP} משלו שעונה מהר יותר מהאמיתי, ומחלק לקורבנות \enLR{default gateway} ו-\enLR{DNS} \underline{שלו} – וכך כל התעבורה שלהם עוברת דרכו \enLR{(MITM).} ההגנה: \enLR{DHCP snooping}, שמסמן פורטים "מהימנים" וחוסם הצעות משרתים לא מורשים.\par

\textbf{ש: "מה זה \enLR{Evil Twin} ואיך מתגוננים?"}\newline ת: נקודת גישה זדונית עם \underline{אותו שם רשת \enLR{(SSID)}} כמו החוקית, בעוצמה חזקה יותר – מכשירים מתחברים אליה אוטומטית והתוקף מבצע \enLR{MITM.} הגנה: \enLR{WPA2/WPA3-Enterprise} עם \enLR{802.1X} (אימות הדדי – המכשיר מוודא שהשרת אמיתי) ו-\enLR{WIPS.}\par

\textbf{ש: "מה ההבדל בין \enLR{port security} ל-\enLR{802.1X}?"}\newline ת: \textbf{\enLR{port security}} מגביל \underline{כמה/אילו כתובות \enLR{MAC}} מותרות בפורט (הגנה מפני \enLR{MAC flooding).} \textbf{\enLR{802.1X}} דורש \underline{אימות משתמש/מכשיר} מלא מול שרת \enLR{RADIUS} לפני שנותנים גישה. \enLR{802.1X} חזק בהרבה אך דורש תשתית שרת.\par

\sectionheading{מילון מונחים – הגדרות}

הגדרות תמציתיות של המונחים המרכזיים בפרק. שימושי גם כדף עזר לבחינה.\par

\begin{xltabular}{\linewidth}{|>{\setRTL\raggedleft\arraybackslash}X|>{\setRTL\raggedleft\arraybackslash}X|}
\hline
\rowcolor{hdrbg}\textbf{הגדרה} & \textbf{מונח} \\
\hline
\endhead
שכבת הקישור; עובדת בכתובות \enLR{MAC}, מתגים ומסגרות \enLR{(frames).} & \textbf{שכבה \enLR{2 (Data Link)}} \\
\hline
מכשיר שמחבר מכשירים ברשת מקומית ולומד כתובות \enLR{MAC} לכל פורט. & \textbf{מתג \enLR{(Switch)}} \\
\hline
טבלה במתג הממפה כתובת \enLR{MAC} לפורט; מוגבלת בגודלה. & \textbf{טבלת \enLR{MAC (CAM)}} \\
\hline
הצפת המתג בכתובות \enLR{MAC} מזויפות עד שהוא משדר לכולם (כמו \enLR{hub).} & \textbf{\enLR{MAC flooding / overflow}} \\
\hline
הגבלת מספר/זהות כתובות \enLR{MAC} בפורט; הגנה מפני \enLR{MAC flooding.} & \textbf{\enLR{Port Security}} \\
\hline
תגובת \enLR{port security} להפרה: התעלמות, התראה, או כיבוי הפורט \enLR{(err-disabled).} & \textbf{\enLR{violation (protect/restrict/shutdown)}} \\
\hline
רשת לוגית מבודדת בתוך מתג פיזי; מעבר בין \enLR{VLAN}ים דורש ניתוב. & \textbf{\enLR{VLAN}} \\
\hline
קישור הנושא כמה \enLR{VLAN}ים בין מתגים, עם תיוג של מספר ה-\enLR{VLAN.} & \textbf{\enLR{Trunk (802.1Q)}} \\
\hline
פרוטוקול שמנהל אוטומטית \enLR{trunk}; יש לכבותו כדי למנוע \enLR{VLAN hopping.} & \textbf{\enLR{DTP}} \\
\hline
קבלת גישה ל-\enLR{VLAN}ים אחרים בניצול \enLR{DTP} או תיוג כפול. & \textbf{\enLR{VLAN hopping}} \\
\hline
פרוטוקול שמונע לולאות בשכבה \enLR{2} על ידי חסימת פורטים. & \textbf{\enLR{STP (Spanning Tree)}} \\
\hline
המתג המרכזי ב-\enLR{STP}; נבחר לפי \enLR{Bridge ID} הנמוך \enLR{(priority}, ואז \enLR{MAC).} & \textbf{\enLR{Root Bridge}} \\
\hline
מכבה פורט קצה שמקבל \enLR{BPDU} (חובר אליו מתג/תוקף). & \textbf{\enLR{BPDU Guard}} \\
\hline
דלדול כתובות ה-\enLR{DHCP}, או הקמת שרת \enLR{DHCP} מזויף ל-\enLR{MITM.} & \textbf{\enLR{DHCP starvation / spoofing}} \\
\hline
סימון פורטים מהימנים וחסימת הצעות \enLR{DHCP} מפורטים לא מורשים. & \textbf{\enLR{DHCP Snooping}} \\
\hline
זיוף הודעות \enLR{ARP} כדי להפנות תעבורה לתוקף \enLR{(MITM).} & \textbf{\enLR{ARP spoofing / poisoning}} \\
\hline
בדיקת הודעות \enLR{ARP} מול טבלת ה-\enLR{DHCP snooping}; חוסם זיופים. & \textbf{\enLR{DAI (Dynamic ARP Inspection)}} \\
\hline
הגבלת אחוז תעבורת \enLR{broadcast/multicast} בפורט; מונע סופות. & \textbf{\enLR{Storm Control}} \\
\hline
בקרת גישה שבודקת 'בריאות' מכשיר לפני מתן גישה (מבוסס \enLR{802.1X).} & \textbf{\enLR{NAC}} \\
\hline
נקודת גישה זדונית עם \enLR{SSID} זהה לחוקית, ל-\enLR{MITM} ברשת אלחוטית. & \textbf{\enLR{Evil Twin}} \\
\hline
\end{xltabular}


\clearpage
\phantomsection\addcontentsline{toc}{section}{פרק \enLR{7} – מערכות הצפנה וקריפטולוגיה}

\sloppy
\chapterheading{פרק \enLR{7} – מערכות הצפנה וקריפטולוגיה}{\small\color{subgray}\enLR{11} שעות עיוני + \enLR{3} מעשי · שבועות \enLR{20}–\enLR{23}\par}\vskip4pt

\begin{metabox}\textbf{מטרות:} קריפטולוגיה · הצפנה, קידוד, \enLR{hashing} · סימטרי מול א-סימטרי · חתימה דיגיטלית · \enLR{PKI}\\ \textbf{בבגרות:} \enLR{Diffie-Hellman}, צופן קיסר ו-\enLR{mod 26}, מפתח ציבורי/פרטי, \enLR{AES} מול \enLR{hash} (סודיות)\end{metabox}

\begin{defbox}\boxtitle{לפני שמתחילים – שלושה מושגים שמבלבלים (למי שאין רקע)}\par  כל הפרק בנוי על הבחנה בין שלושה דברים שנראים דומים אבל שונים לגמרי:  \textbf{קידוד \enLR{(Encoding)}} – שינוי \underline{צורת הכתיבה} של מידע לנוחות, \underline{בלי סוד}. כל אחד יכול להחזיר. דוגמה: מורס, \enLR{Base64.} \textbf{זו אינה אבטחה.} \textbf{הצפנה \enLR{(Encryption)}} – ערבוב המידע בעזרת \underline{מפתח סודי}, כך שרק מי שיש לו המפתח יכול לפענח. הפיך \underline{עם המפתח}. מטרה: \textbf{סודיות}. \textbf{גיבוב \enLR{(Hashing)}} – "טביעת אצבע" של המידע: פונקציה \underline{חד-כיוונית} שאי אפשר להפוך. מטרה: לבדוק ש\textbf{המידע לא שונה (שלמות)}.  \textbf{אנלוגיה:} קידוד = לכתוב מילה הפוך (כל אחד קורא). הצפנה = כספת עם מפתח (רק בעל המפתח פותח). גיבוב = לטחון פרי לשייק (אי אפשר להחזיר את הפרי, אבל אותו פרי תמיד נותן אותו שייק). \newline \textbf{מושג מפתח:} "מפתח" \enLR{(Key)} בהצפנה = מספר/מחרוזת סודית שקובעת איך המידע יעורבב. בלי המפתח הנכון, המידע המוצפן נראה כג'יבריש.\end{defbox}

\sectionheading{\enLR{7.1} מושגי יסוד – ולמה קידוד אינו הצפנה}

\begin{xltabular}{\linewidth}{|>{\setRTL\raggedleft\arraybackslash}X|>{\setRTL\raggedleft\arraybackslash}X|}
\hline
מדע כתיבת ההצפנות (יצירת צפנים) & \textbf{קריפטוגרפיה} \\
\hline
מדע שבירת ההצפנות & \textbf{קריפטואנליזה} \\
\hline
שם כולל לשתיהן & \textbf{קריפטולוגיה} \\
\hline
הטקסט הגלוי, לפני הצפנה & \textbf{\enLR{Plaintext}} \\
\hline
הטקסט המוצפן & \textbf{\enLR{Ciphertext}} \\
\hline
שיטת ההצפנה (קיסר, \enLR{AES)} & \textbf{\enLR{Cipher} / אלגוריתם} \\
\hline
הסוד שמשנה את פעולת האלגוריתם & \textbf{\enLR{Key} (מפתח)} \\
\hline
\end{xltabular}

\begin{defbox}\boxtitle{שלושה מושגים שמבלבלים – חשוב מאוד להבחין}\par  \textbf{קידוד \enLR{(Encoding)}} – המרת ייצוג לנוחות/תאימות, \underline{ללא סוד}. \enLR{Base64, ASCII, Morse.} כל אחד יכול לפענח – זו לא אבטחה!\newline  \textbf{הצפנה \enLR{(Encryption)}} – המרה \underline{הפיכה} בעזרת מפתח סודי; רק מי שיש לו המפתח מפענח. מטרה: \textbf{סודיות}.\newline  \textbf{גיבוב \enLR{(Hashing)}} – פונקציה \underline{חד-כיוונית}, ללא מפתח, פלט קבוע-אורך. אי אפשר להפוך. מטרה: \textbf{שלמות}.\end{defbox}

\begin{mistakebox}\boxtitle{טעות נפוצה מסוכנת}\par  תלמידים (ומתכנתים!) חושבים ש-\enLR{Base64} היא "הצפנה". \code{btoa("password")} אינו מאבטח דבר. אם רואים "==" בסוף מחרוזת – זה כמעט תמיד \enLR{Base64} (קידוד), ופענוח לוקח שנייה. השאלה המבחינה: "האם צריך סוד כדי לחזור?" קידוד – לא. הצפנה – כן.\end{mistakebox}

\subheading{עקרון קרקהופס \enLR{(Kerckhoffs)}}

"מערכת צריכה להיות מאובטחת גם אם \textbf{הכול ידוע חוץ מהמפתח}." המשמעות: לא סומכים על סודיות האלגוריתם ("\enLR{security through obscurity"} זה רע), אלא על סודיות המפתח בלבד. לכן \enLR{AES} פורסם לכל העולם – וזו דווקא הסיבה שסומכים עליו: מיליוני מומחים ניסו לשבור ולא הצליחו.\par

\sectionheading{\enLR{7.2} היסטוריה – מקיסר עד אניגמה}

\subheading{צופן קיסר \enLR{(Caesar)} – נשאל בבגרות}

הזזת כל אות במספר קבוע. יוליוס קיסר השתמש בהזזה של \enLR{3: A}→\enLR{D, B}→\enLR{E}… הנוסחה (עבור אנגלית, \enLR{26} אותיות, \enLR{A=0)}:\par

\begin{defbox}\boxtitle{הצפנה: \enLR{C = (P + K) mod 26} · פענוח: \enLR{P = (C} − \enLR{K) mod 26}}\par  דוגמה עם \enLR{K=10}: האות \enLR{H (=7)} ⇒ \enLR{C = (7+10) mod 26 = 17 = R.} האות \enLR{Z (=25)} ⇒ \enLR{C = (25+10) mod 26 = 35 mod 26 = 9 = J.} ה-\enLR{mod "}מגלגל" חזרה להתחלה.\end{defbox}

\begin{exambox}\boxtitle{בבגרות (שאלה \enLR{5}ג)}\par  "בצופן קיסר \enLR{C=P+10 (mod 26).} מה תפקידו של \enLR{mod 26}?" – \textbf{מספר האותיות (האפשרויות) בשפה האנגלית}; ה-\enLR{mod} מבטיח שהתוצאה תישאר בטווח \enLR{0}–\enLR{25} ו"תתגלגל" מ-\enLR{Z} חזרה ל-\enLR{A.} \emph{לא} מספר הפעולות ולא גרסת האלגוריתם.\end{exambox}

\subheading{צפנים היסטוריים נוספים}

\begin{xltabular}{\linewidth}{|>{\setRTL\raggedleft\arraybackslash}X|>{\setRTL\raggedleft\arraybackslash}X|>{\setRTL\raggedleft\arraybackslash}X|>{\setRTL\raggedleft\arraybackslash}X|}
\hline
\rowcolor{hdrbg}\textbf{חולשה} & \textbf{רעיון} & \textbf{סוג} & \textbf{צופן} \\
\hline
\endhead
\enLR{26} מפתחות בלבד; ניתוח תדירות & הזזה קבועה & החלפה \enLR{(substitution)}, חד-אלפביתי & \enLR{Caesar / ROT13} \\
\hline
ניתוח תדירות אותיות \enLR{(E} הכי נפוצה) & כל אות ↔ אות אחרת \enLR{(26}! מפתחות) & החלפה & \enLR{Monoalphabetic} \\
\hline
נשבר \enLR{(Kasiski)} אך חזק בהרבה & קיסר עם מפתח מתחלף לפי מילה & החלפה רב-אלפביתית & \enLR{Vigen}è\enLR{re} \\
\hline
סדר בלבד – ניתן לניתוח & שינוי סדר האותיות & שחלוף & \enLR{Transposition (Rail Fence)} \\
\hline
נשברה ע"י טיורינג בבלצ'לי פארק & מכונת רוטורים גרמנית \enLR{(WWII)} & רוטורי & \enLR{Enigma} \\
\hline
לא מעשי – מפתח באורך ההודעה, פעם אחת & המפתח היחיד שהוכח מתמטית כבלתי שביר & \enLR{XOR} עם מפתח אקראי באורך ההודעה & \enLR{One-Time Pad} \\
\hline
\end{xltabular}

\begin{storybox}\boxtitle{סיפור מהחיים: אלן טיורינג ואניגמה}\par  בזמן מלחמת העולם השנייה הצי הגרמני הצפין הודעות במכונת \enLR{Enigma} – \enLR{158} מיליון מיליון מיליון הגדרות אפשריות. אלן טיורינג וצוותו בבלצ'לי פארק בנו מכונה חשמלית ("\enLR{The Bombe")} שניצלה חולשה: אות לעולם לא מוצפנת לעצמה, ומילים צפויות ("\enLR{Wetter"}=מזג אוויר) הופיעו כל בוקר. ההערכה: שבירת \enLR{Enigma} קיצרה את המלחמה בכשנתיים. הלקח לכיתה: אלגוריתם עם מרחב מפתחות ענק עדיין נשבר אם יש בו חולשה מבנית או שימוש לא נכון – בדיוק העיקרון של קרקהופס.\end{storybox}

\sectionheading{\enLR{7.3} מטרות האבטחה ואיזה כלי קריפטוגרפי נותן מה}

\begin{xltabular}{\linewidth}{|>{\setRTL\raggedleft\arraybackslash}X|>{\setRTL\raggedleft\arraybackslash}X|}
\hline
\rowcolor{hdrbg}\textbf{הכלי הקריפטוגרפי} & \textbf{מטרה} \\
\hline
\endhead
הצפנה – סימטרית \enLR{(AES)} או א-סימטרית \enLR{(RSA)} & \textbf{סודיות} \enLR{(Confidentiality)} \\
\hline
\enLR{Hash (SHA-256)} & \textbf{שלמות} \enLR{(Integrity)} \\
\hline
\enLR{HMAC} & \textbf{אימות} \enLR{(Authentication)} + שלמות עם מפתח משותף \\
\hline
חתימה דיגיטלית \enLR{(hash} מוצפן במפתח פרטי) & \textbf{אימות + אי-הכחשה} \enLR{(Non-repudiation)} \\
\hline
\end{xltabular}

\begin{exambox}\boxtitle{בבגרות (שאלה \enLR{5}ה)}\par  "איזה אלגוריתם מבטיח סודיות של מידע מוצפן?" \enLR{1. MD5 2. WPA-2 3. AES 4. HASH} – התשובה \textbf{\enLR{3. AES}}. \enLR{MD5} ו-\enLR{HASH} נותנים שלמות (לא סודיות); \enLR{WPA-2} הוא פרוטוקול אבטחת \enLR{Wi-Fi}, לא אלגוריתם הצפנה.\end{exambox}

\sectionheading{\enLR{7.4} גיבוב \enLR{(Hashing)}}

פונקציית גיבוב מקבלת קלט בכל אורך ומחזירה פלט קבוע \enLR{(digest).} תכונות נדרשות: \textbf{חד-כיווניות} (אי אפשר לשחזר קלט), \textbf{דטרמיניסטיות} (אותו קלט → אותו פלט), \textbf{אפקט מפולת} (שינוי ביט אחד → פלט שונה לגמרי), \textbf{עמידות להתנגשויות} (קשה למצוא שני קלטים עם אותו \enLR{hash).}\par

\begin{xltabular}{\linewidth}{|>{\setRTL\raggedleft\arraybackslash}X|>{\setRTL\raggedleft\arraybackslash}X|>{\setRTL\raggedleft\arraybackslash}X|}
\hline
\rowcolor{hdrbg}\textbf{מצב} & \textbf{אורך פלט} & \textbf{אלגוריתם} \\
\hline
\endhead
\textbf{שבור} – התנגשויות ידועות. לא לאבטחה. (מוזכר בתכנית) & \enLR{128} ביט & \textbf{\enLR{MD5}} \\
\hline
\textbf{שבור} \enLR{(2017, "SHATTERED").} לא לשימוש. (מוזכר בתכנית) & \enLR{160} ביט & \textbf{\enLR{SHA-1}} \\
\hline
הסטנדרט הנוכחי & \enLR{256/512} ביט & \textbf{\enLR{SHA-2}} \enLR{(SHA-256/512)} \\
\hline
חלופה מודרנית & משתנה & \textbf{\enLR{SHA-3}} \\
\hline
\end{xltabular}

\subheading{שימושים}

\begin{itemize}
\item \textbf{שלמות קבצים} – מורידים \enLR{ISO} ומשווים \enLR{SHA-256} לזה שבאתר; אם שונה – הקובץ שונה/פגום.
\item \textbf{אחסון סיסמאות} – שומרים \enLR{hash}, לא את הסיסמה. בכניסה מגבבים ומשווים. חובה \textbf{\enLR{salt}} (מלח – ערך אקראי לכל סיסמה) נגד \enLR{rainbow tables}, ופונקציה איטית \enLR{(bcrypt/Argon2/PBKDF2)} נגד \enLR{brute force.}
\item בסיס ל-\enLR{HMAC} ולחתימה דיגיטלית.
\end{itemize}

\begin{faqbox}\boxtitle{שאלת תלמיד: "אם \enLR{hash} חד-כיווני, איך אתרים 'שוברים' \enLR{MD5}?"}\par  הם לא "הופכים" אותו – הם מנחשים. מחשבים מראש \enLR{hash} של מיליארדי סיסמאות נפוצות \enLR{(rainbow table)} ומחפשים התאמה. לכן: \enLR{(1) salt} – הופך כל טבלה מוכנה ללא רלוונטית; \enLR{(2)} פונקציה איטית – אם כל ניחוש לוקח \enLR{0.5} שנייה, מיליארד ניחושים = \enLR{15} שנה. \enLR{MD5} מהיר מדי בדיוק בגלל זה.\end{faqbox}

\sectionheading{\enLR{7.5} הצפנה סימטרית}

\textbf{מפתח אחד} משמש גם להצפנה וגם לפענוח, ומשותף לשני הצדדים.\par

\begin{center}\fbox{\enLR{Plaintext}} \;→ 🔑 →\; \fbox{\enLR{Ciphertext}} \;→ 🔑(אותו) →\; \fbox{\enLR{Plaintext}}\end{center}

\begin{xltabular}{\linewidth}{|>{\setRTL\raggedleft\arraybackslash}X|>{\setRTL\raggedleft\arraybackslash}X|>{\setRTL\raggedleft\arraybackslash}X|}
\hline
\rowcolor{hdrbg}\textbf{הערכה} & \textbf{אורך מפתח} & \textbf{אלגוריתם} \\
\hline
\endhead
שבור (מפתח קצר מדי, נפרץ ב-\enLR{1998} ביום) & \enLR{56} ביט & \textbf{\enLR{DES}} \\
\hline
מיושן אך נסבל; \enLR{DES} שלוש פעמים & \enLR{112/168} ביט & \textbf{\enLR{3DES}} \\
\hline
\textbf{הסטנדרט העולמי} \enLR{(Rijndael, 2001).} מהיר, מאובטח, בחומרה & \enLR{128/192/256} & \textbf{\enLR{AES}} \\
\hline
שבור – היה ב-\enLR{WEP/SSL} & משתנה (זרם) & \enLR{RC4} \\
\hline
מוזכרים בתכנית; פחות בשימוש היום &  & \enLR{SEAL, Blowfish} \\
\hline
\end{xltabular}

שני סוגים: \textbf{\enLR{block cipher}} \enLR{(AES} – מצפין בלוקים של \enLR{128} ביט; מצבי פעולה \enLR{CBC, GCM)} ו-\textbf{\enLR{stream cipher}} \enLR{(RC4} – ביט אחר ביט).\par

\begin{xltabular}{\linewidth}{|>{\setRTL\raggedleft\arraybackslash}X|>{\setRTL\raggedleft\arraybackslash}X|}
\hline
\rowcolor{hdrbg}\textbf{חסרונות סימטרי} & \textbf{יתרונות סימטרי} \\
\hline
\endhead
\textbf{בעיית חלוקת המפתח} – איך מעבירים את המפתח בבטחה? \textbf{בעיית מדרגיות} – ל-\enLR{N} משתמשים צריך \enLR{N(N}−\enLR{1)/2} מפתחות \enLR{(100} איש = \enLR{4,950} מפתחות) & מהיר מאוד (פי \enLR{1000} מא-סימטרי); מתאים לכמויות מידע גדולות; מפתחות קצרים יחסית \\
\hline
\end{xltabular}

\sectionheading{\enLR{7.6} הצפנה א-סימטרית (מפתח ציבורי)}

לכל צד \textbf{זוג מפתחות} מתמטית קשורים: \textbf{ציבורי} (מפרסמים לכולם) ו\textbf{פרטי} (סוד מוחלט). מה שאחד מצפין – רק השני מפענח. זה פותר את בעיית חלוקת המפתח.\par

\begin{defbox}\boxtitle{שני שימושים הפוכים – חובה להבין}\par  \textbf{לסודיות:} מצפינים ב\underline{ציבורי של הנמען} ⇒ רק הנמען (בעל הפרטי) יפענח. (בבגרות \enLR{4}ו + \enLR{5}ד)\newline  \textbf{לחתימה/אימות:} מצפינים ב\underline{פרטי של השולח} ⇒ כל אחד מפענח בציבורי שלו, ובכך מוכיח שהשולח אכן שלח (רק לו יש את הפרטי).\end{defbox}

\begin{exambox}\boxtitle{בבגרות (שאלה \enLR{5}ד)}\par  "האלגוריתם הא-סימטרי משתמש במפתח הציבורי להצפנה – איזה מפתח לפענוח?" – \textbf{\enLR{1.} מפתח פרטי}.\end{exambox}

\begin{xltabular}{\linewidth}{|>{\setRTL\raggedleft\arraybackslash}X|>{\setRTL\raggedleft\arraybackslash}X|}
\hline
\rowcolor{hdrbg}\textbf{שימוש} & \textbf{אלגוריתם} \\
\hline
\endhead
הצפנה + חתימה; הנפוץ ביותר. בנתב סיסקו: מפתחות \enLR{RSA} ל-\enLR{SSH} (פרק \enLR{2}!) & \textbf{\enLR{RSA}} \\
\hline
\textbf{החלפת מפתחות} – מסכימים על מפתח סודי משותף בערוץ ציבורי, בלי לשלוח אותו & \textbf{\enLR{Diffie-Hellman (DH)}} \\
\hline
חתימה דיגיטלית בלבד (מוזכר בתכנית) & \textbf{\enLR{DSA}} \\
\hline
עקומים אליפטיים – אבטחה זהה במפתח קצר בהרבה (מודרני) & \textbf{\enLR{ECC}} \\
\hline
\end{xltabular}

\begin{defbox}\boxtitle{\enLR{Diffie-Hellman} במשל הצבעים (הכי טוב לכיתה)}\par  אליס ובוב רוצים סוד משותף מול מאזין. \enLR{1.} מסכימים בפומבי על צבע בסיס (צהוב). \enLR{2.} כל אחד בוחר צבע סודי ומערבב עם הצהוב. \enLR{3.} מחליפים את התערובות \textbf{בפומבי} (המאזין רואה אותן). \enLR{4.} כל אחד מוסיף את הסוד \emph{שלו} לתערובת שקיבל. שניהם מגיעים לאותו צבע סופי – \textbf{אבל המאזין לא יכול}, כי הפרדת צבעים היא "בלתי הפיכה" (כמו לוגריתם דיסקרטי במתמטיקה). זה בדיוק \enLR{DH.}\end{defbox}

\begin{exambox}\boxtitle{בבגרות (שאלה \enLR{4}ו)}\par  "איזה אלגוריתם מאפשר לצדדים להסכים על מפתח פרטי על גבי ערוץ ציבורי?" – \textbf{\enLR{2. Diffie-Hellman}}. \enLR{(NOBLE, Dijkstra, Caesar} – מסיחים; \enLR{Dijkstra} הוא אלגוריתם מסלולים, לא הצפנה.)\end{exambox}

\sectionheading{\enLR{7.7} הצפנה היברידית – איך זה באמת עובד \enLR{(TLS/HTTPS)}}

המערכות האמיתיות משלבות: א-סימטרי לפתיחה, סימטרי לתעבורה – "טוב משני העולמות".\par

\begin{defbox}\boxtitle{\enLR{HTTPS} – מה קורה כשנכנסים לאתר בנק}\par  \enLR{1.} הדפדפן והשרת מבצעים \textbf{\enLR{DH}} (או שהדפדפן מצפין מפתח ב-\enLR{RSA} הציבורי של השרת) כדי להסכים על \textbf{מפתח סימטרי \enLR{(session key)}}. \enLR{2.} השרת מציג \textbf{תעודה} חתומה על ידי \enLR{CA} שמוכיחה שהוא באמת הבנק \enLR{(7.9). 3.} מכאן כל התעבורה מוצפנת ב-\textbf{\enLR{AES}} (מהיר) עם ה-\enLR{session key.} הא-סימטרי שימש רק לשנייה הראשונה. זו הסיבה שרואים 🔒.\end{defbox}

\sectionheading{\enLR{7.8 HMAC} וחתימה דיגיטלית}

\subheading{\enLR{HMAC (Hash-based Message Authentication Code)}}

\enLR{hash} לבדו נותן שלמות, אבל תוקף \enLR{MITM} יכול לשנות גם את ההודעה וגם את ה-\enLR{hash.} \textbf{\enLR{HMAC}} = \enLR{hash} של (ההודעה + מפתח סודי משותף). עכשיו רק מי שיש לו המפתח יכול לחשב \enLR{HMAC} תקין ⇒ \textbf{שלמות + אימות מקור}. משמש ב-\enLR{IPsec} (פרק \enLR{8)}, ב-\enLR{API}, ו-\enLR{JWT.}\par

\subheading{חתימה דיגיטלית}

\begin{defbox}\boxtitle{התהליך}\par  \textbf{חתימה (השולח):} \enLR{1.} מחשב \enLR{hash} של המסמך. \enLR{2.} מצפין את ה-\enLR{hash} ב\underline{מפתח הפרטי שלו} = החתימה. שולח מסמך + חתימה.\newline  \textbf{אימות (המקבל):} \enLR{1.} מפענח את החתימה ב\underline{מפתח הציבורי של השולח} ⇒ מקבל את ה-\enLR{hash} המקורי. \enLR{2.} מחשב \enLR{hash} של המסמך שקיבל. \enLR{3.} משווה. תואם ⇒ המסמך לא שונה (שלמות) \textbf{וגם} נשלח באמת על ידו (אימות + אי-הכחשה – רק לו יש את הפרטי).\end{defbox}

למה מצפינים את ה-\enLR{hash} ולא את כל המסמך? כי א-סימטרי איטי; ה-\enLR{hash} קטן וקבוע.\par

\sectionheading{\enLR{7.9 PKI} ותעודות דיגיטליות}

בעיית \enLR{DH} ו-\enLR{RSA}: איך אני יודע שהמפתח הציבורי "של הבנק" באמת שייך לבנק ולא לתוקף \enLR{MITM}? הפתרון: \textbf{\enLR{PKI}} \enLR{(Public Key Infrastructure)} – מערכת של גורמים מהימנים שמאשרים "מפתח ציבורי \enLR{X} שייך לישות \enLR{Y".}\par

\begin{xltabular}{\linewidth}{|>{\setRTL\raggedleft\arraybackslash}X|>{\setRTL\raggedleft\arraybackslash}X|}
\hline
מסמך שקושר זהות למפתח ציבורי, \textbf{חתום דיגיטלית על ידי \enLR{CA}}. תקן: \enLR{X.509.} & \textbf{תעודה דיגיטלית \enLR{(Certificate)}} \\
\hline
הגורם המהימן שמנפיק וחותם תעודות \enLR{(DigiCert, Let's Encrypt).} & \textbf{\enLR{CA (Certificate Authority)}} \\
\hline
מאמת את זהות המבקש לפני הנפקה. & \textbf{\enLR{RA (Registration Authority)}} \\
\hline
רשימת/בדיקת תעודות שבוטלו (למשל אם מפתח פרטי נגנב). & \textbf{\enLR{CRL / OCSP}} \\
\hline
ה-\enLR{CA} העליון; התעודה שלו מובנית במערכת ההפעלה/דפדפן ("\enLR{trust anchor").} & \textbf{\enLR{Root CA}} \\
\hline
\end{xltabular}

\begin{defbox}\boxtitle{שרשרת אמון \enLR{(Chain of Trust)} – היררכיה}\par  \enLR{Root CA} (סומכים עליו מראש) → חותם על → \enLR{Intermediate CA} → חותם על → תעודת השרת \enLR{(bank.com).} הדפדפן בודק את השרשרת עד ל-\enLR{root} שהוא מכיר. אם חוליה חסרה או פגה – אזהרה. זו ה"היררכיית סטנדרטים ואישורי סטנדרטים" שבתכנית.\end{defbox}

\begin{storybox}\boxtitle{סיפור מהחיים: \enLR{DigiNotar} – כש-\enLR{CA} נפרץ \enLR{(2011)}}\par  תוקף פרץ ל-\enLR{CA} ההולנדי \enLR{DigiNotar} והנפיק לעצמו תעודה מזויפת ל-*.\enLR{google.com.} עם התעודה הזו, המודיעין האיראני (ככל הנראה) ביצע \enLR{MITM} על \enLR{Gmail} של \enLR{300,000} אזרחים איראנים – הדפדפן שלהם הציג 🔒 ירוק ולא חשד בכלום, כי התעודה הייתה "חתומה כדין". כשהתגלה, הדפדפנים הסירו את \enLR{DigiNotar} מרשימת ה-\enLR{root} המהימנים, והחברה פשטה רגל תוך חודש. הלקח: כל מודל האמון של האינטרנט נשען על כך שה-\enLR{CA}-ים ישרים ומאובטחים – חוליה אחת שבורה מפילה הכול.\end{storybox}

\sectionheading{\enLR{7.10} חוזק מפתחות ומדיניות \enLR{NSA/NIST}}

\begin{itemize}
\item \textbf{מפתח ארוך = מרחב חיפוש גדול יותר}. סימטרי \enLR{128} ביט = \enLR{2}\textsuperscript{\enLR{128}} אפשרויות – מעבר להישג יד גם למחשבי-על. אבל אורך מפתח א-סימטרי אינו בר-השוואה ישירה: \enLR{RSA 2048} ≈ \enLR{AES 112} בערך.
\item המלצות \enLR{(NIST SP 800): AES-128} ומעלה, \enLR{RSA-2048} ומעלה (רצוי \enLR{3072), SHA-256} ומעלה, \enLR{ECC-256.}
\item \textbf{\enLR{Suite B}} של ה-\enLR{NSA} (מוזכר "חוקי \enLR{NSA"): AES, ECDH, ECDSA, SHA-2} – לתקשורת ממשלתית.
\item \textbf{ניהול מפתחות} (הנושא "ניהול מפתחות ע"י שינוי אחסון וריפיקציה"): יצירה, אחסון מאובטח \enLR{(HSM)}, החלפה תקופתית \enLR{(rotation)}, ביטול, גיבוי. מפתח שנחשף = כל מה שהוצפן בו חשוף.
\item איום עתידי: \textbf{מחשוב קוונטי} יאיים על \enLR{RSA/DH} ⇒ קריפטוגרפיה פוסט-קוונטית \enLR{(NIST} בחר אלגוריתמים ב-\enLR{2024).}
\end{itemize}

\sectionheading{\enLR{7.11} תרגילים לתלמידים}

\begin{exbox}\boxtitle{תרגיל \enLR{1} – צופן קיסר}\par  עם \enLR{K=3}: הצפינו את המילה \code{NET}. עם \enLR{K=5}: פענחו את \code{MJQQT}. \enLR{(A=0}…\enLR{Z=25)} \par\vskip2pt\hrule\vskip3pt\textbf{}\enLR{NET: N(13)}→\enLR{Q, E(4)}→\enLR{H, T(19)}→\enLR{W} = \textbf{\enLR{QHW}}. פענוח \enLR{MJQQT} עם \enLR{K=5: M(12)-5=H, J(9)-5=E, Q(16)-5=L, Q}→\enLR{L, T(19)-5=O} = \textbf{\enLR{HELLO}}.\end{exbox}

\begin{exbox}\boxtitle{תרגיל \enLR{2} – איזה כלי לאיזו מטרה}\par  לכל דרישה בחרו: \enLR{hash} / הצפנה סימטרית / א-סימטרית / \enLR{HMAC} / חתימה דיגיטלית.\newline  א. לוודא שקובץ שהורד לא שונה. ב. להצפין \enLR{10GB} גיבוי במהירות. ג. לשלוח מפתח סודי דרך האינטרנט בבטחה. ד. לוודא שעדכון תוכנה הגיע באמת מהיצרן. ה. לאחסן סיסמאות במסד נתונים. \par\vskip2pt\hrule\vskip3pt\textbf{}א. \enLR{hash} ב. סימטרית \enLR{(AES)} ג. א-סימטרית / \enLR{DH} ד. חתימה דיגיטלית ה. \enLR{hash + salt} (פונקציה איטית)\end{exbox}

\begin{exbox}\boxtitle{תרגיל \enLR{3} – ציבורי או פרטי?}\par  דנה רוצה לשלוח ליוסי הודעה סודית, ולחתום עליה. איזה מפתח היא משתמשת בכל שלב?\newline  א. להצפין את ההודעה כך שרק יוסי יקרא. ב. לחתום כדי שיוסי ידע שזה ממנה. ואצל יוסי: ג. לפענח את ההודעה. ד. לאמת את החתימה. \par\vskip2pt\hrule\vskip3pt\textbf{}א. הציבורי של יוסי ב. הפרטי של דנה ג. הפרטי של יוסי ד. הציבורי של דנה\end{exbox}

\begin{exbox}\boxtitle{תרגיל \enLR{4} – קידוד/הצפנה/גיבוב}\par  סווגו כל אחד: א. \code{btoa} ל-\enLR{Base64.} ב. \enLR{SHA-256.} ג. \enLR{AES-256-CBC.} ד. \enLR{Morse.} ה. \enLR{bcrypt.} \par\vskip2pt\hrule\vskip3pt\textbf{}א. קידוד ב. גיבוב ג. הצפנה ד. קידוד ה. גיבוב (איטי, לסיסמאות)\end{exbox}

\begin{exbox}\boxtitle{תרגיל \enLR{5} – דיון: התנגשות}\par  הסבירו במילים שלכם מדוע "עמידות להתנגשויות" חשובה לחתימה דיגיטלית. מה תוקף יכול לעשות אם הוא מוצא שני מסמכים עם אותו \enLR{hash}? \par\vskip2pt\hrule\vskip3pt\textbf{}בחתימה חותמים על ה-\enLR{hash.} אם התוקף מוצא מסמך זדוני עם אותו \enLR{hash} כמו מסמך תמים שנחתם – החתימה "תקפה" גם על הזדוני, והוא יכול להחליף אותם. בדיוק בגלל זה \enLR{MD5} ו-\enLR{SHA-1} פסולים לחתימה.\end{exbox}

\sectionheading{\enLR{7.12} שאלות בסגנון בגרות}

\begin{exambox}\boxtitle{שאלה \enLR{1}}\par  איזה אלגוריתם מבטיח סודיות של מידע מוצפן? \enLR{MD5 / WPA-2 / AES / HASH} \par\vskip2pt\hrule\vskip3pt\textbf{}\textbf{\enLR{AES}}. השאר: \enLR{MD5/HASH} – שלמות; \enLR{WPA-2} – פרוטוקול.\end{exambox}

\begin{exambox}\boxtitle{שאלה \enLR{2}}\par  בצופן \enLR{C=(P+K) mod 26}, מה תפקיד ה-\enLR{mod 26}? \par\vskip2pt\hrule\vskip3pt\textbf{}מספר האותיות בשפה; "מגלגל" את התוצאה חזרה לטווח \enLR{0}–\enLR{25.}\end{exambox}

\begin{exambox}\boxtitle{שאלה \enLR{3}}\par  כאשר האלגוריתם הא-סימטרי מצפין במפתח ציבורי – באיזה מפתח מפענחים? \par\vskip2pt\hrule\vskip3pt\textbf{}במפתח הפרטי (של הנמען).\end{exambox}

\begin{exambox}\boxtitle{שאלה \enLR{4}}\par  איזה אלגוריתם מאפשר לשני צדדים להסכים על מפתח סודי משותף מעל ערוץ ציבורי? \par\vskip2pt\hrule\vskip3pt\textbf{}\enLR{Diffie-Hellman.}\end{exambox}

\begin{exambox}\boxtitle{שאלה \enLR{5}}\par  מהו ההבדל העיקרי בין הצפנה סימטרית לא-סימטרית מבחינת מפתחות, ומהו החיסרון המרכזי של כל אחת? \par\vskip2pt\hrule\vskip3pt\textbf{}סימטרית – מפתח אחד משותף; חיסרון: חלוקת המפתח ומדרגיות. א-סימטרית – זוג מפתחות (ציבורי/פרטי); חיסרון: איטית מאוד.\end{exambox}

\begin{exambox}\boxtitle{שאלה \enLR{6}}\par  מדוע נדרשת \enLR{PKI}, ומה תפקיד ה-\enLR{CA}? \par\vskip2pt\hrule\vskip3pt\textbf{}כדי לוודא שמפתח ציבורי אכן שייך לישות שהוא טוען (נגד \enLR{MITM).} ה-\enLR{CA} הוא גורם מהימן שמנפיק וחותם תעודות שמקשרות זהות למפתח ציבורי.\end{exambox}

\sectionheading{\enLR{7.13} מעבדה \enLR{(3} שעות)}

\begin{enumerate}
\item \textbf{קיסר:} באתר \enLR{cryptii.com} הצפינו ופענחו; בדקו כמה מהר \enLR{brute force} על \enLR{25} מפתחות שובר אותו.
\item \textbf{\enLR{Hash}:} באתר (או \code{echo -n "password" | sha256sum}) חשבו \enLR{SHA-256} של "\enLR{hello"} ושל "\enLR{Hello"} – ראו את אפקט המפולת. שנו ביט אחד וראו פלט שונה לגמרי.
\item \textbf{שלמות:} הורידו קובץ, שנו בו תו אחד, חשבו \enLR{hash} לפני ואחרי – הראו שהשתנה.
\item \textbf{סיסמאות:} הראו \enLR{rainbow table} אונליין ששובר \enLR{MD5} של "\enLR{123456"}, ואיך \enLR{salt} מסכל אותו.
\item \textbf{תעודות:} בדפדפן, לחצו על 🔒 של אתר בנק, הציגו את התעודה, את ה-\enLR{CA} ואת שרשרת האמון ותוקף התעודה.
\item \textbf{\enLR{RSA} ל-\enLR{SSH} (קישור לפרק \enLR{2)}:} ב-\enLR{Packet Tracer} צרו מפתחות \enLR{RSA} בנתב והתחברו ב-\enLR{SSH} – הדגישו שאלה בדיוק הא-סימטרי בפעולה.
\end{enumerate}

\sectionheading{בנק שאלות שתלמידים שואלים – ותשובות מוכנות}

שאלות נפוצות בפרק ההצפנה, עם תשובות מוכנות.\par

\textbf{ש: "אם גיבוב הוא חד-כיווני, איך אתרים 'שוברים' סיסמאות \enLR{MD5}?"}\newline ת: הם לא הופכים את הגיבוב – הם \underline{מנחשים}. מחשבים מראש את הגיבוב של מיליארדי סיסמאות נפוצות (טבלת \enLR{rainbow)} ומחפשים התאמה. לכן: מוסיפים \textbf{\enLR{salt}} (ערך אקראי לכל סיסמה) שהופך טבלה מוכנה לחסרת ערך, ומשתמשים בפונקציה \underline{איטית} \enLR{(bcrypt)} שהופכת ניחוש המוני לבלתי מעשי.\par

\textbf{ש: "אז \enLR{Base64} זו הצפנה?"}\newline ת: לא! \enLR{Base64} הוא \textbf{קידוד} – שינוי צורת כתיבה בלי סוד. כל אחד מפענח אותו בשנייה (רמז: מחרוזת שמסתיימת ב-'=='). קידוד ≠ הצפנה. השאלה המבחינה: "האם צריך מפתח סודי כדי לחזור?".\par

\textbf{ש: "אם אני מצפין הכול ב-\enLR{AES}, אני מוגן לגמרי?"}\newline ת: הצפנה נותנת \underline{סודיות} בלבד. היא לא מונעת מחיקה של הקובץ (זמינות – כופרה דווקא \emph{משתמשת} בהצפנה נגדכם!) ולא בהכרח מוודאת שהמידע לא שונה (שלמות – לזה צריך גיבוב/\enLR{HMAC).} צריך את שלושתם.\par

\textbf{ש: "מה ההבדל בין סימטרי לא-סימטרי, ומתי משתמשים בכל אחד?"}\newline ת: \textbf{סימטרי} \enLR{(AES)} = מפתח אחד משותף, מהיר מאוד – טוב לכמויות מידע גדולות, אבל יש בעיה: איך מעבירים את המפתח בבטחה? \textbf{א-סימטרי} \enLR{(RSA/DH)} = זוג מפתחות ציבורי/פרטי, פותר את בעיית העברת המפתח, אבל איטי. בפועל \underline{משלבים}: א-סימטרי כדי להסכים על מפתח, ואז סימטרי לתעבורה (כך עובד \enLR{HTTPS).}\par

\textbf{ש: "אם המפתח הציבורי גלוי לכולם, איך זה מאובטח?"}\newline ת: כי מה שמוצפן במפתח הציבורי ניתן לפענוח \underline{רק} במפתח הפרטי התואם, שנשאר סודי אצל הבעלים. גם אם כל העולם מכיר את המפתח הציבורי שלכם – רק אתם יכולים לפענח את מה ששלחו אליכם.\par

\textbf{ש: "איך \enLR{Diffie-Hellman} מסכים על מפתח סודי אם המאזין רואה הכול?"}\newline ת: משל הצבעים: כל צד מערבב צבע סודי עם צבע בסיס משותף, מחליפים את התערובות בפומבי, וכל אחד מוסיף שוב את הסוד שלו. שניהם מגיעים לאותו צבע – אבל המאזין לא יכול, כי "הפרדת צבעים" (לוגריתם דיסקרטי) היא בלתי הפיכה מעשית. \enLR{DH} לא מעביר את המפתח – הוא \underline{מייצר} אותו בשני הצדדים.\par

\textbf{ש: "אם \enLR{RSA} כל כך טוב, למה בכלל צריך תעודות ו-\enLR{PKI}?"}\newline ת: כי \enLR{RSA} לבדו לא אומר \underline{למי} שייך המפתח הציבורי. תוקף \enLR{MITM} יכול לתת לכם את המפתח \emph{שלו} ולטעון שהוא הבנק. תעודה חתומה על ידי \enLR{CA} מהימן מוכיחה "המפתח הזה באמת שייך לבנק". זה מה שקורה כשרואים 🔒 בדפדפן.\par

\textbf{ש: "מה ההבדל בין חתימה דיגיטלית להצפנה?"}\newline ת: הצפנה מסתירה תוכן (סודיות), ומשתמשת ב\underline{מפתח הציבורי של הנמען}. חתימה מוכיחה מי שלח ושהמסמך לא שונה (אימות+שלמות), ומשתמשת ב\underline{מפתח הפרטי של השולח}. שני שימושים הפוכים של אותו זוג מפתחות.\par

\textbf{ש: "האם מחשב קוונטי ישבור את כל ההצפנות?"}\newline ת: הוא מאיים בעיקר על ה\underline{א-סימטריות} \enLR{(RSA/DH)}, פחות על סימטריות \enLR{(AES-256} נחשב עמיד). לכן \enLR{NIST} כבר בחר \enLR{(2024)} אלגוריתמים "פוסט-קוונטיים". זה איום עתידי שמתכוננים אליו כבר היום.\par

\sectionheading{מילון מונחים – הגדרות}

הגדרות תמציתיות של המונחים המרכזיים בפרק. שימושי גם כדף עזר לבחינה.\par

\begin{xltabular}{\linewidth}{|>{\setRTL\raggedleft\arraybackslash}X|>{\setRTL\raggedleft\arraybackslash}X|}
\hline
\rowcolor{hdrbg}\textbf{הגדרה} & \textbf{מונח} \\
\hline
\endhead
מדע יצירת ההצפנות. & \textbf{קריפטוגרפיה} \\
\hline
מדע שבירת ההצפנות. & \textbf{קריפטואנליזה} \\
\hline
הטקסט הגלוי (לפני) והטקסט המוצפן (אחרי). & \textbf{\enLR{Plaintext / Ciphertext}} \\
\hline
סוד שקובע את פעולת ההצפנה; בלעדיו הטקסט המוצפן חסר משמעות. & \textbf{מפתח \enLR{(Key)}} \\
\hline
המרת ייצוג ללא סוד \enLR{(Base64}, מורס) – אינו אבטחה. & \textbf{קידוד \enLR{(Encoding)}} \\
\hline
המרה הפיכה בעזרת מפתח סודי; מטרה: סודיות. & \textbf{הצפנה \enLR{(Encryption)}} \\
\hline
פונקציה חד-כיוונית שמפיקה טביעת אצבע קבועה; מטרה: שלמות. & \textbf{גיבוב \enLR{(Hashing)}} \\
\hline
ערך אקראי המתווסף לסיסמה לפני גיבוב, נגד טבלאות \enLR{rainbow.} & \textbf{\enLR{Salt}} \\
\hline
אלגוריתמי גיבוב ישנים ושבורים – לא לשימוש אבטחתי. & \textbf{\enLR{MD5 / SHA-1}} \\
\hline
אלגוריתם הגיבוב הסטנדרטי כיום. & \textbf{\enLR{SHA-256}} \\
\hline
צופן החלפה שמזיז כל אות במספר קבוע: \enLR{C=(P+K) mod 26.} & \textbf{צופן קיסר \enLR{(Caesar)}} \\
\hline
מפתח אחד משותף להצפנה ולפענוח; מהיר \enLR{(AES, 3DES, DES).} & \textbf{הצפנה סימטרית} \\
\hline
זוג מפתחות ציבורי/פרטי; פותר העברת מפתח; איטי \enLR{(RSA).} & \textbf{הצפנה א-סימטרית} \\
\hline
תקן ההצפנה הסימטרי המוביל כיום \enLR{(128/192/256} ביט). & \textbf{\enLR{AES}} \\
\hline
אלגוריתם א-סימטרי נפוץ להצפנה ולחתימה; משמש גם ל-\enLR{SSH.} & \textbf{\enLR{RSA}} \\
\hline
שיטה להסכמה על מפתח סודי משותף מעל ערוץ ציבורי. & \textbf{\enLR{Diffie-Hellman (DH)}} \\
\hline
זוג מפתחות; מה שמוצפן באחד מפוענח רק בשני. & \textbf{מפתח ציבורי / פרטי} \\
\hline
גיבוב עם מפתח סודי משותף; נותן שלמות ואימות מקור. & \textbf{\enLR{HMAC}} \\
\hline
גיבוב המסמך המוצפן במפתח הפרטי של השולח; אימות + אי-הכחשה. & \textbf{חתימה דיגיטלית} \\
\hline
תשתית מפתח ציבורי – מערכת גורמים מהימנים המנפיקים תעודות. & \textbf{\enLR{PKI}} \\
\hline
מסמך חתום ע"י \enLR{CA} הקושר זהות למפתח ציבורי. & \textbf{תעודה דיגיטלית \enLR{(X.509)}} \\
\hline
גורם מהימן המנפיק וחותם תעודות. & \textbf{\enLR{CA (Certificate Authority)}} \\
\hline
היררכיה: \enLR{Root CA} → \enLR{Intermediate} → תעודת שרת. & \textbf{שרשרת אמון \enLR{(Chain of Trust)}} \\
\hline
\end{xltabular}


\clearpage
\phantomsection\addcontentsline{toc}{section}{פרק \enLR{8} – מערכות \enLR{VPN}}

\sloppy
\chapterheading{פרק \enLR{8} – מערכות \enLR{VPN}}{\small\color{subgray}\enLR{10} שעות עיוני + \enLR{3} מעשי · שבועות \enLR{24}–\enLR{26}\par}\vskip4pt

\begin{metabox}\textbf{מטרות:} מערכות \enLR{VPN} · הגדרת \enLR{VPN} · \enLR{IPsec} · הגדרה ב-\enLR{CLI} · \enLR{GRE} · \enLR{SSL VPN}\\ \textbf{בבגרות:} \enLR{IPsec} מאבטח את חבילת \enLR{TCP/IP} (שכבה \enLR{3)}\end{metabox}

\begin{defbox}\boxtitle{לפני שמתחילים – מה זו "מנהרה מוצפנת"? (למי שאין רקע)}\par   \textbf{הבעיה:} כשמידע עובר באינטרנט, הוא עובר דרך נתבים רבים של חברות שונות. כל אחד בדרך יכול, עקרונית, לראות אותו. זה מסוכן למידע רגיש. \textbf{\enLR{VPN}} \enLR{(Virtual Private Network) = "}רשת פרטית וירטואלית". הוא יוצר \textbf{מנהרה מוצפנת} דרך האינטרנט הציבורי: המידע מוצפן בצד אחד, עובר את האינטרנט כג'יבריש, ומפוענח רק בצד השני. \textbf{אנלוגיה:} במקום לשלוח גלויה שכל דוור קורא (אינטרנט רגיל), שולחים מזוודה נעולה שרק לשולח ולנמען יש מפתח אליה \enLR{(VPN).} \textbf{שני שימושים:} חיבור \underline{סניף לסניף} \enLR{(Site-to-Site}, כמו לחבר שני משרדים), וחיבור \underline{עובד מהבית} \enLR{(Remote Access).}  \textbf{חשוב:} פרק זה מסתמך ישירות על פרק \enLR{7} (הצפנה). כל מה שנלמד שם – מפתח סימטרי, מפתח ציבורי/פרטי, גיבוב – חוזר כאן בפעולה. אם משהו לא ברור, חזרו לפרק \enLR{7.}\end{defbox}

\sectionheading{\enLR{8.1} מהו \enLR{VPN} ולמה הוא נחוץ}

\textbf{\enLR{VPN}} \enLR{(Virtual Private Network)} יוצר "מנהרה" מוצפנת דרך רשת ציבורית (האינטרנט), כך שנתונים עוברים כאילו הם על קו פרטי. במקום לשכור קו תקשורת יקר בין סניפים \enLR{(leased line)}, משתמשים באינטרנט הזול – והצפנה מספקת את הפרטיות. \enLR{VPN} מספק את שלוש מטרות ה-\enLR{CIA} לתעבורה: \textbf{סודיות} (הצפנה), \textbf{שלמות} \enLR{(HMAC)}, \textbf{אימות} (מי הצד השני), ולעיתים \textbf{\enLR{anti-replay}} (מניעת שידור חוזר של חבילה שנלכדה).\par

\begin{storybox}\boxtitle{סיפור מהחיים: בית הקפה והמנהל שעבד מרחוק}\par  מנהל מכירות יושב בבית קפה, מתחבר ל-\enLR{Wi-Fi} הפתוח, ופותח את מערכת ה-\enLR{CRM} של החברה. בלי \enLR{VPN}, כל מי שיושב בבית הקפה עם \enLR{Wireshark} (או \enLR{Evil Twin} – פרק \enLR{6)} רואה את הכול. עם \enLR{VPN}: המחשב פותח מנהרת \enLR{IPsec/SSL} מוצפנת אל שער החברה עוד לפני שנשלח ביט אחד של \enLR{CRM.} המאזין רואה רק "תעבורה מוצפנת ל-\enLR{VPN gateway"} – ג'יבריש. זה בדיוק מה שלמדנו בפרק \enLR{7 (AES+HMAC+DH)}, עכשיו ארוז כשירות רשת.\end{storybox}

\sectionheading{\enLR{8.2} סוגי \enLR{VPN}}

\begin{xltabular}{\linewidth}{|>{\setRTL\raggedleft\arraybackslash}X|>{\setRTL\raggedleft\arraybackslash}X|>{\setRTL\raggedleft\arraybackslash}X|}
\hline
\rowcolor{hdrbg}\textbf{דוגמה} & \textbf{מי מתחבר למי} & \textbf{סוג} \\
\hline
\endhead
סניף תל אביב ↔ סניף חיפה & שער (נתב/\enLR{ASA)} לשער – מחבר שתי רשתות שלמות. שקוף למשתמשים & \textbf{\enLR{Site-to-Site}} \\
\hline
עובד מהבית ↔ המשרד \enLR{(Cisco AnyConnect)} & משתמש יחיד (לקוח) ↔ שער החברה & \textbf{\enLR{Remote Access}} \\
\hline
\end{xltabular}

\begin{xltabular}{\linewidth}{|>{\setRTL\raggedleft\arraybackslash}X|>{\setRTL\raggedleft\arraybackslash}X|}
\hline
\rowcolor{hdrbg}\textbf{} & \textbf{לפי מי מפעיל} \\
\hline
\endhead
הארגון מנהל \enLR{(IPsec, SSL)} – מה שנלמד כאן & \textbf{\enLR{Enterprise VPN}} \\
\hline
ספק התקשורת מספק – \textbf{\enLR{MPLS}}, \enLR{Metro Ethernet} (מוזכר בתכנית: \enLR{MPLS, GRE)} & \textbf{\enLR{Service Provider VPN}} \\
\hline
\end{xltabular}

\begin{mistakebox}\boxtitle{הבהרה חשובה}\par  "\enLR{VPN} מסחרי" \enLR{(NordVPN} וכו') נועד להסתרת זהות/מיקום – שימוש שונה מ-\enLR{VPN} ארגוני שמחבר משתמש/סניף לרשת החברה בבטחה. בקורס מדובר ב-\enLR{VPN} ארגוני.\end{mistakebox}

\sectionheading{\enLR{8.3 GRE} – מנהור בסיסי (בלי הצפנה!)}

\textbf{\enLR{GRE}} \enLR{(Generic Routing Encapsulation)} הוא פרוטוקול מנהור של סיסקו שעוטף חבילה בתוך חבילת \enLR{IP} חדשה. יתרון: יכול לשאת \textbf{\enLR{multicast} ו-\enLR{broadcast}} (ולכן פרוטוקולי ניתוב כמו \enLR{OSPF/EIGRP} רצים דרכו) וגם תעבורה שאינה \enLR{IP.} חיסרון קריטי: \textbf{\enLR{GRE} אינו מצפין כלום} – הוא רק עוטף. לכן בפועל: \textbf{\enLR{GRE over IPsec}} – \enLR{GRE} נותן את הגמישות \enLR{(multicast, routing), IPsec} נותן את ההצפנה.\par

\begin{codebox}\begin{english}\codefont\footnotesize\color{cfg}\setlength{\parindent}{0pt}
\textcolor{ccom}{\texthebrew{! מנהרת \enLR{GRE} בין \enLR{R1} ל-\enLR{R2}}}\\
R1(config)\# \textcolor{cbold}{\textbf{interface tunnel 0}}\\
R1(config-if)\# \textcolor{cbold}{\textbf{ip address 172.16.1.1 255.255.255.252}}    \textcolor{ccom}{\texthebrew{! כתובת פנימית של המנהרה}}\\
R1(config-if)\# \textcolor{cbold}{\textbf{tunnel source 200.1.1.1}}                   \textcolor{ccom}{\texthebrew{! ה-\enLR{IP} הציבורי המקומי}}\\
R1(config-if)\# \textcolor{cbold}{\textbf{tunnel destination 200.2.2.2}}              \textcolor{ccom}{\texthebrew{! ה-\enLR{IP} הציבורי המרוחק}}\\
R1(config-if)\# tunnel mode gre ip                            \textcolor{ccom}{\texthebrew{! ברירת מחדל}}\\
\textcolor{ccom}{\texthebrew{! עכשיו מנתבים דרך המנהרה}}\\
R1(config)\# ip route 10.2.2.0 255.255.255.0 172.16.1.2\\
R1\# \textcolor{cbold}{\textbf{show ip interface brief | include Tunnel}}\\
R1\# show interfaces tunnel 0
\end{english}\end{codebox}

\sectionheading{\enLR{8.4 IPsec} – לב הפרק}

\textbf{\enLR{IPsec}} אינו פרוטוקול יחיד אלא \textbf{מסגרת} \enLR{(framework)} של פרוטוקולים שמאבטחת תעבורה ב\textbf{שכבה \enLR{3}} \enLR{(IP).} זו התשובה בבגרות: "\enLR{IPsec} נועד לאבטח את חבילת הפרוטוקולים \enLR{TCP/IP".} כיוון שהוא בשכבה \enLR{3}, הוא מאבטח \emph{כל} תעבורה מעליו \enLR{(TCP, UDP}, כל אפליקציה) בשקיפות – בניגוד ל-\enLR{SSL/TLS} שבשכבה גבוהה יותר.\par

\begin{exambox}\boxtitle{בבגרות (שאלה \enLR{6}א)}\par  "איזה סוג תקשורת נועדה חבילת הפרוטוקולים \enLR{IPSEC} לאבטח?" – \textbf{\enLR{2. TCP/IP}} (שכבה \enLR{3).} לא \enLR{SSH}, לא \enLR{Telnet}, לא \enLR{UDP} ספציפית.\end{exambox}

\subheading{שני פרוטוקולי ההגנה}

\begin{xltabular}{\linewidth}{|>{\setRTL\raggedleft\arraybackslash}X|>{\setRTL\raggedleft\arraybackslash}X|>{\setRTL\raggedleft\arraybackslash}X|}
\hline
\rowcolor{hdrbg}\textbf{\enLR{ESP (Encapsulating Security Payload)}} & \textbf{\enLR{AH (Authentication Header)}} & \textbf{} \\
\hline
\endhead
\enLR{50} & \enLR{51} & פרוטוקול \enLR{IP} \\
\hline
כן (המטען) & כן (כולל הכותרת החיצונית) & שלמות + אימות \\
\hline
\textbf{כן} & \textbf{לא!} & \textbf{סודיות (הצפנה)} \\
\hline
כן (עם \enLR{NAT-T)} & לא (מגן על הכותרת) & עובד עם \enLR{NAT} \\
\hline
\textbf{הבחירה כמעט תמיד} & כמעט לא בשימוש (אין הצפנה) & שימוש \\
\hline
\end{xltabular}

\subheading{שני מצבי פעולה}

\begin{itemize}
\item \textbf{\enLR{Transport mode}} – מצפין רק את המטען \enLR{(payload)}; הכותרת המקורית נשמרת. בין שני מארחים.
\item \textbf{\enLR{Tunnel mode}} – מצפין את כל החבילה המקורית ועוטף בכותרת \enLR{IP} חדשה. בין שערים \enLR{(site-to-site).} \textbf{הנפוץ.}
\end{itemize}

\subheading{ארבעת מרכיבי ה-\enLR{IPsec} (מה שהמנהל צריך לבחור)}

\begin{xltabular}{\linewidth}{|>{\setRTL\raggedleft\arraybackslash}X|>{\setRTL\raggedleft\arraybackslash}X|}
\hline
\rowcolor{hdrbg}\textbf{אפשרויות} & \textbf{תפקיד} \\
\hline
\endhead
\enLR{DES / 3DES} / \textbf{\enLR{AES}} & הצפנה (סודיות) \\
\hline
\enLR{MD5} / \textbf{\enLR{SHA}} \enLR{(HMAC)} & שלמות \enLR{(hash)} \\
\hline
\textbf{\enLR{PSK}} (מפתח משותף מראש) / תעודות \enLR{RSA (PKI} – פרק \enLR{7)} & אימות הצדדים \\
\hline
\textbf{\enLR{DH}} \enLR{(Diffie-Hellman)} – קבוצות \enLR{1,2,5,14}… & החלפת מפתחות \\
\hline
\end{xltabular}

שימו לב: כל מה שלמדנו בפרק \enLR{7} חוזר כאן – \enLR{AES} לסודיות, \enLR{SHA-HMAC} לשלמות, \enLR{DH} להחלפת מפתחות, \enLR{RSA/PKI} לאימות. \enLR{IPsec} הוא "פרק \enLR{7} בפעולה".\par

\sectionheading{\enLR{8.5 IKE} – ניהול המפתחות והמנהרה}

\textbf{\enLR{IKE}} \enLR{(Internet Key Exchange, "}מפתחות אינטרנט משתנים" בתכנית) הוא הפרוטוקול שמקים את מנהרת ה-\enLR{IPsec}: מבצע \enLR{DH}, מאמת את הצדדים, ומסכים על הפרמטרים. פועל בשני שלבים:\par

\begin{defbox}\boxtitle{\enLR{IKE Phase 1} – בונים ערוץ ניהול מאובטח \enLR{(ISAKMP SA)}}\par  הצדדים מסכימים על מדיניות (הצפנה/\enLR{hash/DH}/אימות), מבצעים \enLR{DH}, ומאמתים זה את זה \enLR{(PSK} או תעודות). התוצאה: ערוץ מוצפן \emph{דו-כיווני} אחד שבו ינוהל שלב \enLR{2.} מצבים: \enLR{Main mode (6} הודעות, בטוח) או \enLR{Aggressive mode (3}, מהיר).\end{defbox}

\begin{defbox}\boxtitle{\enLR{IKE Phase 2} – בונים את מנהרת הנתונים \enLR{(IPsec SA)}}\par  בתוך הערוץ המאובטח משלב \enLR{1}, מסכימים על ה-\enLR{transform set (ESP-AES-SHA)} ומקימים שני \textbf{\enLR{SA}} חד-כיווניים (אחד לכל כיוון) שדרכם תעבור התעבורה בפועל. מצב: \enLR{Quick mode.}\end{defbox}

\textbf{\enLR{SA}} \enLR{(Security Association)} – "הסכם" חד-כיווני שמגדיר איך לאבטח זרם תעבורה מסוים, עם מזהה \textbf{\enLR{SPI}}. \enLR{Phase 2} מייצר \enLR{PFS (Perfect Forward Secrecy)} אם מפעילים \enLR{DH} מחדש – כך שדליפת מפתח אחד לא חושפת תעבורה עבר.\par

\sectionheading{\enLR{8.6} הגדרת \enLR{Site-to-Site IPsec VPN} ב-\enLR{CLI (5} שלבים)}

\begin{codebox}\begin{english}\codefont\footnotesize\color{cfg}\setlength{\parindent}{0pt}
\textcolor{ccom}{\texthebrew{! ===== שלב \enLR{0: ACL "}מעניין" – איזו תעבורה להצפין =====}}\\
R1(config)\# \textcolor{cbold}{\textbf{access-list 100 permit ip 10.1.1.0 0.0.0.255 10.2.2.0 0.0.0.255}}\\
~\\
\textcolor{ccom}{\texthebrew{! ===== שלב \enLR{1}: מדיניות \enLR{IKE Phase 1 (ISAKMP)} =====}}\\
R1(config)\# \textcolor{cbold}{\textbf{crypto isakmp policy 10}}\\
R1(config-isakmp)\# \textcolor{cbold}{\textbf{encryption aes 256}}\\
R1(config-isakmp)\# \textcolor{cbold}{\textbf{hash sha}}\\
R1(config-isakmp)\# \textcolor{cbold}{\textbf{authentication pre-share}}\\
R1(config-isakmp)\# \textcolor{cbold}{\textbf{group 14}}                 \textcolor{ccom}{\texthebrew{! \enLR{DH group 14 (2048-bit)}}}\\
R1(config-isakmp)\# lifetime 3600\\
R1(config-isakmp)\# exit\\
\textcolor{ccom}{\texthebrew{! מפתח משותף מראש + כתובת השכן}}\\
R1(config)\# \textcolor{cbold}{\textbf{crypto isakmp key S3cr3tPSK address 200.2.2.2}}\\
~\\
\textcolor{ccom}{\texthebrew{! ===== שלב \enLR{2: transform set} – הגנת \enLR{Phase 2} =====}}\\
R1(config)\# \textcolor{cbold}{\textbf{crypto ipsec transform-set TSET esp-aes 256 esp-sha-hmac}}\\
R1(cfg-crypto-trans)\# mode tunnel\\
R1(cfg-crypto-trans)\# exit\\
~\\
\textcolor{ccom}{\texthebrew{! ===== שלב \enLR{3: crypto map} – מקשר הכול =====}}\\
R1(config)\# \textcolor{cbold}{\textbf{crypto map CMAP 10 ipsec-isakmp}}\\
R1(config-crypto-map)\# \textcolor{cbold}{\textbf{set peer 200.2.2.2}}\\
R1(config-crypto-map)\# \textcolor{cbold}{\textbf{set transform-set TSET}}\\
R1(config-crypto-map)\# \textcolor{cbold}{\textbf{match address 100}}          \textcolor{ccom}{\texthebrew{! ה-\enLR{ACL} המעניין}}\\
R1(config-crypto-map)\# set pfs group14\\
R1(config-crypto-map)\# exit\\
~\\
\textcolor{ccom}{\texthebrew{! ===== שלב \enLR{4}: החלה על הממשק החיצוני =====}}\\
R1(config)\# interface g0/0\\
R1(config-if)\# \textcolor{cbold}{\textbf{crypto map CMAP}}\\
\textcolor{ccom}{\texthebrew{! \enLR{(R2} מוגדר במראה: כתובות ו-\enLR{ACL} הפוכים, אותה מדיניות ואותו \enLR{PSK)}}}
\end{english}\end{codebox}

\begin{mistakebox}\boxtitle{טעויות נפוצות ב-\enLR{IPsec}}\par  • מדיניות \enLR{Phase 1} לא תואמת בין הצדדים \enLR{(encryption/hash/group/auth} שונים) ⇒ המנהרה לא קמה. \textbf{שני הצדדים חייבים להסכים.}\newline  • \enLR{PSK} שונה משני הצדדים.\newline  • ה-\enLR{ACL "}המעניין" לא מראה \enLR{(R1: 10.1}→\enLR{10.2; R2} חייב \enLR{10.2}→\enLR{10.1).} אם לא הפוך – תעבורה לא נכנסת למנהרה.\newline  • שכחו \code{crypto map} על הממשק – הכול מוגדר אבל כלום לא קורה.\newline  • יש \enLR{NAT} שמתרגם את התעבורה המעניינת \emph{לפני} ההצפנה ⇒ ה-\enLR{ACL} לא תואם. פתרון: \enLR{NAT exemption (ACL} עם \enLR{deny} לתעבורת ה-\enLR{VPN).}\newline  • המנהרה קמה רק כשיש "תעבורה מעניינת" – \enLR{ping} אחד לא מספיק לפעמים; שלחו כמה.\end{mistakebox}

\subheading{וריפיקציה ופתרון תקלות}

\begin{codebox}\begin{english}\codefont\footnotesize\color{cfg}\setlength{\parindent}{0pt}
R1\# \textcolor{cbold}{\textbf{show crypto isakmp sa}}          \textcolor{ccom}{\texthebrew{! \enLR{Phase 1} – מצב \enLR{QM\_IDLE} = הצליח}}\\
R1\# \textcolor{cbold}{\textbf{show crypto ipsec sa}}           \textcolor{ccom}{\texthebrew{! \enLR{Phase 2} – ספירת חבילות \enLR{encrypted/decrypted}}}\\
R1\# \textcolor{cbold}{\textbf{show crypto map}}\\
R1\# \textcolor{cbold}{\textbf{show crypto session}}\\
R1\# debug crypto isakmp\\
R1\# debug crypto ipsec
\end{english}\end{codebox}

\begin{faqbox}\boxtitle{שאלת תלמיד: "איך אני יודע שהמנהרה עובדת ולא סתם ה-\enLR{ping} עובר רגיל?"}\par  ב-\code{show crypto ipsec sa} תראו את המונים \code{\#pkts encrypt} ו-\code{\#pkts decrypt} עולים כשאתם שולחים תעבורה בין הרשתות. אם הם עולים – התעבורה עוברת במנהרה מוצפנת. אם ה-\enLR{ping} עובר אבל המונים \enLR{0} – התעבורה עוברת \emph{מחוץ} למנהרה (בעיה ב-\enLR{ACL} המעניין או ב-\enLR{NAT).} זה הכלי הכי חשוב לדיבוג.\end{faqbox}

\sectionheading{\enLR{8.7 SSL/TLS VPN}}

\begin{xltabular}{\linewidth}{|>{\setRTL\raggedleft\arraybackslash}X|>{\setRTL\raggedleft\arraybackslash}X|>{\setRTL\raggedleft\arraybackslash}X|}
\hline
\rowcolor{hdrbg}\textbf{\enLR{SSL VPN}} & \textbf{\enLR{IPsec VPN}} & \textbf{} \\
\hline
\endhead
\enLR{4}–\enLR{7} (מעל \enLR{TCP 443)} & \enLR{3} (רשת) & שכבה \\
\hline
\textbf{\enLR{Clientless}} – רק דפדפן (או לקוח קל \enLR{AnyConnect)} & צריך תוכנת לקוח מותקנת ומוגדרת & לקוח \\
\hline
עובר בקלות – זה "רק \enLR{HTTPS"} על \enLR{443} & לעיתים חסום (פרוטוקולים \enLR{50/51, UDP 500)} & מעבר חומות אש/\enLR{NAT} \\
\hline
אפליקציות \enLR{web} ספציפיות \enLR{(portal)} או מלא \enLR{(tunnel)} & גישת רשת מלאה & גישה \\
\hline
גישה מרחוק ממחשב לא מנוהל, ספקים, ניידים & \enLR{site-to-site}, גישה קבועה & מתאים ל \\
\hline
\end{xltabular}

\enLR{SSL VPN} משתמש באותו \enLR{TLS} מפרק \enLR{7} (תעודות, \enLR{DH, AES).} זו הסיבה שהוא "פשוט עובד" מכל מקום – חומות אש כמעט תמיד מאפשרות \enLR{443.}\par

\sectionheading{\enLR{8.8} ציוד וארכיטקטורה}

\begin{itemize}
\item \textbf{\enLR{Cisco ASA 5500}} – מכשיר האבטחה המרכזי: חומת אש + \enLR{VPN (IPsec} ו-\enLR{SSL) + IPS.} יורשו: \enLR{Firepower.}
\item \textbf{\enLR{Cisco PIX 500}} – חומת האש/\enLR{VPN} ההיסטורית (הוחלפה ב-\enLR{ASA)}; מוזכרת בתכנית.
\item \textbf{\enLR{Cisco VPN 3000 Concentrator}} – ריכוז \enLR{VPN} מרחוק (הופסק); מוזכר בתכנית.
\item \textbf{\enLR{SOHO}} – נתב קטן עם \enLR{VPN} מובנה לעסק/בית.
\item האצת חומרה (בתכנית): \textbf{\enLR{VAC+, SPA, AIM}} – מודולי הצפנה שמורידים עומס מה-\enLR{CPU.} היום מובנה במעבד.
\item \textbf{\enLR{Cisco AnyConnect}} – לקוח ה-\enLR{VPN} המודרני \enLR{(IPsec/SSL)} למחשב ולנייד.
\end{itemize}

\sectionheading{\enLR{8.9} מדריך פקודות מרוכז – פרק \enLR{8}}

\begin{xltabular}{\linewidth}{|>{\setRTL\raggedleft\arraybackslash}X|>{\setRTL\raggedleft\arraybackslash}X|}
\hline
\rowcolor{hdrbg}\textbf{שלב} & \textbf{פקודה} \\
\hline
\endhead
\enLR{GRE} & \code{interface tunnel 0} + \code{tunnel source/destination} \\
\hline
\enLR{IKE Phase 1} & \code{crypto isakmp policy N} \enLR{(encryption/hash/authentication/group)} \\
\hline
\enLR{PSK} & \code{crypto isakmp key K address PEER} \\
\hline
הגנת \enLR{Phase 2} & \code{crypto ipsec transform-set NAME esp-aes esp-sha-hmac} \\
\hline
תעבורה מעניינת & \code{access-list 100 permit ip SRC DST} \\
\hline
קישור & \code{crypto map NAME N ipsec-isakmp} \enLR{(set peer / transform-set / match address)} \\
\hline
החלה & \code{crypto map NAME} (על ממשק) \\
\hline
וריפיקציה \enLR{Phase 1} & \code{show crypto isakmp sa} \\
\hline
וריפיקציה \enLR{Phase 2} (מונים) & \code{show crypto ipsec sa} \\
\hline
\end{xltabular}

\sectionheading{\enLR{8.10} תרגילים לתלמידים}

\begin{exbox}\boxtitle{תרגיל \enLR{1} – מתאים לאיזה \enLR{VPN}?}\par  לכל תרחיש: \enLR{Site-to-Site} או \enLR{Remote Access}; ו-\enLR{IPsec} או \enLR{SSL.}\newline  א. חיבור קבוע בין \enLR{3} סניפים. ב. עובד שמתחבר מבית קפה מהלפטופ. ג. ספק חיצוני שצריך לגשת לאפליקציית \enLR{web} אחת ממחשב לא מנוהל. ד. שרתי שני מוקדי נתונים שצריכים להעביר גיבויים. \par\vskip2pt\hrule\vskip3pt\textbf{}א. \enLR{Site-to-Site / IPsec} ב. \enLR{Remote Access / IPsec} או \enLR{SSL (AnyConnect)} ג. \enLR{Remote Access / SSL clientless} ד. \enLR{Site-to-Site / IPsec}\end{exbox}

\begin{exbox}\boxtitle{תרגיל \enLR{2} – \enLR{AH} או \enLR{ESP}?}\par  א. צריך גם הצפנה. ב. הצדדים מאחורי \enLR{NAT.} ג. רוצים רק לוודא שהחבילה לא שונתה בלי להצפין. מה נבחר בכל מקרה, ולמה בפועל תמיד \enLR{ESP}? \par\vskip2pt\hrule\vskip3pt\textbf{}א. \enLR{ESP} (רק הוא מצפין). ב. \enLR{ESP (AH} נשבר עם \enLR{NAT).} ג. תיאורטית \enLR{AH}, אבל בפועל \enLR{ESP} במצב \enLR{authentication-only.} תמיד \enLR{ESP} כי \enLR{AH} לא מצפין ולא עובד עם \enLR{NAT.}\end{exbox}

\begin{exbox}\boxtitle{תרגיל \enLR{3} – זהו את התקלה}\par  המנהל הגדיר \enLR{IPsec}, \code{show crypto isakmp sa} מראה \enLR{QM\_IDLE}, אבל \enLR{ping} בין הרשתות נכשל ו-\code{show crypto ipsec sa} מראה \enLR{0 encrypted.} מה כנראה הבעיה? \par\vskip2pt\hrule\vskip3pt\textbf{}\enLR{Phase 1} הצליח (יש \enLR{ISAKMP SA)}, אבל \enLR{Phase 2} לא מעביר תעבורה: כנראה ה-\enLR{ACL} המעניין לא תואם/לא הפוך בין הצדדים, או ש-\enLR{NAT} מתרגם את התעבורה לפני ההצפנה. בודקים \code{match address} ו-\enLR{NAT exemption.}\end{exbox}

\begin{exbox}\boxtitle{תרגיל \enLR{4} – מיפוי לפרק \enLR{7}}\par  בקונפיג: \code{encryption aes 256}, \code{hash sha}, \code{group 14}, \code{authentication pre-share}. לכל אחד – איזו מטרת \enLR{CIA}/תפקיד הוא ממלא ואיזה מושג מפרק \enLR{7} הוא? \par\vskip2pt\hrule\vskip3pt\textbf{}\enLR{aes 256} → סודיות (סימטרי). \enLR{sha} → שלמות \enLR{(HMAC). group 14} → החלפת מפתחות \enLR{(Diffie-Hellman). pre-share} → אימות הצדדים (חלופה: תעודות \enLR{RSA/PKI).}\end{exbox}

\sectionheading{\enLR{8.11} שאלות בסגנון בגרות}

\begin{exambox}\boxtitle{שאלה \enLR{1}}\par  איזה סוג תקשורת נועדה חבילת הפרוטוקולים \enLR{IPsec} לאבטח? \enLR{SSH / TCP/IP / TELNET / UDP} \par\vskip2pt\hrule\vskip3pt\textbf{}\textbf{\enLR{TCP/IP}} (שכבה \enLR{3} – ולכן כל התעבורה מעליו).\end{exambox}

\begin{exambox}\boxtitle{שאלה \enLR{2}}\par  מהו ההבדל בין \enLR{AH} ל-\enLR{ESP}, ומדוע בפועל משתמשים ב-\enLR{ESP}? \par\vskip2pt\hrule\vskip3pt\textbf{}\enLR{AH} – שלמות ואימות בלבד, ללא הצפנה. \enLR{ESP} – שלמות, אימות \textbf{והצפנה}. משתמשים ב-\enLR{ESP} כי רק הוא נותן סודיות, וגם עובד עם \enLR{NAT.}\end{exambox}

\begin{exambox}\boxtitle{שאלה \enLR{3}}\par  מה תפקיד \enLR{IKE Phase 1} לעומת \enLR{Phase 2}? \par\vskip2pt\hrule\vskip3pt\textbf{}\enLR{Phase 1} – בונה ערוץ ניהול מאובטח \enLR{(ISAKMP SA): DH}, אימות הצדדים, הסכמה על מדיניות. \enLR{Phase 2} – בונה בתוכו את מנהרות הנתונים \enLR{(IPsec SA)} שדרכן עוברת התעבורה.\end{exambox}

\begin{exambox}\boxtitle{שאלה \enLR{4}}\par  מדוע \enLR{GRE} לבדו אינו מספיק לחיבור מאובטח בין סניפים, וכיצד פותרים? \par\vskip2pt\hrule\vskip3pt\textbf{}\enLR{GRE} עוטף אך אינו מצפין – התעבורה גלויה. פתרון: \enLR{GRE over IPsec} – \enLR{GRE} לגמישות \enLR{(multicast/routing), IPsec} להצפנה.\end{exambox}

\begin{exambox}\boxtitle{שאלה \enLR{5}}\par  מתי נעדיף \enLR{SSL VPN} על פני \enLR{IPsec VPN}? \par\vskip2pt\hrule\vskip3pt\textbf{}גישה מרחוק ממחשב לא מנוהל / דרך דפדפן בלבד \enLR{(clientless)}, וכשצריך לעבור חומות אש/\enLR{NAT} בקלות – \enLR{SSL} על \enLR{443} עובר כמעט תמיד.\end{exambox}

\begin{exambox}\boxtitle{שאלה \enLR{6}}\par  השלימו את הפקודה שמגדירה \enLR{transform set} עם \enLR{AES} ו-\enLR{SHA}: \code{crypto ipsec transform-set TSET \_\_\_\_\_\_ \_\_\_\_\_\_} \par\vskip2pt\hrule\vskip3pt\textbf{}\code{esp-aes esp-sha-hmac}\end{exambox}

\sectionheading{\enLR{8.12} מעבדה \enLR{(3} שעות) – \enLR{IPsec Site-to-Site} ב-\enLR{Packet Tracer}}

\begin{enumerate}
\item טופולוגיה: \enLR{LAN1 (10.1.1.0/24)} — \enLR{R1} — "אינטרנט" \enLR{(R-ISP)} — \enLR{R2} — \enLR{LAN2 (10.2.2.0/24).} ודאו ניתוב מקצה לקצה \enLR{(ping} עובר לפני \enLR{VPN).}
\item ב-\enLR{R1} ו-\enLR{R2} הגדירו את \enLR{5} השלבים מ-\enLR{8.6} (במראה). \enLR{PSK} זהה, מדיניות זהה, \enLR{ACL} הפוך.
\item מ-\enLR{PC} ב-\enLR{LAN1} עשו \enLR{ping} ל-\enLR{PC} ב-\enLR{LAN2} – יופעל "\enLR{interesting traffic".}
\item \code{show crypto isakmp sa} ⇒ \enLR{QM\_IDLE.} \code{show crypto ipsec sa} ⇒ מונים עולים.
\item ב-\enLR{Simulation mode} לכדו חבילה על קו ה"אינטרנט" – ראו שהיא \enLR{ESP} ולא ניתן לקרוא את התוכן.
\item \textbf{\enLR{GRE} (אופציונלי):} הקימו מנהרת \enLR{GRE} והריצו \enLR{OSPF} דרכה; ראו ש-\enLR{multicast} של \enLR{OSPF} עובר – מה שלא היה עובד ב-\enLR{IPsec} לבדו.
\end{enumerate}

\sectionheading{בנק שאלות שתלמידים שואלים – ותשובות מוכנות}

שאלות נפוצות בפרק ה-\enLR{VPN}, עם תשובות מוכנות.\par

\textbf{ש: "מה ההבדל בין ה-\enLR{VPN} שאנחנו לומדים ל-\enLR{VPN} מסחרי כמו \enLR{NordVPN}?"}\newline ת: \enLR{VPN} מסחרי נועד להסתיר את הזהות/המיקום שלכם ולעקוף חסימות. \enLR{VPN} \underline{ארגוני} (מה שאנחנו לומדים) מחבר בבטחה עובד או סניף אל רשת החברה. אותה טכנולוגיית מנהרה מוצפנת, מטרות שונות.\par

\textbf{ש: "אם \enLR{GRE} יוצר מנהרה, למה הוא לא מספיק לחיבור מאובטח?"}\newline ת: כי \enLR{GRE} רק \underline{עוטף} את החבילה – הוא \textbf{לא מצפין} כלום. כל מי שמאזין בדרך רואה את התוכן. לכן משלבים "\enLR{GRE over IPsec": GRE} נותן גמישות (מעביר \enLR{multicast}/ניתוב), \enLR{IPsec} נותן את ההצפנה.\par

\textbf{ש: "\enLR{IPsec} מצפין – אז מה זה \enLR{AH} ו-\enLR{ESP}?"}\newline ת: שני "אופנים" של \enLR{IPsec.} \textbf{\enLR{AH}} נותן שלמות ואימות בלבד – \underline{בלי הצפנה}. \textbf{\enLR{ESP}} נותן שלמות, אימות \underline{וגם הצפנה}. בפועל תמיד משתמשים ב-\enLR{ESP} (רק הוא מצפין, וגם עובד עם \enLR{NAT).}\par

\textbf{ש: "למה \enLR{IKE} עובד בשני שלבים? זה נשמע מסובך."}\newline ת: \enLR{Phase 1} בונה קודם \underline{ערוץ ניהול מאובטח} אחד (מבצע \enLR{Diffie-Hellman} ומאמת את הצדדים). רק בתוך הערוץ המוגן הזה, \enLR{Phase 2} בונה את מנהרות הנתונים בפועל. זה כמו קודם לבנות חדר סגור, ואז לנהל בו את המשא-ומתן.\par

\textbf{ש: "איך אני יודע שהמנהרה באמת עובדת ולא שה-\enLR{ping} סתם עובר רגיל?"}\newline ת: \code{show crypto ipsec sa} – מסתכלים על המונים \code{\#pkts encrypt/decrypt}. אם הם עולים כששולחים תעבורה, היא באמת עוברת מוצפנת. אם ה-\enLR{ping} עובר אבל המונים \enLR{0} – התעבורה עוברת \underline{מחוץ} למנהרה (בעיה ב-\enLR{ACL} המעניין או ב-\enLR{NAT).}\par

\textbf{ש: "מתי עדיף \enLR{SSL VPN} על \enLR{IPsec}?"}\newline ת: כשצריך גישה מרחוק ממחשב \underline{לא מנוהל} או דרך דפדפן בלבד \enLR{(clientless)}, וכשרוצים מעבר קל דרך חומות אש – \enLR{SSL} רץ על פורט \enLR{443} שכמעט תמיד פתוח. \enLR{IPsec} מתאים יותר לחיבורים קבועים (סניף לסניף).\par

\textbf{ש: "המנהרה קמה \enLR{(QM\_IDLE)} אבל ה-\enLR{ping} נכשל. מה קורה?"}\newline ת: \enLR{Phase 1} הצליח (יש \enLR{ISAKMP SA)} אבל \enLR{Phase 2} לא מעביר תעבורה. הסיבות הנפוצות: ה-\enLR{ACL "}המעניין" לא \underline{הפוך} בין שני הצדדים, או ש-\enLR{NAT} מתרגם את התעבורה \emph{לפני} ההצפנה. בודקים \code{match address} ופטור מ-\enLR{NAT.}\par

\textbf{ש: "כל מה שהגדרתי ב-\enLR{IPsec} נראה נכון אבל המנהרה לא קמה בכלל. מה לבדוק?"}\newline ת: מדיניות \enLR{Phase 1} חייבת להיות \underline{זהה} בשני הצדדים \enLR{(encryption/hash/DH group/authentication)} וגם ה-\enLR{PSK} זהה. אי-התאמה של אחד מהם מונעת מהמנהרה לקום. בודקים עם \code{show crypto isakmp sa} ו-\code{debug crypto isakmp}.\par

\sectionheading{מילון מונחים – הגדרות}

הגדרות תמציתיות של המונחים המרכזיים בפרק. שימושי גם כדף עזר לבחינה.\par

\begin{xltabular}{\linewidth}{|>{\setRTL\raggedleft\arraybackslash}X|>{\setRTL\raggedleft\arraybackslash}X|}
\hline
\rowcolor{hdrbg}\textbf{הגדרה} & \textbf{מונח} \\
\hline
\endhead
רשת פרטית וירטואלית – מנהרה מוצפנת דרך רשת ציבורית. & \textbf{\enLR{VPN}} \\
\hline
\enLR{VPN} בין שני שערים (סניפים), שקוף למשתמשים. & \textbf{\enLR{Site-to-Site}} \\
\hline
\enLR{VPN} של משתמש יחיד (עובד מהבית) אל שער החברה. & \textbf{\enLR{Remote Access}} \\
\hline
פרוטוקול מנהור שעוטף חבילות (כולל \enLR{multicast)} – אך אינו מצפין. & \textbf{\enLR{GRE}} \\
\hline
מסגרת פרוטוקולים לאבטחת תעבורה בשכבה \enLR{3 (TCP/IP).} & \textbf{\enLR{IPsec}} \\
\hline
פרוטוקול \enLR{IPsec}: שלמות ואימות בלבד – ללא הצפנה. & \textbf{\enLR{AH}} \\
\hline
פרוטוקול \enLR{IPsec}: שלמות, אימות והצפנה; הבחירה המקובלת. & \textbf{\enLR{ESP}} \\
\hline
\enLR{IPsec} בין מארחים \enLR{(transport)} או בין שערים עם עטיפה \enLR{(tunnel).} & \textbf{\enLR{Transport / Tunnel mode}} \\
\hline
פרוטוקול שמקים את מנהרת ה-\enLR{IPsec} (מבצע \enLR{DH} ואימות). & \textbf{\enLR{IKE}} \\
\hline
\enLR{IKE}: בונה ערוץ ניהול מאובטח \enLR{(1)} ואז את מנהרות הנתונים \enLR{(2).} & \textbf{\enLR{Phase 1 / Phase 2}} \\
\hline
הסכם חד-כיווני המגדיר איך לאבטח זרם תעבורה, עם מזהה \enLR{SPI.} & \textbf{\enLR{SA (Security Association)}} \\
\hline
מפתח משותף מראש לאימות הצדדים ב-\enLR{IKE} (חלופה לתעודות). & \textbf{\enLR{PSK (Pre-Shared Key)}} \\
\hline
צירוף אלגוריתמי ההגנה של \enLR{Phase 2} (למשל \enLR{esp-aes esp-sha-hmac).} & \textbf{\enLR{Transform set}} \\
\hline
אובייקט שמקשר \enLR{peer, transform set} ו-\enLR{ACL '}מעניין', ומוחל על ממשק. & \textbf{\enLR{Crypto map}} \\
\hline
התעבורה (מוגדרת ב-\enLR{ACL)} שאמורה לעבור במנהרה המוצפנת. & \textbf{\enLR{interesting traffic}} \\
\hline
\enLR{VPN} מעל \enLR{TCP 443}; לרוב \enLR{clientless} (דרך דפדפן); עובר חומות אש בקלות. & \textbf{\enLR{SSL/TLS VPN}} \\
\hline
מכשיר אבטחה המשלב חומת אש, \enLR{VPN} ו-\enLR{IPS.} & \textbf{\enLR{Cisco ASA}} \\
\hline
\end{xltabular}


\clearpage
\phantomsection\addcontentsline{toc}{section}{פרק \enLR{9} – ניהול אבטחת רשת מתקדם}

\sloppy
\chapterheading{פרק \enLR{9} – ניהול אבטחת רשת מתקדם}{\small\color{subgray}\enLR{11} שעות עיוני + \enLR{3} מעשי · שבועות \enLR{27}–\enLR{29}\par}\vskip4pt

\begin{metabox}\textbf{מטרות:} עיצוב רשת מאובטחת · ניתוח פגיעות · ניהול סיכונים · הגנה עצמית · \enLR{Nessus/Nmap}\\ \textbf{בבגרות:} בקרות מנהלי/פיזי/לוגי, \enLR{Nmap}, אתיקה וחוק\end{metabox}

\begin{defbox}\boxtitle{לפני שמתחילים – מהניהול הטכני לניהול הארגוני (למי שאין רקע)}\par  עד כה למדנו \underline{כלים} (חומת אש, הצפנה, \enLR{VPN).} פרק זה עוסק ב\underline{ניהול} – ההחלטות והנהלים סביב הכלים:  \textbf{סיכון \enLR{(Risk)}} – ההסתברות שמשהו רע יקרה, כפול הנזק. אבטחה = ניהול סיכונים, לא ביטולם המוחלט (שבלתי אפשרי). \textbf{מדיניות \enLR{(Policy)}} – מסמך שקובע מה מותר ומה אסור בארגון (למשל "אסור להתחבר לרשת הארגון מ-\enLR{Wi-Fi} ציבורי"). \textbf{בקרה \enLR{(Control)}} – כל אמצעי שמקטין סיכון. יש שלושה סוגים: \textbf{מנהלית} (נהלים), \textbf{פיזית} (מנעולים), \textbf{לוגית} (סיסמאות, הצפנה). \textbf{גיבוי \enLR{(Backup)}} – עותק של המידע שנשמר בנפרד, כדי לשחזר אחרי תקלה/מתקפה. \textbf{בדיקת חדירה \enLR{(Penetration Test)}} – לשכור "האקר טוב" שינסה לפרוץ בכוונה, כדי למצוא חולשות לפני שהרעים ימצאו.  \textbf{הרעיון:} הפרק מחבר את כל הקורס – איך משלבים את כל הכלים לכדי תוכנית אבטחה שלמה לארגון, ומה החוק והאתיקה דורשים.\end{defbox}

\sectionheading{\enLR{9.1} מהניתוק הטכני לתמונה הניהולית}

שמונת הפרקים הקודמים היו על \emph{כלים} \enLR{(ACL, IPS, VPN}…). פרק \enLR{9} הוא על \emph{המערכת הניהולית} שמחליטה איזה כלי, איפה, וכמה – כי אבטחה מושלמת עולה אינסוף, ואבטחה אפסית עולה בפריצה. זהו איזון של סיכון מול עלות. זה גם הפרק שקושר את הקורס לעולם המקצועי \enLR{(SOC, CISO}, ביקורת, חוק).\par

\sectionheading{\enLR{9.2} ניהול סיכונים \enLR{(Risk Management)}}

\begin{defbox}\boxtitle{הנוסחה}\par  \textbf{סיכון = איום × פגיעות × ערך הנכס}. אם אחד מהם אפס – אין סיכון. אי אפשר לבטל איומים (הם קיימים), אבל אפשר להקטין פגיעות (בקרות) ולהקטין את ההשפעה.\end{defbox}

\subheading{תהליך ניהול הסיכונים}

\begin{enumerate}
\item \textbf{זיהוי נכסים} – מה יש לנו ומה ערכו (מידע, שרתים, מוניטין).
\item \textbf{הערכת איומים ופגיעויות} – מה עלול לקרות (סריקות, ראיונות, בדיקות – \enLR{9.5).}
\item \textbf{ניתוח סיכונים} – \textbf{כמותי} (בכסף: \enLR{ALE = SLE} × \enLR{ARO} – ההפסד השנתי הצפוי) או \textbf{איכותי} (מטריצה גבוה/בינוני/נמוך).
\item \textbf{טיפול בסיכון} – ראו להלן.
\item \textbf{ניטור וחזרה} – סיכון משתנה עם הזמן.
\end{enumerate}

\subheading{ארבע דרכי טיפול בסיכון (נושא מרכזי)}

\begin{xltabular}{\linewidth}{|>{\setRTL\raggedleft\arraybackslash}X|>{\setRTL\raggedleft\arraybackslash}X|>{\setRTL\raggedleft\arraybackslash}X|}
\hline
\rowcolor{hdrbg}\textbf{דוגמה} & \textbf{משמעות} & \textbf{אסטרטגיה} \\
\hline
\endhead
התקנת \enLR{IPS}; הפסקת שירות ישן ופגיע & בקרות שמקטינות סיכון, או הפסקת הפעילות המסוכנת & \textbf{הפחתה / הימנעות} \enLR{(Mitigate/Avoid)} \\
\hline
ביטוח סייבר; מיקור חוץ לענן & מעבירים את הנטל לגורם אחר & \textbf{העברה} \enLR{(Transfer)} \\
\hline
לא מצפינים מסך פנימי בעל ערך נמוך & הסיכון נמוך/עלות הטיפול גבוהה – חיים איתו במודע & \textbf{קבלה} \enLR{(Accept)} \\
\hline
רשת נפרדת למערכת \enLR{legacy; DMZ; air-gap} & מפרידים את המערכת המסוכנת & \textbf{בידוד} \enLR{(Isolation)} \\
\hline
\end{xltabular}

\sectionheading{\enLR{9.3} עקרונות עיצוב רשת מאובטחת}

\begin{itemize}
\item \textbf{הגנה לעומק} \enLR{(Defense in Depth)} – שכבות עצמאיות; כשל באחת לא מפיל הכול.
\item \textbf{מינימום הרשאות} \enLR{(Least Privilege)} – כל אחד מקבל בדיוק מה שצריך, לא יותר.
\item \textbf{הפרדת תפקידים} \enLR{(Separation of Duties)} – פעולה קריטית דורשת שני אנשים; מי שמאשר תשלום ≠ מי שמבצע. מונע הונאה פנימית.
\item \textbf{סבב תפקידים} \enLR{(Job Rotation)} – חושף הונאות, מפחית תלות באדם יחיד.
\item \textbf{\enLR{Fail-Secure}} – במקרה כשל, ברירת המחדל היא לחסום \enLR{(deny by default).}
\item \textbf{\enLR{Zero Trust}} – "לעולם אל תסמוך, תמיד תאמת"; אין "פנים מהימן".
\item \textbf{\enLR{KISS} / פשטות} – מערכת מורכבת קשה לאבטח; מורכבות היא אויב האבטחה.
\item \textbf{מגוון הגנות} \enLR{(Diversity)} – ספקים שונים בשכבות שונות, כדי שחולשה אחת לא תפיל הכול.
\end{itemize}

\begin{exambox}\boxtitle{בבגרות (שאלה \enLR{6}ב) – שלוש קטגוריות בקרה}\par  \textbf{מנהלית}: כתיבת מדיניות, הנחיות, נהלים, תקנים, איסור עבודה ב-\enLR{Zoom} מרשת לא מאובטחת.\newline  \textbf{פיזית}: הגנה על חדר שרתים וארונות תקשורת.\newline  \textbf{לוגית}: שימוש ברשימות גישה, סיסמאות, מערכות למניעת גישה.\end{exambox}

חלוקה משלימה לפי \emph{מתי} הבקרה פועלת: \textbf{מונעת} \enLR{(Preventive} – מנעול, הצפנה), \textbf{מגלה} \enLR{(Detective} – \enLR{IDS}, מצלמה, לוג), \textbf{מתקנת} \enLR{(Corrective} – גיבוי, \enLR{patch)}, \textbf{מרתיעה} \enLR{(Deterrent} – שלט אזהרה), \textbf{משחזרת} \enLR{(Recovery} – \enLR{DRP).}\par

\sectionheading{\enLR{9.4} מדיניות, נהלים ותקנים}

\begin{xltabular}{\linewidth}{|>{\setRTL\raggedleft\arraybackslash}X|>{\setRTL\raggedleft\arraybackslash}X|}
\hline
מסמך-על: \emph{מה} ולמה. "כל גישה מרחוק דרך \enLR{VPN".} & \textbf{מדיניות \enLR{(Policy)}} \\
\hline
חובה מדידה: "סיסמה ≥ \enLR{12} תווים, \enLR{MFA".} & \textbf{תקן \enLR{(Standard)}} \\
\hline
המלצה גמישה. & \textbf{הנחיה \enLR{(Guideline)}} \\
\hline
צעד-אחר-צעד: \emph{איך} מגדירים \enLR{VPN.} & \textbf{נוהל \enLR{(Procedure)}} \\
\hline
\end{xltabular}

מסמכים חשובים: \textbf{\enLR{AUP}} \enLR{(Acceptable Use Policy} – מה מותר לעובד לעשות עם הציוד), מדיניות סיסמאות, מדיניות גיבוי, מדיניות תגובה לאירועים, סיווג מידע (סודי/פנימי/ציבורי).\par

\sectionheading{\enLR{9.5} בדיקות אבטחה: הערכה, סריקה ובדיקת חדירה}

\begin{xltabular}{\linewidth}{|>{\setRTL\raggedleft\arraybackslash}X|>{\setRTL\raggedleft\arraybackslash}X|}
\hline
\rowcolor{hdrbg}\textbf{מה זה} & \textbf{סוג בדיקה} \\
\hline
\endhead
בדיקת עמידה במדיניות/תקן \enLR{(ISO 27001)} – מסמכים ותצורות & \textbf{\enLR{Security Audit}} \\
\hline
סריקה אוטומטית לזיהוי פגיעויות ידועות \enLR{(Nessus)} & \textbf{\enLR{Vulnerability Assessment}} \\
\hline
ניסיון פריצה מבוקר – "לחשוב כמו תוקף" ולהוכיח ניצול & \textbf{\enLR{Penetration Test}} \\
\hline
בדיקה שיטתית שהבקרות עובדות כמתוכנן (מוזכר בתכנית) & \textbf{\enLR{ST\&E}} \enLR{(Security Test \& Evaluation)} \\
\hline
\end{xltabular}

\subheading{סוגי \enLR{pentest}}

\textbf{\enLR{Black box}} (התוקף לא יודע כלום – כמו תוקף חיצוני), \textbf{\enLR{White box}} (יודע הכול – יעיל, מהיר), \textbf{\enLR{Gray box}} (ידע חלקי – כמו עובד/משתמש). שלבים: תכנון והרשאה → סיור → סריקה → ניצול → דוח.\par

\begin{mistakebox}\boxtitle{חובה: אישור בכתב \enLR{(Scope \& Authorization)}}\par  בדיקת חדירה ללא אישור כתוב = עבירה פלילית. חייבים: היקף מוגדר (אילו כתובות), חלון זמן, "\enLR{get out of jail" letter}, ואיש קשר. זה גם מה שמלמדים לתלמידים לגבי הכלים בהמשך.\end{mistakebox}

\sectionheading{\enLR{9.6} כלי סריקת פגיעויות (נשאל בבגרות)}

\begin{xltabular}{\linewidth}{|>{\setRTL\raggedleft\arraybackslash}X|>{\setRTL\raggedleft\arraybackslash}X|}
\hline
\rowcolor{hdrbg}\textbf{מה עושה} & \textbf{כלי} \\
\hline
\endhead
\textbf{גילוי, ניתוח ובקרה ברשתות תקשורת}: מיפוי מארחים חיים, סריקת פורטים פתוחים, זיהוי שירותים וגרסאות (-\enLR{sV)}, זיהוי מערכת הפעלה (-\enLR{O).} זו התשובה בבגרות. \enLR{Zenmap} = ה-\enLR{GUI.} & \textbf{\enLR{Nmap}} \enLR{(Network Mapper)} \\
\hline
סורק פגיעויות: מריץ אלפי בדיקות מול מארח ומדווח על חולשות ידועות \enLR{(CVE)} עם דירוג חומרה והמלצות תיקון. (התכנית כותבת "\enLR{NESUS".)} & \textbf{\enLR{Nessus}} \\
\hline
סורק פורטים ותיקוני \enLR{TCP/UDP} למערכות \enLR{Windows} (ישן; מוזכר בתכנית). & \textbf{\enLR{SuperScan}} \\
\hline
חלופה חינמית ל-\enLR{Nessus} & \enLR{OpenVAS} \\
\hline
ניתוח תעבורה \enLR{(packet analyzer)} & \enLR{Wireshark} \\
\hline
מסגרת לניצול פגיעויות (בדיקות חדירה) & \enLR{Metasploit} \\
\hline
הפצת לינוקס עם כל כלי הבדיקה מותקנים & \enLR{Kali Linux} \\
\hline
\end{xltabular}

\begin{codebox}\begin{english}\codefont\footnotesize\color{cfg}\setlength{\parindent}{0pt}
\textcolor{ccom}{\texthebrew{! דוגמאות \enLR{Nmap} (רק ברשת מעבדה מבודדת ובאישור!)}}\\
nmap 192.168.1.0/24              \textcolor{ccom}{\texthebrew{\# אילו מארחים חיים ואילו פורטים}}\\
nmap -sV 192.168.1.10            \textcolor{ccom}{\texthebrew{\# גרסאות שירותים}}\\
nmap -O 192.168.1.10             \textcolor{ccom}{\texthebrew{\# ניחוש מערכת הפעלה}}\\
nmap -sS -p 1-1000 192.168.1.10  \textcolor{ccom}{\texthebrew{\# \enLR{SYN scan}}}
\end{english}\end{codebox}

\begin{storybox}\boxtitle{סיפור מהחיים: \enLR{Equifax} – פגיעות אחת שלא תוקנה \enLR{(2017)}}\par  דליפת המידע של \enLR{Equifax} חשפה נתונים של \enLR{147} מיליון אמריקאים (מספרי ביטוח לאומי, תאריכי לידה). הסיבה: פגיעות ידועה ב-\enLR{Apache Struts (CVE-2017-5638)} עם תיקון זמין \emph{חודשיים} קודם. סורק פגיעויות היה מוצא אותה תוך דקות – אבל התהליך של \enLR{Equifax} לא סרק את השרת הזה. הלקח: ניהול פגיעויות הוא לא אירוע חד-פעמי אלא תהליך מתמשך \enLR{(scan} → \enLR{prioritize} → \enLR{patch} → \enLR{verify).} זה הלב של פרק \enLR{9.}\end{storybox}

\sectionheading{\enLR{9.7} הגנה עצמית ותגובה לאירועים}

"מערכות הגנה עצמית" בתכנית: היכולת של הרשת לזהות ולהגיב אוטומטית. סיסקו קראה לזה \textbf{\enLR{SDN}} \enLR{(Self-Defending Network)} בעבר – שילוב \enLR{IPS, NAC}, קורלציה ותגובה אוטומטית (בידוד מחשב נגוע). היום: \textbf{\enLR{SOAR}} \enLR{(Security Orchestration, Automation and Response)} ו-\textbf{\enLR{XDR}}.\par

\subheading{\enLR{Incident Response} – \enLR{6} שלבים \enLR{(SANS/NIST)}}

\begin{center}\fbox{\enLR{Preparation}} \;→\; \fbox{\enLR{Identification}} \;→\; \fbox{\enLR{Containment}} \;→\; \fbox{\enLR{Eradication}} \;→\; \fbox{\enLR{Recovery}} \;→\; \fbox{\enLR{Lessons Learned}}\end{center}

הכנה (נהלים, כלים) → זיהוי \enLR{(SIEM, IDS)} → בלימה (ניתוק מהרשת) → מיגור (הסרת הנוזקה) → שחזור (החזרה לפעילות) → הפקת לקחים. הצוות: \textbf{\enLR{CSIRT/CERT}}.\par

\sectionheading{\enLR{9.8} המשכיות עסקית והתאוששות מאסון}

\begin{xltabular}{\linewidth}{|>{\setRTL\raggedleft\arraybackslash}X|>{\setRTL\raggedleft\arraybackslash}X|>{\setRTL\raggedleft\arraybackslash}X|}
\hline
\rowcolor{hdrbg}\textbf{\enLR{DRP (Disaster Recovery Plan)}} & \textbf{\enLR{BCP (Business Continuity Plan)}} & \textbf{} \\
\hline
\endhead
שחזור מערכות ה-\enLR{IT} אחרי אסון & שהעסק ימשיך לתפקד בזמן אסון & מטרה \\
\hline
טכנולוגי (שרתים, נתונים) & כל העסק (אנשים, תהליכים, מבנים) & היקף \\
\hline
\end{xltabular}

\textbf{\enLR{BIA}} \enLR{(Business Impact Analysis)} מזהה תהליכים קריטיים ומגדיר שני יעדים: \textbf{\enLR{RTO}} \enLR{(Recovery Time Objective} – כמה זמן מותר להיות מושבתים) ו-\textbf{\enLR{RPO}} \enLR{(Recovery Point Objective} – כמה נתונים מותר לאבד = תדירות הגיבוי). \textbf{\enLR{WCS}} \enLR{(Worst Case Scenario} – בתכנית) = תכנון לתרחיש הגרוע ביותר.\par

\subheading{גיבויים ואתרי גיבוי}

\begin{itemize}
\item \textbf{כלל \enLR{3-2-1}}: \enLR{3} עותקים, \enLR{2} מדיות שונות, \enLR{1} מחוץ לאתר \enLR{(off-site).} קריטי נגד כופרה.
\item סוגים: מלא \enLR{(Full)}, הפרשי \enLR{(Differential)}, משלים \enLR{(Incremental).}
\item אתרי גיבוי: \textbf{\enLR{Hot site}} (פעיל, מיידי, יקר), \textbf{\enLR{Warm site}} (חלקי), \textbf{\enLR{Cold site}} (מבנה ריק, זול, איטי).
\item \textbf{בדיקת גיבוי} – גיבוי שלא נבדק שחזור ממנו = לא גיבוי. תרגלו שחזור!
\end{itemize}

\sectionheading{\enLR{9.9 SDLC} – פיתוח מאובטח}

\textbf{\enLR{SDLC}} \enLR{(Software/System Development Life Cycle)} – מחזור חיי פיתוח מערכת. אבטחה חייבת להיכנס \textbf{מההתחלה} ("\enLR{Security by Design")}, לא כטלאי בסוף – זול ויעיל פי כמה. שלבים: תכנון → ניתוח דרישות → עיצוב → מימוש → בדיקות → הטמעה → תחזוקה (וזה חוזר במעגל). התכנית מציינת "תוכנית מערך \enLR{SDLC"} ו"תחזוקה ותפעול \enLR{SDLC".}\par

ב-\textbf{\enLR{DevSecOps}} משלבים בדיקות אבטחה אוטומטיות בכל \enLR{commit.} עקרונות קוד מאובטח: בדיקת קלט (נגד \enLR{injection} ו-\enLR{buffer overflow} – פרק \enLR{1), least privilege}, ניהול סודות נכון, עדכון ספריות \enLR{(Log4Shell).}\par

\sectionheading{\enLR{9.10} אבטחת מערכות מגוונות ושילוב גורמים}

\begin{itemize}
\item \textbf{מייל} – \enLR{ESA/anti-spam, SPF/DKIM/DMARC} (נגד זיוף שולח – פרק \enLR{1 phishing).}
\item \textbf{טלפוניה \enLR{(VoIP)}} – \enLR{SRTP, voice VLAN} (פרק \enLR{6).}
\item \textbf{אפליקציות} – \enLR{WAF, OWASP Top 10.}
\item \textbf{אלחוט} – \enLR{WPA3-Enterprise, WIPS} (פרק \enLR{6).}
\item \textbf{שילוב גורמים} – לקוחות, שותפים, עובדים: לכל אחד רמת גישה ומדיניות \enLR{(least privilege}, פילוח).
\item ניהול מרכזי וחלוקת תפקידים \enLR{(SOC, NOC}, הפרדת תפקידים).
\end{itemize}

\sectionheading{\enLR{9.11} אתיקה, חוק ואחריות מקצועית}

\begin{mistakebox}\boxtitle{חובה ללמד – זו לא רק "המלצה"}\par  בישראל: \textbf{חוק המחשבים, התשנ"ה-\enLR{1995}} אוסר חדירה לחומר מחשב שלא כדין (סעיף \enLR{4)} – עד \enLR{3} שנות מאסר; שיבוש/הפרעה – עד \enLR{5} שנים. \textbf{חוק הגנת הפרטיות} ותקנות אבטחת מידע \enLR{(2017)} מחייבים ארגונים. באירופה: \textbf{\enLR{GDPR}} (קנסות עד \enLR{4\%} מהמחזור). הבהירו לתלמידים: הרצת \enLR{Nmap} או כלי פריצה על רשת שאינה שלכם וללא אישור בכתב – עבירה פלילית, גם "רק כדי לבדוק".\end{mistakebox}

\begin{itemize}
\item \textbf{\enLR{Hacker} אתי \enLR{(White hat)}} פועל תמיד עם אישור והיקף מוגדר.
\item קודים אתיים מקצועיים: \enLR{(ISC)}², \enLR{EC-Council.}
\item \textbf{\enLR{Responsible/Coordinated Disclosure}} – מצאתם פגיעות? מדווחים ליצרן ונותנים זמן לתקן לפני פרסום.
\item דילמות: פרטיות מול ניטור עובדים; חשיפת פגיעות מול ניצולה.
\end{itemize}

\sectionheading{\enLR{9.12} תרגילים לתלמידים}

\begin{exbox}\boxtitle{תרגיל \enLR{1} – סיווג בקרות}\par  סווגו כל בקרה (מנהלית/פיזית/לוגית) וגם (מונעת/מגלה/מתקנת):\newline  א. מצלמת אבטחה בכניסה. ב. נוהל גיבוי יומי. ג. חומת אש. ד. הדרכת מודעות לעובדים. ה. גיבוי ששוחזר אחרי כופרה. ו. מנעול ביומטרי לחדר שרתים. \par\vskip2pt\hrule\vskip3pt\textbf{}א. פיזית/מגלה ב. מנהלית/מתקנת (מכינה לשחזור) ג. לוגית/מונעת ד. מנהלית/מונעת ה. לוגית או מנהלית/מתקנת \enLR{(recovery)} ו. פיזית/מונעת\end{exbox}

\begin{exbox}\boxtitle{תרגיל \enLR{2} – טיפול בסיכון}\par  לכל מצב בחרו: הפחתה / העברה / קבלה / בידוד.\newline  א. שרת ישן הכרחי שאי אפשר לעדכן. ב. רכשו ביטוח סייבר. ג. סיכון נמוך של הדפסה בזבזנית – מתעלמים. ד. מתקינים \enLR{IPS} מול האינטרנט. \par\vskip2pt\hrule\vskip3pt\textbf{}א. בידוד (רשת נפרדת) ב. העברה ג. קבלה ד. הפחתה\end{exbox}

\begin{exbox}\boxtitle{תרגיל \enLR{3} – \enLR{RTO/RPO}}\par  חנות אונליין מגבה כל \enLR{6} שעות ולוקח לה \enLR{2} שעות לשחזר. אירע כשל ב-\enLR{13:00}, הגיבוי האחרון ב-\enLR{12:00.} א. כמה נתונים אבדו \enLR{(RPO} בפועל)? ב. מתי חוזרים לפעול \enLR{(RTO)}? ג. אם החברה דורשת \enLR{RPO} של שעה – מה משנים? \par\vskip2pt\hrule\vskip3pt\textbf{}א. שעה \enLR{(12:00}–\enLR{13:00).} ב. \textasciitilde{}\enLR{15:00 (2} שעות שחזור). ג. מגבים כל שעה (או רפליקציה רציפה) כדי לצמצם את חלון האובדן.\end{exbox}

\begin{exbox}\boxtitle{תרגיל \enLR{4} – דיון אתי}\par  תלמיד גילה שאתר בית הספר חושף ציונים דרך שינוי כתובת ב-\enLR{URL.} מה עליו לעשות (ומה אסור)? נמקו לפי חוק המחשבים. \par\vskip2pt\hrule\vskip3pt\textbf{}מותר/ראוי: לדווח למורה/מנהל הרשת מיד, לתעד בלי לגשת לנתונים של אחרים. אסור: להוריד/לשנות ציונים, לשתף, "לבדוק עד כמה זה גרוע" בגישה נרחבת – זו חדירה לחומר מחשב שלא כדין. \enLR{Responsible disclosure.}\end{exbox}

\begin{exbox}\boxtitle{תרגיל \enLR{5} – פרויקט מסכם}\par  בחרו ארגון דמיוני (מרפאה / חנות / בית ספר). כתבו מסמך של עמוד: \enLR{3} נכסים קריטיים, איום ופגיעות לכל אחד, ובקרה מכל פרק בקורס שתגן עליו. (זה מחבר את כל \enLR{9} הפרקים.) \par\vskip2pt\hrule\vskip3pt\textbf{}דוגמה למרפאה: נכס=תיקים רפואיים → הצפנה \enLR{(7)+ACL (4)+AAA (3)}; נכס=זמינות מערכת → \enLR{IPS (5)}+גיבוי/\enLR{DRP (9)+storm control (6)}; נכס=גישה מרחוק של רופאים → \enLR{VPN (8)+MFA (3)}+נתב מוקשח \enLR{(2).} מודעות עובדים נגד \enLR{phishing (1).}\end{exbox}

\sectionheading{\enLR{9.13} שאלות בסגנון בגרות}

\begin{exambox}\boxtitle{שאלה \enLR{1}}\par  התאימו אמצעי אבטחה (מנהלי/פיזי/לוגי) לתיאור: \enLR{(1)} הגנה על חדר שרתים; \enLR{(2)} כתיבת מדיניות ונהלים; \enLR{(3)} שימוש ברשימות גישה וסיסמאות. \par\vskip2pt\hrule\vskip3pt\textbf{}\enLR{(1)} פיזי \enLR{(2)} מנהלי \enLR{(3)} לוגי\end{exambox}

\begin{exambox}\boxtitle{שאלה \enLR{2}}\par  מה מאפשרת תוכנת \enLR{Nmap}? \enLR{1.} זיהוי וירוסים \enLR{2.} זיהוי סוסים טרויאנים \enLR{3.} גילוי, ניתוח ובקרה ברשתות תקשורת \enLR{4.} חסימת פורטים \par\vskip2pt\hrule\vskip3pt\textbf{}\textbf{\enLR{3}} – גילוי מארחים, סריקת פורטים וזיהוי שירותים. אינה אנטי-וירוס ואינה חוסמת.\end{exambox}

\begin{exambox}\boxtitle{שאלה \enLR{3}}\par  מהן ארבע הדרכים לטיפול בסיכון? תנו דוגמה לכל אחת. \par\vskip2pt\hrule\vskip3pt\textbf{}הפחתה \enLR{(IPS)}, העברה (ביטוח), קבלה (חיים עם סיכון נמוך), בידוד/הימנעות (רשת נפרדת/הפסקת שירות).\end{exambox}

\begin{exambox}\boxtitle{שאלה \enLR{4}}\par  מה ההבדל בין \enLR{BCP} ל-\enLR{DRP}? \par\vskip2pt\hrule\vskip3pt\textbf{}\enLR{BCP} – המשכיות תפקוד כלל העסק בזמן אסון (אנשים, תהליכים). \enLR{DRP} – שחזור מערכות ה-\enLR{IT} והנתונים.\end{exambox}

\begin{exambox}\boxtitle{שאלה \enLR{5}}\par  עובד הריץ \enLR{Nmap} על רשת חברה מתחרה "כדי לבדוק". האם זו עבירה? נמקו. \par\vskip2pt\hrule\vskip3pt\textbf{}כן – סריקה של רשת שאינה שלו וללא אישור מהווה ניסיון חדירה/גישה לפי חוק המחשבים התשנ"ה-\enLR{1995.} גם "רק בדיקה" אסורה ללא היקף ואישור בכתב.\end{exambox}

\begin{exambox}\boxtitle{שאלה \enLR{6}}\par  מדוע עדיף לשלב אבטחה כבר בשלב העיצוב של \enLR{SDLC} ולא בסוף? \par\vskip2pt\hrule\vskip3pt\textbf{}תיקון פגיעות בשלב מוקדם זול ויעיל בהרבה; "\enLR{Security by Design"} מונע פגמים מבניים שקשה וריקם לתקן במערכת גמורה.\end{exambox}

\sectionheading{\enLR{9.14} מעבדה \enLR{(3} שעות)}

\begin{enumerate}
\item \textbf{\enLR{Nmap} במעבדה מבודדת בלבד} (רשת סגורה, לא רשת בית הספר, באישור): מפו מארחים, סרקו פורטים, זהו שירותים וגרסאות. דונו בכל פלט: מה תוקף לומד ומה מגן לומד.
\item סרקו \enLR{VM} עם שירות ישן; זהו פורט/גרסה פגיעים; הציעו טיפול \enLR{(patch}/כיבוי/\enLR{ACL)} – מדגים את מעגל ניהול הפגיעויות.
\item \textbf{תכנון (על נייר):} קבלו טופולוגיה של ארגון וסמנו את כל הבקרות מ-\enLR{9} הפרקים במקומן (נתב מוקשח, \enLR{AAA, ACL/ZPF, IPS}, מתג מאובטח, \enLR{VPN}, גיבוי).
\item \textbf{סימולציית אירוע:} "מחשב עובד נדבק בכופרה" – עברו בכיתה את \enLR{6} שלבי ה-\enLR{Incident Response.}
\item \textbf{פרויקט מסכם} (תרגיל \enLR{5)} כמצגת קבוצתית – חזרה מצוינת לבגרות.
\end{enumerate}

\sectionheading{בנק שאלות שתלמידים שואלים – ותשובות מוכנות}

שאלות נפוצות בפרק ניהול האבטחה, עם תשובות מוכנות.\par

\textbf{ש: "אם יש לנו את כל הכלים (חומת אש, \enLR{IPS, VPN)}, למה בכלל צריך 'ניהול'?"}\newline ת: כי כלים בלי מדיניות ותהליך לא מספיקים. מי מחליט מה לחסום? מי בודק שהגיבוי עובד? מה עושים כשקורה אירוע? אבטחה היא תהליך ניהולי מתמשך, לא מוצר שקונים פעם אחת.\par

\textbf{ש: "אפשר להגיע לאבטחה מושלמת \enLR{(100\%)}?"}\newline ת: לא, וזו לא המטרה. אבטחה מושלמת עולה אינסוף ומשתקת עבודה. המטרה היא \underline{ניהול סיכונים}: להשקיע היכן שהסיכון (הסתברות × נזק) הכי גבוה, ולקבל סיכונים קטנים במודע.\par

\textbf{ש: "מה ההבדל בין בקרה מנהלית, פיזית ולוגית?"}\newline ת: \textbf{מנהלית} = מדיניות, נהלים, הדרכות (מה שאנשים עושים). \textbf{פיזית} = מנעולים, מצלמות, חדר שרתים נעול (הגנה על החומרה). \textbf{לוגית/טכנית} = סיסמאות, הצפנה, חומת אש, \enLR{ACL} (הגנה בתוכנה). שאלה נפוצה בבגרות.\par

\textbf{ש: "מה מאפשרת תוכנת \enLR{Nmap}? האם היא כלי תקיפה או הגנה?"}\newline ת: \enLR{Nmap} מגלה מארחים חיים, סורק פורטים פתוחים ומזהה שירותים/גרסאות. היא \underline{דו-שימושית}: תוקף משתמש בה לסיור (פרק \enLR{1)}, ומגן משתמש בה כדי לגלות מה חשוף ברשת שלו לפני שהתוקף יעשה זאת. היא \textbf{אינה} אנטי-וירוס ואינה חוסמת פורטים.\par

\textbf{ש: "מה ההבדל בין סריקת פגיעויות \enLR{(Nessus)} לבדיקת חדירה \enLR{(pentest)}?"}\newline ת: סורק פגיעויות \enLR{(Nessus)} הוא כלי \underline{אוטומטי} שמדווח על חולשות ידועות. בדיקת חדירה היא \underline{אדם} שמנסה בפועל לנצל אותן ולפרוץ, "לחשוב כמו תוקף", ולהוכיח את הנזק. הסריקה מוצאת דלתות; ה-\enLR{pentest} בודק אילו באמת נפתחות.\par

\textbf{ש: "מה ההבדל בין \enLR{BCP} ל-\enLR{DRP}?"}\newline ת: \textbf{\enLR{BCP}} (המשכיות עסקית) = איך העסק \underline{ממשיך לתפקד} בזמן אסון (אנשים, תהליכים, מבנים). \textbf{\enLR{DRP}} (התאוששות מאסון) = איך \underline{משחזרים את מערכות ה-\enLR{IT}} והנתונים. \enLR{DRP} הוא חלק טכני בתוך ה-\enLR{BCP} הרחב.\par

\textbf{ש: "מהו כלל \enLR{3-2-1} לגיבויים?"}\newline ת: \enLR{3} עותקים של המידע, על \enLR{2} סוגי מדיה שונים, כשלפחות \enLR{1} נמצא \underline{מחוץ לאתר} \enLR{(off-site}/ענן). זה קריטי נגד כופרה: אם המתקפה הצפינה את הרשת, יש עותק מנותק לשחזר ממנו. וזכרו – \underline{גיבוי שלא נבדק שחזור ממנו אינו גיבוי}.\par

\textbf{ש: "תלמיד מצא חולשה באתר בית הספר. מה מותר לו לעשות?"}\newline ת: לדווח מיד למורה/מנהל הרשת, ולתעד בלי לגשת למידע של אחרים. \underline{אסור} להוריד/לשנות נתונים, לשתף, או "לבדוק עד כמה זה חמור" בגישה נרחבת – זו חדירה לחומר מחשב שלא כדין (חוק המחשבים \enLR{1995).} זה נקרא \enLR{Responsible Disclosure.}\par

\textbf{ש: "למה לשלב אבטחה כבר בשלב התכנון \enLR{(SDLC)} ולא בסוף?"}\newline ת: כי תיקון חולשה שהתגלתה בסוף עולה פי כמה, ולפעמים דורש לבנות מחדש. "אבטחה מהעיצוב" \enLR{(Security by Design)} זולה ויעילה בהרבה מטלאי מאוחר.\par

\sectionheading{מילון מונחים – הגדרות}

הגדרות תמציתיות של המונחים המרכזיים בפרק. שימושי גם כדף עזר לבחינה.\par

\begin{xltabular}{\linewidth}{|>{\setRTL\raggedleft\arraybackslash}X|>{\setRTL\raggedleft\arraybackslash}X|}
\hline
\rowcolor{hdrbg}\textbf{הגדרה} & \textbf{מונח} \\
\hline
\endhead
תהליך זיהוי, הערכה וטיפול בסיכוני אבטחה. & \textbf{ניהול סיכונים} \\
\hline
= איום × פגיעות × ערך הנכס. & \textbf{סיכון \enLR{(Risk)}} \\
\hline
הפחתה, העברה (ביטוח), קבלה, או בידוד/הימנעות. & \textbf{טיפול בסיכון} \\
\hline
מסמך הקובע מה מותר ואסור בארגון (לא איך). & \textbf{מדיניות \enLR{(Policy)}} \\
\hline
הוראות צעד-אחר-צעד ליישום המדיניות. & \textbf{נוהל \enLR{(Procedure)}} \\
\hline
מדיניות, נהלים, הדרכות, בדיקות רקע. & \textbf{בקרה מנהלית} \\
\hline
מנעולים, מצלמות, חדר שרתים נעול. & \textbf{בקרה פיזית} \\
\hline
סיסמאות, הצפנה, חומת אש, \enLR{ACL.} & \textbf{בקרה לוגית} \\
\hline
שכבות הגנה עצמאיות; כשל באחת אינו מפיל הכול. & \textbf{הגנה לעומק \enLR{(Defense in Depth)}} \\
\hline
מתן בדיוק ההרשאות הנדרשות, לא יותר. & \textbf{מינימום הרשאות} \\
\hline
פעולה קריטית דורשת שני אנשים; מונע הונאה. & \textbf{הפרדת תפקידים} \\
\hline
'לעולם אל תסמוך, תמיד אמת' – אין אזור פנימי מהימן. & \textbf{\enLR{Zero Trust}} \\
\hline
סריקה אוטומטית לזיהוי פגיעויות ידועות \enLR{(Nessus).} & \textbf{\enLR{Vulnerability Assessment}} \\
\hline
ניסיון פריצה מבוקר, באישור, כדי לאתר חולשות. & \textbf{\enLR{Penetration Test}} \\
\hline
כלי לגילוי מארחים, סריקת פורטים וזיהוי שירותים. & \textbf{\enLR{Nmap}} \\
\hline
סורק פגיעויות המדווח על חולשות ידועות \enLR{(CVE).} & \textbf{\enLR{Nessus}} \\
\hline
בדיקה שיטתית שהבקרות עובדות כמתוכנן. & \textbf{\enLR{ST\&E}} \\
\hline
תוכנית המשכיות עסקית – שהעסק ימשיך לתפקד בזמן אסון. & \textbf{\enLR{BCP}} \\
\hline
תוכנית התאוששות מאסון – שחזור מערכות ה-\enLR{IT} והנתונים. & \textbf{\enLR{DRP}} \\
\hline
זמן ההשבתה המרבי המותר \enLR{(RTO)} וכמות הנתונים המרבית לאיבוד \enLR{(RPO).} & \textbf{\enLR{RTO / RPO}} \\
\hline
\enLR{3} עותקי גיבוי, \enLR{2} מדיות, \enLR{1} מחוץ לאתר. & \textbf{כלל \enLR{3-2-1}} \\
\hline
מחזור חיי פיתוח מערכת; משלבים אבטחה כבר בעיצוב. & \textbf{\enLR{SDLC}} \\
\hline
החוק האוסר חדירה ושיבוש של חומר מחשב שלא כדין. & \textbf{חוק המחשבים התשנ"ה-\enLR{1995}} \\
\hline
\end{xltabular}

\end{document}
